\hypertarget{nace_kartlar_13-15}{%
\chapter{Faz 3, parça 2 --- Tekstil, giyim ve deri: Gümrük Birliği'nden çıkışın sınıf düzeyinde analizi (NACE 13--15)}\label{nace_kartlar_13-15}}

\emph{Durum: Faz 3 / 9, parça 2 / 16; 18 Ekim 2026. Önceki parça: \belgebag{nace_kartlar_10-12}{}{\kod{nace\_\allowbreak{}kartlar\_\allowbreak{}10-\allowbreak{}12.\allowbreak{}md}}. Önceki fazlar: \belgebag{01_teshis}{}{\kod{01\_\allowbreak{}teshis.\allowbreak{}md}}, \belgebag{02a_politika_uzayi}{}{\kod{02a\_\allowbreak{}politika\_\allowbreak{}uzayi.\allowbreak{}md}}, \belgebag{02b_secenekler_hukuk}{}{\kod{02b\_\allowbreak{}secenekler\_\allowbreak{}hukuk.\allowbreak{}md}}, \belgebag{02c_hazirlik}{}{\kod{02c\_\allowbreak{}hazirlik.\allowbreak{}md}}. Kümülatif defter: \belgebag{defter}{}{\kod{defter.\allowbreak{}md}}. Bütün 230 imalat sınıfının ana matrisi: \href{../veri/nace_matris.csv}{\kod{veri/\allowbreak{}nace\_\allowbreak{}matris.\allowbreak{}csv}}.}

\hypertarget{nace_kartlar_13-15-yuxf6netici-uxf6zeti}{%
\section{Yönetici özeti}\label{nace_kartlar_13-15-yuxf6netici-uxf6zeti}}

Tekstil, giyim ve deri, Gümrük Birliği'nin Türkiye'deki en büyük işvereni. Üç bölümün 21 sınıfında 1.238 bin kişi çalışıyor; bu, imalat istihdamının \%24,6 kadarı. Yıllık ihracat 35,9 milyar dolar ve bunun \%48 kadarı AB'ye gidiyor; giyimde bu oran \%59. Gıdanın tersine bu ticaretin tamamı Gümrük Birliği'nin sanayi rejiminde. 2025'te AB'ye satılan 13.132 milyon avroluk tekstil, giyim ve deri ürünü bugün tarifesiz ve menşe belgesi aranmadan, yalnızca A.TR dolaşım belgesiyle giriyor. Anlaşmasız çıkışta AB bu ürünlere ortalama \%10,2 tarife uygular. Giyimde ortalama \%11,8, yani AB'nin giyimdeki MFN tarifesinin tamamı.

Bu yüzden çıkışın imalattaki en büyük istihdam maliyeti bu parçada. Anlaşmasız çıkışta (S4) üç bölüm T+5'te medyan senaryoda 39.397 kişi kaybediyor; dağılımın alt ucunda (P10) kayıp 60.851 kişi. Bu, aynı senaryoda imalatın bütününde beklenen istihdam kaybının \%85 kadarı. Kaybın büyük kısmı giyimde (29.021 kişi) ve dört sınıfta toplanıyor: dış giyim (11.910 kişi), iç giyim ve tişört (9.401), triko (3.773) ve çorap (2.277). AB'ye ihracat yılda 4.104 milyon dolar azalıyor, cari denge 2.634 milyon dolar kötüleşiyor. Önerilen bitiş durumu olan derin ve kapsamlı serbest ticaret anlaşmasında (S2) kayıp belirgin biçimde küçülüyor: AB'ye ihracat 1.854 milyon dolar, istihdam 14.419 kişi azalıyor. Bu parçada S2 ile S4 arasındaki fark, yani iyi bir anlaşmanın değeri, çalışmanın şimdiye kadar ölçtüğü en büyük fark.

İkinci bulgu, Türk tekstilinin AB'ye yalnızca doğrudan ulaşmadığı. Fas, Tunus, Mısır, Sırbistan ve Kuzey Makedonya'daki konfeksiyoncular Türk kumaşını dikip Pan-Avrupa-Akdeniz (PEM) menşe kümülasyonuyla AB'ye gümrüksüz satıyor. 2024'te Türkiye'nin bu ülkelere sattığı kumaş ve teknik tekstil 1.736 milyon dolar. Ortakların giyim ihracatındaki AB payıyla ağırlıklandırıldığında bunun 1.131 milyon doları AB'ye giden giysiye dönüşüyor; Fas tek başına 613 milyon dolar. PEM'de kümülasyon, girdinin menşe ülkesiyle varış ülkesi arasında bir tercihli ticaret anlaşması olmasını şart koşuyor \etiketV{}. Anlaşmasız çıkışta bu şart düşüyor ve Türk kumaşı bu zincirde menşesiz sayılıyor. Fas'taki bir konfeksiyoncu AB'ye tercihli satmak için kumaşı başka yerden almak zorunda kalıyor. Model bu kanalı ayrıca hesaplıyor. S4'te Türkiye'nin AB dışına kumaş ihracatı 673 milyon dolar azalıyor ve 4.750 kişilik ek istihdam kaybı doğuyor. Tekstil bölümünün (13) S4'teki istihdam kaybı 10.042 kişi ve bunun 4.750 kişisi bu kanaldan geliyor. S2'de ve S3'te kanal kapanmıyor, çünkü yeni anlaşma PEM'e atıf yaptığı sürece kümülasyon sürüyor. Bu yüzden menşe protokolünde PEM'e dinamik atıf ve kümülasyon bildiriminin çıkış gününde yayımlanması (E-13-01) bu parçanın ilk müzakere talebi.

Üçüncü bulgu, S2'nin menşe yükü. Gümrük Birliği'nde menşe sorulmuyor. Serbest ticaret anlaşmasında ise giyim, tercihli tarifeden ancak revize PEM kurallarındaki çift dönüşüm şartını karşılarsa yararlanıyor \etiketV{}. Örme giyimde kumaşın örülmesi ve dikilmesi, dokuma giyimde kumaşın dokunması (ya da Türkiye'de baskılanması) ve dikilmesi PEM bölgesinde yapılmalı. Türkiye'nin kumaş kapasitesi büyük olduğu için ihracatın çoğu bu kuralı rahatça karşılıyor. Asya kumaşıyla üretilen giyim ise karşılamıyor; bunun ihracattaki payı \%10--15 dolayında tahmin ediliyor \etiketH{}. Dış giyimde S2'nin AB'ye ihracat kaybı (453 milyon dolar), S4'teki kaybın \%45 kadarı. Menşeye ilişkin ikinci kural, geri ödeme yasağı, yeni bir yük getirmiyor. Dahilde İşleme Rejimi'yle ithal edilen PEM dışı girdiye, ürün AB'ye A.TR ile satıldığında bugün de telafi edici vergi alınıyor \etiketV{}. Önerilen çözüm üç parçalı: kumaş menşei ve izlenebilirlik için ortak bir veri altyapısı (E-13-02), Türkiye'de baskı kapasitesinin genişletilmesi ve S2 müzakeresinde PEM dışı kumaşla üretilen giyim için bir menşe kotası (E-13-06).

Rakiplerin konumu tarifenin ağırlığını artırıyor. AB pazarındaki büyük tedarikçilerin çoğu tercihli erişime sahip. Bangladeş ve Kamboçya AB'nin en az gelişmiş ülkelere tanıdığı ``Silah Dışında Her Şey'' (EBA) rejimiyle, Pakistan GSP+ ile, Fas, Tunus ve Vietnam serbest ticaret anlaşmalarıyla sıfır ya da sıfıra yakın tarifeyle giriyor. İç giyim ve tişörtte AB'nin AB dışından tedarikinin \%32 kadarı Bangladeş'ten. Bangladeş'in en az gelişmiş ülke statüsünden çıkışı Kasım 2026'ya planlanmış olsa da EBA tercihleri üç yıllık geçiş süresince sürüyor \etiketV{}; bu tablo 2029'dan önce değişmiyor. Anlaşmasız çıkışta Türkiye, AB'nin giyim pazarında Çin'le birlikte tam tarife ödeyen büyük tedarikçi konumuna düşer. Elinde kalan üstünlük hız ve yakınlık: kısa termin ve küçük partili hızlı moda siparişleri.

Dördüncü bulgu ithalat tarafında. Faz 2A'da önerilen c2 rejiminde, yerli üretimi güçlü satırlarda ek gümrük vergisi (İGV) AB menşeli ürüne de uygulanıyor. Giyim ve ayakkabı bu satırlar arasında, dolayısıyla AB'den gelen dış giyimin tarifesi S4'te \%37,2 oluyor. Bu, bölümün TÜFE'ye birinci tur etkisini S4'te 15,6 baz puana çıkarıyor; kümülatif 50 baz puanlık bütçenin büyük bir dilimi tek bir parçada harcanıyor. Faz 2A'da tasarlanan vergi takası bu etkiyi tasarım gereği siliyor: KDV giyimde \%10'dan \%8,7'ye, ayakkabıda \%5,6'ya iniyor ve indirim bütün menşelere eşit uygulanıyor (E-13-04). Takasla S4'te parçanın enflasyon yükü 3,6 baz puana iniyor. Takasın net mali maliyeti, AB menşeli ürünlere gelen gümrük geliri düşüldükten sonra yılda yaklaşık 17,5 milyar TL \etiketT{}.

Bu parçada çıkış, yapısal bir daralmanın üstüne geliyor. Giyim eşyası imalatında sigortalı ücretli çalışan sayısı 2025'te bir yılda yaklaşık 86 bin azaldı \etiketV{}. Maliyet baskısı altındaki firmaların bir kısmı üretimini Mısır'a taşıyor \etiketV{}. Bu daralma Gümrük Birliği'nden bağımsız. Ancak anlaşmasız çıkış bu kaymayı hızlandırır, çünkü Mısır'a giden bir Türk firması hem AB--Mısır anlaşmasından hem PEM kümülasyonundan yararlanırken Türkiye'de kalan firma \%12 öder. Kaybın bölgesel dağılımı da bu yüzden önemli. İstanbul en büyük mutlak kaybı yaşıyor (12.209 kişi). Ama Mardin, Batman, Şırnak ve Siirt'te tekstil, giyim ve deri imalat istihdamının \%78 kadarını oluşturuyor; bölgenin S4 kaybı 3.388 kişi ve alternatif sanayi istihdamı sınırlı.

Duruşlar sektörün iki yüzünü yansıtıyor. 21 sınıfın 11 tanesi için ``Geliştir'', 2 tanesi için ``Dönüştür'', 3 tanesi için ``Yönetilen uyum'', 5 tanesi için ``Koru'' öneriliyor. Kumaş, halı, ev tekstili ve teknik tekstil küresel rekabet gücü yüksek ihracatçılar. Dış giyim, tasarım ve hızlı tedarik hizmetine dönüşmesi gereken büyük bir işveren. İç giyim, tişört ve trikoda ise standart ürünlerin bir kısmının Asya'ya kaybedilmesi kaçınılmaz görünüyor ve önerilen politika bu uyumu yönetmek. Paketin yıllık mali maliyeti S2 yolunda 0,82 milyar TL, S4 yolunda 21,07 milyar TL; S4'teki tutarın çoğu vergi takası. Parça kümülatif enflasyon bütçesinin S2 yolunda 1,7, S4 yolunda önlemlerle 3,6 baz puanını kullanıyor.

İki olumsuz bulgu açıkça yazılmalı. Birincisi, S4'te giyimdeki kayıp önlemlerle önlenemiyor, ancak yumuşatılabiliyor. DTÖ kuralları sanayi ürünlerinde ihracat sübvansiyonunu yasaklıyor. Kullanılabilecek araçlar, yani kısa çalışma ödeneği, işgücü uyum programları, dönüşüm desteği ve pazar çeşitlendirme, işin yerini değiştirmiyor; geçişi kolaylaştırıyor. Giyim için S2'nin yerini tutacak bir tek taraflı politika yok. Çıkış kararı, bu sektörün kaderini büyük ölçüde müzakerenin başarısına bağlıyor. İkincisi, kümülasyon kanalının büyüklüğü bir üst sınır tahmini. Türk kumaşının ortak ülkelerde AB'ye giden üretimde orantılı kullanıldığı ve kaybın \%30--90'ının telafi edilemeyeceği varsayıldı \etiketH{}. Gerçek değer, ihracatçı birliklerinin firma verisiyle kesinleşmeli (VB-37).

\Needspace{7\baselineskip}

\textbf{Tablo Ö.1 --- Çıkışın 13--15'e etkisi (T+5'te yıllık; bugüne göre fark)}

\nopagebreak[4]

\begingroup\gbtablo\footnotesize

\begin{longtable}[]{@{}
  >{\raggedright\arraybackslash}p{(\columnwidth - 8\tabcolsep) * \real{0.5036}}
  >{\raggedleft\arraybackslash}p{(\columnwidth - 8\tabcolsep) * \real{0.0797}}
  >{\raggedleft\arraybackslash}p{(\columnwidth - 8\tabcolsep) * \real{0.1602}}
  >{\raggedleft\arraybackslash}p{(\columnwidth - 8\tabcolsep) * \real{0.0797}}
  >{\raggedleft\arraybackslash}p{(\columnwidth - 8\tabcolsep) * \real{0.1768}}@{}}
\toprule\noalign{}
\begin{minipage}[b]{\linewidth}\raggedright
\textbf{Gösterge}
\end{minipage} & \begin{minipage}[b]{\linewidth}\raggedleft
\textbf{S2: P50}
\end{minipage} & \begin{minipage}[b]{\linewidth}\raggedleft
\textbf{S2: P10 / P90}
\end{minipage} & \begin{minipage}[b]{\linewidth}\raggedleft
\textbf{S4: P50}
\end{minipage} & \begin{minipage}[b]{\linewidth}\raggedleft
\textbf{S4: P10 / P90}
\end{minipage} \\
\midrule\noalign{}
\endhead
\bottomrule\noalign{}
\endlastfoot
AB'ye ihracat (\$ mn) & −1.854 & −4.010 / +347 & −4.104 & −6.394 / −1.698 \\
PEM kümülasyonu kaybı: AB dışı ihracat (\$ mn) & 0 & 0 / 0 & −673 & −858 / −483 \\
İstihdam (kişi) & −14.419 & −30.880 / +1.234 & −39.397 & −60.851 / −20.019 \\
Cari denge (\$ mn) & −893 & −1.812 / +56 & −2.634 & −3.841 / −1.473 \\
TÜFE, birinci tur (bp) & 1,7 & 1,4 / 1,9 & 15,6 & 13,4 / 17,7 \\
TÜFE, önerilen önlemlerle (bp) & 1,7 & --- & 3,6 & --- \\
Önlemlerin mali maliyeti (mlr TL/yıl) & 0,82 & --- & 21,07 & --- \\
\end{longtable}

\endgroup

\emph{Kaynak:} \etiketT{} f3e Monte Carlo (N = 2.000; yüzdelikler bölüm toplamının dağılımından). S2: derin ve kapsamlı STA; S4: anlaşmasız çıkış ve tek taraflı paket. PEM kümülasyonu kaybı: Türk kumaşının PEM ortaklarında AB'ye giden giyimde menşe kazandırmaması (yalnız S4; ihracat, istihdam ve cari denge satırlarına dahil). Önlemler ve mali maliyet Bölüm 4'teki eylemlerden.

\hypertarget{nace_kartlar_13-15-1-bu-paruxe7anux131n-sorusu-ve-okuma-kux131lavuzu}{%
\section{1. Bu parçanın sorusu ve okuma kılavuzu}\label{nace_kartlar_13-15-1-bu-paruxe7anux131n-sorusu-ve-okuma-kux131lavuzu}}

Faz 3, Gümrük Birliği'nden çıkışın imalat sanayisine etkisini NACE Rev. 2'nin dört basamaklı sınıfları düzeyinde ölçüyor ve her sınıf için bir strateji öneriyor. Bu ikinci parça tekstil ürünleri (13), giyim eşyaları (14) ile deri ve ilgili ürünleri (15) kapsıyor. Sorular ilk parçadakiyle aynı. Çıkış bu sınıfın ihracatını, üretimini, istihdamını, cari dengeye katkısını ve tüketici fiyatlarını hangi kanaldan ve ne ölçüde etkiler? Hangi yetkinlikler korunmalı, hangileri geliştirilmeli? AB ile ve üçüncü ülkelerle neyin müzakeresi yapılmalı? Hangi önlemi kim, hangi hukuki araçla ve ne zaman almalı?

Senaryolar ilk parçadaki gibi. \textbf{S2}, AB ile derin ve kapsamlı bir serbest ticaret anlaşması; sanayi ürünlerinde tarifeler menşe şartıyla sıfır ve çalışmanın önerdiği bitiş durumu bu. \textbf{S3}, AB'nin üye devletlerin onayı olmadan imzalayabileceği temel bir ticaret anlaşması; S2'ye giden yolda ara durum. \textbf{S4}, anlaşmasız çıkış ve Türkiye'nin tek taraflı paketi: AB, Türk mallarına MFN tarifesini uyguluyor; Türkiye de AB mallarına MFN tarifesini ve seçici ek gümrük vergisini (c2 rejimi) uyguluyor, girdilerde tarifeleri sıfırlıyor ve tüketici fiyatlarını vergi takasıyla dengeliyor. AB'nin tutumu hasmane ve üyelik perspektifi yok. İki takvim var: T2028'de çıkış tarihi T = 1 Ocak 2028, T2029'da T = 1 Ocak 2029. Yöntem, veri kaynakları ve kanıt etiketleri ilk parçanın birinci bölümünde ve Ek B'sinde anlatıldı. Bu parçanın Ek B'si yalnızca yenilikleri veriyor.

Bu parça üç nedenle Faz 3'ün en ağır parçalarından biri. Birincisi ölçek: üç bölüm imalat istihdamının dörtte birini veriyor ve neredeyse bütün sınıflarında emek yoğun. İkincisi bağımlılık: ihracatın yarıya yakını AB'ye gidiyor ve bu ticaret tümüyle Gümrük Birliği'nin sanayi rejiminde. Gıdadaki gibi çıkışın dokunmadığı bir tarım rejimi yok. Üçüncüsü menşe: tekstil ve giyim, PEM menşe kurallarının en ayrıntılı ve en katı olduğu fasıllar. Serbest ticaret anlaşmasının değeri burada tarifeden çok bu kuralların nasıl yazıldığına bağlı.

Bu parçada modele iki ek yapıldı. Birincisi PEM kümülasyonu kanalı. Türkiye'nin Fas, Tunus, Mısır ve Batı Balkanlar'a sattığı kumaş, oralarda dikilip AB'ye tercihli giriyor. Anlaşmasız çıkışta bu zincir kopuyor. Kanalın büyüklüğü ve varsayımları Bölüm 3.2'de ve Ek B'de anlatılıyor; kanal ilk parçanın sınıflarında sıfır. İkincisi bölgesel dağılım. Giyim istihdamı birkaç bölgede yoğunlaştığı ve bunların bazılarında sanayinin neredeyse tek kolu olduğu için S4'ün istihdam etkisi İBBS-2 bölgelerine dağıtılarak gösteriliyor (Bölüm 4.5).

Belge ilk parçayla aynı sırada ilerliyor. İkinci bölüm sektörün Gümrük Birliği içindeki yerini, üçüncü bölüm çıkışın kanallarını anlatıyor. Dördüncü bölüm duruşları, eylem paketini, müzakere pozisyonunu, erken uyarı göstergelerini ve bölgesel etkileri bir araya getiriyor. Beşinci bölüm sınıfları tek tek inceliyor; etkisi sınırlı altı sınıf tek bir tabloda toplanıyor. Altıncı bölüm kümülatif defteri ve kalite kontrollerini veriyor. Ekler ana matris satırlarını (Ek A), yöntemdeki yenilikleri (Ek B), sınıf veri sayfalarını (Ek C), önceki parçalara ve fazlara düzeltmeleri (Ek D) ve varsayım ile veri boşluğu kayıtlarını (Ek E) içeriyor. On beş alanlı sınıf kartı yine üç yerde tamamlanıyor: nitel alanlar Bölüm 5'te, sayısal alanlar Ek C'de, eylemler ve tetikler Bölüm 4'te.

\hypertarget{nace_kartlar_13-15-2-sektuxf6ruxfcn-guxfcmruxfck-birliux11fi-iuxe7indeki-yeri}{%
\section{2. Sektörün Gümrük Birliği içindeki yeri}\label{nace_kartlar_13-15-2-sektuxf6ruxfcn-guxfcmruxfck-birliux11fi-iuxe7indeki-yeri}}

\Needspace{7\baselineskip}

\textbf{Tablo 2.1 --- Bölümlerin büyüklüğü ve dış ticareti (2024)}

\nopagebreak[4]

\begingroup\gbtablo\footnotesize

\begin{longtable}[]{@{}
  >{\raggedright\arraybackslash}p{(\columnwidth - 16\tabcolsep) * \real{0.1400}}
  >{\raggedleft\arraybackslash}p{(\columnwidth - 16\tabcolsep) * \real{0.0784}}
  >{\raggedleft\arraybackslash}p{(\columnwidth - 16\tabcolsep) * \real{0.1043}}
  >{\raggedleft\arraybackslash}p{(\columnwidth - 16\tabcolsep) * \real{0.1278}}
  >{\raggedleft\arraybackslash}p{(\columnwidth - 16\tabcolsep) * \real{0.1177}}
  >{\raggedleft\arraybackslash}p{(\columnwidth - 16\tabcolsep) * \real{0.1356}}
  >{\raggedleft\arraybackslash}p{(\columnwidth - 16\tabcolsep) * \real{0.0801}}
  >{\raggedleft\arraybackslash}p{(\columnwidth - 16\tabcolsep) * \real{0.1136}}
  >{\raggedleft\arraybackslash}p{(\columnwidth - 16\tabcolsep) * \real{0.1024}}@{}}
\toprule\noalign{}
\begin{minipage}[b]{\linewidth}\raggedright
\textbf{Bölüm}
\end{minipage} & \begin{minipage}[b]{\linewidth}\raggedleft
\textbf{Sınıf}
\end{minipage} & \begin{minipage}[b]{\linewidth}\raggedleft
\textbf{Çalışan (bin)}
\end{minipage} & \begin{minipage}[b]{\linewidth}\raggedleft
\textbf{Katma değer (\$ mlr)}
\end{minipage} & \begin{minipage}[b]{\linewidth}\raggedleft
\textbf{İhracat (\$ mlr)}
\end{minipage} & \begin{minipage}[b]{\linewidth}\raggedleft
\textbf{AB'ye ihracat (\$ mlr)}
\end{minipage} & \begin{minipage}[b]{\linewidth}\raggedleft
\textbf{AB payı}
\end{minipage} & \begin{minipage}[b]{\linewidth}\raggedleft
\textbf{İthalat (\$ mlr)}
\end{minipage} & \begin{minipage}[b]{\linewidth}\raggedleft
\textbf{Denge (\$ mlr)}
\end{minipage} \\
\midrule\noalign{}
\endhead
\bottomrule\noalign{}
\endlastfoot
13 Tekstil ürünleri & 10 & 493,4 & 12,84 & 14,08 & 4,85 & \%34 & 4,42 & +9,66 \\
14 Giyim eşyaları & 8 & 668,7 & 8,42 & 19,99 & 11,85 & \%59 & 4,02 & +15,97 \\
15 Deri ve ilgili ürünler & 3 & 75,8 & 1,20 & 1,82 & 0,64 & \%35 & 2,53 & −0,71 \\
\textbf{13--15 toplamı} & \textbf{21} & \textbf{1.237,9} & \textbf{22,46} & \textbf{35,89} & \textbf{17,34} & \textbf{\%48} & \textbf{10,97} & \textbf{+24,92} \\
İmalat (C) içindeki payı & 21 / 230 & \%24,6 & \%13,4 & \%14,3 & \%16,4 & --- & \%4,0 & --- \\
\end{longtable}

\endgroup

\emph{Kaynak:} TÜİK Yıllık Sanayi ve Hizmet İstatistikleri 2024 (çalışan, katma değer); BACI HS17 2024 (dış ticaret; Kıbrıs akımları AB dışı, Ek D). \etiketV{}; HS6 → NACE eşlemesi \etiketT{}.

Üç bölümün istihdamını giyim (14) taşıyor: 669 bin çalışan ve 20,0 milyar dolarlık ihracat. Dış giyim (14.13) tek başına 366 bin kişiyle imalatın en büyük işveren sınıfı. Bu sınıfta 41.798 girişim var ve büyük çoğunluğu en fazla dokuz kişinin çalıştığı atölyeler. Onu iç giyim ve tişört (184 bin kişi) izliyor. Tekstilde (13) istihdam daha dengeli dağılıyor: dokuma (96 bin), boya, baskı ve apre hizmetleri (89 bin), ev tekstili (87 bin) ve iplik (86 bin). Deri bölümünün (15) ağırlığı ayakkabıda (51 bin kişi).

Sektör emek yoğun ve verimlilik farkı büyük. Dış giyimde kişi başı katma değer 11,0 bin dolar; imalat ortalaması 33,2 bin dolar. AB27'deki aynı sınıfa göre verimlilik, AB27 = 1 alındığında dış giyimde 0,31, iç giyimde 0,78, dokumada 0,49. Türk giyiminin AB'deki rekabet gücü bu yüzden ücretten değil, hızdan, esneklikten ve kumaştan dikime uzanan yerli tedarik zincirinden geliyor. Avrupalı perakendeciler sezon içi yeniden siparişlerin önemli kısmını Türkiye'ye veriyor, çünkü Türk tedarikçi küçük partiyi iki ila dört haftada teslim edebiliyor \etiketH{}.

Türkiye birçok üründe dünya ihracatının ön sıralarında. 2024'te dokuma halıda dünya ihracatının \%41,5 kadarını yaparak birinci; sentetik tufting halıda (\%14,0), çorapta (\%9,7) ve örme kumaşta (\%6,7) ikinci; denimde (\%9,9) ve kadın pamuklu pantolonunda (\%7,5) üçüncü. Havlu ve bukle mutfak tekstilinde (\%8,9) ve kürk eşyada (\%8,3) dördüncü sırada. AB pazarında da güçlü: AB'nin AB dışından yaptığı ithalat içinde Türkiye'nin payı örme kumaşta \%31,9, çorapta \%25,8, iplikte \%24,3 ve halıda \%21,8; dördünde de Çin'den sonra ikinci. Giyimin büyük sınıflarında Türkiye üçüncü tedarikçi: dış giyimde \%8,7, iç giyim ve tişörtte \%10,7, trikoda \%10,6.

AB'ye ihracat giyimde yoğunlaşıyor. Dış giyim (5.034 milyon dolar), iç giyim ve tişört (3.381), triko (1.784), ev tekstili (1.257), dokuma kumaş (986), iplik (877) ve çorap (845) bölüm toplamının \%82 kadarını yapıyor. Dış giyimde ilk üç AB pazarı Almanya (\%24), Hollanda (\%20) ve İspanya (\%18); Hollanda ve İspanya'nın payı büyük perakendecilerin dağıtım merkezlerinden geliyor. AB dışında en büyük pazar Birleşik Krallık. Birleşik Krallık, çorabın AB dışı ihracatının \%42, trikonun \%27, iç giyimin \%19 kadarını alıyor. Türkiye--Birleşik Krallık STA'sı Gümrük Birliği'nden bağımsız ve bütün senaryolarda sürüyor; çıkışta en değerli yedek pazar bu.

Tekstilin AB dışı pazarları ise başka bir hikâye anlatıyor. Dokuma kumaşın AB dışı ihracatının \%32 kadarı Fas'a, ipliğin \%15 kadarı Mısır'a, örme kumaşın \%12 kadarı yine Fas'a gidiyor. Bu ülkeler Türk kumaşını ve ipliğini işleyip giysi olarak AB'ye (Mısır ağırlıkla ABD'ye) satıyor. Yani Türk tekstilinin AB'ye bağımlılığı, doğrudan AB'ye ihracat payının (\%34) gösterdiğinden büyük. Bu dolaylı bağın çıkışta nasıl koptuğu Bölüm 3.2'nin konusu. Örme kumaşın öteki büyük AB dışı pazarları Belarus (\%22) ve Rusya (\%12); bu pazarlar çıkıştan bağımsız bir siyasi risk taşıyor.

\Needspace{7\baselineskip}

\textbf{Tablo 2.2 --- AB'ye ihracatın rejimi ve S4'te eklenecek AB tarifesi (2025)}

\nopagebreak[4]

\begingroup\gbtablo\footnotesize

\begin{longtable}[]{@{}
  >{\raggedright\arraybackslash}p{(\columnwidth - 12\tabcolsep) * \real{0.1420}}
  >{\raggedleft\arraybackslash}p{(\columnwidth - 12\tabcolsep) * \real{0.1436}}
  >{\raggedleft\arraybackslash}p{(\columnwidth - 12\tabcolsep) * \real{0.1284}}
  >{\raggedleft\arraybackslash}p{(\columnwidth - 12\tabcolsep) * \real{0.1468}}
  >{\raggedleft\arraybackslash}p{(\columnwidth - 12\tabcolsep) * \real{0.1416}}
  >{\raggedleft\arraybackslash}p{(\columnwidth - 12\tabcolsep) * \real{0.1870}}
  >{\raggedleft\arraybackslash}p{(\columnwidth - 12\tabcolsep) * \real{0.1106}}@{}}
\toprule\noalign{}
\begin{minipage}[b]{\linewidth}\raggedright
\textbf{Bölüm}
\end{minipage} & \begin{minipage}[b]{\linewidth}\raggedleft
\textbf{AB'ye ihracat (€ mn)}
\end{minipage} & \begin{minipage}[b]{\linewidth}\raggedleft
\textbf{GB (sanayi) payı}
\end{minipage} & \begin{minipage}[b]{\linewidth}\raggedleft
\textbf{S4'te eklenecek tarife}
\end{minipage} & \begin{minipage}[b]{\linewidth}\raggedleft
\textbf{AB MFN ortalaması}
\end{minipage} & \begin{minipage}[b]{\linewidth}\raggedleft
\textbf{Tarifesi ≥ \%10 olan satırların payı}
\end{minipage} & \begin{minipage}[b]{\linewidth}\raggedleft
\textbf{KN8 satırı}
\end{minipage} \\
\midrule\noalign{}
\endhead
\bottomrule\noalign{}
\endlastfoot
13 Tekstil ürünleri & 4.190 & \%100 & \%7,3 & \%7,3 & \%16 & 718 \\
14 Giyim eşyaları & 8.433 & \%100 & \%11,8 & \%11,8 & \%98 & 368 \\
15 Deri ve ilgili ürünler & 508 & \%100 & \%7,4 & \%7,4 & \%19 & 162 \\
\textbf{13--15 toplamı} & \textbf{13.132} & \textbf{\%100} & \textbf{\%10,2} & \textbf{\%10,2} & \textbf{\%69} & \textbf{1.248} \\
İmalat (C) & 98.030 & \%91 & \%5,6 & \%6,1 & \%29 & 6.930 \\
\end{longtable}

\endgroup

\emph{Kaynak:} Eurostat Comext 2025 (TR→AB27, KN8), AB Kombine Nomenklatürü 2026; \kod{kod/\allowbreak{}f3d\_\allowbreak{}ab\_\allowbreak{}tarife.\allowbreak{}py}. Tarifeler TR→AB ihracatı ağırlıklı \etiketT{}. Bu parçanın sınıflarında tarım ve işlenmiş tarım ürünü rejimi yok; S4'te AB'nin MFN tarifesi bütün ihracata ekleniyor.

Tablo 2.2'nin söylediği yalın: AB'ye satılan 13.132 milyon avroluk ürünün tamamı Gümrük Birliği'nin sanayi rejiminde. Tarım rejimi ya da işlenmiş tarım ürünü yok; dolayısıyla anlaşmasız çıkışta AB'nin MFN tarifesi bütün ihracata tam olarak ekleniyor. Ortalama \%10,2, imalatın genelindeki \%5,6 ortalamanın yaklaşık iki katı. Giyimde neredeyse bütün satırlar \%12'lik MFN'ye tabi. En büyük kalemler kadın ve erkek denim pantolonu (KN 6204 62 31 ve 6203 42 31; AB'ye sırasıyla €508 milyon ve €361 milyon) ile pamuklu tişört (KN 6109 10 00; €838 milyon). Ayakkabıda tarife satıra göre \%3 ile \%17 arasında değişiyor ve ortalama \%9,9. Tekstilde ise iplikte \%3,8, dokuma ve örme kumaşta \%8, ev tekstilinde \%10,0.

İthalat tarafında bölümün 11,0 milyar dolarlık ithalatının yalnızca \%28 kadarı AB'den. İplik ithalatının ilk kaynakları Özbekistan (\%25) ve Hindistan, dokuma kumaşınki Çin (\%23), Mısır ve İtalya, dış giyiminki Bangladeş (\%16) ve Çin, ayakkabınınki Çin (\%32) ve Vietnam (\%27). Ayakkabı ve çantada Türkiye net ithalatçı; deri bölümünün dış ticaret açığı Tablo 2.1'de görülüyor. AB'den gelen ithalatın ağırlığı İtalyan kumaşı, İspanyol ve İtalyan markalarının giyim ve ayakkabısı ile deri. Çıkışın ithalat tarafındaki etkisi, Bölüm 3.4'te görüleceği gibi, bu ürünlerin c2 tarifesiyle pahalanmasından geliyor.

Bu tablo bir daralma döneminde çekildi. 2024 verileri sektörün güçlü bir yılını gösteriyor; 2025 ise farklı. TEPAV'ın SGK verisine dayanan izleme bültenine göre giyim eşyası imalatında sigortalı ücretli çalışan sayısı 2025'te bir yılda 86 binden fazla azaldı \etiketV{}. Basına yansıyan haberler, maliyet baskısı altındaki firmaların bir kısmının üretimini Mısır'a taşıdığını gösteriyor; Mısır'ın Süveyş Kanalı Ekonomik Bölgesi'nde Türk sermayeli büyük konfeksiyon yatırımları açıldı \etiketV{}. Taşınan firma sayısına ilişkin sendika açıklamaları ise doğrulanmadı \etiketD{}. Bu daralmanın nedeni Gümrük Birliği değil, reel kur ve ücret maliyeti. Ama çıkış senaryolarının etkisini değerlendirirken bu arka planı unutmamak gerekiyor: çıkış, zaten küçülen ve taşınma eğiliminde olan bir sektörü vuruyor. Sınıf düzeyinde 2025 istihdam verisi henüz yok (VB-38); bu yüzden modelin taban yılı 2024.

\hypertarget{nace_kartlar_13-15-3-uxe7ux131kux131ux15f-bu-sektuxf6rleri-hangi-kanallardan-etkiliyor}{%
\section{3. Çıkış bu sektörleri hangi kanallardan etkiliyor?}\label{nace_kartlar_13-15-3-uxe7ux131kux131ux15f-bu-sektuxf6rleri-hangi-kanallardan-etkiliyor}}

Çıkış tekstil, giyim ve deriyi beş kanaldan etkiliyor. En büyüğü AB'nin MFN tarifesi: anlaşmasız çıkışta giyimde \%12'lik bir duvar. İkincisi görünmeyen bir kanal: Türk kumaşının Kuzey Afrika ve Balkanlar üzerinden AB'ye giden zincirinin kopması. Üçüncüsü S2'nin menşe kuralları. Dördüncüsü ithalat tarafı: AB'den gelen giyim ve ayakkabının pahalanması ve bunun enflasyona etkisi. Beşincisi Gümrük Birliği'nden bağımsız düzenlemeler ve sektörün yapısal daralması. Aşağıda kanallar bu sırayla ele alınıyor.

\hypertarget{nace_kartlar_13-15-31-tarife-duvarux131-ve-tercihli-rakipler}{%
\subsection{3.1 Tarife duvarı ve tercihli rakipler}\label{nace_kartlar_13-15-31-tarife-duvarux131-ve-tercihli-rakipler}}

\Needspace{7\baselineskip}

\textbf{Tablo 3.1 --- GB ve İTÜ rejimindeki sınıflar: S4'te AB tarifesi ve S2'de menşe (T+5, P50)}

\nopagebreak[4]

\begingroup\gbtablo\footnotesize

\begin{longtable}[]{@{}
  >{\raggedright\arraybackslash}p{(\columnwidth - 16\tabcolsep) * \real{0.1538}}
  >{\raggedleft\arraybackslash}p{(\columnwidth - 16\tabcolsep) * \real{0.1204}}
  >{\raggedleft\arraybackslash}p{(\columnwidth - 16\tabcolsep) * \real{0.0945}}
  >{\raggedleft\arraybackslash}p{(\columnwidth - 16\tabcolsep) * \real{0.1058}}
  >{\raggedright\arraybackslash}p{(\columnwidth - 16\tabcolsep) * \real{0.0998}}
  >{\raggedleft\arraybackslash}p{(\columnwidth - 16\tabcolsep) * \real{0.1147}}
  >{\raggedleft\arraybackslash}p{(\columnwidth - 16\tabcolsep) * \real{0.1147}}
  >{\raggedleft\arraybackslash}p{(\columnwidth - 16\tabcolsep) * \real{0.0774}}
  >{\raggedleft\arraybackslash}p{(\columnwidth - 16\tabcolsep) * \real{0.1190}}@{}}
\toprule\noalign{}
\begin{minipage}[b]{\linewidth}\raggedright
\textbf{Sınıf}
\end{minipage} & \begin{minipage}[b]{\linewidth}\raggedleft
\textbf{AB'ye ihracat (€ mn, 2025)}
\end{minipage} & \begin{minipage}[b]{\linewidth}\raggedleft
\textbf{GB ve İTÜ payı}
\end{minipage} & \begin{minipage}[b]{\linewidth}\raggedleft
\textbf{S4'te AB tarifesi}
\end{minipage} & \begin{minipage}[b]{\linewidth}\raggedright
\textbf{Menşe riski S2/S3}
\end{minipage} & \begin{minipage}[b]{\linewidth}\raggedleft
\textbf{AB'ye ihracat S2 (\$ mn)}
\end{minipage} & \begin{minipage}[b]{\linewidth}\raggedleft
\textbf{AB'ye ihracat S4 (\$ mn)}
\end{minipage} & \begin{minipage}[b]{\linewidth}\raggedleft
\textbf{S2 ÷ S4}
\end{minipage} & \begin{minipage}[b]{\linewidth}\raggedleft
\textbf{İstihdam S4 (kişi)}
\end{minipage} \\
\midrule\noalign{}
\endhead
\bottomrule\noalign{}
\endlastfoot
14.14 İç giyim ve tişört & 2.246 & \%100 & \%12,0 & D/O & −445 & −1.003 & \%44 & −9.401 \\
14.13 Dış giyim & 3.588 & \%100 & \%12,0 & O/O & −453 & −999 & \%45 & −11.910 \\
14.39 Triko (kazak, hırka) & 1.306 & \%100 & \%12,0 & D/D & −138 & −425 & \%32 & −3.773 \\
13.92 Ev tekstili & 1.068 & \%100 & \%10,0 & O/O & −163 & −351 & \%47 & −2.804 \\
14.31 Çorap & 657 & \%100 & \%12,0 & D/D & −101 & −271 & \%37 & −2.277 \\
13.10 İplik & 834 & \%100 & \%3,8 & D/D & −146 & −229 & \%64 & −1.031 \\
13.20 Dokuma kumaş & 839 & \%100 & \%8,0 & D/D & −83 & −208 & \%40 & −3.966 \\
15.20 Ayakkabı & 303 & \%100 & \%9,9 & D/D & −53 & −128 & \%41 & −412 \\
13.91 Örme kumaş & 437 & \%100 & \%8,0 & D/D & −47 & −118 & \%40 & −1.632 \\
14.19 Diğer giyim ve aksesuar & 515 & \%100 & \%11,2 & O/O & −49 & −107 & \%46 & −1.289 \\
13.93 Halı & 383 & \%100 & \%7,9 & D/O & −27 & −60 & \%45 & −340 \\
13.95 Dokusuz kumaş & 311 & \%100 & \%4,3 & D/D & −29 & −53 & \%55 & −36 \\
13.96 Teknik tekstil & 183 & \%100 & \%6,4 & D/D & −22 & −48 & \%46 & −50 \\
15.11 Deri tabaklama & 89 & \%100 & \%3,7 & D/D & −9 & −16 & \%58 & −13 \\
15.12 Bavul ve çanta & 116 & \%100 & \%3,9 & D/D & −6 & −12 & \%54 & +214 \\
13.99 Diğer tekstil (şerit, dantel, keçe) & 101 & \%100 & \%6,4 & D/O & −6 & −11 & \%50 & −131 \\
14.12 İş giysisi & 25 & \%100 & \%12,0 & O/O & −4 & −8 & \%49 & −62 \\
13.94 Halat ve ağ & 34 & \%100 & \%6,8 & O/O & −3 & −5 & \%57 & −25 \\
14.11 Deri giyim & 66 & \%100 & \%4,0 & D/D & −3 & −5 & \%53 & −27 \\
14.20 Kürk & 30 & \%100 & \%3,7 & D/D & −1 & −2 & \%55 & −8 \\
\textbf{Toplam} & \textbf{13.132} & --- & --- & --- & \textbf{−1.788} & \textbf{−4.060} & \textbf{\%44} & \textbf{−38.973} \\
\end{longtable}

\endgroup

\emph{Kaynak:} \etiketT{} f3d, f3e. Menşe riski D/O/Y: düşük, orta, yüksek (\kod{veri/\allowbreak{}f3\_\allowbreak{}sinif\_\allowbreak{}varsayim.\allowbreak{}csv}, \etiketH{}). Seçim: AB'ye ihracatının en az yarısı GB veya İTÜ rejiminde, €10 mn üzerinde ve S4'te AB tarifesi en az \%1 olan sınıflar.

Tablo 3.1, sınıfları S4'teki AB'ye ihracat kaybına göre sıralıyor. Kayıp giyimde toplanıyor, çünkü giyim hem en yüksek tarifeyle karşılaşıyor hem de AB'ye en çok satan bölüm. İç giyim ve tişört ile dış giyim, S4'te her biri yaklaşık bir milyar dolarlık ihracat kaybediyor; bu, dış giyimin AB'ye ihracatının \%20, iç giyim ve tişörtünkinin \%30 kadarı. Trikoda kayıp 425 milyon dolar, çorapta 271 milyon dolar. Tekstil sınıflarında tarife daha düşük, kayıp da daha küçük: ev tekstilinde 351, iplikte 229, dokuma kumaşta 208 milyon dolar.

Tarifenin etkisini rakiplerin konumu büyütüyor. AB'nin giyim pazarındaki büyük tedarikçilerin çoğu tercihli. Dış giyimde AB'nin AB dışından yaptığı ithalatın \%58 kadarı, iç giyim ve tişörtte \%57 kadarı, Türkiye dışındaki tercihli tedarikçilerden geliyor. Bangladeş ve Kamboçya EBA rejimiyle sıfır tarife ödüyor. Pakistan'ın GSP+ statüsü AB'nin Ocak 2027'de yürürlüğe girecek yeni genel tercihler tüzüğü altında yeniden başvuruya bağlandı \etiketD{}. Vietnam'ın AB anlaşmasındaki giyim tarifeleri aşamalı olarak sıfırlanıyor; Fas ve Tunus PEM kümülasyonuyla sıfır tarife ödüyor. Bangladeş'in en az gelişmiş ülke statüsünden çıkışı 24 Kasım 2026'ya planlanmıştı; ülke ertelenmesini istedi \etiketD{}. Her durumda EBA tercihleri çıkıştan sonra üç yıl sürüyor \etiketV{}, yani 2029'a kadar değişmiyor. S4'te Türkiye bu pazarda Çin'le birlikte tam tarife ödeyen iki büyük tedarikçiden biri olur. Türkiye'nin elinde kalan üstünlük, Asya'nın sunamadığı hız. Avrupalı perakendecilerin sezon içi siparişlerinde iki haftalık bir teslim süresi \%12'lik bir tarifeyi kısmen telafi edebiliyor \etiketH{}. Ama bu, ürünün standartlaştıkça kaybolan bir üstünlük: tişört ve iç çamaşırında Bangladeş'in fiyat avantajı hızdan daha belirleyici.

Model, sınır sürtünmesini bütün Gümrük Birliği satırlarında aynı varsayıyor: S4'te \%0,5--2 arasında bir ek maliyet. Hızlı modada bir günlük gecikmenin değeri ortalama sanayi ürününden yüksek olduğu için gerçek maliyet giyimde bu aralığın üstünde olabilir \etiketH{}. Bu, sonuçların giyimde iyimser tarafta kalabileceğinin bir nedeni (VB-39).

S2'de tablo değişiyor. Tablo 3.1'in son satırına göre S2'nin AB'ye ihracat kaybı, S4 kaybının \%44 kadarı. Kalan kayıp üç kaynaktan geliyor: menşe kuralını karşılayamayan ihracatın ödediği MFN tarifesi, menşe ispatının maliyeti ve yarıya inen sınır sürtünmesi. Giyimde S2'yi S4'ten ayıran şey bu yüzden tarife değil, menşe kuralının nasıl yazıldığı. Bölüm 3.3 bunu ayrıntılı ele alıyor.

\hypertarget{nace_kartlar_13-15-32-guxf6ruxfcnmeyen-maruziyet-pem-kuxfcmuxfclasyonu-ve-kuzey-afrika-zinciri}{%
\subsection{3.2 Görünmeyen maruziyet: PEM kümülasyonu ve Kuzey Afrika zinciri}\label{nace_kartlar_13-15-32-guxf6ruxfcnmeyen-maruziyet-pem-kuxfcmuxfclasyonu-ve-kuzey-afrika-zinciri}}

Türk tekstilinin AB'ye ulaşmasının ikinci bir yolu var. Fas'taki, Tunus'taki ya da Sırbistan'daki bir konfeksiyoncu Türkiye'den aldığı kumaşı kesip dikiyor ve giysiyi AB'ye gümrüksüz satıyor. Bu mümkün, çünkü bu ülkelerin hepsi Pan-Avrupa-Akdeniz menşe sisteminin tarafı. PEM'in çapraz kümülasyon kuralına göre bir PEM tarafının menşeli malzemesi, başka bir tarafta işlendiğinde orada menşeli sayılıyor (PEM Ek I md. 7) \etiketV{}. Giyimde menşe kuralı kumaşın PEM bölgesinde dokunmasını ya da örülmesini istediği için, Fas'taki konfeksiyoncunun AB'ye tercihli satabilmesi Türk kumaşının ``PEM menşeli'' sayılmasına bağlı.

Kümülasyonun bir şartı var: girdinin menşe tarafıyla son ürünün varış tarafı arasında GATT XXIV. maddesine uygun bir tercihli ticaret anlaşması yürürlükte olmalı (PEM Ek I md. 8(1)(a)) \etiketV{}. Bugün bu anlaşma Gümrük Birliği. S2'de ve S3'te bu yeni serbest ticaret anlaşması olur. Anlaşmasız çıkışta ise Türkiye ile AB arasında böyle bir anlaşma kalmıyor ve Türk kumaşı Fas'ın AB'ye ihracatında menşesiz sayılıyor. Fas'taki konfeksiyoncu ya kumaşı AB'den, Mısır'dan ya da kendi ülkesinden almak ya da AB'de \%12 tarife ödemek zorunda kalıyor.

\Needspace{7\baselineskip}

\textbf{Tablo 3.2 --- PEM kümülasyonu: Türk kumaşının AB'ye ortaklar üzerinden ulaşan kısmı (2024)}

\nopagebreak[4]

\begingroup\gbtablo\small

\begin{longtable}[]{@{}
  >{\raggedright\arraybackslash}p{(\columnwidth - 6\tabcolsep) * \real{0.1530}}
  >{\raggedleft\arraybackslash}p{(\columnwidth - 6\tabcolsep) * \real{0.3566}}
  >{\raggedleft\arraybackslash}p{(\columnwidth - 6\tabcolsep) * \real{0.2584}}
  >{\raggedleft\arraybackslash}p{(\columnwidth - 6\tabcolsep) * \real{0.2320}}@{}}
\toprule\noalign{}
\begin{minipage}[b]{\linewidth}\raggedright
\textbf{PEM ortağı}
\end{minipage} & \begin{minipage}[b]{\linewidth}\raggedleft
\textbf{Türkiye'nin kumaş ve teknik tekstil ihracatı (\$ mn)}
\end{minipage} & \begin{minipage}[b]{\linewidth}\raggedleft
\textbf{Ortağın giyim ihracatında AB payı}
\end{minipage} & \begin{minipage}[b]{\linewidth}\raggedleft
\textbf{Kümülasyon maruziyeti (\$ mn)}
\end{minipage} \\
\midrule\noalign{}
\endhead
\bottomrule\noalign{}
\endlastfoot
Fas & 848 & \%74 & 613 \\
Tunus & 224 & \%82 & 179 \\
Sırbistan & 151 & \%67 & 86 \\
Ukrayna & 107 & \%82 & 80 \\
Mısır & 202 & \%31 & 57 \\
Kuzey Makedonya & 38 & \%86 & 30 \\
Arnavutluk & 31 & \%86 & 23 \\
Bosna-Hersek & 34 & \%72 & 22 \\
Diğer 6 ortak & 101 & --- & 40 \\
\textbf{Toplam} & \textbf{1.736} & --- & \textbf{1.131} \\
\end{longtable}

\endgroup

\emph{Kaynak:} BACI HS17 2024 \etiketV{}; \kod{kod/\allowbreak{}f3a\_\allowbreak{}ticaret.\allowbreak{}py} \etiketT{}. Kumaş: κ \textgreater{} 0 sınıflar (13.20, 13.91 κ = 1; 13.96, 13.99 κ = 0,3 \etiketH{}). Maruziyet = κ × ihracat × ortağın Fasıl 61--63 ihracatında AB payı; Türk kumaşının ortağın AB'ye giden üretiminde orantılı kullanıldığı varsayımıyla üst sınır. Ortaklar hem Türkiye hem AB ile tercihli anlaşması olan PEM tarafları.

Tablo 3.2 bu zincirin büyüklüğünü gösteriyor. Türkiye'nin PEM ortaklarına kumaş ve teknik tekstil ihracatı 1.736 milyon dolar. Bunun ne kadarının AB'ye giden giysiye dönüştüğü doğrudan gözlenmiyor. Ortağın giyim ihracatındaki AB payıyla ağırlıklandırınca ortaya 1.131 milyon dolarlık bir maruziyet çıkıyor. Bu bir üst sınır, çünkü Türk kumaşının ortak ülkede AB'ye giden üretimde orantılı kullanıldığı varsayılıyor (VB-37). Fas açık ara en büyük kalem (613 milyon dolar): Fas giyim ihracatının \%74 kadarını AB'ye, çoğunu İspanya'ya yapıyor ve Türk dokuma kumaşının en büyük AB dışı alıcısı. Tunus (\%82) ve Kuzey Makedonya (\%86) giyim ihracatında AB'ye daha da bağımlı. Mısır'ın giyim ihracatı ise ağırlıkla ABD'ye gidiyor (AB payı \%31), bu yüzden maruziyeti ihracatına göre küçük.

\Needspace{7\baselineskip}

\textbf{Tablo 3.3 --- PEM kümülasyonu kanalı: sınıflar (S4, T+5)}

\nopagebreak[4]

\begingroup\gbtablo\footnotesize

\begin{longtable}[]{@{}
  >{\raggedright\arraybackslash}p{(\columnwidth - 12\tabcolsep) * \real{0.1652}}
  >{\raggedleft\arraybackslash}p{(\columnwidth - 12\tabcolsep) * \real{0.1738}}
  >{\raggedleft\arraybackslash}p{(\columnwidth - 12\tabcolsep) * \real{0.1237}}
  >{\raggedleft\arraybackslash}p{(\columnwidth - 12\tabcolsep) * \real{0.0362}}
  >{\raggedleft\arraybackslash}p{(\columnwidth - 12\tabcolsep) * \real{0.1570}}
  >{\raggedleft\arraybackslash}p{(\columnwidth - 12\tabcolsep) * \real{0.1605}}
  >{\raggedleft\arraybackslash}p{(\columnwidth - 12\tabcolsep) * \real{0.1835}}@{}}
\toprule\noalign{}
\begin{minipage}[b]{\linewidth}\raggedright
\textbf{Sınıf}
\end{minipage} & \begin{minipage}[b]{\linewidth}\raggedleft
\textbf{PEM ortaklarına ihracat (\$ mn, 2024)}
\end{minipage} & \begin{minipage}[b]{\linewidth}\raggedleft
\textbf{Maruziyet (\$ mn)}
\end{minipage} & \begin{minipage}[b]{\linewidth}\raggedleft
\textbf{κ}
\end{minipage} & \begin{minipage}[b]{\linewidth}\raggedleft
\textbf{S4 kaybı (\$ mn; P50, P10 / P90)}
\end{minipage} & \begin{minipage}[b]{\linewidth}\raggedleft
\textbf{Kanalın istihdam etkisi (kişi, P50)}
\end{minipage} & \begin{minipage}[b]{\linewidth}\raggedleft
\textbf{Sınıfın S4 istihdam etkisi, kanal dahil (P50)}
\end{minipage} \\
\midrule\noalign{}
\endhead
\bottomrule\noalign{}
\endlastfoot
13.20 Dokuma kumaş & 1.095 & 768 & 1,0 & −457 (−583 / −328) & −3.399 & −3.966 \\
13.91 Örme kumaş & 495 & 334 & 1,0 & −199 (−254 / −143) & −1.197 & −1.632 \\
13.96 Teknik tekstil & 100 & 63 & 0,3 & −11 (−14 / −8) & −86 & −50 \\
13.99 Diğer tekstil (şerit, dantel, keçe) & 46 & 32 & 0,3 & −6 (−7 / −4) & −68 & −131 \\
\end{longtable}

\endgroup

\emph{Kaynak:} \etiketT{} f3a, f3e. κ: Türk girdisinin ortağın AB'ye tercihli ürününde menşe kazandıran adım olma derecesi (\kod{veri/\allowbreak{}f3\_\allowbreak{}sinif\_\allowbreak{}varsayim.\allowbreak{}csv}, \etiketH{}). κ = 0 olan sınıflarda kanal yok: iplikte ortakta yapılan dokuma ya da örme, giyim ve ayakkabıda ortağın kendi üretimi menşe kazandırıyor. Kayıp oranı λ \textasciitilde{} üçgen(0,3; 0,6; 0,9) \etiketH{}. Sınıf toplamında kanalın etkisi iç ikame ve diğer kanallarla birlikte.

Kanal yalnızca kumaşta işliyor. Giyimde menşe kuralı kumaş oluşturmayı (dokuma ya da örme) menşe kazandıran adım saydığı için, Fas'ta dokunan ya da örülen kumaşın ipliği nereden gelirse gelsin giysi menşeli oluyor. Bu yüzden Türk ipliği kümülasyona bağlı değil. Türk giysisi ve ayakkabısı da PEM ortaklarına son ürün olarak gidiyor. Dokuma ve örme kumaşta ise kanal büyük: S4'te dokumanın PEM ortaklarına ihracatı 457 milyon dolar, örme kumaşınki 199 milyon dolar azalıyor. Bu, iki sınıfın S4'teki istihdam kaybının çoğu. Dokumada S4'ün istihdam etkisi 3.966 kişi ve bunun 3.399 kişisi bu kanaldan geliyor. Kaybın oranı (λ) için \%30--90 arası, tepe değeri \%60 olan bir dağılım varsayıldı \etiketH{}. Ortak ülkelerdeki alıcıların bir kısmı AB'ye satmayı sürdürmek için kumaşı başka yerden alacak. Bir kısmı Türk kumaşını AB dışı pazarlara ya da Türkiye'ye satılan ürünlerde kullanmaya devam edecek. Bir kısmı ise PEM kuralının sunduğu bir çıkış yolunu kullanacak: dokuma giyimde baskı ve dikim birlikte menşe kazandırıyor. Ortak ülkede baskılanan Türk ham kumaşı, menşesiz olsa da giysiye menşe kazandırabiliyor \etiketV{}. Bu yol yalnızca baskılı dokuma ürünlerde açık.

Kanalın üç politika sonucu var. Birincisi, S2 ve S3'ün menşe protokolü PEM'e dinamik atıf yapmalı ve AB, kümülasyon için gereken bildirimi (PEM Ek I md. 8(2)) çıkış günü yayımlamalı. Hasmane bir AB'nin bu bildirimi geciktirmesi, anlaşma imzalansa bile zinciri aylarca koparabilir \etiketH{}. Bu yüzden bildirim, anlaşmanın yürürlüğüyle aynı güne bağlanmalı (E-13-01). İkincisi, Türkiye'nin Fas ve Tunus ile ilişkisi revize PEM kurallarına taşınmalı. Kümülasyon ancak aynı kuralları uygulayan taraflar arasında işliyor ve bu iki ülkeyle geçiş henüz tamamlanmadı (Faz 2C, H-M8). Üçüncüsü, kanal S4'ün en az görünen maliyetlerinden biri ve Fas ile Mısır'a Türk tekstil yatırımı kaymasını hızlandırıyor. Kumaşı Fas'ta ya da Mısır'da üreten bir Türk firması, S4'te bu zinciri kendi içinde yeniden kurabiliyor.

\hypertarget{nace_kartlar_13-15-33-s2nin-sux131navux131-uxe7ift-duxf6nuxfcux15fuxfcm-baskux131-ve-geri-uxf6deme-yasaux11fux131}{%
\subsection{3.3 S2'nin sınavı: çift dönüşüm, baskı ve geri ödeme yasağı}\label{nace_kartlar_13-15-33-s2nin-sux131navux131-uxe7ift-duxf6nuxfcux15fuxfcm-baskux131-ve-geri-uxf6deme-yasaux11fux131}}

Gümrük Birliği'nde AB'ye satılan bir Türk giysisinin kumaşının nereden geldiği sorulmuyor. Serbest ticaret anlaşmasında ise giysinin tercihli tarifeden yararlanması menşe kuralına bağlı. AB'nin Türkiye'yle yapacağı her anlaşmada bu kuralın revize PEM kuralları olacağını varsaymak gerçekçi; AB'nin Akdeniz ve Balkan ortaklarıyla bütün anlaşmaları bunlara atıf yapıyor. Revize kurallar 2025'ten beri uygulanıyor ve tekstilde eski kurallara göre belirgin biçimde esnek \etiketV{}. Tablo 3.4 bu parçanın sınıflarına uygulanacak kuralları özetliyor.

\Needspace{7\baselineskip}

\textbf{Tablo 3.4 --- Revize PEM kuralları: tekstil, giyim ve deri (PEM Ortak Komitesi Kararı 1/2023)}

\nopagebreak[4]

\begingroup\gbtablo\small

\begin{longtable}[]{@{}
  >{\raggedright\arraybackslash}p{(\columnwidth - 4\tabcolsep) * \real{0.1684}}
  >{\raggedright\arraybackslash}p{(\columnwidth - 4\tabcolsep) * \real{0.4860}}
  >{\raggedright\arraybackslash}p{(\columnwidth - 4\tabcolsep) * \real{0.3456}}@{}}
\toprule\noalign{}
\begin{minipage}[b]{\linewidth}\raggedright
\textbf{Ürün (fasıl)}
\end{minipage} & \begin{minipage}[b]{\linewidth}\raggedright
\textbf{Menşe kazandıran işlem}
\end{minipage} & \begin{minipage}[b]{\linewidth}\raggedright
\textbf{Türk ihracatı için anlamı}
\end{minipage} \\
\midrule\noalign{}
\endhead
\bottomrule\noalign{}
\endlastfoot
İplik (50--55) & Doğal liflerin eğrilmesi; sentetik lif çekimi + eğirme; büküm + mekanik işlem & İthal pamuk ve POY serbest; risk düşük \\
Dokuma kumaş (50--55) & Dokuma + boyama, baskı ya da kaplama; iplik boyama + dokuma; tek başına baskı & İthal iplik serbest; ithal kumaş Türkiye'de baskılanırsa menşe kazanıyor \\
Örme kumaş (60) & Örme + boyama, baskı ya da kaplama; iplik boyama + örme & İthal iplik serbest; risk düşük \\
Halı (57) & Eğirme ya da filament çekimi + dokuma veya tufting; tufting + boyama ya da baskı & İplik çekimi entegre üreticide risk düşük \\
Ev tekstili (6301--6304) & Dokuma ya da örme + konfeksiyon & Çift dönüşüm; ithal ham bezle üretim tercih dışı \\
Örme giyim (61) & Örme + konfeksiyon (kesim dahil); biçimine göre örülmüş üründe örme ile konfeksiyon tek işlemde & Kumaş PEM bölgesinde örülmeli; ithal iplik serbest \\
Dokuma giyim (62) & Dokuma + konfeksiyon; ya da tek başına baskı + konfeksiyon & Kumaş PEM bölgesinde dokunmalı ya da Türkiye'de baskılanmalı \\
Deri (4104--4107) & Yeniden tabaklama ya da başlık değişikliği & İthal ham ve yarı işlenmiş deri serbest \\
Çanta, deri eşya (42) & Başlık değişikliği ya da menşesiz malzeme ≤ \%50 & Risk düşük \\
Ayakkabı (64) & Her başlıktan imalat; yalnızca 6406'daki hazır saya ve taban birleşimleri kullanılamaz & Risk düşük \\
Toleranslar & Kurala uymayan tekstil malzemesi ≤ fabrika çıkış fiyatının \%15'i (Not 7.1); karışık liflerde ağırlıkça \%15 (Not 6) & Astar, şerit ve aksesuar çoğunlukla tolerans içinde \\
Kümülasyon ve geri ödeme & Tam kümülasyon 50--63'te yok, ikili kümülasyon var (md. 7(3)--(4)); 50--63'te geri ödeme yasağı (md. 16) & AB'de yapılan işlem AB'ye ihracatta sayılıyor; DİR girdisinde vergi ödenir \\
\end{longtable}

\endgroup

\emph{Kaynak:} \etiketV{} PEM Ortak Komitesi Kararı 1/2023 (RG L 2024/390), Ek I md. 7, 8, 16 ve giriş notları 6--7; ürüne özgü kurallar listesi. Sağdaki sütun \etiketH{}.

Revize kurallar Türkiye'nin yapısına uygun. Türkiye kumaşın büyük kısmını kendisi üretiyor; giyim ihracatının çoğu çift dönüşüm şartını zaten karşılıyor. Sorun, Asya'dan, ağırlıkla Çin'den gelen kumaşla üretilen giysilerde. Faz 2C'nin ürün düzeyindeki hesabına göre giyim sanayisinin PEM dışından kumaş ithalatı 2024'te yaklaşık 0,8 milyar dolar. Kumaşın giysi değerindeki payıyla çevrildiğinde bu, giyim ihracatının yaklaşık \%10--15'ine karşılık geliyor (02c, Tablo 4.3) \etiketH{}. Bu ihracat S2'de tercihli tarifeden yararlanamıyor. Model bunu sınıf düzeyinde bir uyum oranıyla yansıtıyor: dış giyimde ihracatın \%86 kadarı, iç giyim ve tişörtte \%88 kadarı menşe kuralını karşılıyor ve uyumlu ihracat da fabrika çıkış fiyatının \%2--2,5'i kadar bir menşe maliyetine katlanıyor \etiketH{}. Dokuma giyimde baskı seçeneği bu kaybın bir kısmını kurtarıyor: Asya'dan gelen ham kumaş Türkiye'de baskılanırsa giysi menşeli oluyor. Örme giyimde bu seçenek yok; yalnızca sütyen ve korsede var. S2 müzakeresinde bunun örme giyime genişletilmesi istenmeli (E-13-06).

Menşeye ilişkin ikinci kural geri ödeme yasağı. PEM'de 50--63. fasılların ürünlerinde, ihraç edilen ürüne giren PEM dışı malzemenin gümrük vergisi geri ödenemiyor ya da askıya alınamıyor (Ek I md. 16) \etiketV{}. Bu ilk bakışta Dahilde İşleme Rejimi'yle (DİR) Asya kumaşı ve ipliği ithal eden ihracatçı için ağır bir yük gibi görünüyor. Ama Türkiye bugün de benzer bir kural uyguluyor: DİR kapsamında üretilen ürün AB'ye A.TR belgesiyle ihraç edildiğinde, PEM dışı menşeli ithal girdiye telafi edici vergi (TEV) alınıyor \etiketV{}. Dolayısıyla S2'de geri ödeme yasağı AB'ye ihracatta yeni bir maliyet getirmiyor; model de buna ek maliyet eklemiyor (V-42). Farkı S4 yaratıyor. S4'te AB'ye ihracatta tercih olmadığı için TEV de yok ve DİR'in vergi muafiyeti tam olarak kullanılabiliyor. Bu, S4'te Asya kumaşına geçişi ucuzlatan, yani yerli kumaş üreticisinin aleyhine işleyen bir etki.

\hypertarget{nace_kartlar_13-15-34-iux307thalat-tarafux131-c2-tarifesi-enflasyon-ve-vergi-takasux131}{%
\subsection{3.4 İthalat tarafı: c2 tarifesi, enflasyon ve vergi takası}\label{nace_kartlar_13-15-34-iux307thalat-tarafux131-c2-tarifesi-enflasyon-ve-vergi-takasux131}}

\Needspace{7\baselineskip}

\textbf{Tablo 3.5 --- İthalat tarafı: TÜFE'ye birinci tur etkisi en büyük sınıflar (bp, P50)}

\nopagebreak[4]

\begingroup\gbtablo\footnotesize

\begin{longtable}[]{@{}
  >{\raggedright\arraybackslash}p{(\columnwidth - 16\tabcolsep) * \real{0.1629}}
  >{\raggedleft\arraybackslash}p{(\columnwidth - 16\tabcolsep) * \real{0.1152}}
  >{\raggedleft\arraybackslash}p{(\columnwidth - 16\tabcolsep) * \real{0.1442}}
  >{\raggedleft\arraybackslash}p{(\columnwidth - 16\tabcolsep) * \real{0.1368}}
  >{\raggedleft\arraybackslash}p{(\columnwidth - 16\tabcolsep) * \real{0.1444}}
  >{\raggedleft\arraybackslash}p{(\columnwidth - 16\tabcolsep) * \real{0.0481}}
  >{\raggedleft\arraybackslash}p{(\columnwidth - 16\tabcolsep) * \real{0.0584}}
  >{\raggedleft\arraybackslash}p{(\columnwidth - 16\tabcolsep) * \real{0.0940}}
  >{\raggedleft\arraybackslash}p{(\columnwidth - 16\tabcolsep) * \real{0.0960}}@{}}
\toprule\noalign{}
\begin{minipage}[b]{\linewidth}\raggedright
\textbf{Sınıf}
\end{minipage} & \begin{minipage}[b]{\linewidth}\raggedleft
\textbf{TÜFE ağırlığı}
\end{minipage} & \begin{minipage}[b]{\linewidth}\raggedleft
\textbf{Tüketimde ithal pay}
\end{minipage} & \begin{minipage}[b]{\linewidth}\raggedleft
\textbf{İthalatta AB payı}
\end{minipage} & \begin{minipage}[b]{\linewidth}\raggedleft
\textbf{S4'te AB'ye tarife (c2)}
\end{minipage} & \begin{minipage}[b]{\linewidth}\raggedleft
\textbf{S2}
\end{minipage} & \begin{minipage}[b]{\linewidth}\raggedleft
\textbf{S4}
\end{minipage} & \begin{minipage}[b]{\linewidth}\raggedleft
\textbf{Önlem (S4)}
\end{minipage} & \begin{minipage}[b]{\linewidth}\raggedleft
\textbf{S4, önlemle}
\end{minipage} \\
\midrule\noalign{}
\endhead
\bottomrule\noalign{}
\endlastfoot
14.13 Dış giyim & \%3,86 & \%11 & \%26 & \%37,2 & 0,45 & 4,33 & −4,33 & 0,00 \\
15.20 Ayakkabı & \%1,11 & \%34 & \%21 & \%48,7 & 0,39 & 3,93 & −3,93 & 0,00 \\
15.12 Bavul ve çanta & \%0,65 & \%36 & \%51 & \%15,4 & 0,25 & 1,81 & --- & 1,81 \\
14.39 Triko (kazak, hırka) & \%1,30 & \%8 & \%27 & \%50,9 & 0,15 & 1,49 & −1,49 & 0,00 \\
14.19 Diğer giyim ve aksesuar & \%0,90 & \%15 & \%23 & \%44,8 & 0,14 & 1,37 & −1,37 & 0,00 \\
14.14 İç giyim ve tişört & \%0,64 & \%8 & \%21 & \%50,7 & 0,08 & 0,84 & −0,84 & 0,00 \\
13.20 Dokuma kumaş & \%0,00 & \%17 & \%31 & \%28,0 & 0,06 & 0,54 & --- & 0,54 \\
13.92 Ev tekstili & \%0,69 & \%6 & \%34 & \%30,7 & 0,05 & 0,46 & --- & 0,46 \\
Diğer 13 sınıf & \%0,37 & --- & --- & --- & 0,10 & 0,80 & −0,06 & 0,75 \\
\textbf{13--15 toplamı} & \textbf{\%9,52} & --- & --- & --- & \textbf{1,67} & \textbf{15,58} & \textbf{−12,01} & \textbf{3,57} \\
\end{longtable}

\endgroup

\emph{Kaynak:} \etiketT{} f3e (doğrudan kanal COICOP 2018 sepeti, dolaylı kanal I-O 2023 Leontief fiyat modeli). c2: S4'te AB menşeli mala uygulanacak MFN + İGV (İGV, AB'nin baskın, yerli üretimin zayıf olduğu satırlarda herkes için kalkar; Faz 2A). Önlemler Bölüm 4'te.

Anlaşmasız çıkışta Türkiye AB menşeli ürüne kendi MFN tarifesini uygulamaya başlıyor. Faz 2A'da önerilen c2 rejiminde ek gümrük vergisi (İGV), AB'nin baskın tedarikçi ve yerli üretimin zayıf olduğu satırlarda bütün menşeler için kaldırılıyor; diğer satırlarda AB'ye de uygulanıyor. Giyim ve ayakkabıda yerli üretim güçlü, bu yüzden İGV AB menşeli ürüne de geliyor. AB'den gelen dış giyimin tarifesi S4'te \%37,2, iç giyim ve tişörtün \%50,7, ayakkabının \%48,7 oluyor.

Bu oranlar Türkiye'nin bugün Çin'e uyguladığı oranlarla aynı; c2 uygulanan MFN tarifesini yükseltmiyor, yalnızca AB'nin tercihini kaldırıyor. DTÖ açısından bu yüzden bağlı oran sorunu doğmuyor. Ayrıca Türkiye giyimde tarife satırlarının yalnızca \%3,7'sini, tekstilde \%23,4'ünü bağlamış durumda \etiketV{} (VB-32 bu parçada bağlayıcı değil).

Etkinin büyüklüğünü TÜFE sepetindeki ağırlık belirliyor. Dış giyimin TÜFE ağırlığı \%3,86 ve tüketimdeki ithal payı küçük (\%11). Buna karşılık sepet o kadar büyük ki AB'den gelen giysinin pahalanması tek başına 4,3 baz puanlık bir etki yaratıyor. Ayakkabıda ithal pay daha yüksek (\%34) ve etki 3,9 baz puan. Parçanın toplam S4 etkisi 15,6 baz puan. Bu, kümülatif 50 baz puanlık bütçenin büyük bir dilimi ve parçanın TÜFE sepetindeki payının üstünde. S2'de etki 1,7 baz puan; tamamı sınır sürtünmesinden ve AB ihracatçısının menşe maliyetinden geliyor.

Bu etki için hazır bir araç var: Faz 2A'da tasarlanan vergi takası (T-A). Takas, tarife ile aynı ürün grubunda bütün menşelere eşit bir KDV indirimini bir arada uyguluyor. Giyimde KDV \%10'dan \%8,7'ye, ayakkabıda \%5,6'ya iniyor ve grubun fiyat endeksi değişmiyor (02a, Tablo 1.5a). İndirim menşeden bağımsız olduğu için DTÖ'nün ulusal muamele kuralıyla (GATT III:2) uyumlu. Yerli üreticiye avantajı ise tarife veriyor, vergi değil. Faz 2A'nın merkez hesabında takasın yıllık net mali maliyeti giyimde 319, ayakkabıda 63 milyon dolar \etiketT{}. Bu, AB menşeli ürünlere gelen gümrük geliri düşüldükten sonraki tutar; 2026 kuruyla yaklaşık 17,5 milyar TL. Takasla birlikte parçanın S4'teki enflasyon yükü 3,6 baz puana iniyor. Kalan etkinin yaklaşık yarısı çantada (1,8 baz puan). Faz 2A'nın takas grupları çantayı içermiyordu; çantanın KDV oranı ve takasa eklenmesi ayrıca değerlendirilmeli \etiketD{} (VB-39).

Takasın istihdam tarafında da etkisi var. Faz 2A'ya göre c2 tarifesi ve KDV indirimi birlikte iç pazarda yerli giyim ve ayakkabının payını artırıyor ve S4'te yaklaşık 20 bin (giyim) ve 3 bin (ayakkabı) kişilik bir istihdam kazancı yaratıyor \etiketT{}. Bu parçanın senaryo modeli c2'nin iç ikame etkisini zaten içeriyor; KDV indiriminin ek talep etkisi ise içermiyor. İki model farklı olduğu için kazançlar toplanmadı. Sonuç tablolarındaki S4 istihdam kaybı bu açıdan temkinli bir tahmin.

Girdi tarafında c2 ters yönde çalışıyor. AB'den gelen dokuma kumaşın tarifesi S4'te \%28,0, örme kumaşınki \%34,7. İhracatçı bu kumaşı DİR'le vergisiz getirebiliyor; iç pazara üreten konfeksiyoncu ise getiremiyor. Etki üst segmentte yoğunlaşıyor: İtalyan yünlü ve ipekli kumaşı kullanan takım elbise ve gömlek üreticileri. Dolaylı kanaldan TÜFE'ye etkisi küçük (0,5 baz puan); bu yüzden ayrı bir önlem önerilmiyor.

\hypertarget{nace_kartlar_13-15-35-guxfcmruxfck-birliux11finden-baux11fux131msux131z-engeller-ve-yapux131sal-daralma}{%
\subsection{3.5 Gümrük Birliği'nden bağımsız engeller ve yapısal daralma}\label{nace_kartlar_13-15-35-guxfcmruxfck-birliux11finden-baux11fux131msux131z-engeller-ve-yapux131sal-daralma}}

AB'nin tekstil ve giyime yönelik yeni düzenlemeleri çıkıştan bağımsız ve AB'ye satan herkese uygulanıyor. Dördü bu parça için önemli. Zorla çalıştırmayla üretilen ürünlerin AB pazarına girişini yasaklayan tüzük (Tüzük 2024/3015) 14 Aralık 2027'den itibaren uygulanıyor \etiketV{}. Pamuk bu düzenlemenin en hassas olduğu ürünlerden biri. Türkiye'nin iplik ithalatının \%25 kadarı Özbekistan'dan geliyor ve Orta Asya pamuğunun izlenebilirliği AB'nin gündeminde \etiketH{}. Sürdürülebilir ürünler için ekotasarım tüzüğü (ESPR) büyük işletmelerin satılmayan giyim ve ayakkabıyı imha etmesini 19 Temmuz 2026'dan beri yasaklıyor \etiketV{}. Tekstil için ekotasarım kurallarının ve dijital ürün pasaportunun 2027'de kabul edilmesi bekleniyor \etiketV{}. Atık Çerçeve Direktifi'nin 2025 değişikliği (Direktif 2025/1892), üye devletlerin Nisan 2028'e kadar tekstil ve ayakkabıda genişletilmiş üretici sorumluluğu sistemleri kurmasını istiyor \etiketV{}. Dördüncüsü olumlu bir gelişme. AB'nin ormansızlaşma tüzüğü (EUDR) 2026/2102 sayılı Yetki Devrine Dayanan Tüzükle değişti ve sığır derisi 18 Eylül 2026'da kapsamdan çıktı \etiketV{}. Türk tabakhanelerinin AB'ye ihracatının çoğu bu yükten kurtuldu.

Bu düzenlemeler Türk ihracatçısını AB'li ve Asyalı rakipleriyle aynı yükün altına sokuyor. Türkiye'nin üstünlüğü, kumaştan dikime kadar zinciri kendi içinde tutması; izlenebilirlik verisi bu yüzden Türk tedarikçi için Asyalı rakibe göre daha kolay üretilebilir. Ama hasmane bir AB'de uygulamanın sıkılığı ve denetimlerin hedefi bir baskı aracına dönüşebilir \etiketH{}. Önerilen yanıt, menşe kanıtı, pamuk izlenebilirliği ve dijital ürün pasaportu için tek bir veri altyapısı (E-13-02). Bu altyapı çıkış hiç gerçekleşmese de gerekli; pişmanlık yaratmayan bir eylem.

Ticaret savunması ikinci risk. AB 1997'de Mısır, Hindistan ve Pakistan menşeli pamuklu nevresime damping vergisi koydu (Tüzük 2398/97). Hindistan'ın DTÖ'ye taşıdığı dava (DS141) AB yöntemine önemli sınırlar getirdi \etiketV{}. Türk ev tekstili bugün bu tür önlemlerin dışında. Hasmane bir AB'de Türk ev tekstili ve sentetik iplik muhtemel hedefler arasında \etiketH{}. Bu yüzden ilk parçada su ürünleri için kurulan ticaret savunması hazırlığının (E-10-04) kapsamı ev tekstilini de içerecek biçimde genişletiliyor (E-13.92-01).

Sektörün kendi dinamiği ise çıkıştan daha güçlü. Bölüm 2'de anlatılan 2025 daralması sürerse, giyimde T+5'e kadar yalnızca maliyet kaynaklı istihdam kaybı, S4'ün çıkış kaynaklı kaybından büyük olabilir \etiketH{}. Bu, çıkışın etkisini küçültmüyor; ama politika açısından iki kaybın aynı yerde birleştiğini gösteriyor. Mısır'a kayan bir firma için S4, kararı öne çeken bir neden olur. S2 ise Türkiye'de kalmanın en güçlü gerekçesi olan AB pazarına tarifesiz erişimi korur. Önerilen dönüşüm programı (E-14.13-01) bu yüzden çıkıştan bağımsız tasarlandı.

\hypertarget{nace_kartlar_13-15-4-strateji-neyi-korumalux131-neyi-geliux15ftirmeli-nasux131l-muxfczakere-etmeli}{%
\section{4. Strateji: neyi korumalı, neyi geliştirmeli, nasıl müzakere etmeli?}\label{nace_kartlar_13-15-4-strateji-neyi-korumalux131-neyi-geliux15ftirmeli-nasux131l-muxfczakere-etmeli}}

\hypertarget{nace_kartlar_13-15-41-duruux15flar}{%
\subsection{4.1 Duruşlar}\label{nace_kartlar_13-15-41-duruux15flar}}

Duruşlar ilk parçadaki yöntemle atandı: önce kabiliyet, AB maruziyeti, ithalat bağımlılığı, enflasyon duyarlılığı ve stratejik önem skorlarına göre kurallarla, sonra her sınıf için ayrı ayrı gözden geçirilerek. Bu parçada dört duruş çıkıyor ve sektörün iki yüzünü yansıtıyor.

\Needspace{7\baselineskip}

\textbf{Tablo 4.1 --- Sınıfların duruşa göre dağılımı}

\nopagebreak[4]

\begingroup\gbtablo\footnotesize

\begin{longtable}[]{@{}
  >{\raggedright\arraybackslash}p{(\columnwidth - 8\tabcolsep) * \real{0.1247}}
  >{\raggedleft\arraybackslash}p{(\columnwidth - 8\tabcolsep) * \real{0.0618}}
  >{\raggedright\arraybackslash}p{(\columnwidth - 8\tabcolsep) * \real{0.4771}}
  >{\raggedleft\arraybackslash}p{(\columnwidth - 8\tabcolsep) * \real{0.1205}}
  >{\raggedleft\arraybackslash}p{(\columnwidth - 8\tabcolsep) * \real{0.2159}}@{}}
\toprule\noalign{}
\begin{minipage}[b]{\linewidth}\raggedright
\textbf{Duruş}
\end{minipage} & \begin{minipage}[b]{\linewidth}\raggedleft
\textbf{Sınıf}
\end{minipage} & \begin{minipage}[b]{\linewidth}\raggedright
\textbf{Sınıflar}
\end{minipage} & \begin{minipage}[b]{\linewidth}\raggedleft
\textbf{Çalışan (bin)}
\end{minipage} & \begin{minipage}[b]{\linewidth}\raggedleft
\textbf{AB'ye ihracat (\$ mn, 2024)}
\end{minipage} \\
\midrule\noalign{}
\endhead
\bottomrule\noalign{}
\endlastfoot
Koru & 5 & 13.30, 13.95, 13.99, 14.12, 15.20 & 181,8 & 863 \\
Geliştir & 11 & 13.10, 13.20, 13.91, 13.92, 13.93, 13.94, 13.96, 14.11, 14.20, 14.31, 15.12 & 423,0 & 5.487 \\
Dönüştür & 2 & 14.13, 15.11 & 374,6 & 5.141 \\
Yönetilen uyum & 3 & 14.14, 14.19, 14.39 & 258,5 & 5.847 \\
\end{longtable}

\endgroup

\emph{Kaynak:} \kod{veri/\allowbreak{}nace\_\allowbreak{}matris.\allowbreak{}csv} (\kod{durus}; kural tabanlı öneri ve uzman düzeltmesi, \kod{uzman\_\allowbreak{}duzeltmesi} sütunu).

``Geliştir'' duruşundaki 11 sınıfın çekirdeği tekstilin küresel rekabetçi kesimi: iplik, dokuma ve örme kumaş, ev tekstili, halı ve teknik tekstil. Bu sınıflar AB'ye entegre ama AB'ye bağımlılıkları giyimden düşük. Kabiliyetleri, yani dünya pazarındaki konumları, yurt içi katma değerleri ve tedarik zinciri bağlantıları güçlü. Strateji, AB pazarını S2'de korurken AB dışı pazarları ve teknik tekstil gibi yüksek katma değerli alanları büyütmek. Kumaşta buna PEM kümülasyonunun korunması ekleniyor. Çorap da bu gruba girdi: çorap sermaye yoğun ve otomasyonu yüksek bir sınıf, AB pazarında Çin'den sonra ikinci tedarikçi. Sorunu kabiliyet değil, \%12'lik tarife. Çanta ve deri giyim gibi küçük sınıflar da bu grupta.

``Dönüştür'' duruşu iki sınıfa öneriliyor. Dış giyim (14.13) imalatın en büyük işvereni ve AB pazarının üçüncü tedarikçisi. Kural tabanlı sınıflama onu ``AB'ye entegre ihracat şampiyonu'' (A) olarak işaretliyor. Ancak kişi başı katma değeri AB'deki muadilinin üçte biri dolayında ve 2025'te maliyet baskısıyla küçülüyor. Bu sınıfın geleceği ucuz emeğe dayalı fason dikimde değil; tasarım, küçük partili hızlı tedarik ve markada. Bu yüzden ikincil arketip olarak ``emek yoğun, rekabet gücü zayıflayan'' (F) eklendi ve duruş ``Dönüştür'' yapıldı. Deri tabaklamada (15.11) dönüşüm yeşil: ormansızlaşma tüzüğünden çıkışla birlikte AB'ye erişimin önündeki asıl engel, AB'li alıcıların çevre ve özen yükümlülükleri. Arıtma, krom azaltımı ve sertifikasyon bu yüzden sınıfın stratejisinin merkezinde.

``Yönetilen uyum'' duruşundaki üç sınıf, iç giyim ve tişört (14.14), diğer giyim ve aksesuar (14.19) ve triko (14.39), standart ürünlerin ağırlıkta olduğu ve Asya ile fiyat rekabetinin belirleyici olduğu sınıflar. S4'te bu üç sınıfın dağılımın alt ucundaki istihdam kaybı \%5'i aşıyor; trikoda \%19. Önerilen politika bu kaybı tarifeyle ya da sübvansiyonla durdurmaya çalışmak değil. Önerilen, kaybın yönetilmesi: işgücünün bölge içinde başka sanayi kollarına geçişi, yeniden eğitim ve bu sınıflardaki yüksek katma değerli üreticilerin dönüşüm programına alınması. ``Yönetilen uyum'' bir vazgeçiş değil. Bu sınıfların yüksek katma değerli kesimi (örme ile konfeksiyonu entegre eden, kendi kumaşını üreten tişört üreticileri) S2'de rekabetçi kalmayı sürdürüyor.

``Koru'' duruşundaki 5 sınıf ya küçük ya da çıkıştan sınırlı etkileniyor: boya, baskı ve apre hizmetleri, dokusuz kumaş, iş giysisi, diğer tekstil ve ayakkabı. Ayakkabı ve dokusuz kumaş için korumanın gerekçesi farklı. Ayakkabıda Türkiye net ithalatçı ve S4'te iç pazarda yerli üreticinin payı artıyor; dokusuz kumaş ise salgın döneminin gösterdiği gibi tıbbi tekstilde arz güvenliğinin girdisi.

\hypertarget{nace_kartlar_13-15-42-uxf6nerilen-paket-ve-sux131ralamasux131}{%
\subsection{4.2 Önerilen paket ve sıralaması}\label{nace_kartlar_13-15-42-uxf6nerilen-paket-ve-sux131ralamasux131}}

Paket, bölüm düzeyinde 6 ortak eylemden ve 8 sınıfa özgü 8 eylemden oluşuyor. Eylemler ilk parçadaki gibi üç katmanda sıralanıyor.

İlk katman hemen başlayacak ve pişmanlık yaratmayacak hazırlık. PEM kümülasyonunun sürekliliği (E-13-01) bu katmanın başında: S2 ve S3'ün menşe protokolüne PEM'e dinamik atıf ve kümülasyon bildirimini çıkış gününe bağlayan bir hüküm konmalı, Fas ve Tunus ile ilişki revize kurallara taşınmalı. Menşe ve izlenebilirlik veri altyapısı (E-13-02) aynı katmanda. Firma bazında kumaş menşe haritası, uzun dönemli tedarikçi beyanı, pamuk izlenebilirliği ve dijital ürün pasaportu hazırlığı tek bir sistemde toplanıyor. Bu altyapı S2'de menşe kuralına uyumu, her senaryoda da AB'nin zorla çalıştırma ve ekotasarım düzenlemelerine uyumu kolaylaştırıyor. Pazar çeşitlendirme (E-13-05) ve S2 hücum gündemi (E-13-06) de hemen başlıyor. Dönüşüm programı (E-14.13-01), pamuk izlenebilirliği (E-13.10-01) ve dokusuz kumaş için kriz kapasite envanteri (E-13.95-01) bu katmanın sınıf eylemleri.

İkinci katman çıkıştan yaklaşık bir yıl önce devreye girecek düzenlemeler: ev tekstilinde ticaret savunması hazırlığı (E-13.92-01), teknik tekstilde savunma ve tıbbi tedarik programı (E-13.96-01) ve tabaklamada yeşil dönüşüm (E-15.11-01). Bunlar kurumsal hazırlık ve yatırım süresi gerektiriyor.

Üçüncü katman yalnızca S4'te ve çıkış gününde devreye giriyor. Bunlar vergi takası (E-13-04), kısa çalışma ve işgücü uyumu (E-13-03, E-14.14-01) ve PEM kümülasyon kaybına karşı kumaş ihracatçılarına yönelik tampon (E-13.20-01). Bu katmanın tasarımında DTÖ'nün Sübvansiyonlar ve Telafi Edici Önlemler Anlaşması (ASCM) belirleyici. Sanayi ürünlerinde ihracata bağlı her destek yasak. Bu yüzden tamponda ihracat performansına bağlı bir ödeme yok. Kısa çalışma ödeneği ve işgücü uyum programı ihracattan bağımsız, sektörel ya da bölgesel kriz gerekçesine dayanan işgücü piyasası araçları (4447 sayılı Kanun ek md. 2). Eximbank'ın kredi sigortası ise piyasa primiyle sunuluyor. Vergi takası bütün menşelere eşit uygulanıyor ve GATT III:2 ile uyumlu.

\begin{landscape}

\textbf{Tablo 4.2 --- Bölüm düzeyi ortak eylemler (13--15)}

\nopagebreak[4]

\begingroup\gbtablo\footnotesize

\begin{longtable}[]{@{}
  >{\raggedright\arraybackslash}p{(\columnwidth - 18\tabcolsep) * \real{0.0556}}
  >{\raggedright\arraybackslash}p{(\columnwidth - 18\tabcolsep) * \real{0.1069}}
  >{\raggedright\arraybackslash}p{(\columnwidth - 18\tabcolsep) * \real{0.2228}}
  >{\raggedright\arraybackslash}p{(\columnwidth - 18\tabcolsep) * \real{0.0962}}
  >{\raggedright\arraybackslash}p{(\columnwidth - 18\tabcolsep) * \real{0.1133}}
  >{\raggedright\arraybackslash}p{(\columnwidth - 18\tabcolsep) * \real{0.0616}}
  >{\raggedright\arraybackslash}p{(\columnwidth - 18\tabcolsep) * \real{0.0698}}
  >{\raggedright\arraybackslash}p{(\columnwidth - 18\tabcolsep) * \real{0.0964}}
  >{\raggedright\arraybackslash}p{(\columnwidth - 18\tabcolsep) * \real{0.0784}}
  >{\raggedright\arraybackslash}p{(\columnwidth - 18\tabcolsep) * \real{0.0991}}@{}}
\toprule\noalign{}
\begin{minipage}[b]{\linewidth}\raggedright
\textbf{Eylem}
\end{minipage} & \begin{minipage}[b]{\linewidth}\raggedright
\textbf{Araç}
\end{minipage} & \begin{minipage}[b]{\linewidth}\raggedright
\textbf{Parametre}
\end{minipage} & \begin{minipage}[b]{\linewidth}\raggedright
\textbf{Sorumlu}
\end{minipage} & \begin{minipage}[b]{\linewidth}\raggedright
\textbf{Hukuki araç}
\end{minipage} & \begin{minipage}[b]{\linewidth}\raggedright
\textbf{Bütçe}
\end{minipage} & \begin{minipage}[b]{\linewidth}\raggedright
\textbf{Başlangıç → bitiş}
\end{minipage} & \begin{minipage}[b]{\linewidth}\raggedright
\textbf{Performans koşulu}
\end{minipage} & \begin{minipage}[b]{\linewidth}\raggedright
\textbf{Sunset}
\end{minipage} & \begin{minipage}[b]{\linewidth}\raggedright
\textbf{KPI}
\end{minipage} \\
\midrule\noalign{}
\endhead
\bottomrule\noalign{}
\endlastfoot
E-13-01 & PEM kümülasyonunun sürekliliği & (i) S2 ve S3 metninde PEM Sözleşmesi'ne dinamik atıf ve AB--Türkiye--PEM ortakları kümülasyon bildiriminin (PEM Ek I md. 8(2)) T'de yayımlanması; (ii) Fas ve Tunus ile STA'ların revize PEM kurallarına taşınması (H-M8); (iii) S4'te kümülasyon kaybına karşı ortakta baskı + konfeksiyon yolu (dokuma giyim) için ihracatçı rehberi & Ticaret Bakanlığı (AB GM, Anlaşmalar GM), Dışişleri Bakanlığı (AB Başkanlığı) & S2 menşe protokolü; TR--Fas ve TR--Tunus Ortak Komite kararları & İdari & Hemen → T & Kümülasyon bildirimi T'de yürürlükte; Fas ve Tunus ile revize kurallar & Yok (anlaşmayla kalıcı) & PEM ortaklarına kumaş ihracatı T+1'de T−12 düzeyinin en az \%95'i (S2) \\
E-13-02 & Menşe ve izlenebilirlik veri altyapısı & Firma bazında kumaş menşe haritası ve uzun dönemli tedarikçi beyanı (H-M6); Asya kumaşıyla yapılan ihracat için PEM kumaşına geçiş ve Türkiye'de baskı kapasitesi; aynı altyapıda pamuk izlenebilirliği (zorla çalıştırma yasağı, Tüzük 2024/3015) ve dijital ürün pasaportu hazırlığı & Ticaret Bakanlığı (İhracat GM, Gümrükler GM), İHKİB, İTHİB, TİM & Bakanlık genelgesi; ihracatçı birlikleri programı & ≈ 0,05 mlr TL/yıl \etiketT{} & Hemen → T+24 & Menşe haritası tamamlanan ihracatçı sayısı & T+24 & S2'de AB'ye giyim ihracatının menşe ispatlı payı ≥ \%90 \\
E-13-03 & S4 geçiş tamponu: kısa çalışma ve işgücü uyumu (yalnız S4) & AB'ye ihracatı yıllık \%15'ten fazla düşen giyim ve ev tekstili firmalarında kısa çalışma ödeneği (en çok 6 ay; bölgesel kriz kararı: İstanbul, Mardin--Batman, Şanlıurfa--Diyarbakır, Tekirdağ); İŞKUR mesleki eğitim ve yeniden yerleştirme; Eximbank'ın piyasa primli kredi sigortası & Çalışma ve Sosyal Güvenlik Bakanlığı, İŞKUR, Türk Eximbank & 4447 sayılı Kanun ek md. 2 (sektörel ve bölgesel kriz); İŞKUR Yönetim Kurulu kararı & ≈ 2,3 mlr TL/yıl × 2 yıl \etiketT{} & T → T+24 & Ödeneğin sürdüğü firmada istihdamın korunması & T+24 & Kapsamdaki firmalarda T+2'de istihdam ≥ T−12 düzeyinin \%90'ı \\
E-13-04 & Giyim ve ayakkabıda KDV takası (yalnız S4) & KDV giyimde \%10 → \%8,7, ayakkabıda \%10 → \%5,6; bütün menşelere eşit (Faz 2A, T-A). AB menşeli ürüne gelen c2 tarifesinin geliri indirimi finanse eder & Hazine ve Maliye Bakanlığı (Gelir Politikaları GM), Ticaret Bakanlığı & Cumhurbaşkanı Kararı (3065 sayılı Kanun md. 28) & Net ≈ 17,5 mlr TL/yıl (c2 gümrük geliri düşülü) \etiketT{} & T → T+36 & Giyim ve ayakkabı TÜFE alt endekslerinin yıllık artışı ≤ genel TÜFE & T+36 (S2'ye geçişte kalkar) & S4'te 13--15'in TÜFE etkisi ≤ 4 bp \\
E-13-05 & Pazar çeşitlendirme & Birleşik Krallık STA'sının yükseltilmesi ve AB--TR--BK kümülasyonu (H-T4); giyimde yakın tedarik ve marka programı (H-Y4): BK, ABD, Körfez ve Orta Asya'da ürüne özgü olmayan pazar geliştirme & Ticaret Bakanlığı (Anlaşmalar GM, İhracat GM), TİM & STA müzakeresi; İhracat Destekleri Kararı & ≈ 0,1 mlr TL/yıl \etiketT{} & Hemen → T+36 & Yeni pazarlarda alıcı sayısı & T+36 & Giyim ihracatında AB dışı pazarların payı T+3'te +5 puan \\
E-13-06 & Tekstil ve giyimde S2 hücum gündemi & (i) PEM dışı kumaşla üretilen giyim için menşe kotası (yıllık AB'ye giyim ihracatının ≈ \%10'u kadar; CETA'daki menşe kotası emsal \etiketD{}); (ii) örme giyimde de baskı + konfeksiyon seçeneği; (iii) ev tekstili ve sentetik iplikte ticaret savunması öncesi istişare mekanizması & Ticaret Bakanlığı & S2 anlaşma metni (menşe protokolü, ticaret savunması bölümü) & İdari & Hemen → T & S2 metninde menşe kotası & --- & Kota kapsamındaki giyimin AB'ye ihracatı T+5'te ≥ T−12 düzeyi \\
\end{longtable}

\endgroup

\end{landscape}

\emph{Kaynak:} Bütçeler \etiketT{}, 2026 fiyatlarıyla; kurumsal sorumluluk mevcut görev dağılımına göre. Mali yükler ve cari denge katkıları kümülatif defterde (Bölüm 6) toplanıyor.

\begin{landscape}

\textbf{Tablo 4.3 --- Sınıf eylemleri}

\nopagebreak[4]

\begingroup\gbtablo\footnotesize

\begin{longtable}[]{@{}
  >{\raggedright\arraybackslash}p{(\columnwidth - 18\tabcolsep) * \real{0.0728}}
  >{\raggedright\arraybackslash}p{(\columnwidth - 18\tabcolsep) * \real{0.0918}}
  >{\raggedright\arraybackslash}p{(\columnwidth - 18\tabcolsep) * \real{0.2308}}
  >{\raggedright\arraybackslash}p{(\columnwidth - 18\tabcolsep) * \real{0.1114}}
  >{\raggedright\arraybackslash}p{(\columnwidth - 18\tabcolsep) * \real{0.0978}}
  >{\raggedright\arraybackslash}p{(\columnwidth - 18\tabcolsep) * \real{0.0599}}
  >{\raggedright\arraybackslash}p{(\columnwidth - 18\tabcolsep) * \real{0.0710}}
  >{\raggedright\arraybackslash}p{(\columnwidth - 18\tabcolsep) * \real{0.1046}}
  >{\raggedleft\arraybackslash}p{(\columnwidth - 18\tabcolsep) * \real{0.0502}}
  >{\raggedright\arraybackslash}p{(\columnwidth - 18\tabcolsep) * \real{0.1096}}@{}}
\toprule\noalign{}
\begin{minipage}[b]{\linewidth}\raggedright
\textbf{Eylem}
\end{minipage} & \begin{minipage}[b]{\linewidth}\raggedright
\textbf{Araç}
\end{minipage} & \begin{minipage}[b]{\linewidth}\raggedright
\textbf{Parametre}
\end{minipage} & \begin{minipage}[b]{\linewidth}\raggedright
\textbf{Sorumlu}
\end{minipage} & \begin{minipage}[b]{\linewidth}\raggedright
\textbf{Hukuki araç}
\end{minipage} & \begin{minipage}[b]{\linewidth}\raggedright
\textbf{Bütçe}
\end{minipage} & \begin{minipage}[b]{\linewidth}\raggedright
\textbf{Başlangıç → bitiş}
\end{minipage} & \begin{minipage}[b]{\linewidth}\raggedright
\textbf{Performans koşulu}
\end{minipage} & \begin{minipage}[b]{\linewidth}\raggedleft
\textbf{Sunset}
\end{minipage} & \begin{minipage}[b]{\linewidth}\raggedright
\textbf{KPI}
\end{minipage} \\
\midrule\noalign{}
\endhead
\bottomrule\noalign{}
\endlastfoot
E-13.10-01 & Pamuk izlenebilirliği & İthal pamuk ve iplikte (Özbekistan, Türkmenistan, Hindistan) kaynak beyanı ve ithalatçı veri havuzu; E-13-02 altyapısında & Ticaret Bakanlığı (İthalat GM), TOB, İTHİB & Bakanlık genelgesi & ≈ 0,02 mlr TL/yıl \etiketT{} & Hemen → T+24 & Kaynak beyanı veren ithalatçı sayısı & T+24 & Tüzük 2024/3015 kapsamında Türk ipliğine yönelik AB soruşturması sıfır \\
E-13.20-01 & Kümülasyon kaybı tamponu (yalnız S4) & PEM ortaklarına kumaş satan ihracatçılar için ortakta baskı + konfeksiyon yolu rehberi; Fas, Tunus ve Mısır'daki alıcılarla menşe danışma masası; Eximbank'ın piyasa primli kredi sigortası & Ticaret Bakanlığı (İhracat GM), Türk Eximbank, İTHİB & Eximbank programı; bakanlık genelgesi & ≈ 0,05 mlr TL/yıl \etiketT{} & T → T+36 & Kapsamdaki ihracatçıların PEM ortaklarına ihracatı & T+36 & PEM ortaklarına kumaş ihracatı T+3'te ≥ T−12 düzeyinin \%70'i \\
E-13.92-01 & Ticaret savunması hazırlığı (ev tekstili) & AB'nin pamuklu nevresim ve havlu ithalatına olası damping soruşturmalarına hazırlık (emsal: AB'nin 1997'de Mısır, Hindistan ve Pakistan nevresimine uyguladığı damping vergisi ve DTÖ'deki DS141 davası \etiketV{}); maliyet ve fiyat dosyası; E-10-04 ekibine katılım & Ticaret Bakanlığı (İthalat GM, İhracat GM), İTHİB & Bakanlık iç talimatı & İdari & T−12 → T+36 & --- & T+36 & Soruşturma açılırsa 30 gün içinde tam yanıt \\
E-13.95-01 & Kriz kapasite envanteri & Tıbbi dokusuz kumaş (meltblown, SMS) kapasitesi ve stok envanteri; salgın ve afet planına bağlanması & Sağlık Bakanlığı (TİTCK), Sanayi ve Teknoloji Bakanlığı & Bakanlık genelgesi & İdari & Hemen → T+12 & Envanter yayımlandı & --- & Krizde 30 günlük maske ve önlük girdisinin yurt içinden karşılanması \\
E-13.96-01 & Savunma ve tıbbi teknik tekstil programı & SSB yerlileştirme listesinde balistik, aramid ve yangına dayanıklı tekstil; TÜBİTAK sanayi Ar-Ge destekleri; kamu alımında performans şartnamesi (menşe şartı olmadan) & Savunma Sanayii Başkanlığı, TÜBİTAK, Sağlık Bakanlığı & SSB yerlileştirme programı; TÜBİTAK çağrıları & ≈ 0,1 mlr TL/yıl \etiketT{} & T−12 → T+60 & Yerli tedarikle karşılanan kamu ihtiyacı & T+60 & Teknik tekstil ihracatı T+5'te +\%25 \\
E-14.13-01 & Dönüşüm programı & Tasarım, ürün geliştirme ve hızlı tedarik hizmeti (küçük parti, kısa termin); kesim ve dikimde otomasyon; dijital ürün pasaportu hazırlığı; yatırım teşvik sisteminde katma değer ölçütlü destek (menşe şartı yok) & Sanayi ve Teknoloji Bakanlığı, Ticaret Bakanlığı, İHKİB & Yatırım teşvik mevzuatı; İhracat Destekleri Kararı & ≈ 0,5 mlr TL/yıl \etiketT{} & Hemen → T+60 & Proje bazında kişi başı katma değer artışı & T+60 & Kişi başı katma değer (AB27 = 1) 0,31 → 0,40 (T+5) \\
E-14.14-01 & Bölgesel işgücü uyum programı & Mardin--Batman, Şanlıurfa--Diyarbakır ve İstanbul'da İŞKUR mesleki eğitimi, işbaşı eğitimi ve yeniden yerleştirme; giyimden ayrılan işçiler için bölge içi sanayi istihdamıyla eşleştirme & Çalışma ve Sosyal Güvenlik Bakanlığı, İŞKUR, kalkınma ajansları (Dicle, Karacadağ, İstanbul) & İŞKUR Yönetim Kurulu kararı & ≈ 0,4 mlr TL/yıl (yalnız S4) \etiketT{} & T → T+36 & Yeniden yerleştirme oranı & T+36 & Programdaki işçilerin \%60'ı 12 ay içinde yeniden istihdamda \\
E-15.11-01 & Yeşil dönüşüm & Organize deri sanayi bölgelerinde ortak arıtma kapasitesinin yenilenmesi; krom azaltımı ve kromsuz tabaklama; LWG sertifikası için danışmanlık & Çevre, Şehircilik ve İklim Değişikliği Bakanlığı; Sanayi ve Teknoloji Bakanlığı; organize sanayi bölgeleri & Kalkınma ajansı ve yeşil dönüşüm destek programları & ≈ 0,05 mlr TL/yıl \etiketT{} & T−12 → T+36 & Sertifikalı tabakhane sayısı & T+36 & AB'ye ihracatta LWG sertifikalı payı ≥ \%60 \\
\end{longtable}

\endgroup

\end{landscape}

\emph{Kaynak:} Bütçeler \etiketT{}, 2026 fiyatlarıyla. Tarihler T'ye göre (T2028'de T = 1.1.2028, T2029'da 1.1.2029).

Paketin yıllık mali maliyeti S2 yolunda 0,82 milyar TL, S4 yolunda 21,07 milyar TL. S4'teki tutarın çoğu vergi takasının net maliyeti (17,5 milyar TL). Kısa çalışma ödeneği (yılda 2,3 milyar TL, iki yıl) ve işgücü uyum programı (0,4 milyar TL) S4'e özgü. Dönüşüm programı (0,5 milyar TL) her iki yolda da uygulanıyor. Paketin S2'de cari dengeye katkısı yılda 150 milyon dolar \etiketH{}. Bunun çoğu menşe kotasının ve menşe hazırlığının S2'de geri kazandıracağı ihracattan geliyor. S4'te paketin cari dengeye kayda değer bir katkısı yok; S4'ün bu sektördeki cari denge maliyeti politika ile telafi edilemiyor.

\hypertarget{nace_kartlar_13-15-43-muxfczakere-pozisyonu}{%
\subsection{4.3 Müzakere pozisyonu}\label{nace_kartlar_13-15-43-muxfczakere-pozisyonu}}

AB ile müzakerede bu sektörün tutumu savunma. Hedef, Gümrük Birliği'ndeki sıfır tarifenin S2'de sürmesi ve menşe kurallarının ihracatçının bugünkü tedarik yapısını cezalandırmaması. Tarife tarafında pazarlık konusu yok: AB, sanayi ürünlerinde bütün serbest ticaret anlaşmalarında sıfır tarifeyi kabul ediyor. Asıl müzakere menşe protokolünde ve dört talep içeriyor. Birincisi PEM'e dinamik atıf ve kümülasyon bildiriminin anlaşmanın yürürlüğüyle aynı gün yayımlanması (E-13-01). İkincisi PEM dışı kumaşla üretilen giyim için bir menşe kotası: yıllık AB'ye giyim ihracatının yaklaşık \%10'u kadar bir hacim, çift dönüşüm şartı aranmadan tercihli tarifeden yararlanmalı. AB'nin Kanada ile anlaşmasında (CETA) giyim için benzer menşe kotaları var \etiketD{}. Üçüncüsü, dokuma giyimde var olan ``baskı + konfeksiyon'' seçeneğinin örme giyime genişletilmesi. Dördüncüsü, ev tekstili ve sentetik iplikte bir ticaret savunması soruşturmasından önce zorunlu istişare (E-13-06).

Bu taleplerin AB'deki karşılığı çeşitli. Menşe kotası, AB'nin Portekiz, İtalya ve İspanya'daki kumaş üreticilerinin itirazıyla karşılaşacaktır \etiketH{}. Buna karşılık AB'nin büyük perakendecileri ve markaları, Türk tedarik zincirinin S2'de kesintisiz sürmesinden yana. Bu, Türkiye'nin müzakerede kullanabileceği en güçlü iç AB müttefiki. Inditex, H\&M ve Alman markalarının Türkiye'deki tedarik ağları, AB içinde Türkiye'nin menşe taleplerinin doğal savunucusu olabilir \etiketH{}. Türkiye'nin karşılığında verebileceği tavizler sınırlı. AB'nin bu sektörde Türkiye'den isteyeceği tek şey, AB menşeli kumaş ve giyimde S2'de sıfır tarife ve kendi ihracatçısının menşe maliyetinin düşük tutulması. Bu da Türkiye'nin zaten vereceği bir şey.

Üçüncü ülkelerle önceliklerin başında Birleşik Krallık geliyor. Türkiye--Birleşik Krallık STA'sı Gümrük Birliği'nden bağımsız ve yükseltme müzakeresi sürüyor (Faz 2C, H-T4). Bu müzakerede menşe faslının AB--Türkiye--Birleşik Krallık üçlü kümülasyonuna hazır yazılması ve giyimde Birleşik Krallık'ın Türk kumaşını menşeli sayması isteniyor (E-13-05). Körfez ülkeleriyle STA ayakkabı, halı ve ev tekstili için, ABD ile bir ticaret anlaşması ise halı, ev tekstili ve denim için büyük fırsat. ABD'nin 2025'ten beri uyguladığı ek tarifeler bu fırsatı daraltıyor \etiketD{}. Mısır, Fas ve Tunus ile ilişkide amaç, PEM kümülasyonunu revize kurallarda sürdürmek ve S4'te zincirin kopmasının maliyetini sınırlamak. Bu ülkeler aynı zamanda Türk tekstil yatırımlarının kaydığı yerler.

\hypertarget{nace_kartlar_13-15-44-erken-uyarux131-ve-tetikler}{%
\subsection{4.4 Erken uyarı ve tetikler}\label{nace_kartlar_13-15-44-erken-uyarux131-ve-tetikler}}

Paketin S4 katmanı koşullu: tampon önlemleri ancak belirli eşikler aşıldığında devreye giriyor. En kritik iki gösterge AB'ye giyim ihracatı (aylık Comext verisi; yıllık \%15'lik düşüş) ve PEM ortaklarına kumaş ihracatı (aylık TÜİK verisi; yıllık \%20'lik düşüş). İlki kısa çalışma ve pazar çeşitlendirme önlemlerini, ikincisi kumaş ihracatçılarına yönelik tamponu başlatıyor. Bölgesel istihdam göstergesi olarak SGK'nın aylık sigortalı verisi, giyimin yoğunlaştığı bölgelerde izleniyor. TÜFE tarafında ayakkabı alt endeksinin genel TÜFE'den belirgin biçimde ayrışması, takasın yeniden ayarlanması için bir tetik.

\Needspace{7\baselineskip}

\textbf{Tablo 4.4 --- Erken uyarı göstergeleri ve tetikler}

\nopagebreak[4]

\begingroup\gbtablo\small

\begin{longtable}[]{@{}
  >{\raggedright\arraybackslash}p{(\columnwidth - 6\tabcolsep) * \real{0.0679}}
  >{\raggedright\arraybackslash}p{(\columnwidth - 6\tabcolsep) * \real{0.3528}}
  >{\raggedright\arraybackslash}p{(\columnwidth - 6\tabcolsep) * \real{0.2164}}
  >{\raggedright\arraybackslash}p{(\columnwidth - 6\tabcolsep) * \real{0.3628}}@{}}
\toprule\noalign{}
\begin{minipage}[b]{\linewidth}\raggedright
\textbf{Sınıf}
\end{minipage} & \begin{minipage}[b]{\linewidth}\raggedright
\textbf{Gösterge}
\end{minipage} & \begin{minipage}[b]{\linewidth}\raggedright
\textbf{Eşik}
\end{minipage} & \begin{minipage}[b]{\linewidth}\raggedright
\textbf{Tetiklenecek önlem}
\end{minipage} \\
\midrule\noalign{}
\endhead
\bottomrule\noalign{}
\endlastfoot
13.10 & AB'nin zorla çalıştırma risk veritabanında Türk tedarik zincirine ilişkin bildirim & Bildirim & E-13-02: izlenebilirlik paketi; AB ithalatçılarıyla ortak savunma \\
13.20 & PEM ortaklarına dokuma ve örme kumaş ihracatı (aylık) & Yıllık −\%20 & E-13.20-01'in genişletilmesi; S3 ya da S2 müzakeresinde kümülasyonun öne alınması \\
13.91 & Belarus ve Rusya'ya örme kumaş ihracatı (aylık) & Yıllık −\%25 & Pazar çeşitlendirme (E-13-05) \\
13.92 & AB'ye nevresim ve havlu ihracatı (KN 6302; aylık) & Yıllık −\%15 & E-13-03 ve E-13-05; AB Resmî Gazetesi'nde soruşturma ilanı izlenir \\
14.13 & AB'ye dış giyim ihracatı (Comext; aylık) & Yıllık −\%15 & E-13-03'ün devreye girmesi; E-13-05 hızlandırılır \\
14.13 & Bangladeş'in LDC'den çıkış tarihi ve AB'nin EBA geçiş kararı & Karar & S2 müzakeresinde menşe kotası talebinin (E-13-06) güçlendirilmesi \\
14.14 & Giyimde SGK sigortalı sayısı (aylık; Mardin--Batman ve İstanbul) & Yıllık −\%10 & E-14.14-01'in kapsamının genişletilmesi \\
14.31 & AB'ye çorap ihracatı (KN 6115; aylık) & Yıllık −\%15 & E-13-03; BK ve ABD'de pazar geliştirme (E-13-05) \\
14.39 & AB'ye triko ihracatı (KN 6110; aylık) & Yıllık −\%20 & E-13-03 ve E-14.14-01 \\
15.20 & Ayakkabı TÜFE alt endeksi (aylık) & 3 ay üst üste genel TÜFE'nin 5 puan üstünde & E-13-04'teki KDV indiriminin genişletilmesi (takas tavanı içinde) \\
\end{longtable}

\endgroup

\emph{Kaynak:} Göstergeler aylık izlenir: TÜİK dış ticaret ve TÜFE alt endeksleri, Eurostat Comext, AB Resmî Gazetesi.

\hypertarget{nace_kartlar_13-15-45-kazananlar-kaybedenler-ve-buxf6lgeler}{%
\subsection{4.5 Kazananlar, kaybedenler ve bölgeler}\label{nace_kartlar_13-15-45-kazananlar-kaybedenler-ve-buxf6lgeler}}

Çıkışın bu sektördeki dağılım etkileri ilk parçadakinden çok daha büyük ve bölgesel olarak yoğun. S4'te kaybedenler belli: AB'ye satan giyim üreticileri, özellikle fason çalışan ve tek bir AB alıcısına bağlı KOBİ'ler; PEM ortaklarına kumaş satan dokuma ve örme üreticileri; ve bu işletmelerde çalışan, çoğunluğu kadın ve genç, düşük ücretli işgücü \etiketH{}. Kazananlar sınırlı: iç pazara satan ayakkabı ve çanta üreticileri, c2 tarifesiyle AB ürününe karşı iç pazarda pay kazanan markalar ve takasla ucuzlayan giyimi alan tüketiciler.

\Needspace{7\baselineskip}

\textbf{Tablo 4.5 --- Bölgeler: S4'te istihdam kaybının en büyük olduğu 10 İBBS-2 bölgesi (T+5, P50)}

\nopagebreak[4]

\begingroup\gbtablo\footnotesize

\begin{longtable}[]{@{}
  >{\raggedright\arraybackslash}p{(\columnwidth - 12\tabcolsep) * \real{0.2095}}
  >{\raggedleft\arraybackslash}p{(\columnwidth - 12\tabcolsep) * \real{0.1451}}
  >{\raggedleft\arraybackslash}p{(\columnwidth - 12\tabcolsep) * \real{0.1643}}
  >{\raggedleft\arraybackslash}p{(\columnwidth - 12\tabcolsep) * \real{0.0921}}
  >{\raggedleft\arraybackslash}p{(\columnwidth - 12\tabcolsep) * \real{0.0921}}
  >{\raggedleft\arraybackslash}p{(\columnwidth - 12\tabcolsep) * \real{0.1139}}
  >{\raggedright\arraybackslash}p{(\columnwidth - 12\tabcolsep) * \real{0.1831}}@{}}
\toprule\noalign{}
\begin{minipage}[b]{\linewidth}\raggedright
\textbf{Bölge (İBBS-2)}
\end{minipage} & \begin{minipage}[b]{\linewidth}\raggedleft
\textbf{13--15 çalışanı (bin)}
\end{minipage} & \begin{minipage}[b]{\linewidth}\raggedleft
\textbf{Bölge imalatındaki payı}
\end{minipage} & \begin{minipage}[b]{\linewidth}\raggedleft
\textbf{S2 (kişi)}
\end{minipage} & \begin{minipage}[b]{\linewidth}\raggedleft
\textbf{S4 (kişi)}
\end{minipage} & \begin{minipage}[b]{\linewidth}\raggedleft
\textbf{S4 ÷ çalışan}
\end{minipage} & \begin{minipage}[b]{\linewidth}\raggedright
\textbf{En çok etkilenen sınıflar (S4)}
\end{minipage} \\
\midrule\noalign{}
\endhead
\bottomrule\noalign{}
\endlastfoot
İstanbul & 333,7 & \%31 & −4.475 & −12.209 & −\%3,7 & 14.14 (−3.845), 14.13 (−3.687) \\
Mardin, Batman, Şırnak, Siirt & 87,0 & \%78 & −1.365 & −3.388 & −\%3,9 & 14.13 (−1.663), 14.14 (−1.289) \\
Bursa, Eskişehir, Bilecik & 106,3 & \%22 & −647 & −2.873 & −\%2,7 & 13.20 (−1.249), 14.13 (−426) \\
Tekirdağ, Edirne, Kırklareli & 72,9 & \%32 & −775 & −2.341 & −\%3,2 & 14.31 (−727), 14.13 (−394) \\
Gaziantep, Adıyaman, Kilis & 119,7 & \%55 & −855 & −2.228 & −\%1,9 & 14.13 (−473), 13.92 (−428) \\
Aydın, Denizli, Muğla & 58,2 & \%35 & −490 & −1.662 & −\%2,9 & 13.20 (−510), 13.92 (−493) \\
Hatay, Kahramanmaraş, Osmaniye & 47,1 & \%34 & −333 & −1.506 & −\%3,2 & 13.20 (−557), 14.39 (−346) \\
İzmir & 45,9 & \%15 & −561 & −1.446 & −\%3,2 & 14.14 (−659), 14.13 (−470) \\
Adana, Mersin & 42,1 & \%23 & −433 & −1.333 & −\%3,2 & 14.13 (−525), 13.20 (−347) \\
Şanlıurfa, Diyarbakır & 40,8 & \%47 & −524 & −1.280 & −\%3,1 & 14.13 (−581), 14.14 (−468) \\
\textbf{13--15 toplamı (bütün bölgeler)} & \textbf{1.237,9} & \textbf{\%25} & \textbf{−13.862} & \textbf{−38.973} & \textbf{−\%3,1} & --- \\
\end{longtable}

\endgroup

\emph{Kaynak:} \etiketT{} f3b (YSHİ 2024 sınıf çalışanı, SGK 2025 bölüm × il sigortalısıyla RAS) ve f3e. Bölge değerleri sınıf P50'lerinin bölge paylarıyla dağıtılmasından; toplam, sınıf P50'lerinin toplamıdır (grup yüzdeliği değil).

Bölgesel tablo iki farklı yoğunlaşma gösteriyor. Mutlak kayıpta İstanbul açık ara önde (12.209 kişi): giyimin tasarım, ticaret ve üretim merkezi burası. İstanbul'un sanayisi ise çeşitli: tekstil, giyim ve deri bölgenin imalat istihdamının \%31 kadarı ve işini kaybedenin başka bir sanayi işine geçme olanağı görece geniş. Asıl kırılganlık güneydoğuda. Mardin, Batman, Şırnak ve Siirt bölgesinde (TRC3) bu üç bölüm imalat istihdamının \%78 kadarını oluşturuyor. Son on yılda bölgesel teşviklerle kurulan büyük konfeksiyon tesisleri bölgenin neredeyse tek sanayi kolu \etiketH{}. Şanlıurfa ve Diyarbakır'da (TRC2) bu oran \%47, Gaziantep, Adıyaman ve Kilis'te (TRC1) \%55; Gaziantep'te ağırlık halı, iplik ve dokusuz kumaşta. Bu bölgelerde S4'ün istihdam kaybı mutlak olarak İstanbul'dan küçük, ama işini kaybedenin bölge içinde başka bir sanayi işine geçme olasılığı çok daha düşük.

Bu yüzden S4 katmanının işgücü önlemleri bölgesel olarak hedeflendi. Kısa çalışma ödeneği için ``bölgesel kriz'' kararı İstanbul, Mardin--Batman, Şanlıurfa--Diyarbakır ve Tekirdağ için öngörülüyor (E-13-03). İşgücü uyum programı ise güneydoğuda kalkınma ajanslarıyla birlikte yürütülüyor (E-14.14-01). Ürüne özgü gelir ya da ihracat desteği önerilmiyor. Böyle bir destek hem DTÖ'ye aykırı hem mali açıdan pahalı. Üstelik hasmane bir AB'nin eline telafi edici vergi gerekçesi verir. Bölgesel teşvik sisteminin sunduğu yatırım ve sigorta primi destekleri zaten yürürlükte. Burada önerilen, bu desteklerin dönüşüm programıyla (E-14.13-01) ilişkilendirilmesi: düşük katma değerli fason dikime değil, entegre ve yüksek katma değerli üretime yönlendirilmesi.

\hypertarget{nace_kartlar_13-15-5-sux131nux131f-incelemeleri}{%
\section{5. Sınıf incelemeleri}\label{nace_kartlar_13-15-5-sux131nux131f-incelemeleri}}

Bu bölüm sınıfları tek tek ele alıyor. Her incelemenin başındaki künye satırı sınıfın Gümrük Birliği statüsünü, arketipini, duruşunu, AB ile müzakere tutumunu ve değerlendirmenin güven düzeyini veriyor; sayısal göstergelerin tamamı Ek C'deki veri sayfasında. Etkisi büyük sınıflar için senaryo tablosu metnin içinde; diğerlerinde veri sayfasının altında. İnceleme sırası NACE kodunu izliyor. Etkisi sınırlı altı sınıf Bölüm 5.16'da tek tabloda toplanıyor.

\hypertarget{nace_kartlar_13-15-51-1310-iux307plik}{%
\subsection{5.1 13.10 İplik}\label{nace_kartlar_13-15-51-1310-iux307plik}}

\emph{GB statüsü: GB · Arketip: A (AB'ye entegre ihracat şampiyonu) · Duruş: Geliştir · AB ile müzakere: savunma · Güven: Yüksek · Veri sayfası: Ek C.1}

Türkiye dünyanın büyük iplik üreticilerinden biri. Sınıfta 86.468 kişi çalışıyor; üretim Gaziantep--Adıyaman--Kilis (\%20), Tekirdağ--Edirne--Kırklareli, Hatay--Kahramanmaraş--Osmaniye ve Ege'de yoğunlaşıyor. Büyük işletmelerin istihdam payı yüksek (\%65) ve sınıfın yurt içi tedarikçi bağlantıları imalatın en güçlüleri arasında. İhracat 2024'te 2.056 milyon dolar ve 2020'den bu yana \%41 arttı. AB'ye ihracatın ilk pazarları İtalya (\%20), Portekiz (\%17) ve Almanya; yani Türk ipliği AB'nin örme ve dokuma sanayisinin girdisi. AB pazarında Türkiye Çin'den sonra ikinci tedarikçi (\%24,3).

AB tarifesi düşük (S4'te \%3,8), ama iplik homojen bir ürün ve alıcı fiyat farkına çok duyarlı: sınıfın ikame esnekliği 15,0, imalat medyanının iki katı. Bu yüzden düşük bir tarife bile kayda değer bir hacim kaybı yaratıyor. S4'te AB'ye ihracat 229 milyon dolar, istihdam 1.031 kişi azalıyor. Aynı neden S2'nin kaybını da büyütüyor. Menşe kuralı iplikte kolay (doğal liflerin eğrilmesi menşe kazandırıyor ve ithal pamuk serbest) ama menşe ispatının ve sınır sürtünmesinin küçük maliyeti yüksek esneklikle çarpılınca S2'nin kaybı (146 milyon dolar) S4'ünkine yaklaşıyor. İplik, S2'de ``maliyetsiz'' görünen menşe belgesinin bile önemli olduğu bir sınıf. Belgelemenin dijital ve ucuz olması (E-13-02) burada doğrudan ihracata dönüşüyor.

İthalat tarafında sınıf iki kaynağa bağlı. İplik ithalatının (1.286 milyon dolar) ilk kaynağı Özbekistan (\%25); onu Hindistan, Çin, Bangladeş ve Mısır izliyor. Ham pamukta ilk iki kaynak Brezilya (ithalatın \%31 kadarı) ve ABD (\%26). Bu tedarik yapısı, AB'nin zorla çalıştırma yasağı (Tüzük 2024/3015) açısından sınıfın en hassas noktası. Özbekistan'da pamuk hasadında sistematik zorla çalıştırmanın sona erdiği ILO'nun izlemesiyle 2022'de raporlandı \etiketD{}; ama Orta Asya pamuğu AB'nin risk değerlendirmelerinde öne çıkmaya devam edebilir \etiketH{}. Önerilen eylem, ithal pamuk ve iplikte kaynak beyanı ve ithalatçılarla ortak bir veri havuzu (E-13.10-01). Bu eylem E-13-02'nin altyapısını kullanıyor.

Önerilen duruş ``Geliştir''. Sınıfın kırmızı çizgi yetkinliği, iplikten kumaşa uzanan yerli zincirin ilk halkası olması; Türk giyiminin menşe kuralını rahatça karşılamasının nedeni bu zincir. Hedefler, AB dışı pazarların (Mısır, Birleşik Krallık, ABD, Fas) payını korumak ve geri dönüştürülmüş lif ile yüksek performanslı sentetik iplikte payı artırmak. AB ile müzakerede tutum savunma. Değerlendirmenin güveni yüksek.

\Needspace{7\baselineskip}

\textbf{Tablo 5.1 --- 13.10 İplik: senaryo etkileri (T+5'te yıllık; P50, parantez içinde P10 / P90)}

\nopagebreak[4]

\begingroup\gbtablo\small

\begin{longtable}[]{@{}
  >{\raggedright\arraybackslash}p{(\columnwidth - 6\tabcolsep) * \real{0.2746}}
  >{\raggedleft\arraybackslash}p{(\columnwidth - 6\tabcolsep) * \real{0.2273}}
  >{\raggedleft\arraybackslash}p{(\columnwidth - 6\tabcolsep) * \real{0.2379}}
  >{\raggedleft\arraybackslash}p{(\columnwidth - 6\tabcolsep) * \real{0.2602}}@{}}
\toprule\noalign{}
\begin{minipage}[b]{\linewidth}\raggedright
\textbf{Değişken}
\end{minipage} & \begin{minipage}[b]{\linewidth}\raggedleft
\textbf{S2}
\end{minipage} & \begin{minipage}[b]{\linewidth}\raggedleft
\textbf{S3}
\end{minipage} & \begin{minipage}[b]{\linewidth}\raggedleft
\textbf{S4}
\end{minipage} \\
\midrule\noalign{}
\endhead
\bottomrule\noalign{}
\endlastfoot
AB'ye ihracat (\$ mn) & −146 (−338 / +18) & −183 (−350 / −36) & −229 (−378 / −97) \\
AB'ye ihracat (\%) & −16,6 (−38,6 / +2,0) & −20,8 (−39,9 / −4,1) & −26,1 (−43,1 / −11,1) \\
Toplam ihracat (\$ mn) & −104 (−246 / +12) & −129 (−249 / −25) & −183 (−306 / −78) \\
Nihai mal ithalatı (\$ mn) & −11 (−38 / +10) & −16 (−42 / +6) & −28 (−52 / −6) \\
Üretim (\$ mn) & −93 (−210 / +2) & −112 (−210 / −30) & −154 (−253 / −72) \\
İstihdam (kişi) & −621 (−1.405 / +14) & −744 (−1.413 / −198) & −1.031 (−1.712 / −466) \\
İstihdam (\%) & −0,7 (−1,6 / 0,0) & −0,9 (−1,6 / −0,2) & −1,2 (−2,0 / −0,5) \\
Cari denge (\$ mn) & −72 (−162 / +2) & −87 (−162 / −23) & −119 (−196 / −56) \\
TÜFE (bp) & +0,01 (+0,01 / +0,01) & +0,02 (+0,02 / +0,02) & +0,06 (+0,05 / +0,06) \\
\end{longtable}

\endgroup

\emph{Kaynak:} \etiketT{} f3e; toplamlar Faz 2B genel denge dağılımına kalibre. Yüzdeler 2024 düzeyine göre.

\hypertarget{nace_kartlar_13-15-52-1320-dokuma-kumaux15f}{%
\subsection{5.2 13.20 Dokuma kumaş}\label{nace_kartlar_13-15-52-1320-dokuma-kumaux15f}}

\emph{GB statüsü: GB · Arketip: A (AB'ye entegre ihracat şampiyonu) / B (küresel rekabetçi, AB bağımlılığı düşük) · Duruş: Geliştir · AB ile müzakere: savunma · Güven: Orta · Veri sayfası: Ek C.2}

Dokuma, Türk tekstilinin bel kemiği: 96.394 çalışan, 3.088 milyon dolar ihracat ve imalatın en güçlü yurt içi tedarik bağlantılarından biri. Üretim Bursa--Eskişehir--Bilecik (\%32), Hatay--Kahramanmaraş--Osmaniye ve Denizli'de yoğunlaşıyor. Denim bu sınıfın yıldızı: Türkiye dünya denim ihracatında \%9,9 payla üçüncü. AB'ye ihracat 986 milyon dolar ve ilk pazar İspanya (\%28); İspanyol perakendecilerinin tedarik zinciri burada belirleyici. AB pazarında Türkiye Çin'den sonra ikinci tedarikçi.

Sınıfın asıl özelliği AB dışı pazarlarında. Dokuma kumaşın AB dışı ihracatının \%32 kadarı Fas'a, \%8 kadarı Tunus'a gidiyor. Bu kumaşın büyük kısmı orada giysiye dönüşüp yine İspanya'ya ve Fransa'ya satılıyor. Bu yüzden S4'te sınıf iki kez kaybediyor. Birincisi AB'ye doğrudan ihracatta: \%8'lik tarifeyle 208 milyon dolar. İkincisi PEM kümülasyonunun kopmasıyla AB dışı ihracatta: 457 milyon dolar. Toplam ihracat kaybı 623 milyon dolar. İstihdam kaybı 3.966 kişi (P10'da 5.185); mevcut model sonuçlarında imalatın S4'te en çok istihdam kaybeden üçüncü sınıfı. S2'de kümülasyon sürdüğü ve menşe kuralı kolay olduğu için kayıp küçük (216 kişi). Dokuma, S2 ile S4 arasındaki farkın en büyük olduğu tekstil sınıfı.

İthalat tarafında sınıf hem rakip hem tedarikçi. Dokuma kumaş ithalatının (1.322 milyon dolar) ilk kaynakları Çin (\%23), Mısır (\%18) ve İtalya (\%16). S4'te AB'den gelen kumaşın tarifesi \%28,0 oluyor. Bu, İtalyan kumaşını kullanan üst segment konfeksiyoncunun maliyetini artırıyor; ihracatçı ise DİR'le vergisiz getirebiliyor. Yerli dokumacı için bu bir pay kazanma fırsatı, iç pazara üreten konfeksiyoncu için bir maliyet.

Önerilen duruş ``Geliştir''; strateji iki ayaklı. Birinci ayak PEM kümülasyonunun S2 ve S3'te kesintisiz sürmesi (E-13-01); sınıf için müzakerenin en önemli maddesi bu. İkinci ayak S4'e hazırlık: PEM ortaklarına satan ihracatçılara ortakta baskı ve dikim yolunu anlatan bir rehber, Fas, Tunus ve Mısır'daki alıcılarla bir menşe danışma masası ve piyasa primli kredi sigortası (E-13.20-01). Kırmızı çizgi yetkinlikler denim ve gömleklik kumaşta tasarım, boya ve apre bilgisi ile hızlı numune. Kazananlar S4'te iç pazara satan dokumacılar, kaybedenler PEM ortaklarına satanlar. Değerlendirmenin güveni orta: kümülasyon kanalının parametreleri (κ ve λ) uzman yargısı \etiketH{}.

\Needspace{7\baselineskip}

\textbf{Tablo 5.2 --- 13.20 Dokuma kumaş: senaryo etkileri (T+5'te yıllık; P50, parantez içinde P10 / P90)}

\nopagebreak[4]

\begingroup\gbtablo\small

\begin{longtable}[]{@{}
  >{\raggedright\arraybackslash}p{(\columnwidth - 6\tabcolsep) * \real{0.3572}}
  >{\raggedleft\arraybackslash}p{(\columnwidth - 6\tabcolsep) * \real{0.2003}}
  >{\raggedleft\arraybackslash}p{(\columnwidth - 6\tabcolsep) * \real{0.2003}}
  >{\raggedleft\arraybackslash}p{(\columnwidth - 6\tabcolsep) * \real{0.2422}}@{}}
\toprule\noalign{}
\begin{minipage}[b]{\linewidth}\raggedright
\textbf{Değişken}
\end{minipage} & \begin{minipage}[b]{\linewidth}\raggedleft
\textbf{S2}
\end{minipage} & \begin{minipage}[b]{\linewidth}\raggedleft
\textbf{S3}
\end{minipage} & \begin{minipage}[b]{\linewidth}\raggedleft
\textbf{S4}
\end{minipage} \\
\midrule\noalign{}
\endhead
\bottomrule\noalign{}
\endlastfoot
AB'ye ihracat (\$ mn) & −83 (−193 / +20) & −107 (−208 / −20) & −208 (−337 / −88) \\
AB'ye ihracat (\%) & −8,4 (−19,6 / +2,0) & −10,9 (−21,1 / −2,0) & −21,1 (−34,2 / −8,9) \\
Toplam ihracat (\$ mn) & −60 (−141 / +13) & −76 (−148 / −14) & −623 (−784 / −466) \\
Bunun içinde PEM kümülasyonu (\$ mn) & 0 (0 / 0) & 0 (0 / 0) & −457 (−583 / −328) \\
Nihai mal ithalatı (\$ mn) & −27 (−89 / +26) & −39 (−97 / +15) & −88 (−158 / −18) \\
Üretim (\$ mn) & −29 (−60 / −8) & −36 (−62 / −19) & −536 (−667 / −401) \\
İstihdam (kişi) & −216 (−460 / −59) & −266 (−474 / −140) & −3.966 (−5.185 / −2.895) \\
İstihdam (\%) & −0,2 (−0,5 / −0,1) & −0,3 (−0,5 / −0,1) & −4,1 (−5,4 / −3,0) \\
Cari denge (\$ mn) & −23 (−47 / −6) & −28 (−48 / −15) & −415 (−516 / −310) \\
TÜFE (bp) & +0,06 (+0,05 / +0,07) & +0,11 (+0,10 / +0,13) & +0,54 (+0,46 / +0,61) \\
\end{longtable}

\endgroup

\emph{Kaynak:} \etiketT{} f3e; toplamlar Faz 2B genel denge dağılımına kalibre. Yüzdeler 2024 düzeyine göre.

\hypertarget{nace_kartlar_13-15-53-1391-uxf6rme-kumaux15f}{%
\subsection{5.3 13.91 Örme kumaş}\label{nace_kartlar_13-15-53-1391-uxf6rme-kumaux15f}}

\emph{GB statüsü: GB · Arketip: A (AB'ye entegre ihracat şampiyonu) / B (küresel rekabetçi, AB bağımlılığı düşük) · Duruş: Geliştir · AB ile müzakere: savunma · Güven: Orta · Veri sayfası: Ek C.4}

Örme kumaş, Türk tişört, iç giyim ve spor giyiminin girdisi. Sınıfta 36.092 kişi çalışıyor; İstanbul (\%31), Bursa ve Hatay--Kahramanmaraş öne çıkıyor. Türkiye dünya örme kumaş ihracatında ikinci (\%6,7) ve AB pazarında AB dışı tedarikin \%31,9 kadarını sağlıyor. Ancak bu pay 2020'de \%39,3 idi; sınıf AB pazarında Çin'e pay kaptırıyor. AB'deki ilk pazarlar İtalya (\%34) ve Bulgaristan (\%14). Bulgaristan'daki konfeksiyoncular Türk örme kumaşını dikip AB içinde satıyor \etiketH{}.

Çıkışın sınıfa etkisi dokumayla aynı yapıda. S4'te AB'ye doğrudan ihracat 118 milyon dolar, PEM ortaklarına (Fas, Sırbistan, Ukrayna) ihracat 199 milyon dolar azalıyor. İstihdam kaybı 1.632 kişi ve dağılımın alt ucunda sınıf istihdamının \%5,9 kadarı. S2'de kayıp 154 kişi. Sınıfın AB dışındaki en büyük iki pazarı Belarus (\%22) ve Rusya (\%12). Bu pazarlar Gümrük Birliği'yle ilgili değil, ama hasmane bir AB'nin yaptırım politikası onları da riske sokabilir \etiketH{}. Bunun için ayrı bir tetik tanımlandı.

Önerilen duruş ``Geliştir'': sınıfın stratejisi dokumayla ortak (E-13-01, E-13-02). Kırmızı çizgi yetkinlikler, pamuklu ve elastanlı örme kumaşta hız ve boya--apre kalitesi. Hedef, AB pazarındaki pay kaybını durdurmak ve yakın tedarik talebini (Bulgaristan, Romanya, Yunanistan ve Kuzey Afrika konfeksiyonu) Türk kumaşına bağlamak. Değerlendirmenin güveni orta.

\Needspace{7\baselineskip}

\textbf{Tablo 5.3 --- 13.91 Örme kumaş: senaryo etkileri (T+5'te yıllık; P50, parantez içinde P10 / P90)}

\nopagebreak[4]

\begingroup\gbtablo\small

\begin{longtable}[]{@{}
  >{\raggedright\arraybackslash}p{(\columnwidth - 6\tabcolsep) * \real{0.3572}}
  >{\raggedleft\arraybackslash}p{(\columnwidth - 6\tabcolsep) * \real{0.2003}}
  >{\raggedleft\arraybackslash}p{(\columnwidth - 6\tabcolsep) * \real{0.2003}}
  >{\raggedleft\arraybackslash}p{(\columnwidth - 6\tabcolsep) * \real{0.2422}}@{}}
\toprule\noalign{}
\begin{minipage}[b]{\linewidth}\raggedright
\textbf{Değişken}
\end{minipage} & \begin{minipage}[b]{\linewidth}\raggedleft
\textbf{S2}
\end{minipage} & \begin{minipage}[b]{\linewidth}\raggedleft
\textbf{S3}
\end{minipage} & \begin{minipage}[b]{\linewidth}\raggedleft
\textbf{S4}
\end{minipage} \\
\midrule\noalign{}
\endhead
\bottomrule\noalign{}
\endlastfoot
AB'ye ihracat (\$ mn) & −47 (−109 / +11) & −60 (−118 / −11) & −118 (−193 / −49) \\
AB'ye ihracat (\%) & −8,3 (−19,4 / +2,0) & −10,8 (−21,1 / −2,0) & −21,1 (−34,3 / −8,8) \\
Toplam ihracat (\$ mn) & −34 (−79 / +8) & −43 (−85 / −8) & −293 (−372 / −214) \\
Bunun içinde PEM kümülasyonu (\$ mn) & 0 (0 / 0) & 0 (0 / 0) & −199 (−254 / −143) \\
Nihai mal ithalatı (\$ mn) & −7 (−24 / +6) & −10 (−25 / +4) & −21 (−39 / −5) \\
Üretim (\$ mn) & −25 (−58 / +2) & −32 (−61 / −11) & −271 (−340 / −204) \\
İstihdam (kişi) & −154 (−350 / +10) & −193 (−375 / −69) & −1.632 (−2.141 / −1.188) \\
İstihdam (\%) & −0,4 (−1,0 / 0,0) & −0,5 (−1,0 / −0,2) & −4,5 (−5,9 / −3,3) \\
Cari denge (\$ mn) & −20 (−44 / +1) & −25 (−47 / −9) & −210 (−263 / −157) \\
TÜFE (bp) & +0,02 (+0,01 / +0,02) & +0,03 (+0,03 / +0,04) & +0,15 (+0,13 / +0,17) \\
\end{longtable}

\endgroup

\emph{Kaynak:} \etiketT{} f3e; toplamlar Faz 2B genel denge dağılımına kalibre. Yüzdeler 2024 düzeyine göre.

\hypertarget{nace_kartlar_13-15-54-1392-ev-tekstili}{%
\subsection{5.4 13.92 Ev tekstili}\label{nace_kartlar_13-15-54-1392-ev-tekstili}}

\emph{GB statüsü: GB · Arketip: A (AB'ye entegre ihracat şampiyonu) · Duruş: Geliştir · AB ile müzakere: savunma · Güven: Orta · Veri sayfası: Ek C.5}

Havlu, bornoz, nevresim ve perde Türkiye'nin AB'de güçlü olduğu ürünler. Sınıfta 87.344 kişi ve 9.199 girişim var; Denizli--Aydın--Muğla (\%18), Gaziantep, Bursa ve İstanbul öne çıkıyor. İhracat 2.346 milyon dolar ve \%54 kadarı AB'ye; Almanya tek başına AB'ye ihracatın \%33 kadarını alıyor. Türkiye AB pazarındaki payını 2020--2024'te neredeyse iki katına çıkardı (\%4,4 → \%8,2). Bu pazarın en güçlü rakipleri Çin ve GSP+ ile tarifesiz giren Pakistan (\%18).

S4'te AB tarifesi ortalama \%10,0; havlu ve nevresimde \%12. AB'ye ihracat 351 milyon dolar, istihdam 2.804 kişi azalıyor. Dağılımın alt ucunda sınıf istihdamının \%5,3 kadarı. Pakistan'la rekabet S4'te ağırlaşıyor: Pakistan havluda AB pazarının büyük tedarikçisi ve GSP+ statüsü sürdükçe sıfır tarife ödüyor. S2'de menşe kuralı çift dönüşüm (dokuma ya da örme ile birlikte konfeksiyon). Denizli ve Bursa'daki üreticilerin çoğu kumaşı kendisi dokuduğu için bu kural sorun değil. İthal ham bezle çalışan üreticiler için ise sorun \etiketH{}. S2'de kayıp 1.186 kişi.

Ev tekstili AB ticaret savunmasının geçmişte kullanıldığı bir alan: AB 1997'de Mısır, Hindistan ve Pakistan'ın pamuklu nevresimine damping vergisi uyguladı ve önlem DTÖ'de dava konusu oldu (DS141) \etiketV{}. Hasmane bir AB'de Türk ev tekstili olası bir hedef. Bu yüzden ilk parçadaki ticaret savunması hazırlığı ev tekstilini de kapsayacak biçimde genişletiliyor (E-13.92-01). Önerilen duruş ``Geliştir''. Kırmızı çizgi yetkinlikler entegre dokuma ve konfeksiyon, organik ve geri dönüştürülmüş pamukta sertifika ve büyük perakendecilerle uzun dönemli ilişkiler. Hedef ABD'deki payı (AB dışı ihracatın \%27 kadarı) büyütmek. Değerlendirmenin güveni orta.

\Needspace{7\baselineskip}

\textbf{Tablo 5.4 --- 13.92 Ev tekstili: senaryo etkileri (T+5'te yıllık; P50, parantez içinde P10 / P90)}

\nopagebreak[4]

\begingroup\gbtablo\small

\begin{longtable}[]{@{}
  >{\raggedright\arraybackslash}p{(\columnwidth - 6\tabcolsep) * \real{0.2555}}
  >{\raggedleft\arraybackslash}p{(\columnwidth - 6\tabcolsep) * \real{0.2393}}
  >{\raggedleft\arraybackslash}p{(\columnwidth - 6\tabcolsep) * \real{0.2428}}
  >{\raggedleft\arraybackslash}p{(\columnwidth - 6\tabcolsep) * \real{0.2624}}@{}}
\toprule\noalign{}
\begin{minipage}[b]{\linewidth}\raggedright
\textbf{Değişken}
\end{minipage} & \begin{minipage}[b]{\linewidth}\raggedleft
\textbf{S2}
\end{minipage} & \begin{minipage}[b]{\linewidth}\raggedleft
\textbf{S3}
\end{minipage} & \begin{minipage}[b]{\linewidth}\raggedleft
\textbf{S4}
\end{minipage} \\
\midrule\noalign{}
\endhead
\bottomrule\noalign{}
\endlastfoot
AB'ye ihracat (\$ mn) & −163 (−377 / +25) & −194 (−371 / −38) & −351 (−565 / −149) \\
AB'ye ihracat (\%) & −13,0 (−30,0 / +2,0) & −15,4 (−29,5 / −3,0) & −27,9 (−45,0 / −11,9) \\
Toplam ihracat (\$ mn) & −118 (−272 / +17) & −137 (−266 / −25) & −280 (−454 / −116) \\
Nihai mal ithalatı (\$ mn) & −8 (−27 / +6) & −11 (−28 / +4) & −23 (−41 / −5) \\
Üretim (\$ mn) & −109 (−249 / +11) & −125 (−239 / −29) & −257 (−415 / −112) \\
İstihdam (kişi) & −1.186 (−2.799 / +128) & −1.358 (−2.675 / −323) & −2.804 (−4.607 / −1.236) \\
İstihdam (\%) & −1,4 (−3,2 / +0,1) & −1,6 (−3,1 / −0,4) & −3,2 (−5,3 / −1,4) \\
Cari denge (\$ mn) & −84 (−192 / +8) & −96 (−185 / −23) & −199 (−321 / −87) \\
TÜFE (bp) & +0,05 (+0,04 / +0,06) & +0,10 (+0,08 / +0,11) & +0,46 (+0,40 / +0,53) \\
\end{longtable}

\endgroup

\emph{Kaynak:} \etiketT{} f3e; toplamlar Faz 2B genel denge dağılımına kalibre. Yüzdeler 2024 düzeyine göre.

\hypertarget{nace_kartlar_13-15-55-1393-halux131}{%
\subsection{5.5 13.93 Halı}\label{nace_kartlar_13-15-55-1393-halux131}}

\emph{GB statüsü: GB · Arketip: B (küresel rekabetçi, AB bağımlılığı düşük) · Duruş: Geliştir · AB ile müzakere: nötr · Güven: Yüksek · Veri sayfası: Ek C.6}

Halı, Türkiye'nin dünya ölçeğinde lider olduğu nadir ürünlerden biri. Türkiye 2024'te dünya dokuma halı ihracatının \%41,5 kadarını yaparak birinci. Sınıfın dünya ihracatındaki açıklanmış karşılaştırmalı üstünlüğü 15,8, imalatın en yükseklerinden. Üretim neredeyse tümüyle Gaziantep'te (\%86): büyük ve entegre tesisler polipropilen ve polyester ipliği kendileri çekiyor, dokuyor ve bitiriyor. Sınıfta 46.903 kişi çalışıyor.

Halı, AB'ye bağımlılığı düşük, küresel rekabetçi bir sınıf: ihracatın yalnızca \%16 kadarı AB'ye gidiyor. İlk pazar ABD (AB dışı ihracatın \%33 kadarı); onu Suudi Arabistan ve Irak izliyor. Bu yüzden çıkışın etkisi küçük. S4'te AB'ye ihracat 60 milyon dolar, istihdam 340 kişi azalıyor. Menşe kuralı (iplik çekimi ya da eğirme ile birlikte dokuma veya tufting) entegre üreticiler için sorun değil. Sınıfın asıl riskleri Gümrük Birliği'nin dışında: ABD'nin Türkiye'ye uyguladığı ek tarifeler \etiketD{} ve Körfez talebindeki dalgalanma.

Önerilen duruş ``Geliştir''. Strateji ABD ve Körfez'de pazar geliştirme (E-13-05) ve geri dönüştürülmüş PET iplikle üretimde AB'nin ekotasarım kurallarına hazırlık. Kırmızı çizgi yetkinlikler entegre iplik ve dokuma kapasitesi ile tasarım. AB ile müzakerede tutum nötr. Değerlendirmenin güveni yüksek.

\hypertarget{nace_kartlar_13-15-56-1395-dokusuz-kumaux15f}{%
\subsection{5.6 13.95 Dokusuz kumaş}\label{nace_kartlar_13-15-56-1395-dokusuz-kumaux15f}}

\emph{GB statüsü: GB · Arketip: A (AB'ye entegre ihracat şampiyonu) · Duruş: Koru · AB ile müzakere: savunma · Güven: Orta · Veri sayfası: Ek C.8}

Dokusuz kumaş, bu parçanın kabiliyet skoru en yüksek sınıfı (KS 61,7; 230 sınıf içinde 19. sırada) ve ürün karmaşıklığı pozitif olan iki tekstil sınıfından biri (PCI 0,47; diğeri teknik tekstil). Sınıfta 8.474 kişi çalışıyor; üretim Gaziantep (\%41) ve Tekirdağ'da sermaye yoğun spunbond ve spunlace hatlarında yapılıyor. Türkiye AB pazarında AB dışı tedarikin \%18,1 kadarını sağlıyor ve bu pay artıyor.

Çıkışın sınıfa etkisi küçük. S4'te AB tarifesi \%4,3, AB'ye ihracat 53 milyon dolar azalıyor; istihdam etkisi sınırlı (36 kişi), çünkü sınıf sermaye yoğun. Menşe kuralı polimerden ya da liften bağlama işlemini menşe kazandırıcı saydığı için ithal polipropilen serbest. Sınıfın stratejik önemi tıbbi tekstilden geliyor (SS 37): maske, önlük, hijyen ürünleri ve filtrelerin girdisi. 2020'deki salgın, bu kapasitenin bir kriz anında ne kadar değerli olduğunu gösterdi.

Önerilen duruş ``Koru''. Tek eylem, tıbbi dokusuz kumaş kapasitesinin ve stoklarının bir envanterini çıkarmak ve bunu salgın ve afet planına bağlamak (E-13.95-01). Bu eylem idari ve çıkıştan bağımsız. Değerlendirmenin güveni orta: AB'ye göre verimlilik verisi eksik.

\hypertarget{nace_kartlar_13-15-57-1396-teknik-tekstil}{%
\subsection{5.7 13.96 Teknik tekstil}\label{nace_kartlar_13-15-57-1396-teknik-tekstil}}

\emph{GB statüsü: GB · Arketip: A (AB'ye entegre ihracat şampiyonu) · Duruş: Geliştir · AB ile müzakere: savunma · Güven: Yüksek · Veri sayfası: Ek C.9}

Teknik tekstil (lastik kordu, kaplamalı kumaş, endüstriyel filtre ve konveyör bandı tekstili, balistik ve yangına dayanıklı kumaş) bu parçanın gelecek vaat eden kesimi. Sınıfta 19.387 kişi çalışıyor ve üretimin \%40 kadarı İstanbul'da. Lastik kordunda Kordsa dünya ölçeğinde bir üretici. Sınıfın ürün karmaşıklığı tekstil ortalamasının üstünde (PCI 0,31) ve stratejik önemi savunma, sağlık ve çift kullanımdan geliyor (SS 33). Türkiye bu sınıfta net ithalatçı: ithalat 589 milyon dolar ve \%44 kadarı AB'den.

Çıkışın sınıfa net etkisi küçük. S4'te AB'ye ihracat 48 milyon dolar azalıyor. Ama AB'den gelen teknik tekstilin pahalanması iç pazarda yerli üreticiye pay kazandırıyor ve istihdam etkisi −50 kişide kalıyor. PEM kümülasyonu kanalı bu sınıfta da var ama küçük: kaplamalı kumaşın bir kısmı ortak ülkelerde giyim ve çantaya giriyor. Önerilen duruş ``Geliştir''. Ana eylem, Savunma Sanayii Başkanlığı'nın yerlileştirme programında balistik, aramid ve yangına dayanıklı tekstile yer verilmesi ve kamu alımında menşe şartı olmadan performans şartnamesi kullanılması (E-13.96-01). Tarife tarafında aramid ve cam elyafı gibi yurt içinde üretilmeyen girdilerde sıfır tarife bir aday; bu, sonraki parçalarda kimya ve cam sınıflarıyla birlikte değerlendirilecek. Değerlendirmenin güveni yüksek.

\hypertarget{nace_kartlar_13-15-58-1413-dux131ux15f-giyim}{%
\subsection{5.8 14.13 Dış giyim}\label{nace_kartlar_13-15-58-1413-dux131ux15f-giyim}}

\emph{GB statüsü: GB · Arketip: A (AB'ye entegre ihracat şampiyonu) / F (emek yoğun, rekabet gücü zayıflayan) · Duruş: Dönüştür · AB ile müzakere: savunma · Güven: Orta · Veri sayfası: Ek C.13}

Dış giyim (pantolon, ceket, elbise, gömlek, etek ve takım) imalatın en büyük işveren sınıfı. Sınıfta 365.708 kişi ve 41.798 girişim var; istihdamın \%31 kadarı İstanbul'da, \%14 kadarı Mardin--Batman--Şırnak--Siirt'te. Şanlıurfa--Diyarbakır, Samsun--Tokat--Çorum--Amasya ve Van--Muş--Bitlis--Hakkari, son on yılda bölgesel teşviklerle kurulan konfeksiyon tesisleriyle öne çıkan diğer bölgeler. İhracat 2024'te 8.577 milyon dolar ve \%59 kadarı AB'ye. Kadın pamuklu pantolonda Türkiye dünya ihracatının \%7,5 kadarını yaparak üçüncü; AB'ye en büyük iki kalem kadın ve erkek denim pantolonu.

Sınıfın kabiliyet profili ikili. Bir yanda geniş bir firma tabanı, hızlı ve esnek bir tedarik zinciri ve dünya ölçeğinde markalar (LC Waikiki, Koton, DeFacto, Mavi) var. Öbür yanda düşük karmaşıklıklı ürünler (PCI −1,73) ve düşük verimlilik var: kişi başı katma değer AB'deki muadilinin 0,31 katı. Kabiliyet skoru 46,5 ve sınıf 230 sınıf arasında 154. sırada. Türk dış giyiminin AB'deki yeri fiyat değil hız. Ürünün tasarımdan rafa iki ila dört haftada ulaşması Avrupalı perakendecinin sezon içi riskini azaltıyor \etiketH{}.

Çıkışın etkisi bu sınıfta en büyük. S4'te AB tarifesi bütün satırlarda \%12 ve AB'nin AB dışından yaptığı dış giyim ithalatının \%58 kadarı tercihli tedarikçilerden geliyor. Bangladeş tek başına bu ithalatın \%20 kadarını sağlıyor ve sıfır tarife ödüyor. AB'ye ihracat 999 milyon dolar azalıyor. İstihdam kaybı 11.910 kişi, dağılımın alt ucunda 20.028 kişi. Mevcut model sonuçlarında bu, imalatın bütününde en büyük sınıf kaybı. S2'de kayıp 5.143 kişiye iniyor ve kalanın çoğu menşe kaynaklı. Asya kumaşıyla üretilen ihracat tercih dışı kalıyor ve uyumlu ihracat da menşe ispatının maliyetine katlanıyor. Bu kayıp, Türkiye'de baskı kapasitesi ve PEM kumaşına geçişle (E-13-02) ve müzakerede menşe kotasıyla (E-13-06) azaltılabilir.

İthalat tarafında sınıfın önemi tüketici fiyatlarında. Dış giyimin TÜFE ağırlığı (\%3,86) bu parçanın en büyüğü. S4'te AB'den gelen giysinin tarifesi \%37,2 oluyor ve TÜFE'ye etkisi 4,3 baz puan. Vergi takası (E-13-04) bu etkiyi siliyor. Aynı tarife iç pazarda yerli üreticiye pay kazandırıyor ve S4 kaybının bir kısmını dengeliyor.

Önerilen duruş ``Dönüştür''. Sınıfın geleceği ucuz emeğe dayalı fason dikimde değil. Asıl yön tasarım, ürün geliştirme, küçük partili hızlı tedarik ve marka. Dönüşüm programı (E-14.13-01) yatırım teşvik sisteminde katma değer ölçütlü bir destek, kesim ve dikimde otomasyon ve dijital ürün pasaportu hazırlığını içeriyor; menşe ya da ihracat şartı taşımıyor. Kırmızı çizgi yetkinlikler, kumaştan dikime entegre tedarik zinciri, hızlı numune ve küçük parti üretim kapasitesi ile perakendecilerle kurulmuş tasarım ortaklıkları. Hedefler, AB27'ye göre verimliliği T+5'te 0,40'a çıkarmak ve AB dışı pazarların payını beş puan artırmak. S4'te kısa çalışma (E-13-03) ve vergi takası (E-13-04) devreye giriyor. Kazananlar marka sahibi, tasarım ve hız yetkinliği olan büyük üreticiler; kaybedenler tek bir AB alıcısına fason çalışan küçük atölyeler. Değerlendirmenin güveni orta. S2'deki menşe uyum oranı \etiketH{} ve firmaların Mısır'a taşınma eğilimi sonucu belirleyen iki belirsizlik.

\Needspace{7\baselineskip}

\textbf{Tablo 5.5 --- 14.13 Dış giyim: senaryo etkileri (T+5'te yıllık; P50, parantez içinde P10 / P90)}

\nopagebreak[4]

\begingroup\gbtablo\small

\begin{longtable}[]{@{}
  >{\raggedright\arraybackslash}p{(\columnwidth - 6\tabcolsep) * \real{0.2393}}
  >{\raggedleft\arraybackslash}p{(\columnwidth - 6\tabcolsep) * \real{0.2362}}
  >{\raggedleft\arraybackslash}p{(\columnwidth - 6\tabcolsep) * \real{0.2566}}
  >{\raggedleft\arraybackslash}p{(\columnwidth - 6\tabcolsep) * \real{0.2679}}@{}}
\toprule\noalign{}
\begin{minipage}[b]{\linewidth}\raggedright
\textbf{Değişken}
\end{minipage} & \begin{minipage}[b]{\linewidth}\raggedleft
\textbf{S2}
\end{minipage} & \begin{minipage}[b]{\linewidth}\raggedleft
\textbf{S3}
\end{minipage} & \begin{minipage}[b]{\linewidth}\raggedleft
\textbf{S4}
\end{minipage} \\
\midrule\noalign{}
\endhead
\bottomrule\noalign{}
\endlastfoot
AB'ye ihracat (\$ mn) & −453 (−1.064 / +101) & −528 (−1.030 / −100) & −999 (−1.613 / −412) \\
AB'ye ihracat (\%) & −9,0 (−21,1 / +2,0) & −10,5 (−20,5 / −2,0) & −19,8 (−32,1 / −8,2) \\
Toplam ihracat (\$ mn) & −325 (−770 / +68) & −374 (−733 / −71) & −785 (−1.306 / −327) \\
Nihai mal ithalatı (\$ mn) & −27 (−87 / +31) & −38 (−96 / +18) & −97 (−174 / −19) \\
Üretim (\$ mn) & −294 (−688 / +37) & −333 (−640 / −87) & −689 (−1.133 / −307) \\
İstihdam (kişi) & −5.143 (−11.876 / +688) & −5.744 (−11.181 / −1.508) & −11.910 (−20.028 / −5.328) \\
İstihdam (\%) & −1,4 (−3,2 / +0,2) & −1,6 (−3,1 / −0,4) & −3,3 (−5,5 / −1,5) \\
Cari denge (\$ mn) & −227 (−532 / +29) & −258 (−495 / −68) & −533 (−877 / −237) \\
TÜFE (bp) & +0,45 (+0,39 / +0,52) & +0,86 (+0,73 / +0,98) & +4,33 (+3,72 / +4,91) \\
\end{longtable}

\endgroup

\emph{Kaynak:} \etiketT{} f3e; toplamlar Faz 2B genel denge dağılımına kalibre. Yüzdeler 2024 düzeyine göre.

\hypertarget{nace_kartlar_13-15-59-1414-iux307uxe7-giyim-ve-tiux15fuxf6rt}{%
\subsection{5.9 14.14 İç giyim ve tişört}\label{nace_kartlar_13-15-59-1414-iux307uxe7-giyim-ve-tiux15fuxf6rt}}

\emph{GB statüsü: GB · Arketip: F (emek yoğun, rekabet gücü zayıflayan) / A (AB'ye entegre ihracat şampiyonu) · Duruş: Yönetilen uyum · AB ile müzakere: savunma · Güven: Orta · Veri sayfası: Ek C.14}

İç giyim, tişört ve gömlek sınıfında 183.980 kişi çalışıyor; istihdamın \%41 kadarı İstanbul'da, \%14 kadarı Mardin--Batman--Şırnak--Siirt'te. İhracat 5.647 milyon dolar ve \%60 kadarı AB'ye. En büyük kalem pamuklu tişört (KN 6109 10 00; AB'ye €838 milyon). Sınıfın verimliliği dış giyimden belirgin biçimde yüksek (AB27'ye göre 0,78) ve büyük işletmelerin istihdam payı \%33. Bu, örme kumaştan dikime uzanan entegre üretimin işareti.

Rekabet bu sınıfta en sert. Bangladeş, AB'nin AB dışından yaptığı ithalatın \%32 kadarını EBA rejimiyle sıfır tarifeyle sağlıyor; Türkiye Bangladeş ve Çin'den sonra üçüncü. Standart tişört ve iç çamaşırında fiyat hızdan daha belirleyici. Bu yüzden S4'te kayıp büyük: AB'ye ihracat 1.003 milyon dolar, istihdam 9.401 kişi azalıyor. Dağılımın alt ucunda sınıf istihdamının \%8,5 kadarı kayboluyor. S2'de menşe kuralı tişört için kolay, çünkü örme kumaş büyük ölçüde yerli. Dantel ve elastik kumaşlı iç çamaşırında ise Asya kumaşı riski var. S2'de kayıp 3.732 kişi.

Önerilen duruş ``Yönetilen uyum''. Kural tabanlı sınıflama bu sınıfı ``emek yoğun, rekabet gücü zayıflayan'' (F) olarak işaretliyor ve bu değerlendirme korundu. Standart ürünlerin bir kısmının Asya'ya kaybı S2'de de sürecek bir eğilim. Politika bu kaybı durdurmaya değil, yönetmeye yönelik. İŞKUR'un bölgesel işgücü uyum programı (E-14.14-01) Mardin--Batman, Şanlıurfa--Diyarbakır ve İstanbul'da mesleki eğitim ve yeniden yerleştirme sunuyor. Entegre ve yüksek katma değerli üreticiler ise dönüşüm programına (E-14.13-01) alınabiliyor. Kırmızı çizgi yetkinlik, örme ve konfeksiyonun aynı tesiste yapılması; bu, hem S2'de menşe kuralına hem hızlı teslimata uyumun temeli. Birleşik Krallık bu sınıfın AB dışındaki ilk pazarı (AB dışı ihracatın \%19 kadarı) ve STA yükseltmesinin önceliklerinden. Değerlendirmenin güveni orta.

\Needspace{7\baselineskip}

\textbf{Tablo 5.6 --- 14.14 İç giyim ve tişört: senaryo etkileri (T+5'te yıllık; P50, parantez içinde P10 / P90)}

\nopagebreak[4]

\begingroup\gbtablo\small

\begin{longtable}[]{@{}
  >{\raggedright\arraybackslash}p{(\columnwidth - 6\tabcolsep) * \real{0.2523}}
  >{\raggedleft\arraybackslash}p{(\columnwidth - 6\tabcolsep) * \real{0.2365}}
  >{\raggedleft\arraybackslash}p{(\columnwidth - 6\tabcolsep) * \real{0.2398}}
  >{\raggedleft\arraybackslash}p{(\columnwidth - 6\tabcolsep) * \real{0.2715}}@{}}
\toprule\noalign{}
\begin{minipage}[b]{\linewidth}\raggedright
\textbf{Değişken}
\end{minipage} & \begin{minipage}[b]{\linewidth}\raggedleft
\textbf{S2}
\end{minipage} & \begin{minipage}[b]{\linewidth}\raggedleft
\textbf{S3}
\end{minipage} & \begin{minipage}[b]{\linewidth}\raggedleft
\textbf{S4}
\end{minipage} \\
\midrule\noalign{}
\endhead
\bottomrule\noalign{}
\endlastfoot
AB'ye ihracat (\$ mn) & −445 (−1.018 / +68) & −526 (−987 / −101) & −1.003 (−1.588 / −419) \\
AB'ye ihracat (\%) & −13,2 (−30,1 / +2,0) & −15,5 (−29,2 / −3,0) & −29,7 (−47,0 / −12,4) \\
Toplam ihracat (\$ mn) & −316 (−743 / +46) & −368 (−711 / −71) & −802 (−1.283 / −335) \\
Nihai mal ithalatı (\$ mn) & −19 (−63 / +11) & −26 (−65 / +7) & −46 (−84 / −10) \\
Üretim (\$ mn) & −297 (−681 / +34) & −341 (−650 / −78) & −753 (−1.200 / −323) \\
İstihdam (kişi) & −3.732 (−8.652 / +465) & −4.314 (−8.257 / −983) & −9.401 (−15.548 / −4.121) \\
İstihdam (\%) & −2,0 (−4,7 / +0,3) & −2,3 (−4,5 / −0,5) & −5,1 (−8,5 / −2,2) \\
Cari denge (\$ mn) & −230 (−527 / +27) & −264 (−503 / −60) & −582 (−928 / −250) \\
TÜFE (bp) & +0,08 (+0,07 / +0,09) & +0,15 (+0,13 / +0,18) & +0,84 (+0,72 / +0,95) \\
\end{longtable}

\endgroup

\emph{Kaynak:} \etiketT{} f3e; toplamlar Faz 2B genel denge dağılımına kalibre. Yüzdeler 2024 düzeyine göre.

\hypertarget{nace_kartlar_13-15-510-1419-diux11fer-giyim-ve-aksesuar}{%
\subsection{5.10 14.19 Diğer giyim ve aksesuar}\label{nace_kartlar_13-15-510-1419-diux11fer-giyim-ve-aksesuar}}

\emph{GB statüsü: GB · Arketip: F (emek yoğun, rekabet gücü zayıflayan) / A (AB'ye entegre ihracat şampiyonu) · Duruş: Yönetilen uyum · AB ile müzakere: savunma · Güven: Orta · Veri sayfası: Ek C.15}

Bu sınıf bebek giyimi, eşofman ve spor giyimi, eldiven, atkı, şal, şapka ve diğer giyim aksesuarlarını kapsıyor. Sınıfta 41.506 kişi çalışıyor. Çıkıştan önce de geriliyor: ihracat 2020'de 2.324 milyon dolardan 2024'te 1.461 milyon dolara düştü. AB pazarındaki payı \%6,6 iken \%3,9 oldu. İç pazarda da ithalat güçlü: ithalatın iç talepteki payı \%61 ve ilk kaynak Çin.

S4'te AB'ye ihracat 107 milyon dolar, istihdam 1.289 kişi azalıyor. TÜFE'ye etkisi 1,4 baz puan ve vergi takasının kapsamında. Önerilen duruş ``Yönetilen uyum''; sınıf işgücü uyum programına (E-14.14-01) dahil. Bebek giyimi bu sınıfın korunması gereken kesimi: organik pamuk ve sertifikalı üretim AB'de pazar buluyor. Değerlendirmenin güveni orta.

\hypertarget{nace_kartlar_13-15-511-1431-uxe7orap}{%
\subsection{5.11 14.31 Çorap}\label{nace_kartlar_13-15-511-1431-uxe7orap}}

\emph{GB statüsü: GB · Arketip: A (AB'ye entegre ihracat şampiyonu) / F (emek yoğun, rekabet gücü zayıflayan) · Duruş: Geliştir · AB ile müzakere: savunma · Güven: Yüksek · Veri sayfası: Ek C.17}

Çorap, giyimin sermaye yoğun ve otomasyonu en yüksek sınıfı. Sınıfta 25.205 kişi çalışıyor; üretim İstanbul (\%39) ile Tekirdağ--Edirne--Kırklareli (\%32) arasında, Çorlu ve Çerkezköy çevresinde yoğunlaşıyor. Türkiye dünya çorap ihracatında \%9,7 payla ikinci ve AB pazarında AB dışı tedarikin \%25,8 kadarını sağlıyor. İhracat 1.296 milyon dolar ve \%65 kadarı AB'ye; bu, parçanın en yüksek AB payı. Birleşik Krallık AB dışı ihracatın \%42 kadarını alıyor.

AB'ye bağımlılık yüksek olduğu için S4'ün etkisi büyük. AB'ye ihracat 271 milyon dolar, istihdam 2.277 kişi azalıyor; dağılımın alt ucunda sınıf istihdamının \%14,9 kadarı. S2'de menşe kuralı kolay: çorap biçimine göre örülüyor ve ithal naylon ve elastan iplik serbest. Bu yüzden S2'de kayıp (750 kişi) S4'ünkinin küçük bir kısmı.

Önerilen duruş ``Geliştir''. Çorapta sorun kabiliyet değil tarife; strateji AB pazarını S2'de korumak ve Birleşik Krallık ile ABD'de büyümek. STA yükseltmesinde AB--Türkiye--Birleşik Krallık kümülasyonu bu sınıf için öncelikli (E-13-05). Kırmızı çizgi yetkinlikler yüksek hızlı örme makinesi parkı, tasarım ve büyük Avrupalı markalarla uzun dönemli tedarik ilişkileri. S4'te kısa çalışma (E-13-03) devreye girebilir. Değerlendirmenin güveni yüksek.

\Needspace{7\baselineskip}

\textbf{Tablo 5.7 --- 14.31 Çorap: senaryo etkileri (T+5'te yıllık; P50, parantez içinde P10 / P90)}

\nopagebreak[4]

\begingroup\gbtablo\small

\begin{longtable}[]{@{}
  >{\raggedright\arraybackslash}p{(\columnwidth - 6\tabcolsep) * \real{0.2727}}
  >{\raggedleft\arraybackslash}p{(\columnwidth - 6\tabcolsep) * \real{0.2325}}
  >{\raggedleft\arraybackslash}p{(\columnwidth - 6\tabcolsep) * \real{0.2364}}
  >{\raggedleft\arraybackslash}p{(\columnwidth - 6\tabcolsep) * \real{0.2584}}@{}}
\toprule\noalign{}
\begin{minipage}[b]{\linewidth}\raggedright
\textbf{Değişken}
\end{minipage} & \begin{minipage}[b]{\linewidth}\raggedleft
\textbf{S2}
\end{minipage} & \begin{minipage}[b]{\linewidth}\raggedleft
\textbf{S3}
\end{minipage} & \begin{minipage}[b]{\linewidth}\raggedleft
\textbf{S4}
\end{minipage} \\
\midrule\noalign{}
\endhead
\bottomrule\noalign{}
\endlastfoot
AB'ye ihracat (\$ mn) & −101 (−237 / +17) & −128 (−248 / −23) & −271 (−434 / −113) \\
AB'ye ihracat (\%) & −11,9 (−28,0 / +2,0) & −15,1 (−29,3 / −2,8) & −32,0 (−51,4 / −13,4) \\
Toplam ihracat (\$ mn) & −73 (−171 / +11) & −92 (−176 / −16) & −216 (−350 / −89) \\
Nihai mal ithalatı (\$ mn) & −1 (−4 / +1) & −2 (−4 / 0) & −3 (−5 / −1) \\
Üretim (\$ mn) & −71 (−168 / +11) & −90 (−172 / −17) & −213 (−345 / −88) \\
İstihdam (kişi) & −750 (−1.808 / +123) & −942 (−1.877 / −185) & −2.277 (−3.752 / −963) \\
İstihdam (\%) & −3,0 (−7,2 / +0,5) & −3,7 (−7,4 / −0,7) & −9,0 (−14,9 / −3,8) \\
Cari denge (\$ mn) & −55 (−130 / +8) & −69 (−133 / −13) & −164 (−267 / −68) \\
TÜFE (bp) & 0,00 (0,00 / 0,00) & 0,00 (0,00 / +0,01) & +0,02 (+0,02 / +0,03) \\
\end{longtable}

\endgroup

\emph{Kaynak:} \etiketT{} f3e; toplamlar Faz 2B genel denge dağılımına kalibre. Yüzdeler 2024 düzeyine göre.

\hypertarget{nace_kartlar_13-15-512-1439-triko}{%
\subsection{5.12 14.39 Triko}\label{nace_kartlar_13-15-512-1439-triko}}

\emph{GB statüsü: GB · Arketip: F (emek yoğun, rekabet gücü zayıflayan) / A (AB'ye entegre ihracat şampiyonu) · Duruş: Yönetilen uyum · AB ile müzakere: savunma · Güven: Orta · Veri sayfası: Ek C.18}

Triko (kazak, hırka ve süveter) sınıfında 32.975 kişi çalışıyor ve üretimin \%59 kadarı İstanbul'da. İhracat 2.719 milyon dolar ve AB payı \%66; çorapla birlikte parçanın AB'ye en bağımlı sınıfı. AB pazarında Türkiye Çin ve Bangladeş'ten sonra üçüncü. Bangladeş (\%27) ve Kamboçya EBA rejimiyle tarifesiz giriyor.

Trikoda S4 kaybı göreli olarak bu parçanın en büyüğü. AB'ye ihracat 425 milyon dolar, istihdam 3.773 kişi azalıyor. Dağılımın alt ucunda sınıf istihdamının \%19,2 kadarı kayboluyor; bu, imalatın bütününde göreli kaybı en büyük sınıflardan biri. İç pazarda da ithalat güçlü (ithalatın iç talepteki payı \%55). S2'de menşe kuralı kolay: tam kalıp örme ya da örme ile birlikte konfeksiyon menşe kazandırıyor ve ithal iplik serbest. S2'de kayıp 1.098 kişi.

Önerilen duruş ``Yönetilen uyum''. Trikoda standart ürün ile moda ürünü arasındaki ayrım belirleyici. Tam kalıp örme makinesiyle çalışan, tasarım ve hızlı teslimat sunan üreticiler S2'de rekabetçi kalıyor. Fason örme ve dikimle çalışan küçük atölyeler ise S4'te ağır kayıp yaşıyor. Önerilen paket, ilk grubu dönüşüm programına, ikinci grubun işgücünü işgücü uyum programına (E-14.14-01) yönlendiriyor. Değerlendirmenin güveni orta.

\Needspace{7\baselineskip}

\textbf{Tablo 5.8 --- 14.39 Triko (kazak, hırka): senaryo etkileri (T+5'te yıllık; P50, parantez içinde P10 / P90)}

\nopagebreak[4]

\begingroup\gbtablo\small

\begin{longtable}[]{@{}
  >{\raggedright\arraybackslash}p{(\columnwidth - 6\tabcolsep) * \real{0.2555}}
  >{\raggedleft\arraybackslash}p{(\columnwidth - 6\tabcolsep) * \real{0.2393}}
  >{\raggedleft\arraybackslash}p{(\columnwidth - 6\tabcolsep) * \real{0.2428}}
  >{\raggedleft\arraybackslash}p{(\columnwidth - 6\tabcolsep) * \real{0.2624}}@{}}
\toprule\noalign{}
\begin{minipage}[b]{\linewidth}\raggedright
\textbf{Değişken}
\end{minipage} & \begin{minipage}[b]{\linewidth}\raggedleft
\textbf{S2}
\end{minipage} & \begin{minipage}[b]{\linewidth}\raggedleft
\textbf{S3}
\end{minipage} & \begin{minipage}[b]{\linewidth}\raggedleft
\textbf{S4}
\end{minipage} \\
\midrule\noalign{}
\endhead
\bottomrule\noalign{}
\endlastfoot
AB'ye ihracat (\$ mn) & −138 (−322 / +36) & −179 (−345 / −33) & −425 (−684 / −175) \\
AB'ye ihracat (\%) & −7,7 (−18,1 / +2,0) & −10,1 (−19,4 / −1,8) & −23,8 (−38,3 / −9,8) \\
Toplam ihracat (\$ mn) & −98 (−236 / +24) & −127 (−250 / −23) & −337 (−556 / −141) \\
Nihai mal ithalatı (\$ mn) & −6 (−21 / +5) & −9 (−22 / +3) & −18 (−32 / −4) \\
Üretim (\$ mn) & −91 (−217 / +19) & −118 (−230 / −26) & −319 (−525 / −136) \\
İstihdam (kişi) & −1.098 (−2.653 / +249) & −1.412 (−2.793 / −317) & −3.773 (−6.347 / −1.609) \\
İstihdam (\%) & −3,3 (−8,0 / +0,8) & −4,3 (−8,5 / −1,0) & −11,4 (−19,2 / −4,9) \\
Cari denge (\$ mn) & −71 (−168 / +15) & −91 (−178 / −20) & −246 (−406 / −106) \\
TÜFE (bp) & +0,15 (+0,13 / +0,17) & +0,28 (+0,24 / +0,31) & +1,49 (+1,28 / +1,69) \\
\end{longtable}

\endgroup

\emph{Kaynak:} \etiketT{} f3e; toplamlar Faz 2B genel denge dağılımına kalibre. Yüzdeler 2024 düzeyine göre.

\hypertarget{nace_kartlar_13-15-513-1511-deri-tabaklama}{%
\subsection{5.13 15.11 Deri tabaklama}\label{nace_kartlar_13-15-513-1511-deri-tabaklama}}

\emph{GB statüsü: GB · Arketip: A (AB'ye entegre ihracat şampiyonu) / G (SKDM ve Yeşil Mutabakat maruz) · Duruş: Dönüştür · AB ile müzakere: nötr · Güven: Orta · Veri sayfası: Ek C.19}

Tabaklama sınıfında 8.899 kişi çalışıyor ve üretim organize deri sanayi bölgelerinde: Kocaeli--Sakarya bölgesi (\%28), İstanbul (Tuzla), Tekirdağ (Çorlu) ve İzmir (Menemen). İhracat 268 milyon dolar; AB'deki ilk pazar İtalya (\%26), yani AB'nin lüks deri eşya ve ayakkabı sanayisi. AB dışında Hindistan, Güney Kore ve Çin öne çıkıyor. İthalatın da ilk kaynağı İtalya (\%36).

Çıkışın sınıfa doğrudan etkisi küçük. S4'te AB tarifesi \%3,7, AB'ye ihracat 16 milyon dolar azalıyor ve istihdam etkisi yalnızca 13 kişi. Menşe kuralı ithal ham deriyi ve yarı işlenmiş deriyi serbest bırakıyor. Sınıf için Eylül 2026'nın asıl haberi başka: 2026/2102 sayılı Yetki Devrine Dayanan Tüzük sığır derisini ormansızlaşma tüzüğünün (EUDR) kapsamından çıkardı \etiketV{}. Bu değişiklikten önce sınıfın AB'ye ihracatının yaklaşık beşte üçü coğrafi konum verisiyle özen yükümlülüğüne girecekti.

Önerilen duruş ``Dönüştür''. Kural tabanlı sınıflama, EUDR'den çıkışla birlikte sınıfı ``AB'ye entegre ihracatçı'' (A) olarak işaretliyor; ikincil arketip olarak ``Yeşil Mutabakat'a maruz'' (G) eklendi. Gerekçe, AB'li alıcıların çevre ve tedarik zinciri özen yükümlülükleri: tabakhanenin arıtma performansı, kromun yönetimi ve sertifika. Bunlar AB pazarına erişimin fiili ölçütleri \etiketH{}. Önerilen eylem organize deri sanayi bölgelerinde ortak arıtmanın yenilenmesi, krom azaltımı ve sertifikasyon (E-15.11-01). Değerlendirmenin güveni orta.

\hypertarget{nace_kartlar_13-15-514-1512-bavul-ve-uxe7anta}{%
\subsection{5.14 15.12 Bavul ve çanta}\label{nace_kartlar_13-15-514-1512-bavul-ve-uxe7anta}}

\emph{GB statüsü: GB · Arketip: C (ithalata bağımlı, yerlileştirme potansiyeli) · Duruş: Geliştir · AB ile müzakere: savunma · Güven: Yüksek · Veri sayfası: Ek C.20}

Çanta ve bavul sınıfında 16.310 kişi çalışıyor; üretimin \%67 kadarı İstanbul'da. İhracat 2020--2024'te \%82 artarak 377 milyon dolara çıktı, ama ithalat (668 milyon dolar) hâlâ ihracatın üstünde. İthalatta Çin'in yanında İtalya ve Fransa (lüks markalar) var. Sınıfın kabiliyet skoru düşük (37,6) ve kural tabanlı sınıflama onu ``ithalata bağımlı, yerlileştirme potansiyeli olan'' (C) sınıf olarak işaretliyor.

Bu parçada S4'te istihdamı artan tek sınıf bu: AB'den gelen çantanın pahalanması iç pazarda yerli üreticiye pay kazandırıyor ve istihdam 214 kişi artıyor. Bedeli tüketici ödüyor: TÜFE'ye etkisi 1,8 baz puan. Çanta vergi takasının kapsamında değil, çünkü takas Faz 2A'da giyim ve ayakkabı grupları için tasarlandı ve çantanın KDV oranı bu gruplardan farklı olabilir \etiketD{} (VB-39). Önerilen duruş ``Geliştir'': tasarım ve markayla iç pazardaki payı kalıcı kılmak. Ürün uzayında yakın ve daha karmaşık ürünler (diğer deri eşya, KN 4205) bir büyüme yolu sunuyor. Yeni bir tarife koruması önerilmiyor. Değerlendirmenin güveni yüksek.

\hypertarget{nace_kartlar_13-15-515-1520-ayakkabux131}{%
\subsection{5.15 15.20 Ayakkabı}\label{nace_kartlar_13-15-515-1520-ayakkabux131}}

\emph{GB statüsü: GB · Arketip: --- (belirgin arketip yok) · Duruş: Koru · AB ile müzakere: nötr · Güven: Yüksek · Veri sayfası: Ek C.21}

Ayakkabı sınıfında 50.596 kişi ve 7.162 girişim var; istihdamın \%42 kadarı İstanbul'da, kalanı Gaziantep, İzmir, Şanlıurfa--Diyarbakır ve Konya'da. Türkiye ayakkabıda net ithalatçı. İhracat 1.176 milyon dolar, ithalat 1.720 milyon dolar; ithalatın ilk kaynakları Çin (\%32) ve Vietnam (\%27). AB'ye ihracat 392 milyon dolar ve AB pazarında Türkiye küçük bir tedarikçi (\%1,2). İlk AB dışı pazarlar Suudi Arabistan ve Irak.

S4'te AB tarifesi ortalama \%9,9 ve bazı satırlarda \%17'ye yaklaşıyor. AB'ye ihracat 128 milyon dolar, istihdam 412 kişi azalıyor. Sınıfın asıl etkisi ithalat tarafında. AB'den gelen ayakkabının tarifesi S4'te \%48,7 oluyor ve TÜFE'ye etkisi 3,9 baz puan; bu parçanın ikinci büyük enflasyon kalemi. Vergi takası ayakkabıda KDV'yi \%10'dan \%5,6'ya indirerek bu etkiyi siliyor (E-13-04). Menşe kuralı kolay: hazır saya ve taban birleşimleri dışında her girdi serbest ve 64. fasılda geri ödeme yasağı yok.

Önerilen duruş ``Koru''. Ayakkabıda çıkışın net etkisi küçük; strateji iç pazardaki yerli payını vergi takasıyla birlikte korumak ve Körfez ile Irak'taki konumu güçlendirmek. Askerî bot, sınıfa küçük bir stratejik önem veriyor (SS 8). Ayakkabı TÜFE alt endeksi için bir tetik tanımlandı. Değerlendirmenin güveni yüksek.

\Needspace{7\baselineskip}

\textbf{Tablo 5.9 --- 15.20 Ayakkabı: senaryo etkileri (T+5'te yıllık; P50, parantez içinde P10 / P90)}

\nopagebreak[4]

\begingroup\gbtablo\small

\begin{longtable}[]{@{}lrrr@{}}
\toprule\noalign{}
\textbf{Değişken} & \textbf{S2} & \textbf{S3} & \textbf{S4} \\
\midrule\noalign{}
\endhead
\bottomrule\noalign{}
\endlastfoot
AB'ye ihracat (\$ mn) & −53 (−120 / +8) & −66 (−125 / −13) & −128 (−204 / −53) \\
AB'ye ihracat (\%) & −13,4 (−30,7 / +2,0) & −16,9 (−31,9 / −3,2) & −32,8 (−51,9 / −13,4) \\
Toplam ihracat (\$ mn) & −38 (−87 / +5) & −47 (−90 / −9) & −102 (−164 / −42) \\
Nihai mal ithalatı (\$ mn) & −33 (−108 / +15) & −42 (−107 / +9) & −63 (−117 / −13) \\
Üretim (\$ mn) & −7 (−18 / +27) & −8 (−22 / +23) & −36 (−60 / −21) \\
İstihdam (kişi) & −77 (−215 / +301) & −94 (−254 / +256) & −412 (−700 / −234) \\
İstihdam (\%) & −0,2 (−0,4 / +0,6) & −0,2 (−0,5 / +0,5) & −0,8 (−1,4 / −0,5) \\
Cari denge (\$ mn) & −5 (−14 / +21) & −7 (−17 / +18) & −28 (−46 / −16) \\
TÜFE (bp) & +0,39 (+0,33 / +0,44) & +0,73 (+0,63 / +0,83) & +3,93 (+3,37 / +4,45) \\
\end{longtable}

\endgroup

\emph{Kaynak:} \etiketT{} f3e; toplamlar Faz 2B genel denge dağılımına kalibre. Yüzdeler 2024 düzeyine göre.

\hypertarget{nace_kartlar_13-15-516-etkisi-sux131nux131rlux131-sux131nux131flar}{%
\subsection{5.16 Etkisi sınırlı sınıflar}\label{nace_kartlar_13-15-516-etkisi-sux131nux131rlux131-sux131nux131flar}}

Altı sınıfta çıkışın etkisi küçük ya da ölçülemiyor. Hepsinin veri sayfası ve senaryo tablosu Ek C'de.

\Needspace{7\baselineskip}

\textbf{Tablo 5.10 --- Etkisi sınırlı sınıflar: temel göstergeler}

\nopagebreak[4]

\begingroup\gbtablo\footnotesize

\begin{longtable}[]{@{}
  >{\raggedright\arraybackslash}p{(\columnwidth - 16\tabcolsep) * \real{0.1598}}
  >{\raggedleft\arraybackslash}p{(\columnwidth - 16\tabcolsep) * \real{0.0830}}
  >{\raggedleft\arraybackslash}p{(\columnwidth - 16\tabcolsep) * \real{0.1073}}
  >{\raggedleft\arraybackslash}p{(\columnwidth - 16\tabcolsep) * \real{0.0788}}
  >{\raggedright\arraybackslash}p{(\columnwidth - 16\tabcolsep) * \real{0.0928}}
  >{\raggedleft\arraybackslash}p{(\columnwidth - 16\tabcolsep) * \real{0.1040}}
  >{\raggedleft\arraybackslash}p{(\columnwidth - 16\tabcolsep) * \real{0.1395}}
  >{\raggedleft\arraybackslash}p{(\columnwidth - 16\tabcolsep) * \real{0.1381}}
  >{\raggedleft\arraybackslash}p{(\columnwidth - 16\tabcolsep) * \real{0.0966}}@{}}
\toprule\noalign{}
\begin{minipage}[b]{\linewidth}\raggedright
\textbf{Sınıf}
\end{minipage} & \begin{minipage}[b]{\linewidth}\raggedleft
\textbf{Çalışan}
\end{minipage} & \begin{minipage}[b]{\linewidth}\raggedleft
\textbf{İhracat (\$ mn)}
\end{minipage} & \begin{minipage}[b]{\linewidth}\raggedleft
\textbf{AB payı}
\end{minipage} & \begin{minipage}[b]{\linewidth}\raggedright
\textbf{GB statüsü}
\end{minipage} & \begin{minipage}[b]{\linewidth}\raggedleft
\textbf{S4 AB tarifesi}
\end{minipage} & \begin{minipage}[b]{\linewidth}\raggedleft
\textbf{AB'ye ihracat S4 P50 (\$ mn)}
\end{minipage} & \begin{minipage}[b]{\linewidth}\raggedleft
\textbf{İstihdam S4 P50 (kişi)}
\end{minipage} & \begin{minipage}[b]{\linewidth}\raggedleft
\textbf{TÜFE S4 (bp)}
\end{minipage} \\
\midrule\noalign{}
\endhead
\bottomrule\noalign{}
\endlastfoot
13.30 Bitirme (boya, baskı, apre) & 88.781 & 0 & \%0 & GB & --- & 0,0 & 0 & 0,00 \\
13.94 Halat ve ağ & 2.223 & 99 & \%38 & GB & \%6,8 & −4,9 & −25 & 0,02 \\
13.99 Diğer tekstil (şerit, dantel, keçe) & 21.343 & 297 & \%39 & GB & \%6,4 & −11,5 & −131 & 0,03 \\
14.11 Deri giyim & 5.103 & 128 & \%41 & GB & \%4,0 & −4,9 & −27 & 0,02 \\
14.12 İş giysisi & 12.634 & 69 & \%53 & GB & \%12,0 & −8,1 & −62 & 0,01 \\
14.20 Kürk & 1.602 & 94 & \%37 & GB & \%3,7 & −2,0 & −8 & 0,01 \\
\end{longtable}

\endgroup

\emph{Kaynak:} TÜİK YSHİ 2024, BACI 2024 \etiketV{}; senaryolar \etiketT{} f3e. Veri sayfaları Ek C'de.

Boya, baskı ve apre hizmetleri (13.30) 88.781 kişiyle büyük bir sınıf, ama ticareti yok: fason hizmet olarak çalışıyor ve etkisi kumaş ve giyim sınıflarının üzerinden geliyor. Revize PEM kurallarında dokuma kumaşa tek başına baskının menşe kazandırması bu sınıfa S2'de yeni bir rol veriyor. Asya'dan gelen ham kumaşı Türkiye'de baskılamak, giysiyi menşeli kılmanın en ucuz yolu. Değerlendirmenin güveni düşük, çünkü sınıfın ticaret verisi yok. Halat ve balık ağında (13.94) AB'nin payı küçük, menşe kuralı ise iplik çekimini şart koşuyor. Şerit, dantel, etiket ve keçede (13.99) kayıp küçük. Bu ürünlerin bir kısmı PEM ortaklarında giysiye süsleme olarak giriyor, ama çoğu \%15'lik toleransın içinde kalıyor. Deri giyimde (14.11) ilk pazar ABD. İş giysisinde (14.12) sınıfın değeri ihracattan çok güvenlik güçlerinin tedarikinden geliyor (SS 23). Kürkte (14.20) Türkiye dünya dördüncüsü; ilk pazarlar Fransa, İtalya, ABD ve Kazakistan.

\hypertarget{nace_kartlar_13-15-6-defter-kux131sux131tlar-ve-kalite-kontroluxfc}{%
\section{6. Defter, kısıtlar ve kalite kontrolü}\label{nace_kartlar_13-15-6-defter-kux131sux131tlar-ve-kalite-kontroluxfc}}

Kümülatif defter, tamamlanan parçaların şoklarını ve önerilen önlemlerini bir araya getiriyor ve Faz 0'da konan makro kısıtlara karşı bir ara hesap tutuyor. Kısıtlar şunlar: birinci tur TÜFE etkisi yılda 30, kümülatif olarak 50 baz puanı aşmamalı; cari denge T+5'te en az sıfır olmalı; önlemlerin mali maliyeti yılda GSYH'nin \%0,5'ini geçmemeli; istihdam ve kabiliyet kayıpları sınırlanmalı. Bu parçadan itibaren defter kümülatif: aşağıdaki tablo 10--15 bölümlerinin, yani ilk iki parçanın 54 sınıfının toplamı. Toplamların yüzdelikleri yine sınıf yüzdeliklerinin toplanmasıyla değil, her çekilişteki 10--15 toplamının dağılımından alınıyor. Tablonun ``bu parçanın katkısı'' satırı 13--15'i ayrıca gösteriyor.

\Needspace{7\baselineskip}

\textbf{Tablo 6.1 --- Kümülatif defter (Faz 3, parça 2 sonrası)}

\nopagebreak[4]

\begingroup\gbtablo\small

\begin{longtable}[]{@{}
  >{\raggedright\arraybackslash}p{(\columnwidth - 2\tabcolsep) * \real{0.1314}}
  >{\raggedright\arraybackslash}p{(\columnwidth - 2\tabcolsep) * \real{0.8686}}@{}}
\toprule\noalign{}
\begin{minipage}[b]{\linewidth}\raggedright
\textbf{Kalem}
\end{minipage} & \begin{minipage}[b]{\linewidth}\raggedright
\textbf{Değer}
\end{minipage} \\
\midrule\noalign{}
\endhead
\bottomrule\noalign{}
\endlastfoot
Tamamlanan / kalan sınıf & 54 / 176 C sınıfı; ek kapsam (B, D, E, hizmet) 0 / 29 \\
Kullanılan / kalan enflasyon bütçesi (bp; birinci tur, P50) & S2 yolu: şok 2,6 + önlem −0,11 = 2,4 bp; S4 yolu: şok 21,8 + önlem −13,56 = 8,2 bp. Kümülatif 50 bp bütçeden kalan: S2 47,6, S4 41,8 bp \\
Kümülatif cari denge katkısı (\$ mlr/yıl, T+5) & Senaryo P50: S2 −0,926, S4 −2,753 (S4 P10 −4,15; 1/98 askısında ortalama −3,46); önlemler: S2 +0,221, S4 +0,010 \etiketT{}/\etiketH{} \\
Kümülatif mali maliyet (mlr TL/yıl; GSYH \%) & S2 yolu 1,65 (\%0,002); S4 yolu 22,76 (\%0,031) \\
Riskteki istihdam (kişi, T+5) & S2 P50 −14.630; S4 P50 −40.241, P10 −62.512; 1/98 askısında ortalama −42.977. S4 P10'da kaybı \%5'i aşan sınıf: 13.20 (−\%5,4), 13.91 (−\%5,9), 13.92 (−\%5,3), 14.13 (−\%5,5), 14.14 (−\%8,5), 14.19 (−\%5,4), 14.31 (−\%14,9), 14.39 (−\%19,2) \\
Bu parçanın (13--15) katkısı & TÜFE şoku S2 1,7, S4 15,6 bp; önlemler S2 0,00, S4 −12,01 bp. İstihdam S2 P50 −14.419, S4 P50 −39.397 (P10 −60.851). Mali maliyet S2 0,82, S4 21,07 mlr TL/yıl. Cari denge önlem katkısı S2 +0,150, S4 0,000 \$ mlr; S4'te PEM kümülasyonu kanalı istihdamda −4.750 \\
Açık kalan veri boşlukları & VB-30: sınıf düzeyinde yabancı sermaye payı (TÜİK yabancı kontrollü girişim istatistikleri yalnız bölüm düzeyinde); VB-31: Türkiye'ye özgü sınıf düzeyi Ar-Ge yoğunluğu (OECD sınıflaması vekil); VB-32 (sürüyor): DTÖ bağlı oranları, satır düzeyi. 13--15'te bağlayıcı değil: c2 uygulanan MFN oranlarını değiştirmiyor, yalnız AB tercihini kaldırıyor (giyimde satırların \%3,7'si bağlı); VB-33: TMO yerli buğday fiyatı ile ithal paritesi farkı (menşe maliyeti c\_RoO'nun kalibrasyonu); VB-34: EUDR kapsamındaki ihracatçıların tedarik zinciri verisi; VB-35: YSHİ ile BACI sınıf eşlemesindeki tutarsızlık (10.86; 11.03, 11.04, 11.06'da YSHİ verisi yok); VB-06 (sürüyor): 1/98 tercih oranları satır düzeyinde; 1/98 askısı varyantı bu yüzden üst sınır; VB-36: firma düzeyinde kumaş menşei (giyim ihracatında PEM dışı kumaşın payı; p\_uyum'un kalibrasyonu); VB-37: PEM ortaklarında Türk kumaşının AB'ye giden üretimdeki gerçek payı (kümülasyon maruziyeti üst sınır; κ ve λ \etiketH{}); VB-38: tekstil ve giyimde 2025 istihdam daralmasının sınıf düzeyinde dağılımı (SGK yalnız bölüm düzeyinde); VB-39: deri eşya ve çantada KDV oranı (takasın kapsamı); hızlı modada sınır gecikmesinin maliyeti \\
Varsayım kaydına eklenen yeni maddeler & V-30: Faz 3 senaryoları --- toplam genel dengeden, sınıf dağılımı kısmi dengeden (ortak z, N = 2.000); V-31: S4'te 1/98'in askıya alınma olasılığı 0,25 (hasmane AB) \etiketH{}; V-32: ticaret esnekliği σ --- ProTEE HS6, sınıf içi AB'ye ihracat ağırlıklı medyan, {[}2, 15{]} aralığında; V-33: sınır sürtünmesi \%0,5--2 (yalnız GB ve İTÜ), S2'de × 0,5; AB ihracatçısının menşe maliyeti S2'de \%0,5, S3'te \%1; V-34: arketip eşikleri (KS, AME, ithalat penetrasyonu, EDE ≥ 0,3 bp/puan, SS ≥ 60) \etiketH{}; V-35 (güncellendi): SS boyut puanları ve ağırlıkları \etiketH{} (veri/f3\_ss\_boyutlar.csv; 10--15 gözden geçirildi, diğerleri ilk tahmin); V-36: CKE tanımı --- 1 puan tarife ya da \%1 üretim desteği, 5 yıl, kapasite yatırımının ithalatı (sermaye/hasıla 1, ithal payı 0,6) düşülü; V-37: KN8 → NACE eşlemesinde CPA 2.2 → 2.1 kimlik bağlantısı önceliği (D-1); V-38: Grup toplamlarının yüzdelikleri çekilişlerin toplamından hesaplanır (sınıf yüzdeliklerinin toplamı değil); 1/98 askısının etkisi askı çekilişlerinin koşullu ortalamasıyla raporlanır; V-39: Eylem takvimi T'ye göre; T2028'de T−18 geçmişte kaldığından hazırlık eylemleri 2026'nın son çeyreğinde başlar; V-40: PEM kümülasyonu kanalı yalnız S4'te; kayıp = λ · κ · maruziyet, λ \textasciitilde{} üçgen(0,3; 0,6; 0,9), κ: kumaş 1, teknik ve diğer tekstil 0,3, diğerleri 0 \etiketH{}; genel denge toplamına eklenir; V-41: tekstil, giyim ve deride menşe parametreleri revize PEM kurallarından (Karar 1/2023): giyimde p\_uyum 0,85--0,88 (Asya kumaşı payı \%10--15), kumaşta 0,92 \etiketH{}; V-42: geri ödeme yasağı (PEM Ek I md. 16) bugün AB'ye A.TR ihracatında telafi edici vergi olarak uygulandığından S2'de ek maliyet sayılmadı; V-43: KDV takası (E-13-04) giyim ve ayakkabının S4 TÜFE şokunu tasarım gereği siler; net mali maliyet Faz 2A merkez değeri (laspeyres, c2) \\
\end{longtable}

\endgroup

\emph{Kaynak:} \kod{kod/\allowbreak{}f3g\_\allowbreak{}kartlar.\allowbreak{}py}; matris, grup senaryoları (\kod{veri/\allowbreak{}f3\_\allowbreak{}senaryo\_\allowbreak{}grup.\allowbreak{}csv}) ve eylem kayıtlarından. Senaryo şokları politika seçimi değil, çıkışın kendisi; önlemler Bölüm 4'teki eylemlerden. Makro mutabakat Faz 6'da.

Defterin söylediği şu. Enflasyon tarafında bu parça bütçe hesabını değiştirdi. 13--15'in S4'teki TÜFE şoku (15,6 baz puan) ilk parçanınkinin (6,2 baz puan) iki katından fazla. İlk iki parçanın birlikte S4 şoku 21,8 baz puan; önlemler olmasaydı kümülatif 50 baz puanlık bütçenin büyük bir kısmı iki parçada tükenmiş olacaktı. Vergi takası ve ilk parçanın tarife indirimleriyle S4 yolunda kümülatif kullanım 8,2 baz puana iniyor; S2 yolunda 2,4 baz puan. Bu, kalan on dört parça için bir uyarı: otomotiv, beyaz eşya ve kozmetik gibi c2 tarifesinin AB ürününü pahalandırdığı tüketim mallarında takas, bütçenin korunmasında ana araç olacak. Faz 2A'nın on ürün grubu için yaptığı takas hesabı, bu parçaların defterine aynı yöntemle taşınacak.

Mali tarafta bu parçanın önlemleri S4 yolunda yılda 21,07 milyar TL tutuyor; çoğu vergi takası. Kümülatif mali yük ilk iki parçada S2 yolunda 1,65, S4 yolunda 22,76 milyar TL ve yıllık tavanın (GSYH'nin \%0,5'i) çok altında. Cari denge tarafında tablo daha kötü. 13--15, S4'te cari dengeyi yılda 2.634 milyon dolar kötüleştiriyor ve bu kayıp politika ile telafi edilemiyor. Önlemler istihdamı ve fiyatları koruyor, ihracatı geri getirmiyor. S2'de kayıp 893 milyon dolar. Bunun bir kısmını menşe hazırlığı ve menşe kotası geri kazandırıyor (150 milyon dolar \etiketH{}). T+5'te cari dengenin sıfırın üstünde kalması kısıtı imalatın bütününde, Faz 6'da değerlendirilecek. Bu parçanın katkısı, kısıtın S4 yolunda ne kadar zor olacağını gösteriyor.

İstihdam tarafında bu parça çalışmanın en ağır bulgusunu üretiyor. S4'te ilk iki parçanın toplam istihdam kaybı 40.241 kişi (P10'da 62.512) ve bunun neredeyse tamamı 13--15'ten geliyor. S4'te dağılımın alt ucunda istihdam kaybı \%5'i aşan sınıfların hepsi bu parçada. Önerilen önlemler bu kaybı yumuşatıyor ama ortadan kaldırmıyor. Kısa çalışma ödeneği geçişi zamana yayıyor. İşgücü uyum programı işini kaybedenlerin bir kısmını başka işlere yerleştiriyor. Vergi takasının iç pazarda yarattığı ek talep, Faz 2A'nın hesabına göre giyim ve ayakkabıda birkaç bin kişilik ek istihdam getiriyor \etiketT{}; bu etki tabloda yok. S2'nin bu sektördeki değeri en açık biçimde istihdamda görülüyor: S2'de kayıp 14.419 kişi, yani S4'teki kaybın \%37 kadarı.

Belgenin kalite kontrolü şu noktaları kapsıyor. Her sayının birimi ve dönemi belirtildi: sınıf verileri 2024, AB'ye ihracatın rejim dağılımı ve AB tarifeleri 2025--2026, senaryolar T+5'teki yıllık fark. PEM kümülasyonu kanalı modele yeni eklendi (V-40). Bu kanal yalnız S4'te çalışıyor, genel denge toplamının üstüne ekleniyor ve ilk parçanın sınıflarında sıfır. Ekleme, imalatın S4 istihdam kaybını değiştirdi; ilk parçanın S2 sonuçları ise menşe varsayımlarındaki güncelleme nedeniyle çok az değişti (Ek D). Bütün tarife önerileri MFN değişikliği değil. c2, uygulanan MFN oranlarını yükseltmiyor, yalnızca AB'nin tercihini kaldırıyor. Bu yüzden DTÖ bağlı oranlarıyla bir çatışma doğmuyor. Vergi takası bütün menşelere eşit uygulanıyor ve GATT III:2 ile uyumlu. Paket sanayi ürünlerinde ihracata bağlı hiçbir destek içermiyor (ASCM md. 3). Kısa çalışma ve işgücü uyum araçları ihracattan bağımsız ve sektörel ya da bölgesel kriz gerekçesine dayanıyor; kredi sigortası piyasa primiyle sunuluyor.

Önceki parçada tanımlanan iki tipik hatadan bu parçada da kaçınıldı. Gümrük Birliği dışında kalan bir tarifeye çıkışa bağlı bir artış eklenmedi. Bu parçada tarım rejimi ya da işlenmiş tarım ürünü yok; bütün satırlar sanayi rejiminde ve S4'te tam MFN'yle karşılaşıyor. Olumsuz bulgular gizlenmedi. S4'te giyimdeki kaybın politika ile önlenemeyeceği ve kümülasyon kanalının bir üst sınır tahmini olduğu, yönetici özetinde ve ilgili sınıf incelemelerinde açıkça yazıldı. Belirsiz olgular \etiketD{} olarak işaretlendi ve veri boşlukları Ek E'de listelendi.

\hypertarget{nace_kartlar_13-15-ek-a--ana-matris-satux131rlarux131}{%
\section{Ek A --- Ana matris satırları}\label{nace_kartlar_13-15-ek-a--ana-matris-satux131rlarux131}}

Ana matrisin bütün sütunları ve 230 imalat sınıfının tamamı \kod{veri/\allowbreak{}nace\_\allowbreak{}matris.\allowbreak{}csv} dosyasında. Aşağıdaki iki tablo çalışma çerçevesinin öngördüğü matris sütunlarını 13--15 sınıfları için veriyor. Skorlar 0--100 arasında, C sınıfları içindeki yüzdelik sıra olarak ifade ediliyor.

\begin{landscape}

\textbf{Tablo A.1 --- Ana matris: kimlik, ticaret ve maruziyet}

\nopagebreak[4]

\begingroup\gbtablo\footnotesize

\begin{longtable}[]{@{}
  >{\raggedright\arraybackslash}p{(\columnwidth - 26\tabcolsep) * \real{0.0402}}
  >{\raggedright\arraybackslash}p{(\columnwidth - 26\tabcolsep) * \real{0.1844}}
  >{\raggedright\arraybackslash}p{(\columnwidth - 26\tabcolsep) * \real{0.0770}}
  >{\raggedleft\arraybackslash}p{(\columnwidth - 26\tabcolsep) * \real{0.0402}}
  >{\raggedleft\arraybackslash}p{(\columnwidth - 26\tabcolsep) * \real{0.0684}}
  >{\raggedleft\arraybackslash}p{(\columnwidth - 26\tabcolsep) * \real{0.0600}}
  >{\raggedleft\arraybackslash}p{(\columnwidth - 26\tabcolsep) * \real{0.0586}}
  >{\raggedleft\arraybackslash}p{(\columnwidth - 26\tabcolsep) * \real{0.0572}}
  >{\raggedleft\arraybackslash}p{(\columnwidth - 26\tabcolsep) * \real{0.0980}}
  >{\raggedleft\arraybackslash}p{(\columnwidth - 26\tabcolsep) * \real{0.0735}}
  >{\raggedright\arraybackslash}p{(\columnwidth - 26\tabcolsep) * \real{0.0884}}
  >{\raggedleft\arraybackslash}p{(\columnwidth - 26\tabcolsep) * \real{0.0723}}
  >{\raggedleft\arraybackslash}p{(\columnwidth - 26\tabcolsep) * \real{0.0402}}
  >{\raggedleft\arraybackslash}p{(\columnwidth - 26\tabcolsep) * \real{0.0415}}@{}}
\toprule\noalign{}
\begin{minipage}[b]{\linewidth}\raggedright
\textbf{Kod}
\end{minipage} & \begin{minipage}[b]{\linewidth}\raggedright
\textbf{Ad}
\end{minipage} & \begin{minipage}[b]{\linewidth}\raggedright
\textbf{GB statüsü}
\end{minipage} & \begin{minipage}[b]{\linewidth}\raggedleft
\textbf{KD}
\end{minipage} & \begin{minipage}[b]{\linewidth}\raggedleft
\textbf{İstihdam}
\end{minipage} & \begin{minipage}[b]{\linewidth}\raggedleft
\textbf{İhracat}
\end{minipage} & \begin{minipage}[b]{\linewidth}\raggedleft
\textbf{AB payı}
\end{minipage} & \begin{minipage}[b]{\linewidth}\raggedleft
\textbf{İthalat}
\end{minipage} & \begin{minipage}[b]{\linewidth}\raggedleft
\textbf{AB ithalat payı}
\end{minipage} & \begin{minipage}[b]{\linewidth}\raggedleft
\textbf{τ\_AB (S4)}
\end{minipage} & \begin{minipage}[b]{\linewidth}\raggedright
\textbf{Menşe S2/S3}
\end{minipage} & \begin{minipage}[b]{\linewidth}\raggedleft
\textbf{İthal girdi}
\end{minipage} & \begin{minipage}[b]{\linewidth}\raggedleft
\textbf{RCA}
\end{minipage} & \begin{minipage}[b]{\linewidth}\raggedleft
\textbf{PCI}
\end{minipage} \\
\midrule\noalign{}
\endhead
\bottomrule\noalign{}
\endlastfoot
13.10 & İplik & GB & 2.741 & 86.468 & 2.056 & \%43 & 1.286 & \%12 & \%3,8 & D/D & \%23 & 5,29 & −1,24 \\
13.20 & Dokuma kumaş & GB & 2.815 & 96.394 & 3.088 & \%32 & 1.322 & \%31 & \%8,0 & D/D & \%23 & 3,53 & −0,52 \\
13.30 & Bitirme (boya, baskı, apre) & GB & 1.738 & 88.781 & 0 & \%0 & 0 & \%0 & \%--- & D/D & \%23 & --- & --- \\
13.91 & Örme kumaş & GB & 1.060 & 36.092 & 1.806 & \%31 & 487 & \%20 & \%8,0 & D/D & \%23 & 4,13 & −0,80 \\
13.92 & Ev tekstili & GB & 1.580 & 87.344 & 2.346 & \%54 & 261 & \%34 & \%10,0 & O/O & \%23 & 2,50 & −1,30 \\
13.93 & Halı & GB & 1.493 & 46.903 & 3.056 & \%16 & 93 & \%49 & \%7,9 & D/O & \%23 & 15,79 & −0,89 \\
13.94 & Halat ve ağ & GB & 47 & 2.223 & 99 & \%38 & 53 & \%32 & \%6,8 & O/O & \%23 & 1,30 & −0,68 \\
13.95 & Dokusuz kumaş & GB & 458 & 8.474 & 804 & \%39 & 209 & \%67 & \%4,3 & D/D & \%23 & 3,90 & 0,47 \\
13.96 & Teknik tekstil & GB & 472 & 19.387 & 528 & \%42 & 589 & \%44 & \%6,4 & D/D & \%23 & 1,42 & 0,31 \\
13.99 & Diğer tekstil (şerit, dantel, keçe) & GB & 436 & 21.343 & 297 & \%39 & 119 & \%33 & \%6,4 & D/O & \%23 & 4,00 & −0,31 \\
14.11 & Deri giyim & GB & 61 & 5.103 & 128 & \%41 & 17 & \%57 & \%4,0 & D/D & \%23 & 4,22 & −0,97 \\
14.12 & İş giysisi & GB & 189 & 12.634 & 69 & \%53 & 11 & \%26 & \%12,0 & O/O & \%23 & 1,97 & −1,70 \\
14.13 & Dış giyim & GB & 4.024 & 365.708 & 8.577 & \%59 & 2.081 & \%26 & \%12,0 & O/O & \%23 & 3,36 & −1,73 \\
14.14 & İç giyim ve tişört & GB & 2.482 & 183.980 & 5.647 & \%60 & 883 & \%21 & \%12,0 & D/O & \%23 & 3,77 & −1,90 \\
14.19 & Diğer giyim ve aksesuar & GB & 466 & 41.506 & 1.461 & \%47 & 507 & \%23 & \%11,2 & O/O & \%23 & 1,59 & −1,66 \\
14.20 & Kürk & GB & 37 & 1.602 & 94 & \%37 & 5 & \%91 & \%3,7 & D/D & \%23 & 7,07 & −0,59 \\
14.31 & Çorap & GB & 607 & 25.205 & 1.296 & \%65 & 59 & \%19 & \%12,0 & D/D & \%23 & 7,53 & −1,44 \\
14.39 & Triko (kazak, hırka) & GB & 552 & 32.975 & 2.719 & \%66 & 459 & \%27 & \%12,0 & D/D & \%23 & 3,60 & −1,98 \\
15.11 & Deri tabaklama & GB & 201 & 8.899 & 268 & \%40 & 138 & \%51 & \%3,7 & D/D & \%23 & 1,58 & −0,86 \\
15.12 & Bavul ve çanta & GB & 234 & 16.310 & 377 & \%37 & 668 & \%51 & \%3,9 & D/D & \%23 & 0,37 & −1,39 \\
15.20 & Ayakkabı & GB & 765 & 50.596 & 1.176 & \%33 & 1.720 & \%21 & \%9,9 & D/D & \%23 & 0,62 & −1,46 \\
\end{longtable}

\endgroup

\end{landscape}

\emph{Kaynak:} \kod{veri/\allowbreak{}nace\_\allowbreak{}matris.\allowbreak{}csv}. KD, ihracat ve ithalat \$ mn (2024). τ\_AB: S4'te eklenecek AB tarifesi (1/98 sürer varsayımı).

\begin{landscape}

\textbf{Tablo A.2 --- Ana matris: skorlar, strateji ve senaryo}

\nopagebreak[4]

\begingroup\gbtablo\footnotesize

\begin{longtable}[]{@{}
  >{\raggedright\arraybackslash}p{(\columnwidth - 28\tabcolsep) * \real{0.0408}}
  >{\raggedleft\arraybackslash}p{(\columnwidth - 28\tabcolsep) * \real{0.0285}}
  >{\raggedleft\arraybackslash}p{(\columnwidth - 28\tabcolsep) * \real{0.0285}}
  >{\raggedleft\arraybackslash}p{(\columnwidth - 28\tabcolsep) * \real{0.0420}}
  >{\raggedleft\arraybackslash}p{(\columnwidth - 28\tabcolsep) * \real{0.0388}}
  >{\raggedleft\arraybackslash}p{(\columnwidth - 28\tabcolsep) * \real{0.0388}}
  >{\raggedright\arraybackslash}p{(\columnwidth - 28\tabcolsep) * \real{0.0580}}
  >{\raggedright\arraybackslash}p{(\columnwidth - 28\tabcolsep) * \real{0.0695}}
  >{\raggedright\arraybackslash}p{(\columnwidth - 28\tabcolsep) * \real{0.1421}}
  >{\raggedright\arraybackslash}p{(\columnwidth - 28\tabcolsep) * \real{0.1997}}
  >{\raggedright\arraybackslash}p{(\columnwidth - 28\tabcolsep) * \real{0.0724}}
  >{\raggedleft\arraybackslash}p{(\columnwidth - 28\tabcolsep) * \real{0.0572}}
  >{\raggedleft\arraybackslash}p{(\columnwidth - 28\tabcolsep) * \real{0.0593}}
  >{\raggedright\arraybackslash}p{(\columnwidth - 28\tabcolsep) * \real{0.0740}}
  >{\raggedright\arraybackslash}p{(\columnwidth - 28\tabcolsep) * \real{0.0503}}@{}}
\toprule\noalign{}
\begin{minipage}[b]{\linewidth}\raggedright
\textbf{Kod}
\end{minipage} & \begin{minipage}[b]{\linewidth}\raggedleft
\textbf{KS}
\end{minipage} & \begin{minipage}[b]{\linewidth}\raggedleft
\textbf{SS}
\end{minipage} & \begin{minipage}[b]{\linewidth}\raggedleft
\textbf{AME}
\end{minipage} & \begin{minipage}[b]{\linewidth}\raggedleft
\textbf{EDE}
\end{minipage} & \begin{minipage}[b]{\linewidth}\raggedleft
\textbf{CKE}
\end{minipage} & \begin{minipage}[b]{\linewidth}\raggedright
\textbf{Arketip}
\end{minipage} & \begin{minipage}[b]{\linewidth}\raggedright
\textbf{Duruş}
\end{minipage} & \begin{minipage}[b]{\linewidth}\raggedright
\textbf{Tarife yolu}
\end{minipage} & \begin{minipage}[b]{\linewidth}\raggedright
\textbf{Ana araçlar}
\end{minipage} & \begin{minipage}[b]{\linewidth}\raggedright
\textbf{AB müzakere}
\end{minipage} & \begin{minipage}[b]{\linewidth}\raggedleft
\textbf{S2 cari P50}
\end{minipage} & \begin{minipage}[b]{\linewidth}\raggedleft
\textbf{S2 TÜFE P50}
\end{minipage} & \begin{minipage}[b]{\linewidth}\raggedright
\textbf{İstihdam riski}
\end{minipage} & \begin{minipage}[b]{\linewidth}\raggedright
\textbf{Güven}
\end{minipage} \\
\midrule\noalign{}
\endhead
\bottomrule\noalign{}
\endlastfoot
13.10 & 54 & 0 & 80 & 0,07 & 73 & A & Geliştir & AB ile karşılıklı \%0 (S2); ithal pamuk ve lifte MFN \%0 sürer; S4'te AB'ye c2 & Pamuk izlenebilirliği (Tüzük 2024/3015); menşe verisi (E-13-02); pazar çeşitlendirme & savunma & −72,1 & 0,01 & Düşük & Yüksek \\
13.20 & 55 & 8 & 87 & 0,07 & 50 & A/B & Geliştir & AB ile karşılıklı \%0 (S2); S4'te AB'ye c2; PEM kümülasyonu S2 metninde & PEM kümülasyonunun sürekliliği (E-13-01); ortakta baskı + konfeksiyon yolu; pazar çeşitlendirme & savunma & −22,5 & 0,06 & Yüksek & Orta \\
13.30 & 58 & 0 & 5 & 0,00 & 5 & --- & Koru & Değişiklik yok (hizmet; ticareti yok) & Baskı kapasitesi (tek başına baskı menşe kazandırıyor; E-13-02); enerji ve su verimliliği & nötr & 0,0 & 0,00 & Düşük & Düşük \\
13.91 & 51 & 0 & 85 & 0,03 & 49 & A/B & Geliştir & AB ile karşılıklı \%0 (S2); S4'te AB'ye c2; PEM kümülasyonu S2 metninde & PEM kümülasyonunun sürekliliği (E-13-01); menşe verisi (E-13-02); Belarus ve Rusya pazar riskinin izlenmesi & savunma & −19,6 & 0,02 & Yüksek & Orta \\
13.92 & 53 & 0 & 96 & 0,05 & 31 & A & Geliştir & AB ile karşılıklı \%0 (S2); S4'te AB'ye c2 & Dokuma + konfeksiyon entegrasyonunda menşe; ticaret savunması hazırlığı; ABD ve BK pazarları (E-13-05) & savunma & −84,1 & 0,05 & Yüksek & Orta \\
13.93 & 53 & 0 & 60 & 0,02 & 23 & B & Geliştir & AB ile karşılıklı \%0 (S2); S4'te AB'ye c2; iplik girdisinde mevcut rejim & ABD ve Körfez pazarları (E-13-05); iplik çekimi entegrasyonuyla menşe uyumu & nötr & −14,4 & 0,01 & Düşük & Yüksek \\
13.94 & 48 & 7 & 89 & 0,00 & 49 & C & Geliştir & AB ile karşılıklı \%0 (S2); S4'te AB'ye c2; yeni koruma yok & İzleme & nötr & −1,2 & 0,00 & Düşük & Yüksek \\
13.95 & 62 & 37 & 76 & 0,01 & 59 & A & Koru & AB ile karşılıklı \%0 (S2); S4'te AB'ye c2; polimer girdisinde mevcut rejim & Sağlık arz güvenliği (kriz kapasite envanteri); hijyen ürünlerinde pazar çeşitlendirme & savunma & −6,6 & 0,02 & Düşük & Orta \\
13.96 & 54 & 33 & 90 & 0,03 & 63 & A & Geliştir & AB ile karşılıklı \%0 (S2); S4'te AB'ye c2; aramid ve cam elyafı girdisinde \%0 adayı & SSB yerlileştirme ve kamu alımı; Ar-Ge; otomotiv tedarik zincirinde menşe & savunma & −2,9 & 0,03 & Düşük & Yüksek \\
13.99 & 60 & 0 & 88 & 0,01 & 34 & A & Koru & AB ile karşılıklı \%0 (S2); girdilerde artış yok; S4'te AB'ye MFN (c2) & S2'de \%0 ve elverişli menşe kuralı; ACAA; çapa yatırımcı tutma & savunma & −2,4 & 0,00 & Düşük & Yüksek \\
14.11 & 45 & 0 & 76 & 0,00 & 41 & A & Geliştir & AB ile karşılıklı \%0 (S2); girdilerde artış yok; S4'te AB'ye MFN (c2) & S2'de \%0 ve elverişli menşe kuralı; ACAA; çapa yatırımcı tutma & savunma & −0,9 & 0,00 & Düşük & Yüksek \\
14.12 & 43 & 23 & 97 & 0,00 & 34 & --- & Koru & Değişiklik yok & İzleme & nötr & −2,0 & 0,00 & Düşük & Yüksek \\
14.13 & 46 & 0 & 97 & 0,54 & 30 & A/F & Dönüştür & AB ile karşılıklı \%0 (S2); S4'te AB'ye c2 ve KDV takası (E-13-04) & Dönüşüm programı (tasarım, hızlı tedarik, otomasyon); menşe verisi (E-13-02); S4'te kısa çalışma (E-13-03) ve KDV takası & savunma & −227,4 & 0,45 & Yüksek & Orta \\
14.14 & 51 & 0 & 97 & 0,10 & 44 & F/A & Yönetilen uyum & AB ile karşılıklı \%0 (S2); S4'te AB'ye c2 ve KDV takası (E-13-04) & Yönetilen uyum (işgücü uyumu, bölgesel program); örme ve konfeksiyon entegrasyonunda menşe; BK pazarı & savunma & −229,5 & 0,08 & Yüksek & Orta \\
14.19 & 43 & 0 & 96 & 0,16 & 35 & F/A & Yönetilen uyum & AB ile karşılıklı \%0 (S2); S4'te AB'ye c2 ve KDV takası (E-13-04) & Yönetilen uyum (E-14.14-01); bebek giyiminde menşe; pazar çeşitlendirme & savunma & −24,8 & 0,14 & Yüksek & Orta \\
14.20 & 61 & 0 & 68 & 0,00 & 32 & B & Geliştir & Girdilerde \%0 (P1); nihai üründe MFN korunur & Yeni STA'lar; ihracat finansmanı; girdi tarifesi indirimi & nötr & −0,5 & 0,00 & Düşük & Yüksek \\
14.31 & 51 & 0 & 99 & 0,00 & 56 & A/F & Geliştir & AB ile karşılıklı \%0 (S2); S4'te AB'ye c2 ve KDV takası (E-13-04) & BK pazarı ve AB--TR--BK kümülasyonu (E-13-05); otomasyon; S4'te kısa çalışma (E-13-03) & savunma & −55,1 & 0,00 & Yüksek & Yüksek \\
14.39 & 46 & 0 & 99 & 0,13 & 44 & F/A & Yönetilen uyum & AB ile karşılıklı \%0 (S2); S4'te AB'ye c2 ve KDV takası (E-13-04) & Yönetilen uyum (E-14.14-01); tam kalıp örme otomasyonu; BK pazarı & savunma & −70,5 & 0,15 & Yüksek & Orta \\
15.11 & 45 & 0 & 73 & 0,01 & 56 & A/G & Dönüştür & AB ile karşılıklı \%0 (S2); ham deride MFN \%0 sürer; S4'te AB'ye c2 & Organize deri sanayi bölgelerinde arıtma ve krom azaltımı; LWG sertifikası; AB alıcılarının özen yükümlülüğüne uyum & nötr & −2,5 & 0,01 & Düşük & Orta \\
15.12 & 38 & 0 & 69 & 0,27 & 39 & C & Geliştir & AB ile karşılıklı \%0 (S2); S4'te AB'ye c2 (iç ikame) & Tasarım ve marka; S4'te iç ikame; TÜFE izleme & savunma & +0,7 & 0,25 & Düşük & Yüksek \\
15.20 & 39 & 8 & 92 & 0,47 & 57 & --- & Koru & AB ile karşılıklı \%0 (S2); S4'te AB'ye c2 ve KDV takası (E-13-04) & KDV takası (E-13-04); S2'de menşe (saya kuralı); Körfez ve Irak pazarları & nötr & −5,1 & 0,39 & Düşük & Yüksek \\
\end{longtable}

\endgroup

\end{landscape}

\emph{Kaynak:} Skorlar 0--100 (C sınıfları içinde yüzdelik); EDE bp/puan; S2 cari denge \$ mn; TÜFE bp. Arketip ve duruş kural tabanlı öneri ve uzman düzeltmesi.

\hypertarget{nace_kartlar_13-15-ek-b--yuxf6ntem-bu-paruxe7adaki-yenilikler}{%
\section{Ek B --- Yöntem: bu parçadaki yenilikler}\label{nace_kartlar_13-15-ek-b--yuxf6ntem-bu-paruxe7adaki-yenilikler}}

Hesap zinciri ilk parçanın Ek B'sinde anlatıldı ve değişmedi: f3a--f3f bütün 230 sınıfı birlikte hesaplıyor, f3g her parçanın belgesini üretiyor. Bu parçada zincire dört ekleme yapıldı ve parçaya özgü varsayımlar derlendi. Tablo B.1 değişen betikleri özetliyor.

\Needspace{7\baselineskip}

\textbf{Tablo B.1 --- Bu parçada değişen betikler ve tablolar}

\nopagebreak[4]

\begingroup\gbtablo\small

\begin{longtable}[]{@{}
  >{\raggedright\arraybackslash}p{(\columnwidth - 4\tabcolsep) * \real{0.1654}}
  >{\raggedright\arraybackslash}p{(\columnwidth - 4\tabcolsep) * \real{0.6759}}
  >{\raggedright\arraybackslash}p{(\columnwidth - 4\tabcolsep) * \real{0.1586}}@{}}
\toprule\noalign{}
\begin{minipage}[b]{\linewidth}\raggedright
\textbf{Betik ya da tablo}
\end{minipage} & \begin{minipage}[b]{\linewidth}\raggedright
\textbf{Değişiklik}
\end{minipage} & \begin{minipage}[b]{\linewidth}\raggedright
\textbf{Etiket}
\end{minipage} \\
\midrule\noalign{}
\endhead
\bottomrule\noalign{}
\endlastfoot
\kod{kod/\allowbreak{}f3a\_\allowbreak{}ticaret.\allowbreak{}py} & PEM kümülasyonu maruziyeti: PEM ortaklarına ihracat (\kod{x\_\allowbreak{}pem}) ve ortağın AB payıyla ağırlıklı maruziyet (\kod{x\_\allowbreak{}kum}); \kod{veri/\allowbreak{}f3\_\allowbreak{}pem\_\allowbreak{}ortak\_\allowbreak{}ab\_\allowbreak{}payi.\allowbreak{}csv} & \etiketV{} veri, \etiketT{} eşleme \\
\kod{kod/\allowbreak{}f3d\_\allowbreak{}ab\_\allowbreak{}tarife.\allowbreak{}py} & EUDR kapsamı Tüzük 2026/2102'ye göre güncellendi; 30.12.2027'den eklenen satırlar ayrı bayrakla (\kod{eudr\_\allowbreak{}2027}) & \etiketV{} \\
\kod{kod/\allowbreak{}f3e\_\allowbreak{}senaryo.\allowbreak{}py} & S4'te PEM kümülasyonu kanalı (\kod{dx\_\allowbreak{}kum}, \kod{dl\_\allowbreak{}kum}, \kod{dca\_\allowbreak{}kum}) & \etiketT{}; parametreler \etiketH{} \\
\kod{kod/\allowbreak{}f3f\_\allowbreak{}skor\_\allowbreak{}matris.\allowbreak{}py} & Arketip düzeltmesinin kaydı ikincil arketipi de gösteriyor & --- \\
\kod{kod/\allowbreak{}f3g\_\allowbreak{}kartlar.\allowbreak{}py} & Parça numarası ve etiketi YAML'dan; kümülatif defter bütün parçaların kayıtlarından; bölge ve kümülasyon tabloları; ``−şok'' kaydıyla TÜFE önleminin veriden çözülmesi & --- \\
\kod{veri/\allowbreak{}f3\_\allowbreak{}sinif\_\allowbreak{}varsayim.\allowbreak{}csv} & 13--15 için menşe kuralları, c\_RoO, p\_uyum, değer zinciri rolü, DİR bağımlılığı ve κ & \etiketH{}; kurallar \etiketV{} \\
\kod{veri/\allowbreak{}f3\_\allowbreak{}ss\_\allowbreak{}boyutlar.\allowbreak{}csv} & 13--15 için stratejik önem boyutları gözden geçirildi & \etiketH{} \\
\end{longtable}

\endgroup

\textbf{PEM kümülasyonu kanalı.} Kanal, Türk girdisinin bir PEM ortağında işlenip AB'ye tercihli satıldığı zincirleri yakalıyor. Ortaklar, hem Türkiye hem AB ile yürürlükte tercihli anlaşması olan PEM taraflarıdır: Fas, Tunus, Mısır, Batı Balkan ülkeleri, Gürcistan, Moldova, Ukrayna, EFTA ve Faroe Adaları. İsrail ile ticaret askıda olduğu için, Ürdün, Cezayir ve Lübnan ile de Türkiye'nin yürürlükte bir anlaşması olmadığı için dışarıda bırakıldı. Sınıf i için maruziyet, Türkiye'nin 2024'te her ortağa sınıf i'de yaptığı ihracatın, ortağın aşağı yöndeki ürünlerde AB'ye ihracat payıyla çarpımlarının toplamı. Aşağı yön ürünleri tekstil (13) için Fasıl 61--63, diğer bölümler için ortağın toplam ihracatı. S4'teki kayıp λ · κ\_i · maruziyet\_i. κ\_i, Türk girdisinin ortağın ürününde menşe kazandıran adım olma derecesi: dokuma ve örme kumaşta 1, teknik ve diğer tekstilde 0,3, diğer bütün sınıflarda 0 \etiketH{}. Kayıp oranı λ, \%30--90 arasında tepe değeri \%60 olan üçgen bir dağılımdan, ortak bir çekilişle alınıyor \etiketH{}. Kayıp genel denge toplamının üstüne ekleniyor, çünkü Faz 2B'nin genel denge modeli üçüncü ülkeler üzerinden kurulan menşe zincirlerini içermiyor. S2 ve S3'te kanal yok, çünkü yeni anlaşma PEM'e atıf yaptığı sürece kümülasyon sürüyor. λ çekilişi, önceki çekilişlerden sonra yapıldığı için ilk parçanın sonuçlarını değiştirmiyor.

\textbf{Menşe parametreleri.} Tekstil, giyim ve deri sınıflarının menşe kuralı, revize PEM kurallarının ürüne özgü listesinden alındı (Tablo 3.4) \etiketV{}. Uyum oranı (p\_uyum) ve uyumlu ihracatın menşe maliyeti (c\_RoO) bu kurallara ve Faz 2C'nin Asya kumaşı tahminine göre belirlendi \etiketH{}. Kumaş, iplik, halı, deri eşya ve ayakkabıda risk düşük (p\_uyum 0,92, c\_RoO \%1,5). Giyimde risk orta: dış giyimde p\_uyum 0,86 ve c\_RoO \%2,5, iç giyimde 0,88 ve \%2,0. Geri ödeme yasağı için ek maliyet eklenmedi (V-42).

\textbf{Bölgesel dağılım.} Bölge tablosundaki değerler, her sınıfın S2 ve S4 istihdam etkisinin (P50) o sınıfın İBBS-2 bölge paylarıyla dağıtılmasından geliyor. Bölge payları f3b'de TÜİK sınıf çalışanı ile SGK'nın bölüm × il sigortalı verisinin RAS yöntemiyle uzlaştırılmasından türetildi. Bölgedeki imalat payı aynı yöntemle bütün C sınıfları üzerinden hesaplandı. Değerler sınıf medyanlarının toplamıdır, grup dağılımının medyanı değil.

\textbf{Kümülatif defter ve vergi takası.} Defter bu parçadan itibaren kümülatif. f3g, \kod{kartlar/\allowbreak{}} klasöründeki bütün parça dosyalarını okuyor ve son bölümü bu parçanınkinden büyük olmayan parçaların eylem kayıtlarını topluyor. Senaryo toplamları 10--{[}son bölüm{]} grubunun çekiliş dağılımından alınıyor. Vergi takasının TÜFE etkisi sınıf defterlerinde ``−şok'' olarak kayıtlı. f3g bu kaydı ilgili sınıfın S4 TÜFE şokuyla (P50) değiştiriyor; böylece takas, tasarımına uygun olarak şoku siliyor ve model değiştiğinde kayıt kendiliğinden güncelleniyor. Takasın net mali maliyeti Faz 2A'nın merkez değeri (Laspeyres kuralı, c2 rejimi; giyim 319, ayakkabı 63 milyon dolar; 2026 ortalama kuru 45,84 TL) \etiketT{}.

\hypertarget{nace_kartlar_13-15-ek-c--sux131nux131f-veri-sayfalarux131}{%
\section{Ek C --- Sınıf veri sayfaları}\label{nace_kartlar_13-15-ek-c--sux131nux131f-veri-sayfalarux131}}

Her veri sayfası sınıf kartının sayısal alanlarını bir arada veriyor. Senaryo tablosu Bölüm 5'te verilen sınıflar için sayfada ilgili tabloya atıf var; diğer sınıfların senaryo tablosu sayfanın hemen altında. Veriler \etiketV{}; eşleme, bölge payları ve skorlar \etiketT{}; menşe parametreleri, kümülasyon katsayısı ve stratejik skor \etiketH{}.

\hypertarget{nace_kartlar_13-15-c1-1310-tekstil-elyafux131nux131n-hazux131rlanmasux131-ve-buxfckuxfclmesi}{%
\subsection{C.1 13.10 Tekstil elyafının hazırlanması ve bükülmesi}\label{nace_kartlar_13-15-c1-1310-tekstil-elyafux131nux131n-hazux131rlanmasux131-ve-buxfckuxfclmesi}}

\Needspace{7\baselineskip}

\textbf{Tablo C.1 --- 13.10: veri sayfası}

\nopagebreak[4]

\begingroup\gbtablo\small

\begin{longtable}[]{@{}
  >{\raggedright\arraybackslash}p{(\columnwidth - 4\tabcolsep) * \real{0.1355}}
  >{\raggedright\arraybackslash}p{(\columnwidth - 4\tabcolsep) * \real{0.3580}}
  >{\raggedright\arraybackslash}p{(\columnwidth - 4\tabcolsep) * \real{0.5065}}@{}}
\toprule\noalign{}
\begin{minipage}[b]{\linewidth}\raggedright
\textbf{Alan}
\end{minipage} & \begin{minipage}[b]{\linewidth}\raggedright
\textbf{Gösterge}
\end{minipage} & \begin{minipage}[b]{\linewidth}\raggedright
\textbf{Değer}
\end{minipage} \\
\midrule\noalign{}
\endhead
\bottomrule\noalign{}
\endlastfoot
1 Kimlik & Rev. 2.1 karşılığı; CPA & 13.10 (1:1); CPA 13.10 \\
& Başlıca HS4 (ihracat, 2024) & 5205 \%31, 5509 \%18, 5402 \%15, 5511 \%12 \\
& BEC payları: ara / sermaye / tüketim & \%89 / \%0 / \%11 \\
2 Ölçek & Üretim; katma değer (GSYH payı) & \$8.999 mn; \$2.741 mn (\%0,20) \\
& Çalışan; girişim & 86.468; 1.761 \\
& Girişim ölçeği: 0 / 1--9 / 10--49 / 50--249 / 250+ çalışan & 137 / 1.130 / 334 / 181 / 89 \\
& 250+ girişimlerin istihdam payı; kişi başı KD (AB27 = 1) & \%65; 0,63 \\
& İstihdam değişimi, 2015--2024 & \%3 \\
& İlk üç bölge (İBBS-2) & Gaziantep, Adıyaman, Kilis \%20; Tekirdağ, Edirne, Kırklareli \%12; Hatay, Kahramanmaraş, Osmaniye \%11 \\
& Çapa yatırımcılar & Sanko, Kipaş, İskur (pamuk ipliği; Gaziantep ve Kahramanmaraş), Korteks (polyester tekstüre; Bursa) \etiketH{} \\
3 Dış ticaret & İhracat 2024; 2020--2024 değişimi & \$2.056 mn; \%41 \\
& AB payı; ilk beş AB pazarı & \%43; İtalya \%20; Portekiz \%17; Almanya \%14; Hollanda \%9; İspanya \%8 \\
& İlk beş AB dışı pazar & Mısır \%15; Birleşik Krallık \%14; ABD \%13; Fas \%8; Pakistan \%7 \\
& İthalat 2024; AB ve Çin payı & \$1.286 mn; AB \%12, Çin \%12 \\
& İlk beş ithalat kaynağı & Özbekistan \%25; Hindistan \%15; Çin \%12; Bangladeş \%10; Mısır \%10 \\
& Dış ticaret dengesi & +\$769 mn \\
& İthal girdi içeriği (I-O); KD ÷ üretim & \%23; \%30 \\
& RCA: dünya 2020 → 2024; AB pazarı 2024 & 4,65 → 5,29; 6,09 \\
& Birim değer ÷ AB pazarı ortalaması & 1,07 \\
& AB pazarındaki pay (AB dışı tedarik içinde) 2020 → 2024; Türkiye'nin sırası & \%22,6 → \%24,3; 2. \\
4 AB tarafı & AB'ye ihracat 2025; rejim payları (GB / İTÜ / tarım) & €834,3 mn; \%100 / \%0 / \%0 \\
& S4'te AB tarifesi: 1/98 sürerse / askıda / tam MFN ortalaması & \%3,8 / \%3,8 / \%3,8 \\
& Tepe satırlar (S4 artışı ≥ \%10) & yok \\
& AB pazarında AB dışı rakipler (AB tercihi) & Çin \%24 (MFN/GSP); Hindistan \%15 (MFN/GSP); Mısır \%4 (STA/EPA); Birleşik Krallık \%4 (TCA); Pakistan \%4 (GSP+) \\
& Tercihli rakiplerin payı; AME (puan; ham) & \%26; 80 (2,94) \\
& Menşe kuralı (PEM) & 50--55 iplik: doğal liflerin eğrilmesi, sentetik lif çekimi + eğirme ya da büküm + mekanik işlem (ithal pamuk ve POY serbest) \\
& Menşe riski S2 / S3; c\_RoO (\%); p\_uyum & D/D; 1,5 / 2,0; 0,92 / 0,88 \etiketH{} \\
& Değer zincirindeki rol; DİR bağımlılığı & ara; orta \\
& EUDR ve SKDM kapsamındaki AB'ye ihracat payı & \%0; \%0 \\
& PEM ortaklarına ihracat; kümülasyon maruziyeti; κ & \$380 mn; \$199 mn; 0,0 \etiketH{} \\
5 Türkiye tarafı & Uygulanan tarife 2026 (ithalat ağırlıklı): MFN; İGV dahil & \%3,3; \%11,4 \\
& S4'te AB menşeline tarife (c2) & \%6,9 \\
& Kritik girdiler: 2024 ithalatı, AB payı, ilk kaynaklar & HS4 5201, 5402, 5503: \$2.985 mn (AB \%13; Çin \%23, Brezilya \%15, ABD \%13) \\
6 Kabiliyet & KS; sıra (ağırlık duyarlılığında P10--P90) & 54,2; 68 / 230 (16--146) \\
& KS bileşenleri (yüzdelik): RCA, PCI, yurt içi KD, Ar-Ge, tedarikçi, verimlilik, firma & 85, 6, 64, 25, 96, 79, 53 \\
& PCI (sınıf) & −1,24 \\
& Atlanabilir ürünler (yoğunluk üst çeyrek, PCI \textgreater{} sınıf) & yok \\
7 Stratejik önem & SS; en yüksek boyut; ağırlıklı ortalama & 0,0; --- 0; 0,0 \etiketH{} \\
8 Makro duyarlılık & TÜFE ağırlığı (doğrudan); tüketimde ithal pay & \%0,00; \%22 \\
& ÜFE vekili (imalat üretimindeki pay) & \%1,35 \\
& EDE, 1 puan tarife: tüm menşeler / yalnız AB (bp; ρ = 1) & 0,072 / 0,009 \\
& CKE: tarife (\$ mn / bp); mali (\$ mn / mlr TL); puan & 7.285,3; 27,88; 73 \\
9 Senaryolar & Senaryo tablosu & Tablo 5.1 \\
\end{longtable}

\endgroup

\hypertarget{nace_kartlar_13-15-c2-1320-dokuma}{%
\subsection{C.2 13.20 Dokuma}\label{nace_kartlar_13-15-c2-1320-dokuma}}

\Needspace{7\baselineskip}

\textbf{Tablo C.2 --- 13.20: veri sayfası}

\nopagebreak[4]

\begingroup\gbtablo\small

\begin{longtable}[]{@{}
  >{\raggedright\arraybackslash}p{(\columnwidth - 4\tabcolsep) * \real{0.1351}}
  >{\raggedright\arraybackslash}p{(\columnwidth - 4\tabcolsep) * \real{0.3544}}
  >{\raggedright\arraybackslash}p{(\columnwidth - 4\tabcolsep) * \real{0.5104}}@{}}
\toprule\noalign{}
\begin{minipage}[b]{\linewidth}\raggedright
\textbf{Alan}
\end{minipage} & \begin{minipage}[b]{\linewidth}\raggedright
\textbf{Gösterge}
\end{minipage} & \begin{minipage}[b]{\linewidth}\raggedright
\textbf{Değer}
\end{minipage} \\
\midrule\noalign{}
\endhead
\bottomrule\noalign{}
\endlastfoot
1 Kimlik & Rev. 2.1 karşılığı; CPA & 13.20, 13.96, 23.99 (1:3); CPA 13.20 \\
& Başlıca HS4 (ihracat, 2024) & 5407 \%27, 5209 \%16, 5516 \%12, 5208 \%10 \\
& BEC payları: ara / sermaye / tüketim & \%100 / \%0 / \%0 \\
2 Ölçek & Üretim; katma değer (GSYH payı) & \$8.962 mn; \$2.815 mn (\%0,21) \\
& Çalışan; girişim & 96.394; 3.607 \\
& Girişim ölçeği: 0 / 1--9 / 10--49 / 50--249 / 250+ çalışan & 181 / 2.869 / 498 / 205 / 70 \\
& 250+ girişimlerin istihdam payı; kişi başı KD (AB27 = 1) & \%57; 0,49 \\
& İstihdam değişimi, 2015--2024 & \%1 \\
& İlk üç bölge (İBBS-2) & Bursa, Eskişehir, Bilecik \%32; Hatay, Kahramanmaraş, Osmaniye \%14; Aydın, Denizli, Muğla \%13 \\
& Çapa yatırımcılar & İsko (Sanko), Bossa, Çalık Denim, Kipaş Denim (denim); Bursa ve Denizli dokuma kümeleri \etiketH{} \\
3 Dış ticaret & İhracat 2024; 2020--2024 değişimi & \$3.088 mn; \%35 \\
& AB payı; ilk beş AB pazarı & \%32; İspanya \%28; İtalya \%16; Almanya \%12; Romanya \%6; Polonya \%6 \\
& İlk beş AB dışı pazar & Fas \%32; Tunus \%8; ABD \%6; Mısır \%6; Belarus \%4 \\
& İthalat 2024; AB ve Çin payı & \$1.322 mn; AB \%31, Çin \%23 \\
& İlk beş ithalat kaynağı & Çin \%23; Mısır \%18; İtalya \%16; Pakistan \%7; Malezya \%5 \\
& Dış ticaret dengesi & +\$1.766 mn \\
& İthal girdi içeriği (I-O); KD ÷ üretim & \%23; \%31 \\
& RCA: dünya 2020 → 2024; AB pazarı 2024 & 3,50 → 3,53; 4,31 \\
& Birim değer ÷ AB pazarı ortalaması & 1,44 \\
& AB pazarındaki pay (AB dışı tedarik içinde) 2020 → 2024; Türkiye'nin sırası & \%16,5 → \%17,2; 2. \\
4 AB tarafı & AB'ye ihracat 2025; rejim payları (GB / İTÜ / tarım) & €838,6 mn; \%100 / \%0 / \%0 \\
& S4'te AB tarifesi: 1/98 sürerse / askıda / tam MFN ortalaması & \%8,0 / \%8,0 / \%8,0 \\
& Tepe satırlar (S4 artışı ≥ \%10) & yok \\
& AB pazarında AB dışı rakipler (AB tercihi) & Çin \%42 (MFN/GSP); Pakistan \%12 (GSP+); Birleşik Krallık \%7 (TCA); Hindistan \%5 (MFN/GSP); Güney Kore \%3 (STA/EPA) \\
& Tercihli rakiplerin payı; AME (puan; ham) & \%31; 87 (3,68) \\
& Menşe kuralı (PEM) & 50--55 dokuma kumaş: dokuma + boyama, baskı ya da kaplama; iplik boyama + dokuma; tek başına baskı (ithal iplik serbest) \\
& Menşe riski S2 / S3; c\_RoO (\%); p\_uyum & D/D; 1,5 / 2,0; 0,92 / 0,88 \etiketH{} \\
& Değer zincirindeki rol; DİR bağımlılığı & ara; orta \\
& EUDR ve SKDM kapsamındaki AB'ye ihracat payı & \%0; \%0 \\
& PEM ortaklarına ihracat; kümülasyon maruziyeti; κ & \$1.095 mn; \$768 mn; 1,0 \etiketH{} \\
5 Türkiye tarafı & Uygulanan tarife 2026 (ithalat ağırlıklı): MFN; İGV dahil & \%7,7; \%29,9 \\
& S4'te AB menşeline tarife (c2) & \%28,0 \\
& Kritik girdiler: 2024 ithalatı, AB payı, ilk kaynaklar & HS4 5205, 5402, 5509: \$1.842 mn (AB \%7; Çin \%34, Hindistan \%17, Özbekistan \%16) \\
6 Kabiliyet & KS; sıra (ağırlık duyarlılığında P10--P90) & 54,8; 66 / 230 (21--123) \\
& KS bileşenleri (yüzdelik): RCA, PCI, yurt içi KD, Ar-Ge, tedarikçi, verimlilik, firma & 75, 24, 66, 25, 96, 58, 59 \\
& PCI (sınıf) & −0,52 \\
& Atlanabilir ürünler (yoğunluk üst çeyrek, PCI \textgreater{} sınıf) & yok \\
7 Stratejik önem & SS; en yüksek boyut; ağırlıklı ortalama & 7,6; savunma 10; 2,0 \etiketH{} \\
8 Makro duyarlılık & TÜFE ağırlığı (doğrudan); tüketimde ithal pay & \%0,00; \%17 \\
& ÜFE vekili (imalat üretimindeki pay) & \%1,34 \\
& EDE, 1 puan tarife: tüm menşeler / yalnız AB (bp; ρ = 1) & 0,074 / 0,023 \\
& CKE: tarife (\$ mn / bp); mali (\$ mn / mlr TL); puan & 3.101,1; 13,26; 50 \\
9 Senaryolar & Senaryo tablosu & Tablo 5.2 \\
\end{longtable}

\endgroup

\hypertarget{nace_kartlar_13-15-c3-1330-tekstil-uxfcruxfcnlerinin-bitirilmesi}{%
\subsection{C.3 13.30 Tekstil ürünlerinin bitirilmesi}\label{nace_kartlar_13-15-c3-1330-tekstil-uxfcruxfcnlerinin-bitirilmesi}}

\Needspace{7\baselineskip}

\textbf{Tablo C.3 --- 13.30: veri sayfası}

\nopagebreak[4]

\begingroup\gbtablo\small

\begin{longtable}[]{@{}
  >{\raggedright\arraybackslash}p{(\columnwidth - 4\tabcolsep) * \real{0.1377}}
  >{\raggedright\arraybackslash}p{(\columnwidth - 4\tabcolsep) * \real{0.3799}}
  >{\raggedright\arraybackslash}p{(\columnwidth - 4\tabcolsep) * \real{0.4823}}@{}}
\toprule\noalign{}
\begin{minipage}[b]{\linewidth}\raggedright
\textbf{Alan}
\end{minipage} & \begin{minipage}[b]{\linewidth}\raggedright
\textbf{Gösterge}
\end{minipage} & \begin{minipage}[b]{\linewidth}\raggedright
\textbf{Değer}
\end{minipage} \\
\midrule\noalign{}
\endhead
\bottomrule\noalign{}
\endlastfoot
1 Kimlik & Rev. 2.1 karşılığı; CPA & 13.30 (n:1); CPA 13.30 \\
& Başlıca HS4 (ihracat, 2024) & --- \\
& BEC payları: ara / sermaye / tüketim & --- \\
2 Ölçek & Üretim; katma değer (GSYH payı) & \$5.593 mn; \$1.738 mn (\%0,13) \\
& Çalışan; girişim & 88.781; 5.596 \\
& Girişim ölçeği: 0 / 1--9 / 10--49 / 50--249 / 250+ çalışan & 216 / 5.021 / 610 / 281 / 52 \\
& 250+ girişimlerin istihdam payı; kişi başı KD (AB27 = 1) & \%34; 0,44 \\
& İstihdam değişimi, 2015--2024 & \%18 \\
& İlk üç bölge (İBBS-2) & İstanbul \%23; Bursa, Eskişehir, Bilecik \%19; Tekirdağ, Edirne, Kırklareli \%11 \\
& Çapa yatırımcılar & İstanbul, Bursa ve Çorlu'daki fason boya, baskı ve apre işletmeleri; büyük çoğunluğu KOBİ \etiketH{} \\
3 Dış ticaret & İhracat 2024; 2020--2024 değişimi & \$--- mn; --- \\
& AB payı; ilk beş AB pazarı & ---; --- \\
& İlk beş AB dışı pazar & --- \\
& İthalat 2024; AB ve Çin payı & \$--- mn; AB ---, Çin --- \\
& İlk beş ithalat kaynağı & --- \\
& Dış ticaret dengesi & +\$--- mn \\
& İthal girdi içeriği (I-O); KD ÷ üretim & \%23; \%31 \\
& RCA: dünya 2020 → 2024; AB pazarı 2024 & --- → ---; --- \\
& Birim değer ÷ AB pazarı ortalaması & --- \\
& AB pazarındaki pay (AB dışı tedarik içinde) 2020 → 2024; Türkiye'nin sırası & --- → ---; ---. \\
4 AB tarafı & AB'ye ihracat 2025; rejim payları (GB / İTÜ / tarım) & €--- mn; --- / --- / --- \\
& S4'te AB tarifesi: 1/98 sürerse / askıda / tam MFN ortalaması & --- / --- / --- \\
& Tepe satırlar (S4 artışı ≥ \%10) & yok \\
& AB pazarında AB dışı rakipler (AB tercihi) & --- \\
& Tercihli rakiplerin payı; AME (puan; ham) & ---; 5 (0,00) \\
& Menşe kuralı (PEM) & Hizmet (fason boya, baskı, apre); ticareti yok. Revize PEM'de tek başına baskı dokuma kumaşa menşe kazandırıyor \\
& Menşe riski S2 / S3; c\_RoO (\%); p\_uyum & D/D; 1,5 / 2,0; 0,90 / 0,85 \etiketH{} \\
& Değer zincirindeki rol; DİR bağımlılığı & ara; düşük \\
& EUDR ve SKDM kapsamındaki AB'ye ihracat payı & ---; --- \\
& PEM ortaklarına ihracat; kümülasyon maruziyeti; κ & \$0 mn; \$0 mn; 0,0 \etiketH{} \\
5 Türkiye tarafı & Uygulanan tarife 2026 (ithalat ağırlıklı): MFN; İGV dahil & ---; --- \\
& S4'te AB menşeline tarife (c2) & \%0,0 \\
& Kritik girdiler: 2024 ithalatı, AB payı, ilk kaynaklar & --- \\
6 Kabiliyet & KS; sıra (ağırlık duyarlılığında P10--P90) & 57,7; 42 / 230 (8--119) \\
& KS bileşenleri (yüzdelik): RCA, PCI, yurt içi KD, Ar-Ge, tedarikçi, verimlilik, firma & ---, ---, 65, 25, 96, 45, 70 \\
& PCI (sınıf) & --- \\
& Atlanabilir ürünler (yoğunluk üst çeyrek, PCI \textgreater{} sınıf) & yok \\
7 Stratejik önem & SS; en yüksek boyut; ağırlıklı ortalama & 0,0; --- 0; 0,0 \etiketH{} \\
8 Makro duyarlılık & TÜFE ağırlığı (doğrudan); tüketimde ithal pay & \%0,00; \%0 \\
& ÜFE vekili (imalat üretimindeki pay) & \%0,84 \\
& EDE, 1 puan tarife: tüm menşeler / yalnız AB (bp; ρ = 1) & 0,000 / 0,000 \\
& CKE: tarife (\$ mn / bp); mali (\$ mn / mlr TL); puan & ---; 0,00; 5 \\
9 Senaryolar & Senaryo tablosu & Tablo C.3b (aşağıda) \\
\end{longtable}

\endgroup

\Needspace{7\baselineskip}

\textbf{Tablo C.3b --- 13.30 Bitirme (boya, baskı, apre): senaryo etkileri (T+5'te yıllık; P50, parantez içinde P10 / P90)}

\nopagebreak[4]

\begingroup\gbtablo\small

\begin{longtable}[]{@{}lrrr@{}}
\toprule\noalign{}
\textbf{Değişken} & \textbf{S2} & \textbf{S3} & \textbf{S4} \\
\midrule\noalign{}
\endhead
\bottomrule\noalign{}
\endlastfoot
AB'ye ihracat (\$ mn) & 0 (0 / 0) & 0 (0 / 0) & 0 (0 / 0) \\
AB'ye ihracat (\%) & --- & --- & --- \\
Toplam ihracat (\$ mn) & 0 (0 / 0) & 0 (0 / 0) & 0 (0 / 0) \\
Nihai mal ithalatı (\$ mn) & 0 (0 / 0) & 0 (0 / 0) & 0 (0 / 0) \\
Üretim (\$ mn) & 0 (0 / 0) & 0 (0 / 0) & 0 (0 / 0) \\
İstihdam (kişi) & 0 (0 / 0) & 0 (0 / 0) & 0 (0 / 0) \\
İstihdam (\%) & 0,0 (0,0 / 0,0) & 0,0 (0,0 / 0,0) & 0,0 (0,0 / 0,0) \\
Cari denge (\$ mn) & 0 (0 / 0) & 0 (0 / 0) & 0 (0 / 0) \\
TÜFE (bp) & 0,00 (0,00 / 0,00) & 0,00 (0,00 / 0,00) & 0,00 (0,00 / 0,00) \\
\end{longtable}

\endgroup

\emph{Kaynak:} \etiketT{} f3e; toplamlar Faz 2B genel denge dağılımına kalibre. Yüzdeler 2024 düzeyine göre.

\hypertarget{nace_kartlar_13-15-c4-1391-uxf6rguxfc-triko-veya-tux131ux11f-iux15fi-kroux15fe-kumaux15flarux131n-imalatux131}{%
\subsection{C.4 13.91 Örgü (triko) veya tığ işi (kroşe) kumaşların imalatı}\label{nace_kartlar_13-15-c4-1391-uxf6rguxfc-triko-veya-tux131ux11f-iux15fi-kroux15fe-kumaux15flarux131n-imalatux131}}

\Needspace{7\baselineskip}

\textbf{Tablo C.4 --- 13.91: veri sayfası}

\nopagebreak[4]

\begingroup\gbtablo\small

\begin{longtable}[]{@{}
  >{\raggedright\arraybackslash}p{(\columnwidth - 4\tabcolsep) * \real{0.1348}}
  >{\raggedright\arraybackslash}p{(\columnwidth - 4\tabcolsep) * \real{0.3511}}
  >{\raggedright\arraybackslash}p{(\columnwidth - 4\tabcolsep) * \real{0.5140}}@{}}
\toprule\noalign{}
\begin{minipage}[b]{\linewidth}\raggedright
\textbf{Alan}
\end{minipage} & \begin{minipage}[b]{\linewidth}\raggedright
\textbf{Gösterge}
\end{minipage} & \begin{minipage}[b]{\linewidth}\raggedright
\textbf{Değer}
\end{minipage} \\
\midrule\noalign{}
\endhead
\bottomrule\noalign{}
\endlastfoot
1 Kimlik & Rev. 2.1 karşılığı; CPA & 13.91 (1:1); CPA 13.91 \\
& Başlıca HS4 (ihracat, 2024) & 6006 \%65, 6004 \%24, 6005 \%5, 6001 \%4 \\
& BEC payları: ara / sermaye / tüketim & \%100 / \%0 / \%0 \\
2 Ölçek & Üretim; katma değer (GSYH payı) & \$4.148 mn; \$1.060 mn (\%0,08) \\
& Çalışan; girişim & 36.092; 1.408 \\
& Girişim ölçeği: 0 / 1--9 / 10--49 / 50--249 / 250+ çalışan & 76 / 1.036 / 274 / 78 / 36 \\
& 250+ girişimlerin istihdam payı; kişi başı KD (AB27 = 1) & \%52; 0,55 \\
& İstihdam değişimi, 2015--2024 & \%29 \\
& İlk üç bölge (İBBS-2) & İstanbul \%31; Bursa, Eskişehir, Bilecik \%23; Hatay, Kahramanmaraş, Osmaniye \%14 \\
& Çapa yatırımcılar & İstanbul, Bursa ve Kahramanmaraş örme kumaş üreticileri \etiketH{} \\
3 Dış ticaret & İhracat 2024; 2020--2024 değişimi & \$1.806 mn; \%18 \\
& AB payı; ilk beş AB pazarı & \%31; İtalya \%34; Bulgaristan \%14; Polonya \%8; Yunanistan \%8; Romanya \%7 \\
& İlk beş AB dışı pazar & Belarus \%22; Rusya \%12; Fas \%12; Sırbistan \%6; Ukrayna \%6 \\
& İthalat 2024; AB ve Çin payı & \$487 mn; AB \%20, Çin \%40 \\
& İlk beş ithalat kaynağı & Çin \%40; Malezya \%13; Güney Kore \%13; Mısır \%10; İtalya \%8 \\
& Dış ticaret dengesi & +\$1.319 mn \\
& İthal girdi içeriği (I-O); KD ÷ üretim & \%23; \%26 \\
& RCA: dünya 2020 → 2024; AB pazarı 2024 & 4,84 → 4,13; 8,01 \\
& Birim değer ÷ AB pazarı ortalaması & 1,10 \\
& AB pazarındaki pay (AB dışı tedarik içinde) 2020 → 2024; Türkiye'nin sırası & \%39,3 → \%31,9; 2. \\
4 AB tarafı & AB'ye ihracat 2025; rejim payları (GB / İTÜ / tarım) & €437,5 mn; \%100 / \%0 / \%0 \\
& S4'te AB tarifesi: 1/98 sürerse / askıda / tam MFN ortalaması & \%8,0 / \%8,0 / \%8,0 \\
& Tepe satırlar (S4 artışı ≥ \%10) & yok \\
& AB pazarında AB dışı rakipler (AB tercihi) & Çin \%44 (MFN/GSP); Birleşik Krallık \%4 (TCA); Güney Kore \%4 (STA/EPA); İsviçre \%2 (AEA/STA); Sırbistan \%2 (MFN/GSP) \\
& Tercihli rakiplerin payı; AME (puan; ham) & \%18; 85 (3,59) \\
& Menşe kuralı (PEM) & Fasıl 60: örme + boyama, baskı, kaplama ya da laminasyon; iplik boyama + örme (ithal iplik serbest) \\
& Menşe riski S2 / S3; c\_RoO (\%); p\_uyum & D/D; 1,5 / 2,0; 0,92 / 0,88 \etiketH{} \\
& Değer zincirindeki rol; DİR bağımlılığı & ara; orta \\
& EUDR ve SKDM kapsamındaki AB'ye ihracat payı & \%0; \%0 \\
& PEM ortaklarına ihracat; kümülasyon maruziyeti; κ & \$495 mn; \$334 mn; 1,0 \etiketH{} \\
5 Türkiye tarafı & Uygulanan tarife 2026 (ithalat ağırlıklı): MFN; İGV dahil & \%8,0; \%35,0 \\
& S4'te AB menşeline tarife (c2) & \%34,7 \\
& Kritik girdiler: 2024 ithalatı, AB payı, ilk kaynaklar & HS4 5205, 5402: \$1.653 mn (AB \%7; Çin \%36, Özbekistan \%18, Hindistan \%15) \\
6 Kabiliyet & KS; sıra (ağırlık duyarlılığında P10--P90) & 50,8; 108 / 230 (46--157) \\
& KS bileşenleri (yüzdelik): RCA, PCI, yurt içi KD, Ar-Ge, tedarikçi, verimlilik, firma & 70, 17, 52, 25, 96, 69, 54 \\
& PCI (sınıf) & −0,80 \\
& Atlanabilir ürünler (yoğunluk üst çeyrek, PCI \textgreater{} sınıf) & yok \\
7 Stratejik önem & SS; en yüksek boyut; ağırlıklı ortalama & 0,0; --- 0; 0,0 \etiketH{} \\
8 Makro duyarlılık & TÜFE ağırlığı (doğrudan); tüketimde ithal pay & \%0,00; \%12 \\
& ÜFE vekili (imalat üretimindeki pay) & \%0,62 \\
& EDE, 1 puan tarife: tüm menşeler / yalnız AB (bp; ρ = 1) & 0,027 / 0,005 \\
& CKE: tarife (\$ mn / bp); mali (\$ mn / mlr TL); puan & 3.039,6; 12,30; 49 \\
9 Senaryolar & Senaryo tablosu & Tablo 5.3 \\
\end{longtable}

\endgroup

\hypertarget{nace_kartlar_13-15-c5-1392-giyim-eux15fyasux131-dux131ux15fux131ndaki-tamamlanmux131ux15f-tekstil-uxfcruxfcnlerinin-imalatux131}{%
\subsection{C.5 13.92 Giyim eşyası dışındaki tamamlanmış tekstil ürünlerinin imalatı}\label{nace_kartlar_13-15-c5-1392-giyim-eux15fyasux131-dux131ux15fux131ndaki-tamamlanmux131ux15f-tekstil-uxfcruxfcnlerinin-imalatux131}}

\Needspace{7\baselineskip}

\textbf{Tablo C.5 --- 13.92: veri sayfası}

\nopagebreak[4]

\begingroup\gbtablo\small

\begin{longtable}[]{@{}
  >{\raggedright\arraybackslash}p{(\columnwidth - 4\tabcolsep) * \real{0.1351}}
  >{\raggedright\arraybackslash}p{(\columnwidth - 4\tabcolsep) * \real{0.3537}}
  >{\raggedright\arraybackslash}p{(\columnwidth - 4\tabcolsep) * \real{0.5112}}@{}}
\toprule\noalign{}
\begin{minipage}[b]{\linewidth}\raggedright
\textbf{Alan}
\end{minipage} & \begin{minipage}[b]{\linewidth}\raggedright
\textbf{Gösterge}
\end{minipage} & \begin{minipage}[b]{\linewidth}\raggedright
\textbf{Değer}
\end{minipage} \\
\midrule\noalign{}
\endhead
\bottomrule\noalign{}
\endlastfoot
1 Kimlik & Rev. 2.1 karşılığı; CPA & 13.92, 13.96 (1:2); CPA 13.92 \\
& Başlıca HS4 (ihracat, 2024) & 6302 \%48, 6305 \%15, 6307 \%11, 6304 \%8 \\
& BEC payları: ara / sermaye / tüketim & \%15 / \%0 / \%85 \\
2 Ölçek & Üretim; katma değer (GSYH payı) & \$5.544 mn; \$1.580 mn (\%0,12) \\
& Çalışan; girişim & 87.344; 9.199 \\
& Girişim ölçeği: 0 / 1--9 / 10--49 / 50--249 / 250+ çalışan & 252 / 8.760 / 687 / 183 / 47 \\
& 250+ girişimlerin istihdam payı; kişi başı KD (AB27 = 1) & \%30; 0,51 \\
& İstihdam değişimi, 2015--2024 & \%12 \\
& İlk üç bölge (İBBS-2) & Aydın, Denizli, Muğla \%18; Gaziantep, Adıyaman, Kilis \%15; Bursa, Eskişehir, Bilecik \%10 \\
& Çapa yatırımcılar & Zorlu Tekstil (Taç), Özdilek; Denizli, Bursa ve Uşak havlu ve nevresim üreticileri \etiketH{} \\
3 Dış ticaret & İhracat 2024; 2020--2024 değişimi & \$2.346 mn; −\%12 \\
& AB payı; ilk beş AB pazarı & \%54; Almanya \%33; Fransa \%13; İtalya \%8; Hollanda \%8; İspanya \%5 \\
& İlk beş AB dışı pazar & ABD \%27; Birleşik Krallık \%8; İsviçre \%5; Gürcistan \%4; Irak \%4 \\
& İthalat 2024; AB ve Çin payı & \$261 mn; AB \%34, Çin \%40 \\
& İlk beş ithalat kaynağı & Çin \%40; Almanya \%7; Bulgaristan \%6; Mısır \%4; İtalya \%4 \\
& Dış ticaret dengesi & +\$2.085 mn \\
& İthal girdi içeriği (I-O); KD ÷ üretim & \%23; \%28 \\
& RCA: dünya 2020 → 2024; AB pazarı 2024 & 1,94 → 2,50; 2,06 \\
& Birim değer ÷ AB pazarı ortalaması & 1,44 \\
& AB pazarındaki pay (AB dışı tedarik içinde) 2020 → 2024; Türkiye'nin sırası & \%4,4 → \%8,2; 4. \\
4 AB tarafı & AB'ye ihracat 2025; rejim payları (GB / İTÜ / tarım) & €1.068,1 mn; \%100 / \%0 / \%0 \\
& S4'te AB tarifesi: 1/98 sürerse / askıda / tam MFN ortalaması & \%10,0 / \%10,0 / \%10,0 \\
& Tepe satırlar (S4 artışı ≥ \%10) & 52 satır, AB'ye ihracatın \%63'i: KN 6302 60 00 \%12,0 (€224,5 mn); KN 6302 21 00 \%12,0 (€89,7 mn); KN 6302 10 00 \%12,0 (€85,5 mn) \\
& AB pazarında AB dışı rakipler (AB tercihi) & Çin \%44 (MFN/GSP); Pakistan \%18 (GSP+); Hindistan \%9 (MFN/GSP); Bangladeş \%3 (EBA); Vietnam \%3 (STA/EPA) \\
& Tercihli rakiplerin payı; AME (puan; ham) & \%34; 96 (6,96) \\
& Menşe kuralı (PEM) & 6301--6304: dokuma ya da örme + konfeksiyon (çift dönüşüm); 6307: menşesiz malzeme ≤ \%40 \\
& Menşe riski S2 / S3; c\_RoO (\%); p\_uyum & O/O; 2,0 / 2,5; 0,88 / 0,85 \etiketH{} \\
& Değer zincirindeki rol; DİR bağımlılığı & nihai; orta \\
& EUDR ve SKDM kapsamındaki AB'ye ihracat payı & \%0; \%0 \\
& PEM ortaklarına ihracat; kümülasyon maruziyeti; κ & \$260 mn; \$167 mn; 0,0 \etiketH{} \\
5 Türkiye tarafı & Uygulanan tarife 2026 (ithalat ağırlıklı): MFN; İGV dahil & \%8,5; \%31,4 \\
& S4'te AB menşeline tarife (c2) & \%30,7 \\
& Kritik girdiler: 2024 ithalatı, AB payı, ilk kaynaklar & HS4 5205, 5208, 5209: \$735 mn (AB \%8; Özbekistan \%41, Pakistan \%13, Türkmenistan \%11) \\
6 Kabiliyet & KS; sıra (ağırlık duyarlılığında P10--P90) & 52,9; 81 / 230 (21--154) \\
& KS bileşenleri (yüzdelik): RCA, PCI, yurt içi KD, Ar-Ge, tedarikçi, verimlilik, firma & 81, 6, 58, 25, 96, 63, 72 \\
& PCI (sınıf) & −1,30 \\
& Atlanabilir ürünler (yoğunluk üst çeyrek, PCI \textgreater{} sınıf) & yok \\
7 Stratejik önem & SS; en yüksek boyut; ağırlıklı ortalama & 0,0; --- 0; 0,0 \etiketH{} \\
8 Makro duyarlılık & TÜFE ağırlığı (doğrudan); tüketimde ithal pay & \%0,69; \%6 \\
& ÜFE vekili (imalat üretimindeki pay) & \%0,83 \\
& EDE, 1 puan tarife: tüm menşeler / yalnız AB (bp; ρ = 1) & 0,053 / 0,018 \\
& CKE: tarife (\$ mn / bp); mali (\$ mn / mlr TL); puan & 1.127,5; 6,67; 31 \\
9 Senaryolar & Senaryo tablosu & Tablo 5.4 \\
\end{longtable}

\endgroup

\hypertarget{nace_kartlar_13-15-c6-1393-halux131-ve-kilim-imalatux131}{%
\subsection{C.6 13.93 Halı ve kilim imalatı}\label{nace_kartlar_13-15-c6-1393-halux131-ve-kilim-imalatux131}}

\Needspace{7\baselineskip}

\textbf{Tablo C.6 --- 13.93: veri sayfası}

\nopagebreak[4]

\begingroup\gbtablo\small

\begin{longtable}[]{@{}
  >{\raggedright\arraybackslash}p{(\columnwidth - 4\tabcolsep) * \real{0.1358}}
  >{\raggedright\arraybackslash}p{(\columnwidth - 4\tabcolsep) * \real{0.3606}}
  >{\raggedright\arraybackslash}p{(\columnwidth - 4\tabcolsep) * \real{0.5036}}@{}}
\toprule\noalign{}
\begin{minipage}[b]{\linewidth}\raggedright
\textbf{Alan}
\end{minipage} & \begin{minipage}[b]{\linewidth}\raggedright
\textbf{Gösterge}
\end{minipage} & \begin{minipage}[b]{\linewidth}\raggedright
\textbf{Değer}
\end{minipage} \\
\midrule\noalign{}
\endhead
\bottomrule\noalign{}
\endlastfoot
1 Kimlik & Rev. 2.1 karşılığı; CPA & 13.93 (1:1); CPA 13.93 \\
& Başlıca HS4 (ihracat, 2024) & 5702 \%75, 5703 \%20, 5705 \%3, 5701 \%2 \\
& BEC payları: ara / sermaye / tüketim & \%0 / \%0 / \%100 \\
2 Ölçek & Üretim; katma değer (GSYH payı) & \$4.165 mn; \$1.493 mn (\%0,11) \\
& Çalışan; girişim & 46.903; 1.143 \\
& Girişim ölçeği: 0 / 1--9 / 10--49 / 50--249 / 250+ çalışan & 82 / 847 / 197 / 79 / 49 \\
& 250+ girişimlerin istihdam payı; kişi başı KD (AB27 = 1) & \%67; 0,50 \\
& İstihdam değişimi, 2015--2024 & \%47 \\
& İlk üç bölge (İBBS-2) & Gaziantep, Adıyaman, Kilis \%86; Manisa, Afyon, Kütahya, Uşak \%3; Bursa, Eskişehir, Bilecik \%2 \\
& Çapa yatırımcılar & Merinos, Royal Halı, Atlas Halı, Dinarsu, Sanko Halı (Gaziantep kümesi) \etiketH{} \\
3 Dış ticaret & İhracat 2024; 2020--2024 değişimi & \$3.056 mn; \%15 \\
& AB payı; ilk beş AB pazarı & \%16; Almanya \%30; Polonya \%8; Fransa \%7; Romanya \%7; Belçika \%7 \\
& İlk beş AB dışı pazar & ABD \%33; Suudi Arabistan \%13; Irak \%10; Birleşik Krallık \%8; BAE \%4 \\
& İthalat 2024; AB ve Çin payı & \$93 mn; AB \%49, Çin \%15 \\
& İlk beş ithalat kaynağı & Hollanda \%24; Çin \%15; Hindistan \%12; Belçika \%6; Birleşik Krallık \%5 \\
& Dış ticaret dengesi & +\$2.963 mn \\
& İthal girdi içeriği (I-O); KD ÷ üretim & \%23; \%36 \\
& RCA: dünya 2020 → 2024; AB pazarı 2024 & 17,45 → 15,79; 5,46 \\
& Birim değer ÷ AB pazarı ortalaması & 0,96 \\
& AB pazarındaki pay (AB dışı tedarik içinde) 2020 → 2024; Türkiye'nin sırası & \%22,8 → \%21,8; 2. \\
4 AB tarafı & AB'ye ihracat 2025; rejim payları (GB / İTÜ / tarım) & €382,9 mn; \%100 / \%0 / \%0 \\
& S4'te AB tarifesi: 1/98 sürerse / askıda / tam MFN ortalaması & \%7,9 / \%7,9 / \%7,9 \\
& Tepe satırlar (S4 artışı ≥ \%10) & yok \\
& AB pazarında AB dışı rakipler (AB tercihi) & Çin \%25 (MFN/GSP); Hindistan \%21 (MFN/GSP); Birleşik Krallık \%9 (TCA); Mısır \%6 (STA/EPA); Vietnam \%4 (STA/EPA) \\
& Tercihli rakiplerin payı; AME (puan; ham) & \%27; 60 (1,70) \\
& Menşe kuralı (PEM) & Fasıl 57: eğirme ya da filament çekimi + dokuma veya tufting; tufting + boyama ya da baskı; jüt astar serbest \\
& Menşe riski S2 / S3; c\_RoO (\%); p\_uyum & D/O; 2,0 / 2,5; 0,88 / 0,85 \etiketH{} \\
& Değer zincirindeki rol; DİR bağımlılığı & nihai; orta \\
& EUDR ve SKDM kapsamındaki AB'ye ihracat payı & \%0; \%0 \\
& PEM ortaklarına ihracat; kümülasyon maruziyeti; κ & \$114 mn; \$77 mn; 0,0 \etiketH{} \\
5 Türkiye tarafı & Uygulanan tarife 2026 (ithalat ağırlıklı): MFN; İGV dahil & \%4,6; \%17,7 \\
& S4'te AB menşeline tarife (c2) & \%11,1 \\
& Kritik girdiler: 2024 ithalatı, AB payı, ilk kaynaklar & HS4 3902, 5402, 5503: \$4.446 mn (AB \%19; Suudi Arabistan \%20, Çin \%16, Güney Kore \%11) \\
6 Kabiliyet & KS; sıra (ağırlık duyarlılığında P10--P90) & 53,1; 77 / 230 (24--146) \\
& KS bileşenleri (yüzdelik): RCA, PCI, yurt içi KD, Ar-Ge, tedarikçi, verimlilik, firma & 76, 14, 75, 25, 96, 59, 47 \\
& PCI (sınıf) & −0,89 \\
& Atlanabilir ürünler (yoğunluk üst çeyrek, PCI \textgreater{} sınıf) & yok \\
7 Stratejik önem & SS; en yüksek boyut; ağırlıklı ortalama & 0,0; --- 0; 0,0 \etiketH{} \\
8 Makro duyarlılık & TÜFE ağırlığı (doğrudan); tüketimde ithal pay & \%0,37; \%3 \\
& ÜFE vekili (imalat üretimindeki pay) & \%0,62 \\
& EDE, 1 puan tarife: tüm menşeler / yalnız AB (bp; ρ = 1) & 0,018 / 0,009 \\
& CKE: tarife (\$ mn / bp); mali (\$ mn / mlr TL); puan & 555,8; 3,02; 23 \\
9 Senaryolar & Senaryo tablosu & Tablo C.6b (aşağıda) \\
\end{longtable}

\endgroup

\Needspace{7\baselineskip}

\textbf{Tablo C.6b --- 13.93 Halı: senaryo etkileri (T+5'te yıllık; P50, parantez içinde P10 / P90)}

\nopagebreak[4]

\begingroup\gbtablo\small

\begin{longtable}[]{@{}lrrr@{}}
\toprule\noalign{}
\textbf{Değişken} & \textbf{S2} & \textbf{S3} & \textbf{S4} \\
\midrule\noalign{}
\endhead
\bottomrule\noalign{}
\endlastfoot
AB'ye ihracat (\$ mn) & −27 (−66 / +10) & −34 (−66 / −6) & −60 (−101 / −25) \\
AB'ye ihracat (\%) & −5,7 (−13,8 / +2,0) & −7,0 (−13,9 / −1,3) & −12,7 (−21,1 / −5,3) \\
Toplam ihracat (\$ mn) & −20 (−48 / +6) & −24 (−48 / −4) & −48 (−82 / −20) \\
Nihai mal ithalatı (\$ mn) & −1 (−4 / +3) & −2 (−4 / +2) & −4 (−8 / −1) \\
Üretim (\$ mn) & −19 (−44 / +3) & −22 (−43 / −6) & −44 (−74 / −19) \\
İstihdam (kişi) & −146 (−344 / +27) & −172 (−341 / −47) & −340 (−579 / −145) \\
İstihdam (\%) & −0,3 (−0,7 / +0,1) & −0,4 (−0,7 / −0,1) & −0,7 (−1,2 / −0,3) \\
Cari denge (\$ mn) & −14 (−34 / +3) & −17 (−34 / −5) & −34 (−57 / −15) \\
TÜFE (bp) & +0,01 (+0,01 / +0,02) & +0,03 (+0,02 / +0,03) & +0,08 (+0,07 / +0,10) \\
\end{longtable}

\endgroup

\emph{Kaynak:} \etiketT{} f3e; toplamlar Faz 2B genel denge dağılımına kalibre. Yüzdeler 2024 düzeyine göre.

\hypertarget{nace_kartlar_13-15-c7-1394-halat-urgan-kux131nnap-ve-aux11f-imalatux131}{%
\subsection{C.7 13.94 Halat, urgan, kınnap ve ağ imalatı}\label{nace_kartlar_13-15-c7-1394-halat-urgan-kux131nnap-ve-aux11f-imalatux131}}

\Needspace{7\baselineskip}

\textbf{Tablo C.7 --- 13.94: veri sayfası}

\nopagebreak[4]

\begingroup\gbtablo\small

\begin{longtable}[]{@{}
  >{\raggedright\arraybackslash}p{(\columnwidth - 4\tabcolsep) * \real{0.1357}}
  >{\raggedright\arraybackslash}p{(\columnwidth - 4\tabcolsep) * \real{0.3598}}
  >{\raggedright\arraybackslash}p{(\columnwidth - 4\tabcolsep) * \real{0.5045}}@{}}
\toprule\noalign{}
\begin{minipage}[b]{\linewidth}\raggedright
\textbf{Alan}
\end{minipage} & \begin{minipage}[b]{\linewidth}\raggedright
\textbf{Gösterge}
\end{minipage} & \begin{minipage}[b]{\linewidth}\raggedright
\textbf{Değer}
\end{minipage} \\
\midrule\noalign{}
\endhead
\bottomrule\noalign{}
\endlastfoot
1 Kimlik & Rev. 2.1 karşılığı; CPA & 13.94 (1:1); CPA 13.94 \\
& Başlıca HS4 (ihracat, 2024) & 5607 \%46, 5608 \%31, 5609 \%12, 6310 \%10 \\
& BEC payları: ara / sermaye / tüketim & \%100 / \%0 / \%0 \\
2 Ölçek & Üretim; katma değer (GSYH payı) & \$144 mn; \$47 mn (\%0,00) \\
& Çalışan; girişim & 2.223; 270 \\
& Girişim ölçeği: 0 / 1--9 / 10--49 / 50--249 / 250+ çalışan & 14 / 231 / 30 / 10 / 0 \\
& 250+ girişimlerin istihdam payı; kişi başı KD (AB27 = 1) & \%0; 0,45 \\
& İstihdam değişimi, 2015--2024 & \%41 \\
& İlk üç bölge (İBBS-2) & Balıkesir, Çanakkale \%22; İstanbul \%18; Manisa, Afyon, Kütahya, Uşak \%10 \\
& Çapa yatırımcılar & Balıkesir ve İstanbul'daki halat ve balık ağı üreticileri \etiketH{} \\
3 Dış ticaret & İhracat 2024; 2020--2024 değişimi & \$99 mn; \%30 \\
& AB payı; ilk beş AB pazarı & \%38; İtalya \%25; Almanya \%15; Avusturya \%8; Romanya \%7; Hollanda \%7 \\
& İlk beş AB dışı pazar & ABD \%25; Mısır \%9; Birleşik Krallık \%7; Irak \%4; Kazakistan \%3 \\
& İthalat 2024; AB ve Çin payı & \$53 mn; AB \%32, Çin \%24 \\
& İlk beş ithalat kaynağı & Çin \%24; İtalya \%10; Hindistan \%8; Almanya \%5; Portekiz \%5 \\
& Dış ticaret dengesi & +\$46 mn \\
& İthal girdi içeriği (I-O); KD ÷ üretim & \%23; \%33 \\
& RCA: dünya 2020 → 2024; AB pazarı 2024 & 1,27 → 1,30; 1,32 \\
& Birim değer ÷ AB pazarı ortalaması & 0,85 \\
& AB pazarındaki pay (AB dışı tedarik içinde) 2020 → 2024; Türkiye'nin sırası & \%5,5 → \%5,3; 3. \\
4 AB tarafı & AB'ye ihracat 2025; rejim payları (GB / İTÜ / tarım) & €33,6 mn; \%100 / \%0 / \%0 \\
& S4'te AB tarifesi: 1/98 sürerse / askıda / tam MFN ortalaması & \%6,8 / \%6,8 / \%6,8 \\
& Tepe satırlar (S4 artışı ≥ \%10) & yok (AB'ye ihracatın \%1'inden az) \\
& AB pazarında AB dışı rakipler (AB tercihi) & Çin \%49 (MFN/GSP); Hindistan \%9 (MFN/GSP); ABD \%5 (MFN/GSP); Birleşik Krallık \%4 (TCA); Güney Kore \%2 (STA/EPA) \\
& Tercihli rakiplerin payı; AME (puan; ham) & \%24; 89 (3,92) \\
& Menşe kuralı (PEM) & Fasıl 56: eğirme ya da sentetik lif çekimi + eğirme (ithal iplikten halat menşe kazanmaz) \\
& Menşe riski S2 / S3; c\_RoO (\%); p\_uyum & O/O; 2,5 / 3,0; 0,85 / 0,80 \etiketH{} \\
& Değer zincirindeki rol; DİR bağımlılığı & ara; orta \\
& EUDR ve SKDM kapsamındaki AB'ye ihracat payı & \%0; \%0 \\
& PEM ortaklarına ihracat; kümülasyon maruziyeti; κ & \$15 mn; \$8 mn; 0,0 \etiketH{} \\
5 Türkiye tarafı & Uygulanan tarife 2026 (ithalat ağırlıklı): MFN; İGV dahil & \%6,3; \%18,4 \\
& S4'te AB menşeline tarife (c2) & \%18,8 \\
& Kritik girdiler: 2024 ithalatı, AB payı, ilk kaynaklar & --- \\
6 Kabiliyet & KS; sıra (ağırlık duyarlılığında P10--P90) & 47,8; 140 / 230 (75--176) \\
& KS bileşenleri (yüzdelik): RCA, PCI, yurt içi KD, Ar-Ge, tedarikçi, verimlilik, firma & 48, 18, 70, 25, 96, 47, 61 \\
& PCI (sınıf) & −0,68 \\
& Atlanabilir ürünler (yoğunluk üst çeyrek, PCI \textgreater{} sınıf) & yok \\
7 Stratejik önem & SS; en yüksek boyut; ağırlıklı ortalama & 7,5; gıda güvenliği 10; 1,5 \etiketH{} \\
8 Makro duyarlılık & TÜFE ağırlığı (doğrudan); tüketimde ithal pay & \%0,00; \%20 \\
& ÜFE vekili (imalat üretimindeki pay) & \%0,02 \\
& EDE, 1 puan tarife: tüm menşeler / yalnız AB (bp; ρ = 1) & 0,003 / 0,001 \\
& CKE: tarife (\$ mn / bp); mali (\$ mn / mlr TL); puan & 1.278,3; 24,46; 49 \\
9 Senaryolar & Senaryo tablosu & Tablo C.7b (aşağıda) \\
\end{longtable}

\endgroup

\Needspace{7\baselineskip}

\textbf{Tablo C.7b --- 13.94 Halat ve ağ: senaryo etkileri (T+5'te yıllık; P50, parantez içinde P10 / P90)}

\nopagebreak[4]

\begingroup\gbtablo\small

\begin{longtable}[]{@{}lrrr@{}}
\toprule\noalign{}
\textbf{Değişken} & \textbf{S2} & \textbf{S3} & \textbf{S4} \\
\midrule\noalign{}
\endhead
\bottomrule\noalign{}
\endlastfoot
AB'ye ihracat (\$ mn) & −3 (−7 / +1) & −3 (−7 / −1) & −5 (−8 / −2) \\
AB'ye ihracat (\%) & −7,4 (−18,1 / +2,0) & −8,8 (−17,6 / −1,7) & −13,0 (−21,5 / −5,6) \\
Toplam ihracat (\$ mn) & −2 (−5 / +1) & −2 (−5 / 0) & −4 (−7 / −2) \\
Nihai mal ithalatı (\$ mn) & 0 (−1 / +1) & −1 (−1 / 0) & −1 (−3 / 0) \\
Üretim (\$ mn) & −2 (−4 / 0) & −2 (−3 / −1) & −2 (−4 / −1) \\
İstihdam (kişi) & −17 (−41 / −2) & −19 (−37 / −9) & −25 (−43 / −14) \\
İstihdam (\%) & −0,8 (−1,8 / −0,1) & −0,8 (−1,7 / −0,4) & −1,1 (−1,9 / −0,6) \\
Cari denge (\$ mn) & −1 (−3 / 0) & −1 (−3 / −1) & −2 (−3 / −1) \\
TÜFE (bp) & 0,00 (0,00 / 0,00) & 0,00 (0,00 / 0,00) & +0,02 (+0,01 / +0,02) \\
\end{longtable}

\endgroup

\emph{Kaynak:} \etiketT{} f3e; toplamlar Faz 2B genel denge dağılımına kalibre. Yüzdeler 2024 düzeyine göre.

\hypertarget{nace_kartlar_13-15-c8-1395-dokusuz-kumaux15flarux131n-ve-dokusuz-kumaux15ftan-uxfcruxfcnlerin-imalatux131-giyim-eux15fyasux131-hariuxe7}{%
\subsection{C.8 13.95 Dokusuz kumaşların ve dokusuz kumaştan ürünlerin imalatı (giyim eşyası hariç)}\label{nace_kartlar_13-15-c8-1395-dokusuz-kumaux15flarux131n-ve-dokusuz-kumaux15ftan-uxfcruxfcnlerin-imalatux131-giyim-eux15fyasux131-hariuxe7}}

\Needspace{7\baselineskip}

\textbf{Tablo C.8 --- 13.95: veri sayfası}

\nopagebreak[4]

\begingroup\gbtablo\small

\begin{longtable}[]{@{}
  >{\raggedright\arraybackslash}p{(\columnwidth - 4\tabcolsep) * \real{0.1362}}
  >{\raggedright\arraybackslash}p{(\columnwidth - 4\tabcolsep) * \real{0.3653}}
  >{\raggedright\arraybackslash}p{(\columnwidth - 4\tabcolsep) * \real{0.4984}}@{}}
\toprule\noalign{}
\begin{minipage}[b]{\linewidth}\raggedright
\textbf{Alan}
\end{minipage} & \begin{minipage}[b]{\linewidth}\raggedright
\textbf{Gösterge}
\end{minipage} & \begin{minipage}[b]{\linewidth}\raggedright
\textbf{Değer}
\end{minipage} \\
\midrule\noalign{}
\endhead
\bottomrule\noalign{}
\endlastfoot
1 Kimlik & Rev. 2.1 karşılığı; CPA & 13.95 (1:1); CPA 13.95 \\
& Başlıca HS4 (ihracat, 2024) & 5603 \%100 \\
& BEC payları: ara / sermaye / tüketim & \%100 / \%0 / \%0 \\
2 Ölçek & Üretim; katma değer (GSYH payı) & \$1.536 mn; \$458 mn (\%0,03) \\
& Çalışan; girişim & 8.474; 112 \\
& Girişim ölçeği: 0 / 1--9 / 10--49 / 50--249 / 250+ çalışan & 4 / 82 / 10 / 18 / 10 \\
& 250+ girişimlerin istihdam payı; kişi başı KD (AB27 = 1) & \%72; --- \\
& İstihdam değişimi, 2015--2024 & \%39 \\
& İlk üç bölge (İBBS-2) & Gaziantep, Adıyaman, Kilis \%41; Tekirdağ, Edirne, Kırklareli \%34; Kayseri, Sivas, Yozgat \%11 \\
& Çapa yatırımcılar & Gaziantep ve Tekirdağ'daki spunbond ve spunlace üreticileri (Mogul, Gülsan) \etiketH{} \\
3 Dış ticaret & İhracat 2024; 2020--2024 değişimi & \$804 mn; \%12 \\
& AB payı; ilk beş AB pazarı & \%39; İtalya \%23; Çekya \%16; Almanya \%10; İspanya \%8; Yunanistan \%8 \\
& İlk beş AB dışı pazar & ABD \%19; Birleşik Krallık \%12; İran \%8; Mısır \%6; Cezayir \%6 \\
& İthalat 2024; AB ve Çin payı & \$209 mn; AB \%67, Çin \%11 \\
& İlk beş ithalat kaynağı & İtalya \%26; Almanya \%18; Çin \%11; Mısır \%10; ABD \%4 \\
& Dış ticaret dengesi & +\$595 mn \\
& İthal girdi içeriği (I-O); KD ÷ üretim & \%23; \%30 \\
& RCA: dünya 2020 → 2024; AB pazarı 2024 & 3,67 → 3,90; 4,55 \\
& Birim değer ÷ AB pazarı ortalaması & 0,72 \\
& AB pazarındaki pay (AB dışı tedarik içinde) 2020 → 2024; Türkiye'nin sırası & \%14,8 → \%18,1; 2. \\
4 AB tarafı & AB'ye ihracat 2025; rejim payları (GB / İTÜ / tarım) & €310,8 mn; \%100 / \%0 / \%0 \\
& S4'te AB tarifesi: 1/98 sürerse / askıda / tam MFN ortalaması & \%4,3 / \%4,3 / \%4,3 \\
& Tepe satırlar (S4 artışı ≥ \%10) & yok \\
& AB pazarında AB dışı rakipler (AB tercihi) & Çin \%27 (MFN/GSP); ABD \%16 (MFN/GSP); Japonya \%12 (STA/EPA); İsrail \%6 (STA/EPA); Birleşik Krallık \%6 (TCA) \\
& Tercihli rakiplerin payı; AME (puan; ham) & \%33; 76 (2,68) \\
& Menşe kuralı (PEM) & 5603: lif ya da polimerden bağlama (spunbond, meltblown; ithal polimer serbest) \\
& Menşe riski S2 / S3; c\_RoO (\%); p\_uyum & D/D; 1,5 / 2,0; 0,92 / 0,88 \etiketH{} \\
& Değer zincirindeki rol; DİR bağımlılığı & ara; düşük \\
& EUDR ve SKDM kapsamındaki AB'ye ihracat payı & \%0; \%0 \\
& PEM ortaklarına ihracat; kümülasyon maruziyeti; κ & \$117 mn; \$75 mn; 0,0 \etiketH{} \\
5 Türkiye tarafı & Uygulanan tarife 2026 (ithalat ağırlıklı): MFN; İGV dahil & \%4,3; \%26,3 \\
& S4'te AB menşeline tarife (c2) & \%26,3 \\
& Kritik girdiler: 2024 ithalatı, AB payı, ilk kaynaklar & HS4 3902, 5503: \$3.272 mn (AB \%22; Suudi Arabistan \%27, Güney Kore \%13, Rusya \%9) \\
6 Kabiliyet & KS; sıra (ağırlık duyarlılığında P10--P90) & 61,7; 19 / 230 (4--80) \\
& KS bileşenleri (yüzdelik): RCA, PCI, yurt içi KD, Ar-Ge, tedarikçi, verimlilik, firma & 79, 76, 61, 25, 96, ---, 20 \\
& PCI (sınıf) & 0,47 \\
& Atlanabilir ürünler (yoğunluk üst çeyrek, PCI \textgreater{} sınıf) & yok \\
7 Stratejik önem & SS; en yüksek boyut; ağırlıklı ortalama & 37,2; sağlık 50; 7,5 \etiketH{} \\
8 Makro duyarlılık & TÜFE ağırlığı (doğrudan); tüketimde ithal pay & \%0,00; \%12 \\
& ÜFE vekili (imalat üretimindeki pay) & \%0,23 \\
& EDE, 1 puan tarife: tüm menşeler / yalnız AB (bp; ρ = 1) & 0,012 / 0,008 \\
& CKE: tarife (\$ mn / bp); mali (\$ mn / mlr TL); puan & 4.027,6; 19,84; 59 \\
9 Senaryolar & Senaryo tablosu & Tablo C.8b (aşağıda) \\
\end{longtable}

\endgroup

\Needspace{7\baselineskip}

\textbf{Tablo C.8b --- 13.95 Dokusuz kumaş: senaryo etkileri (T+5'te yıllık; P50, parantez içinde P10 / P90)}

\nopagebreak[4]

\begingroup\gbtablo\small

\begin{longtable}[]{@{}lrrr@{}}
\toprule\noalign{}
\textbf{Değişken} & \textbf{S2} & \textbf{S3} & \textbf{S4} \\
\midrule\noalign{}
\endhead
\bottomrule\noalign{}
\endlastfoot
AB'ye ihracat (\$ mn) & −29 (−68 / +6) & −37 (−74 / −7) & −53 (−88 / −23) \\
AB'ye ihracat (\%) & −9,3 (−21,5 / +2,0) & −11,8 (−23,3 / −2,3) & −16,8 (−27,6 / −7,2) \\
Toplam ihracat (\$ mn) & −21 (−50 / +4) & −27 (−53 / −5) & −42 (−71 / −18) \\
Nihai mal ithalatı (\$ mn) & −11 (−37 / +9) & −16 (−40 / +5) & −33 (−60 / −7) \\
Üretim (\$ mn) & −9 (−17 / −1) & −11 (−18 / −3) & −10 (−18 / −2) \\
İstihdam (kişi) & −33 (−66 / −3) & −41 (−70 / −11) & −36 (−69 / −6) \\
İstihdam (\%) & −0,4 (−0,8 / 0,0) & −0,5 (−0,8 / −0,1) & −0,4 (−0,8 / −0,1) \\
Cari denge (\$ mn) & −7 (−13 / −1) & −8 (−14 / −2) & −7 (−14 / −1) \\
TÜFE (bp) & +0,02 (+0,02 / +0,02) & +0,04 (+0,03 / +0,04) & +0,17 (+0,15 / +0,20) \\
\end{longtable}

\endgroup

\emph{Kaynak:} \etiketT{} f3e; toplamlar Faz 2B genel denge dağılımına kalibre. Yüzdeler 2024 düzeyine göre.

\hypertarget{nace_kartlar_13-15-c9-1396-diux11fer-teknik-ve-enduxfcstriyel-tekstillerin-imalatux131}{%
\subsection{C.9 13.96 Diğer teknik ve endüstriyel tekstillerin imalatı}\label{nace_kartlar_13-15-c9-1396-diux11fer-teknik-ve-enduxfcstriyel-tekstillerin-imalatux131}}

\Needspace{7\baselineskip}

\textbf{Tablo C.9 --- 13.96: veri sayfası}

\nopagebreak[4]

\begingroup\gbtablo\small

\begin{longtable}[]{@{}
  >{\raggedright\arraybackslash}p{(\columnwidth - 4\tabcolsep) * \real{0.1363}}
  >{\raggedright\arraybackslash}p{(\columnwidth - 4\tabcolsep) * \real{0.3661}}
  >{\raggedright\arraybackslash}p{(\columnwidth - 4\tabcolsep) * \real{0.4975}}@{}}
\toprule\noalign{}
\begin{minipage}[b]{\linewidth}\raggedright
\textbf{Alan}
\end{minipage} & \begin{minipage}[b]{\linewidth}\raggedright
\textbf{Gösterge}
\end{minipage} & \begin{minipage}[b]{\linewidth}\raggedright
\textbf{Değer}
\end{minipage} \\
\midrule\noalign{}
\endhead
\bottomrule\noalign{}
\endlastfoot
1 Kimlik & Rev. 2.1 karşılığı; CPA & 13.96 (n:1); CPA 13.96 \\
& Başlıca HS4 (ihracat, 2024) & 5902 \%28, 5903 \%27, 5806 \%12, 5911 \%10 \\
& BEC payları: ara / sermaye / tüketim & \%100 / \%0 / \%0 \\
2 Ölçek & Üretim; katma değer (GSYH payı) & \$1.749 mn; \$472 mn (\%0,03) \\
& Çalışan; girişim & 19.387; 2.218 \\
& Girişim ölçeği: 0 / 1--9 / 10--49 / 50--249 / 250+ çalışan & 96 / 1.952 / 265 / 57 / 5 \\
& 250+ girişimlerin istihdam payı; kişi başı KD (AB27 = 1) & \%8; 0,34 \\
& İstihdam değişimi, 2015--2024 & \%6 \\
& İlk üç bölge (İBBS-2) & İstanbul \%40; Gaziantep, Adıyaman, Kilis \%11; Bursa, Eskişehir, Bilecik \%9 \\
& Çapa yatırımcılar & Kordsa (Sabancı; lastik kordu ve kompozit); İstanbul ve Bursa'daki balistik ve endüstriyel tekstil üreticileri \etiketH{} \\
3 Dış ticaret & İhracat 2024; 2020--2024 değişimi & \$528 mn; \%33 \\
& AB payı; ilk beş AB pazarı & \%42; Almanya \%18; İtalya \%13; Polonya \%9; Romanya \%8; Çekya \%5 \\
& İlk beş AB dışı pazar & ABD \%15; Birleşik Krallık \%12; Sırbistan \%9; Fas \%6; Mısır \%6 \\
& İthalat 2024; AB ve Çin payı & \$589 mn; AB \%44, Çin \%35 \\
& İlk beş ithalat kaynağı & Çin \%35; Almanya \%14; İtalya \%11; Malezya \%6; Romanya \%5 \\
& Dış ticaret dengesi & −\$60 mn \\
& İthal girdi içeriği (I-O); KD ÷ üretim & \%23; \%27 \\
& RCA: dünya 2020 → 2024; AB pazarı 2024 & 1,45 → 1,42; 2,32 \\
& Birim değer ÷ AB pazarı ortalaması & 0,97 \\
& AB pazarındaki pay (AB dışı tedarik içinde) 2020 → 2024; Türkiye'nin sırası & \%8,4 → \%9,2; 2. \\
4 AB tarafı & AB'ye ihracat 2025; rejim payları (GB / İTÜ / tarım) & €183,4 mn; \%100 / \%0 / \%0 \\
& S4'te AB tarifesi: 1/98 sürerse / askıda / tam MFN ortalaması & \%6,4 / \%6,4 / \%6,4 \\
& Tepe satırlar (S4 artışı ≥ \%10) & yok \\
& AB pazarında AB dışı rakipler (AB tercihi) & Çin \%43 (MFN/GSP); Birleşik Krallık \%8 (TCA); ABD \%7 (MFN/GSP); İsviçre \%6 (AEA/STA); Vietnam \%5 (STA/EPA) \\
& Tercihli rakiplerin payı; AME (puan; ham) & \%33; 90 (4,09) \\
& Menşe kuralı (PEM) & 5902--5911: dokuma + emdirme, kaplama ya da laminasyon; lastik kordu (\textgreater{} \%90 tekstil): lif çekimi + dokuma \\
& Menşe riski S2 / S3; c\_RoO (\%); p\_uyum & D/D; 1,5 / 2,0; 0,90 / 0,85 \etiketH{} \\
& Değer zincirindeki rol; DİR bağımlılığı & ara; orta \\
& EUDR ve SKDM kapsamındaki AB'ye ihracat payı & \%0; \%0 \\
& PEM ortaklarına ihracat; kümülasyon maruziyeti; κ & \$100 mn; \$63 mn; 0,3 \etiketH{} \\
5 Türkiye tarafı & Uygulanan tarife 2026 (ithalat ağırlıklı): MFN; İGV dahil & \%7,0; \%16,4 \\
& S4'te AB menşeline tarife (c2) & \%16,0 \\
& Kritik girdiler: 2024 ithalatı, AB payı, ilk kaynaklar & HS4 5402, 7019: \$1.518 mn (AB \%15; Çin \%48, Hindistan \%12, Güney Kore \%6) \\
6 Kabiliyet & KS; sıra (ağırlık duyarlılığında P10--P90) & 54,2; 67 / 230 (31--131) \\
& KS bileşenleri (yüzdelik): RCA, PCI, yurt içi KD, Ar-Ge, tedarikçi, verimlilik, firma & 48, 66, 55, 25, 96, 22, 75 \\
& PCI (sınıf) & 0,31 \\
& Atlanabilir ürünler (yoğunluk üst çeyrek, PCI \textgreater{} sınıf) & yok \\
7 Stratejik önem & SS; en yüksek boyut; ağırlıklı ortalama & 33,4; savunma 40; 18,0 \etiketH{} \\
8 Makro duyarlılık & TÜFE ağırlığı (doğrudan); tüketimde ithal pay & \%0,00; \%30 \\
& ÜFE vekili (imalat üretimindeki pay) & \%0,26 \\
& EDE, 1 puan tarife: tüm menşeler / yalnız AB (bp; ρ = 1) & 0,033 / 0,014 \\
& CKE: tarife (\$ mn / bp); mali (\$ mn / mlr TL); puan & 3.445,8; 28,95; 63 \\
9 Senaryolar & Senaryo tablosu & Tablo C.9b (aşağıda) \\
\end{longtable}

\endgroup

\Needspace{7\baselineskip}

\textbf{Tablo C.9b --- 13.96 Teknik tekstil: senaryo etkileri (T+5'te yıllık; P50, parantez içinde P10 / P90)}

\nopagebreak[4]

\begingroup\gbtablo\small

\begin{longtable}[]{@{}
  >{\raggedright\arraybackslash}p{(\columnwidth - 6\tabcolsep) * \real{0.3774}}
  >{\raggedleft\arraybackslash}p{(\columnwidth - 6\tabcolsep) * \real{0.2075}}
  >{\raggedleft\arraybackslash}p{(\columnwidth - 6\tabcolsep) * \real{0.2075}}
  >{\raggedleft\arraybackslash}p{(\columnwidth - 6\tabcolsep) * \real{0.2075}}@{}}
\toprule\noalign{}
\begin{minipage}[b]{\linewidth}\raggedright
\textbf{Değişken}
\end{minipage} & \begin{minipage}[b]{\linewidth}\raggedleft
\textbf{S2}
\end{minipage} & \begin{minipage}[b]{\linewidth}\raggedleft
\textbf{S3}
\end{minipage} & \begin{minipage}[b]{\linewidth}\raggedleft
\textbf{S4}
\end{minipage} \\
\midrule\noalign{}
\endhead
\bottomrule\noalign{}
\endlastfoot
AB'ye ihracat (\$ mn) & −22 (−51 / +4) & −28 (−54 / −6) & −48 (−77 / −20) \\
AB'ye ihracat (\%) & −10,0 (−23,0 / +2,0) & −12,8 (−24,5 / −2,5) & −21,6 (−34,7 / −9,2) \\
Toplam ihracat (\$ mn) & −16 (−37 / +3) & −20 (−39 / −4) & −49 (−74 / −27) \\
Bunun içinde PEM kümülasyonu (\$ mn) & 0 (0 / 0) & 0 (0 / 0) & −11 (−14 / −8) \\
Nihai mal ithalatı (\$ mn) & −14 (−44 / +15) & −19 (−49 / +9) & −43 (−78 / −9) \\
Üretim (\$ mn) & −4 (−12 / +10) & −2 (−13 / +13) & −7 (−18 / +8) \\
İstihdam (kişi) & −28 (−93 / +77) & −14 (−100 / +98) & −50 (−144 / +58) \\
İstihdam (\%) & −0,1 (−0,5 / +0,4) & −0,1 (−0,5 / +0,5) & −0,3 (−0,7 / +0,3) \\
Cari denge (\$ mn) & −3 (−9 / +8) & −1 (−10 / +10) & −5 (−14 / +6) \\
TÜFE (bp) & +0,03 (+0,02 / +0,03) & +0,05 (+0,04 / +0,06) & +0,20 (+0,17 / +0,22) \\
\end{longtable}

\endgroup

\emph{Kaynak:} \etiketT{} f3e; toplamlar Faz 2B genel denge dağılımına kalibre. Yüzdeler 2024 düzeyine göre.

\hypertarget{nace_kartlar_13-15-c10-1399-baux15fka-yerde-sux131nux131flandux131rux131lmamux131ux15f-diux11fer-tekstillerin-imalatux131}{%
\subsection{C.10 13.99 Başka yerde sınıflandırılmamış diğer tekstillerin imalatı}\label{nace_kartlar_13-15-c10-1399-baux15fka-yerde-sux131nux131flandux131rux131lmamux131ux15f-diux11fer-tekstillerin-imalatux131}}

\Needspace{7\baselineskip}

\textbf{Tablo C.10 --- 13.99: veri sayfası}

\nopagebreak[4]

\begingroup\gbtablo\small

\begin{longtable}[]{@{}
  >{\raggedright\arraybackslash}p{(\columnwidth - 4\tabcolsep) * \real{0.1356}}
  >{\raggedright\arraybackslash}p{(\columnwidth - 4\tabcolsep) * \real{0.3594}}
  >{\raggedright\arraybackslash}p{(\columnwidth - 4\tabcolsep) * \real{0.5049}}@{}}
\toprule\noalign{}
\begin{minipage}[b]{\linewidth}\raggedright
\textbf{Alan}
\end{minipage} & \begin{minipage}[b]{\linewidth}\raggedright
\textbf{Gösterge}
\end{minipage} & \begin{minipage}[b]{\linewidth}\raggedright
\textbf{Değer}
\end{minipage} \\
\midrule\noalign{}
\endhead
\bottomrule\noalign{}
\endlastfoot
1 Kimlik & Rev. 2.1 karşılığı; CPA & 13.99 (1:1); CPA 13.99 \\
& Başlıca HS4 (ihracat, 2024) & 5602 \%33, 5804 \%24, 5606 \%19, 5810 \%19 \\
& BEC payları: ara / sermaye / tüketim & \%99 / \%0 / \%1 \\
2 Ölçek & Üretim; katma değer (GSYH payı) & \$1.230 mn; \$436 mn (\%0,03) \\
& Çalışan; girişim & 21.343; 2.887 \\
& Girişim ölçeği: 0 / 1--9 / 10--49 / 50--249 / 250+ çalışan & 68 / 2.809 / 217 / 62 / 7 \\
& 250+ girişimlerin istihdam payı; kişi başı KD (AB27 = 1) & \%8; 0,58 \\
& İstihdam değişimi, 2015--2024 & −\%8 \\
& İlk üç bölge (İBBS-2) & Bursa, Eskişehir, Bilecik \%24; İstanbul \%19; Gaziantep, Adıyaman, Kilis \%16 \\
& Çapa yatırımcılar & Bursa ve İstanbul'daki şerit, etiket ve dantel üreticileri \etiketH{} \\
3 Dış ticaret & İhracat 2024; 2020--2024 değişimi & \$297 mn; \%38 \\
& AB payı; ilk beş AB pazarı & \%39; Romanya \%21; Almanya \%14; Polonya \%11; İspanya \%9; İtalya \%8 \\
& İlk beş AB dışı pazar & Belarus \%17; ABD \%14; Fas \%7; Birleşik Krallık \%5; Ukrayna \%4 \\
& İthalat 2024; AB ve Çin payı & \$119 mn; AB \%33, Çin \%53 \\
& İlk beş ithalat kaynağı & Çin \%53; İtalya \%10; Almanya \%7; Bulgaristan \%4; Hindistan \%4 \\
& Dış ticaret dengesi & +\$178 mn \\
& İthal girdi içeriği (I-O); KD ÷ üretim & \%23; \%35 \\
& RCA: dünya 2020 → 2024; AB pazarı 2024 & 4,24 → 4,00; 4,06 \\
& Birim değer ÷ AB pazarı ortalaması & 0,85 \\
& AB pazarındaki pay (AB dışı tedarik içinde) 2020 → 2024; Türkiye'nin sırası & \%16,7 → \%16,2; 2. \\
4 AB tarafı & AB'ye ihracat 2025; rejim payları (GB / İTÜ / tarım) & €101,0 mn; \%100 / \%0 / \%0 \\
& S4'te AB tarifesi: 1/98 sürerse / askıda / tam MFN ortalaması & \%6,4 / \%6,4 / \%6,4 \\
& Tepe satırlar (S4 artışı ≥ \%10) & yok \\
& AB pazarında AB dışı rakipler (AB tercihi) & Çin \%42 (MFN/GSP); Hindistan \%11 (MFN/GSP); Birleşik Krallık \%8 (TCA); İsviçre \%5 (AEA/STA); Japonya \%3 (STA/EPA) \\
& Tercihli rakiplerin payı; AME (puan; ham) & \%23; 88 (3,71) \\
& Menşe kuralı (PEM) & Fasıl 58: dokuma veya tufting + boyama, kaplama; nakış: menşesiz malzeme ≤ \%50; keçe: lif çekimi + kumaş oluşumu \\
& Menşe riski S2 / S3; c\_RoO (\%); p\_uyum & D/O; 2,0 / 2,5; 0,88 / 0,85 \etiketH{} \\
& Değer zincirindeki rol; DİR bağımlılığı & ara; orta \\
& EUDR ve SKDM kapsamındaki AB'ye ihracat payı & \%0; \%0 \\
& PEM ortaklarına ihracat; kümülasyon maruziyeti; κ & \$46 mn; \$32 mn; 0,3 \etiketH{} \\
5 Türkiye tarafı & Uygulanan tarife 2026 (ithalat ağırlıklı): MFN; İGV dahil & \%6,3; \%15,9 \\
& S4'te AB menşeline tarife (c2) & \%14,8 \\
& Kritik girdiler: 2024 ithalatı, AB payı, ilk kaynaklar & --- \\
6 Kabiliyet & KS; sıra (ağırlık duyarlılığında P10--P90) & 60,3; 27 / 230 (5--83) \\
& KS bileşenleri (yüzdelik): RCA, PCI, yurt içi KD, Ar-Ge, tedarikçi, verimlilik, firma & 74, 32, 74, 25, 96, 73, 74 \\
& PCI (sınıf) & −0,31 \\
& Atlanabilir ürünler (yoğunluk üst çeyrek, PCI \textgreater{} sınıf) & yok \\
7 Stratejik önem & SS; en yüksek boyut; ağırlıklı ortalama & 0,0; --- 0; 0,0 \etiketH{} \\
8 Makro duyarlılık & TÜFE ağırlığı (doğrudan); tüketimde ithal pay & \%0,00; \%16 \\
& ÜFE vekili (imalat üretimindeki pay) & \%0,18 \\
& EDE, 1 puan tarife: tüm menşeler / yalnız AB (bp; ρ = 1) & 0,007 / 0,002 \\
& CKE: tarife (\$ mn / bp); mali (\$ mn / mlr TL); puan & 1.601,4; 3,91; 34 \\
9 Senaryolar & Senaryo tablosu & Tablo C.10b (aşağıda) \\
\end{longtable}

\endgroup

\Needspace{7\baselineskip}

\textbf{Tablo C.10b --- 13.99 Diğer tekstil (şerit, dantel, keçe): senaryo etkileri (T+5'te yıllık; P50, parantez içinde P10 / P90)}

\nopagebreak[4]

\begingroup\gbtablo\small

\begin{longtable}[]{@{}
  >{\raggedright\arraybackslash}p{(\columnwidth - 6\tabcolsep) * \real{0.3860}}
  >{\raggedleft\arraybackslash}p{(\columnwidth - 6\tabcolsep) * \real{0.1927}}
  >{\raggedleft\arraybackslash}p{(\columnwidth - 6\tabcolsep) * \real{0.2106}}
  >{\raggedleft\arraybackslash}p{(\columnwidth - 6\tabcolsep) * \real{0.2106}}@{}}
\toprule\noalign{}
\begin{minipage}[b]{\linewidth}\raggedright
\textbf{Değişken}
\end{minipage} & \begin{minipage}[b]{\linewidth}\raggedleft
\textbf{S2}
\end{minipage} & \begin{minipage}[b]{\linewidth}\raggedleft
\textbf{S3}
\end{minipage} & \begin{minipage}[b]{\linewidth}\raggedleft
\textbf{S4}
\end{minipage} \\
\midrule\noalign{}
\endhead
\bottomrule\noalign{}
\endlastfoot
AB'ye ihracat (\$ mn) & −6 (−14 / +2) & −7 (−14 / −1) & −11 (−19 / −5) \\
AB'ye ihracat (\%) & −4,9 (−11,7 / +2,0) & −6,0 (−11,9 / −1,2) & −9,8 (−16,4 / −4,2) \\
Toplam ihracat (\$ mn) & −4 (−10 / +2) & −5 (−10 / −1) & −15 (−21 / −9) \\
Bunun içinde PEM kümülasyonu (\$ mn) & 0 (0 / 0) & 0 (0 / 0) & −6 (−7 / −4) \\
Nihai mal ithalatı (\$ mn) & −1 (−3 / +3) & −1 (−4 / +2) & −4 (−7 / −1) \\
Üretim (\$ mn) & −3 (−7 / −1) & −3 (−6 / −2) & −11 (−14 / −8) \\
İstihdam (kişi) & −37 (−84 / −11) & −42 (−79 / −24) & −131 (−180 / −94) \\
İstihdam (\%) & −0,2 (−0,4 / −0,1) & −0,2 (−0,4 / −0,1) & −0,6 (−0,8 / −0,4) \\
Cari denge (\$ mn) & −2 (−5 / −1) & −3 (−5 / −2) & −8 (−11 / −6) \\
TÜFE (bp) & 0,00 (0,00 / 0,00) & +0,01 (+0,01 / +0,01) & +0,03 (+0,02 / +0,03) \\
\end{longtable}

\endgroup

\emph{Kaynak:} \etiketT{} f3e; toplamlar Faz 2B genel denge dağılımına kalibre. Yüzdeler 2024 düzeyine göre.

\hypertarget{nace_kartlar_13-15-c11-1411-deri-giyim-eux15fyasux131-imalatux131}{%
\subsection{C.11 14.11 Deri giyim eşyası imalatı}\label{nace_kartlar_13-15-c11-1411-deri-giyim-eux15fyasux131-imalatux131}}

\Needspace{7\baselineskip}

\textbf{Tablo C.11 --- 14.11: veri sayfası}

\nopagebreak[4]

\begingroup\gbtablo\small

\begin{longtable}[]{@{}
  >{\raggedright\arraybackslash}p{(\columnwidth - 4\tabcolsep) * \real{0.1363}}
  >{\raggedright\arraybackslash}p{(\columnwidth - 4\tabcolsep) * \real{0.3655}}
  >{\raggedright\arraybackslash}p{(\columnwidth - 4\tabcolsep) * \real{0.4982}}@{}}
\toprule\noalign{}
\begin{minipage}[b]{\linewidth}\raggedright
\textbf{Alan}
\end{minipage} & \begin{minipage}[b]{\linewidth}\raggedright
\textbf{Gösterge}
\end{minipage} & \begin{minipage}[b]{\linewidth}\raggedright
\textbf{Değer}
\end{minipage} \\
\midrule\noalign{}
\endhead
\bottomrule\noalign{}
\endlastfoot
1 Kimlik & Rev. 2.1 karşılığı; CPA & 14.24 (n:1); CPA 14.11 \\
& Başlıca HS4 (ihracat, 2024) & 4203 \%100 \\
& BEC payları: ara / sermaye / tüketim & \%0 / \%0 / \%100 \\
2 Ölçek & Üretim; katma değer (GSYH payı) & \$223 mn; \$61 mn (\%0,00) \\
& Çalışan; girişim & 5.103; 726 \\
& Girişim ölçeği: 0 / 1--9 / 10--49 / 50--249 / 250+ çalışan & 33 / 620 / 111 / 12 / 4 \\
& 250+ girişimlerin istihdam payı; kişi başı KD (AB27 = 1) & \%20; 0,30 \\
& İstihdam değişimi, 2015--2024 & −\%38 \\
& İlk üç bölge (İBBS-2) & İstanbul \%59; İzmir \%21; Bursa, Eskişehir, Bilecik \%3 \\
& Çapa yatırımcılar & İstanbul (Zeytinburnu) ve İzmir deri konfeksiyonu \etiketH{} \\
3 Dış ticaret & İhracat 2024; 2020--2024 değişimi & \$128 mn; \%17 \\
& AB payı; ilk beş AB pazarı & \%41; Fransa \%26; Hollanda \%23; Almanya \%15; İtalya \%12; İspanya \%5 \\
& İlk beş AB dışı pazar & ABD \%36; Kazakistan \%11; Rusya \%9; Birleşik Krallık \%6; İsviçre \%5 \\
& İthalat 2024; AB ve Çin payı & \$17 mn; AB \%57, Çin \%1 \\
& İlk beş ithalat kaynağı & İtalya \%34; Hindistan \%28; İspanya \%9; Fransa \%9; Pakistan \%3 \\
& Dış ticaret dengesi & +\$111 mn \\
& İthal girdi içeriği (I-O); KD ÷ üretim & \%23; \%27 \\
& RCA: dünya 2020 → 2024; AB pazarı 2024 & 4,62 → 4,22; 1,94 \\
& Birim değer ÷ AB pazarı ortalaması & 2,25 \\
& AB pazarındaki pay (AB dışı tedarik içinde) 2020 → 2024; Türkiye'nin sırası & \%8,3 → \%7,7; 4. \\
4 AB tarafı & AB'ye ihracat 2025; rejim payları (GB / İTÜ / tarım) & €66,0 mn; \%100 / \%0 / \%0 \\
& S4'te AB tarifesi: 1/98 sürerse / askıda / tam MFN ortalaması & \%4,0 / \%4,0 / \%4,0 \\
& Tepe satırlar (S4 artışı ≥ \%10) & yok \\
& AB pazarında AB dışı rakipler (AB tercihi) & Hindistan \%36 (MFN/GSP); Pakistan \%27 (GSP+); Çin \%9 (MFN/GSP); Vietnam \%4 (STA/EPA); İsviçre \%3 (AEA/STA) \\
& Tercihli rakiplerin payı; AME (puan; ham) & \%43; 76 (2,68) \\
& Menşe kuralı (PEM) & Fasıl 42: başlık değişikliği ya da menşesiz malzeme ≤ \%50 \\
& Menşe riski S2 / S3; c\_RoO (\%); p\_uyum & D/D; 1,5 / 2,0; 0,92 / 0,88 \etiketH{} \\
& Değer zincirindeki rol; DİR bağımlılığı & nihai; düşük \\
& EUDR ve SKDM kapsamındaki AB'ye ihracat payı & \%0; \%0 \\
& PEM ortaklarına ihracat; kümülasyon maruziyeti; κ & \$7 mn; \$3 mn; 0,0 \etiketH{} \\
5 Türkiye tarafı & Uygulanan tarife 2026 (ithalat ağırlıklı): MFN; İGV dahil & \%4,0; \%43,0 \\
& S4'te AB menşeline tarife (c2) & \%43,0 \\
& Kritik girdiler: 2024 ithalatı, AB payı, ilk kaynaklar & --- \\
6 Kabiliyet & KS; sıra (ağırlık duyarlılığında P10--P90) & 45,4; 165 / 230 (86--200) \\
& KS bileşenleri (yüzdelik): RCA, PCI, yurt içi KD, Ar-Ge, tedarikçi, verimlilik, firma & 73, 10, 57, 25, 96, 14, 56 \\
& PCI (sınıf) & −0,97 \\
& Atlanabilir ürünler (yoğunluk üst çeyrek, PCI \textgreater{} sınıf) & yok \\
7 Stratejik önem & SS; en yüksek boyut; ağırlıklı ortalama & 0,0; --- 0; 0,0 \etiketH{} \\
8 Makro duyarlılık & TÜFE ağırlığı (doğrudan); tüketimde ithal pay & \%0,00; \%7 \\
& ÜFE vekili (imalat üretimindeki pay) & \%0,03 \\
& EDE, 1 puan tarife: tüm menşeler / yalnız AB (bp; ρ = 1) & 0,001 / 0,000 \\
& CKE: tarife (\$ mn / bp); mali (\$ mn / mlr TL); puan & 2.159,8; 7,08; 41 \\
9 Senaryolar & Senaryo tablosu & Tablo C.11b (aşağıda) \\
\end{longtable}

\endgroup

\Needspace{7\baselineskip}

\textbf{Tablo C.11b --- 14.11 Deri giyim: senaryo etkileri (T+5'te yıllık; P50, parantez içinde P10 / P90)}

\nopagebreak[4]

\begingroup\gbtablo\small

\begin{longtable}[]{@{}lrrr@{}}
\toprule\noalign{}
\textbf{Değişken} & \textbf{S2} & \textbf{S3} & \textbf{S4} \\
\midrule\noalign{}
\endhead
\bottomrule\noalign{}
\endlastfoot
AB'ye ihracat (\$ mn) & −3 (−6 / +1) & −3 (−7 / −1) & −5 (−8 / −2) \\
AB'ye ihracat (\%) & −4,8 (−11,3 / +2,0) & −6,4 (−12,5 / −1,2) & −9,2 (−14,9 / −3,9) \\
Toplam ihracat (\$ mn) & −2 (−4 / +1) & −2 (−5 / 0) & −4 (−6 / −2) \\
Nihai mal ithalatı (\$ mn) & −1 (−2 / +1) & −1 (−2 / 0) & −2 (−4 / 0) \\
Üretim (\$ mn) & −1 (−3 / 0) & −1 (−3 / −1) & −2 (−3 / −1) \\
İstihdam (kişi) & −19 (−41 / +1) & −24 (−44 / −12) & −27 (−46 / −14) \\
İstihdam (\%) & −0,4 (−0,8 / 0,0) & −0,5 (−0,9 / −0,2) & −0,5 (−0,9 / −0,3) \\
Cari denge (\$ mn) & −1 (−2 / 0) & −1 (−2 / −1) & −1 (−2 / −1) \\
TÜFE (bp) & 0,00 (0,00 / 0,00) & 0,00 (0,00 / 0,00) & +0,02 (+0,02 / +0,02) \\
\end{longtable}

\endgroup

\emph{Kaynak:} \etiketT{} f3e; toplamlar Faz 2B genel denge dağılımına kalibre. Yüzdeler 2024 düzeyine göre.

\hypertarget{nace_kartlar_13-15-c12-1412-iux307ux15f-giysisi-imalatux131}{%
\subsection{C.12 14.12 İş giysisi imalatı}\label{nace_kartlar_13-15-c12-1412-iux307ux15f-giysisi-imalatux131}}

\Needspace{7\baselineskip}

\textbf{Tablo C.12 --- 14.12: veri sayfası}

\nopagebreak[4]

\begingroup\gbtablo\small

\begin{longtable}[]{@{}
  >{\raggedright\arraybackslash}p{(\columnwidth - 4\tabcolsep) * \real{0.1352}}
  >{\raggedright\arraybackslash}p{(\columnwidth - 4\tabcolsep) * \real{0.3548}}
  >{\raggedright\arraybackslash}p{(\columnwidth - 4\tabcolsep) * \real{0.5100}}@{}}
\toprule\noalign{}
\begin{minipage}[b]{\linewidth}\raggedright
\textbf{Alan}
\end{minipage} & \begin{minipage}[b]{\linewidth}\raggedright
\textbf{Gösterge}
\end{minipage} & \begin{minipage}[b]{\linewidth}\raggedright
\textbf{Değer}
\end{minipage} \\
\midrule\noalign{}
\endhead
\bottomrule\noalign{}
\endlastfoot
1 Kimlik & Rev. 2.1 karşılığı; CPA & 14.23 (1:1); CPA 14.12 \\
& Başlıca HS4 (ihracat, 2024) & 6211 \%57, 6203 \%35, 6204 \%9 \\
& BEC payları: ara / sermaye / tüketim & \%0 / \%0 / \%100 \\
2 Ölçek & Üretim; katma değer (GSYH payı) & \$806 mn; \$189 mn (\%0,01) \\
& Çalışan; girişim & 12.634; 1.697 \\
& Girişim ölçeği: 0 / 1--9 / 10--49 / 50--249 / 250+ çalışan & 79 / 1.512 / 180 / 43 / 10 \\
& 250+ girişimlerin istihdam payı; kişi başı KD (AB27 = 1) & \%19; 0,46 \\
& İstihdam değişimi, 2015--2024 & \%59 \\
& İlk üç bölge (İBBS-2) & İstanbul \%18; Kastamonu, Çankırı, Sinop \%13; İzmir \%11 \\
& Çapa yatırımcılar & Kamu ve güvenlik güçlerinin tedarikçileri; Kastamonu ve İstanbul üreticileri \etiketH{} \\
3 Dış ticaret & İhracat 2024; 2020--2024 değişimi & \$69 mn; −\%25 \\
& AB payı; ilk beş AB pazarı & \%53; Almanya \%25; Hollanda \%20; Polonya \%14; İspanya \%8; Fransa \%7 \\
& İlk beş AB dışı pazar & Kazakistan \%10; Irak \%8; Birleşik Krallık \%8; Rusya \%8; Kırgızistan \%6 \\
& İthalat 2024; AB ve Çin payı & \$11 mn; AB \%26, Çin \%20 \\
& İlk beş ithalat kaynağı & Çin \%20; Bangladeş \%15; Vietnam \%7; Mısır \%7; İtalya \%5 \\
& Dış ticaret dengesi & +\$58 mn \\
& İthal girdi içeriği (I-O); KD ÷ üretim & \%23; \%23 \\
& RCA: dünya 2020 → 2024; AB pazarı 2024 & 3,67 → 1,97; 1,53 \\
& Birim değer ÷ AB pazarı ortalaması & 1,01 \\
& AB pazarındaki pay (AB dışı tedarik içinde) 2020 → 2024; Türkiye'nin sırası & \%10,2 → \%6,1; 6. \\
4 AB tarafı & AB'ye ihracat 2025; rejim payları (GB / İTÜ / tarım) & €24,9 mn; \%100 / \%0 / \%0 \\
& S4'te AB tarifesi: 1/98 sürerse / askıda / tam MFN ortalaması & \%12,0 / \%12,0 / \%12,0 \\
& Tepe satırlar (S4 artışı ≥ \%10) & 28 satır, AB'ye ihracatın \%100'i: KN 6211 33 10 \%12,0 (€8,4 mn); KN 6203 43 11 \%12,0 (€3,4 mn); KN 6211 42 10 \%12,0 (€2,2 mn) \\
& AB pazarında AB dışı rakipler (AB tercihi) & Çin \%29 (MFN/GSP); Bangladeş \%12 (EBA); Tunus \%12 (STA/EPA); Vietnam \%8 (STA/EPA); Fas \%7 (STA/EPA) \\
& Tercihli rakiplerin payı; AME (puan; ham) & \%56; 97 (8,28) \\
& Menşe kuralı (PEM) & Fasıl 62: dokuma + konfeksiyon ya da baskı + konfeksiyon; yangına dayanıklı alüminyum kaplı giysi: kaplama + konfeksiyon (kumaş ≤ \%40) \\
& Menşe riski S2 / S3; c\_RoO (\%); p\_uyum & O/O; 2,5 / 3,0; 0,85 / 0,80 \etiketH{} \\
& Değer zincirindeki rol; DİR bağımlılığı & nihai; orta \\
& EUDR ve SKDM kapsamındaki AB'ye ihracat payı & \%0; \%0 \\
& PEM ortaklarına ihracat; kümülasyon maruziyeti; κ & \$6 mn; \$3 mn; 0,0 \etiketH{} \\
5 Türkiye tarafı & Uygulanan tarife 2026 (ithalat ağırlıklı): MFN; İGV dahil & \%12,0; \%51,0 \\
& S4'te AB menşeline tarife (c2) & \%51,0 \\
& Kritik girdiler: 2024 ithalatı, AB payı, ilk kaynaklar & --- \\
6 Kabiliyet & KS; sıra (ağırlık duyarlılığında P10--P90) & 43,2; 179 / 230 (114--207) \\
& KS bileşenleri (yüzdelik): RCA, PCI, yurt içi KD, Ar-Ge, tedarikçi, verimlilik, firma & 54, 2, 48, 25, 96, 50, 64 \\
& PCI (sınıf) & −1,70 \\
& Atlanabilir ürünler (yoğunluk üst çeyrek, PCI \textgreater{} sınıf) & yok \\
7 Stratejik önem & SS; en yüksek boyut; ağırlıklı ortalama & 22,8; savunma 30; 6,0 \etiketH{} \\
8 Makro duyarlılık & TÜFE ağırlığı (doğrudan); tüketimde ithal pay & \%0,00; \%8 \\
& ÜFE vekili (imalat üretimindeki pay) & \%0,12 \\
& EDE, 1 puan tarife: tüm menşeler / yalnız AB (bp; ρ = 1) & 0,001 / 0,000 \\
& CKE: tarife (\$ mn / bp); mali (\$ mn / mlr TL); puan & 2.695,5; 0,80; 34 \\
9 Senaryolar & Senaryo tablosu & Tablo C.12b (aşağıda) \\
\end{longtable}

\endgroup

\Needspace{7\baselineskip}

\textbf{Tablo C.12b --- 14.12 İş giysisi: senaryo etkileri (T+5'te yıllık; P50, parantez içinde P10 / P90)}

\nopagebreak[4]

\begingroup\gbtablo\small

\begin{longtable}[]{@{}lrrr@{}}
\toprule\noalign{}
\textbf{Değişken} & \textbf{S2} & \textbf{S3} & \textbf{S4} \\
\midrule\noalign{}
\endhead
\bottomrule\noalign{}
\endlastfoot
AB'ye ihracat (\$ mn) & −4 (−9 / +1) & −5 (−9 / −1) & −8 (−13 / −3) \\
AB'ye ihracat (\%) & −10,8 (−24,4 / +2,0) & −12,6 (−23,4 / −2,3) & −22,1 (−36,4 / −9,2) \\
Toplam ihracat (\$ mn) & −3 (−7 / 0) & −3 (−6 / −1) & −6 (−11 / −3) \\
Nihai mal ithalatı (\$ mn) & 0 (−1 / 0) & 0 (−1 / 0) & −1 (−1 / 0) \\
Üretim (\$ mn) & −3 (−6 / 0) & −3 (−5 / −1) & −6 (−9 / −3) \\
İstihdam (kişi) & −27 (−63 / +3) & −31 (−59 / −8) & −62 (−104 / −27) \\
İstihdam (\%) & −0,2 (−0,5 / 0,0) & −0,2 (−0,5 / −0,1) & −0,5 (−0,8 / −0,2) \\
Cari denge (\$ mn) & −2 (−4 / 0) & −2 (−4 / −1) & −4 (−7 / −2) \\
TÜFE (bp) & 0,00 (0,00 / 0,00) & 0,00 (0,00 / 0,00) & +0,01 (+0,01 / +0,01) \\
\end{longtable}

\endgroup

\emph{Kaynak:} \etiketT{} f3e; toplamlar Faz 2B genel denge dağılımına kalibre. Yüzdeler 2024 düzeyine göre.

\hypertarget{nace_kartlar_13-15-c13-1413-diux11fer-dux131ux15f-giyim-eux15fyasux131-imalatux131}{%
\subsection{C.13 14.13 Diğer dış giyim eşyası imalatı}\label{nace_kartlar_13-15-c13-1413-diux11fer-dux131ux15f-giyim-eux15fyasux131-imalatux131}}

\Needspace{7\baselineskip}

\textbf{Tablo C.13 --- 14.13: veri sayfası}

\nopagebreak[4]

\begingroup\gbtablo\small

\begin{longtable}[]{@{}
  >{\raggedright\arraybackslash}p{(\columnwidth - 4\tabcolsep) * \real{0.1346}}
  >{\raggedright\arraybackslash}p{(\columnwidth - 4\tabcolsep) * \real{0.3486}}
  >{\raggedright\arraybackslash}p{(\columnwidth - 4\tabcolsep) * \real{0.5169}}@{}}
\toprule\noalign{}
\begin{minipage}[b]{\linewidth}\raggedright
\textbf{Alan}
\end{minipage} & \begin{minipage}[b]{\linewidth}\raggedright
\textbf{Gösterge}
\end{minipage} & \begin{minipage}[b]{\linewidth}\raggedright
\textbf{Değer}
\end{minipage} \\
\midrule\noalign{}
\endhead
\bottomrule\noalign{}
\endlastfoot
1 Kimlik & Rev. 2.1 karşılığı; CPA & 14.10, 14.21 (1:2); CPA 14.13 \\
& Başlıca HS4 (ihracat, 2024) & 6204 \%45, 6203 \%23, 6104 \%19, 6103 \%7 \\
& BEC payları: ara / sermaye / tüketim & \%1 / \%0 / \%99 \\
2 Ölçek & Üretim; katma değer (GSYH payı) & \$14.680 mn; \$4.024 mn (\%0,30) \\
& Çalışan; girişim & 365.708; 41.798 \\
& Girişim ölçeği: 0 / 1--9 / 10--49 / 50--249 / 250+ çalışan & 1.279 / 41.605 / 3.465 / 1.190 / 143 \\
& 250+ girişimlerin istihdam payı; kişi başı KD (AB27 = 1) & \%11; 0,31 \\
& İstihdam değişimi, 2015--2024 & \%22 \\
& İlk üç bölge (İBBS-2) & İstanbul \%31; Mardin, Batman, Şırnak, Siirt \%14; Şanlıurfa, Diyarbakır \%5 \\
& Çapa yatırımcılar & LC Waikiki, Koton, DeFacto, Mavi (marka); Inditex, H\&M, Primark ve Alman markalarının fason tedarik ağları \etiketH{} \\
3 Dış ticaret & İhracat 2024; 2020--2024 değişimi & \$8.577 mn; \%32 \\
& AB payı; ilk beş AB pazarı & \%59; Almanya \%24; Hollanda \%20; İspanya \%18; Polonya \%10; Fransa \%8 \\
& İlk beş AB dışı pazar & Birleşik Krallık \%15; ABD \%11; Irak \%8; İsviçre \%6; Kazakistan \%6 \\
& İthalat 2024; AB ve Çin payı & \$2.081 mn; AB \%26, Çin \%14 \\
& İlk beş ithalat kaynağı & Bangladeş \%16; Çin \%14; İspanya \%9; İtalya \%8; Kamboçya \%7 \\
& Dış ticaret dengesi & +\$6.497 mn \\
& İthal girdi içeriği (I-O); KD ÷ üretim & \%23; \%27 \\
& RCA: dünya 2020 → 2024; AB pazarı 2024 & 3,62 → 3,36; 2,19 \\
& Birim değer ÷ AB pazarı ortalaması & 1,16 \\
& AB pazarındaki pay (AB dışı tedarik içinde) 2020 → 2024; Türkiye'nin sırası & \%8,9 → \%8,7; 3. \\
4 AB tarafı & AB'ye ihracat 2025; rejim payları (GB / İTÜ / tarım) & €3.588,4 mn; \%100 / \%0 / \%0 \\
& S4'te AB tarifesi: 1/98 sürerse / askıda / tam MFN ortalaması & \%12,0 / \%12,0 / \%12,0 \\
& Tepe satırlar (S4 artışı ≥ \%10) & 133 satır, AB'ye ihracatın \%100'i: KN 6204 62 31 \%12,0 (€508,1 mn); KN 6203 42 31 \%12,0 (€360,6 mn); KN 6204 63 18 \%12,0 (€190,0 mn) \\
& AB pazarında AB dışı rakipler (AB tercihi) & Çin \%26 (MFN/GSP); Bangladeş \%20 (EBA); Kamboçya \%6 (EBA); Vietnam \%6 (STA/EPA); Fas \%6 (STA/EPA) \\
& Tercihli rakiplerin payı; AME (puan; ham) & \%58; 97 (9,09) \\
& Menşe kuralı (PEM) & Fasıl 61: örme + konfeksiyon; Fasıl 62: dokuma + konfeksiyon ya da baskı + konfeksiyon; kurala uymayan tekstil ≤ \%15 (Not 7.1) \\
& Menşe riski S2 / S3; c\_RoO (\%); p\_uyum & O/O; 2,5 / 3,0; 0,86 / 0,82 \etiketH{} \\
& Değer zincirindeki rol; DİR bağımlılığı & nihai; orta \\
& EUDR ve SKDM kapsamındaki AB'ye ihracat payı & \%0; \%0 \\
& PEM ortaklarına ihracat; kümülasyon maruziyeti; κ & \$652 mn; \$310 mn; 0,0 \etiketH{} \\
5 Türkiye tarafı & Uygulanan tarife 2026 (ithalat ağırlıklı): MFN; İGV dahil & \%8,7; \%37,0 \\
& S4'te AB menşeline tarife (c2) & \%37,2 \\
& Kritik girdiler: 2024 ithalatı, AB payı, ilk kaynaklar & HS4 5208, 5209, 5407, 6006: \$892 mn (AB \%20; Mısır \%25, Çin \%20, Pakistan \%9) \\
6 Kabiliyet & KS; sıra (ağırlık duyarlılığında P10--P90) & 46,5; 154 / 230 (63--200) \\
& KS bileşenleri (yüzdelik): RCA, PCI, yurt içi KD, Ar-Ge, tedarikçi, verimlilik, firma & 70, 2, 56, 25, 96, 17, 85 \\
& PCI (sınıf) & −1,73 \\
& Atlanabilir ürünler (yoğunluk üst çeyrek, PCI \textgreater{} sınıf) & yok \\
7 Stratejik önem & SS; en yüksek boyut; ağırlıklı ortalama & 0,0; --- 0; 0,0 \etiketH{} \\
8 Makro duyarlılık & TÜFE ağırlığı (doğrudan); tüketimde ithal pay & \%3,86; \%11 \\
& ÜFE vekili (imalat üretimindeki pay) & \%2,20 \\
& EDE, 1 puan tarife: tüm menşeler / yalnız AB (bp; ρ = 1) & 0,544 / 0,141 \\
& CKE: tarife (\$ mn / bp); mali (\$ mn / mlr TL); puan & 398,9; 11,87; 30 \\
9 Senaryolar & Senaryo tablosu & Tablo 5.5 \\
\end{longtable}

\endgroup

\hypertarget{nace_kartlar_13-15-c14-1414-iux307uxe7-giyim-eux15fyasux131-imalatux131}{%
\subsection{C.14 14.14 İç giyim eşyası imalatı}\label{nace_kartlar_13-15-c14-1414-iux307uxe7-giyim-eux15fyasux131-imalatux131}}

\Needspace{7\baselineskip}

\textbf{Tablo C.14 --- 14.14: veri sayfası}

\nopagebreak[4]

\begingroup\gbtablo\small

\begin{longtable}[]{@{}
  >{\raggedright\arraybackslash}p{(\columnwidth - 4\tabcolsep) * \real{0.1348}}
  >{\raggedright\arraybackslash}p{(\columnwidth - 4\tabcolsep) * \real{0.3511}}
  >{\raggedright\arraybackslash}p{(\columnwidth - 4\tabcolsep) * \real{0.5141}}@{}}
\toprule\noalign{}
\begin{minipage}[b]{\linewidth}\raggedright
\textbf{Alan}
\end{minipage} & \begin{minipage}[b]{\linewidth}\raggedright
\textbf{Gösterge}
\end{minipage} & \begin{minipage}[b]{\linewidth}\raggedright
\textbf{Değer}
\end{minipage} \\
\midrule\noalign{}
\endhead
\bottomrule\noalign{}
\endlastfoot
1 Kimlik & Rev. 2.1 karşılığı; CPA & 14.10, 14.21, 14.22 (1:3); CPA 14.14 \\
& Başlıca HS4 (ihracat, 2024) & 6109 \%51, 6206 \%12, 6205 \%10, 6106 \%7 \\
& BEC payları: ara / sermaye / tüketim & \%0 / \%0 / \%100 \\
2 Ölçek & Üretim; katma değer (GSYH payı) & \$10.149 mn; \$2.482 mn (\%0,18) \\
& Çalışan; girişim & 183.980; 8.473 \\
& Girişim ölçeği: 0 / 1--9 / 10--49 / 50--249 / 250+ çalışan & 256 / 6.658 / 1.482 / 671 / 94 \\
& 250+ girişimlerin istihdam payı; kişi başı KD (AB27 = 1) & \%33; 0,78 \\
& İstihdam değişimi, 2015--2024 & \%21 \\
& İlk üç bölge (İBBS-2) & İstanbul \%41; Mardin, Batman, Şırnak, Siirt \%14; İzmir \%7 \\
& Çapa yatırımcılar & İzmir, İstanbul ve Denizli'deki entegre örme ve konfeksiyon üreticileri; Penti, Koton (marka) \etiketH{} \\
3 Dış ticaret & İhracat 2024; 2020--2024 değişimi & \$5.647 mn; \%22 \\
& AB payı; ilk beş AB pazarı & \%60; Almanya \%28; İspanya \%16; Hollanda \%15; Fransa \%10; Polonya \%7 \\
& İlk beş AB dışı pazar & Birleşik Krallık \%19; ABD \%9; Irak \%8; Kazakistan \%6; Rusya \%4 \\
& İthalat 2024; AB ve Çin payı & \$883 mn; AB \%21, Çin \%11 \\
& İlk beş ithalat kaynağı & Bangladeş \%26; Çin \%11; Gürcistan \%9; İtalya \%8; Mısır \%7 \\
& Dış ticaret dengesi & +\$4.765 mn \\
& İthal girdi içeriği (I-O); KD ÷ üretim & \%23; \%24 \\
& RCA: dünya 2020 → 2024; AB pazarı 2024 & 4,41 → 3,77; 2,69 \\
& Birim değer ÷ AB pazarı ortalaması & 1,29 \\
& AB pazarındaki pay (AB dışı tedarik içinde) 2020 → 2024; Türkiye'nin sırası & \%11,6 → \%10,7; 3. \\
4 AB tarafı & AB'ye ihracat 2025; rejim payları (GB / İTÜ / tarım) & €2.246,2 mn; \%100 / \%0 / \%0 \\
& S4'te AB tarifesi: 1/98 sürerse / askıda / tam MFN ortalaması & \%12,0 / \%12,0 / \%12,0 \\
& Tepe satırlar (S4 artışı ≥ \%10) & 61 satır, AB'ye ihracatın \%99'i: KN 6109 10 00 \%12,0 (€837,8 mn); KN 6109 90 20 \%12,0 (€385,2 mn); KN 6205 20 00 \%12,0 (€180,7 mn) \\
& AB pazarında AB dışı rakipler (AB tercihi) & Bangladeş \%32 (EBA); Çin \%19 (MFN/GSP); Hindistan \%8 (MFN/GSP); Vietnam \%5 (STA/EPA); Fas \%4 (STA/EPA) \\
& Tercihli rakiplerin payı; AME (puan; ham) & \%57; 97 (8,95) \\
& Menşe kuralı (PEM) & Fasıl 61: örme + konfeksiyon; 6212 (sütyen, korse): örme + konfeksiyon ya da baskı + konfeksiyon \\
& Menşe riski S2 / S3; c\_RoO (\%); p\_uyum & D/O; 2,0 / 2,5; 0,88 / 0,85 \etiketH{} \\
& Değer zincirindeki rol; DİR bağımlılığı & nihai; orta \\
& EUDR ve SKDM kapsamındaki AB'ye ihracat payı & \%0; \%0 \\
& PEM ortaklarına ihracat; kümülasyon maruziyeti; κ & \$388 mn; \$192 mn; 0,0 \etiketH{} \\
5 Türkiye tarafı & Uygulanan tarife 2026 (ithalat ağırlıklı): MFN; İGV dahil & \%11,5; \%50,5 \\
& S4'te AB menşeline tarife (c2) & \%50,7 \\
& Kritik girdiler: 2024 ithalatı, AB payı, ilk kaynaklar & HS4 6004, 6006, 5205: \$753 mn (AB \%10; Özbekistan \%40, Çin \%10, Türkmenistan \%8) \\
6 Kabiliyet & KS; sıra (ağırlık duyarlılığında P10--P90) & 51,0; 106 / 230 (31--171) \\
& KS bileşenleri (yüzdelik): RCA, PCI, yurt içi KD, Ar-Ge, tedarikçi, verimlilik, firma & 69, 1, 49, 25, 96, 90, 73 \\
& PCI (sınıf) & −1,90 \\
& Atlanabilir ürünler (yoğunluk üst çeyrek, PCI \textgreater{} sınıf) & yok \\
7 Stratejik önem & SS; en yüksek boyut; ağırlıklı ortalama & 0,0; --- 0; 0,0 \etiketH{} \\
8 Makro duyarlılık & TÜFE ağırlığı (doğrudan); tüketimde ithal pay & \%0,64; \%8 \\
& ÜFE vekili (imalat üretimindeki pay) & \%1,52 \\
& EDE, 1 puan tarife: tüm menşeler / yalnız AB (bp; ρ = 1) & 0,098 / 0,020 \\
& CKE: tarife (\$ mn / bp); mali (\$ mn / mlr TL); puan & 1.786,1; 13,65; 44 \\
9 Senaryolar & Senaryo tablosu & Tablo 5.6 \\
\end{longtable}

\endgroup

\hypertarget{nace_kartlar_13-15-c15-1419-diux11fer-giyim-eux15fyasux131-ve-aksesuarlarux131nux131n-imalatux131}{%
\subsection{C.15 14.19 Diğer giyim eşyası ve aksesuarlarının imalatı}\label{nace_kartlar_13-15-c15-1419-diux11fer-giyim-eux15fyasux131-ve-aksesuarlarux131nux131n-imalatux131}}

\Needspace{7\baselineskip}

\textbf{Tablo C.15 --- 14.19: veri sayfası}

\nopagebreak[4]

\begingroup\gbtablo\small

\begin{longtable}[]{@{}
  >{\raggedright\arraybackslash}p{(\columnwidth - 4\tabcolsep) * \real{0.1351}}
  >{\raggedright\arraybackslash}p{(\columnwidth - 4\tabcolsep) * \real{0.3537}}
  >{\raggedright\arraybackslash}p{(\columnwidth - 4\tabcolsep) * \real{0.5112}}@{}}
\toprule\noalign{}
\begin{minipage}[b]{\linewidth}\raggedright
\textbf{Alan}
\end{minipage} & \begin{minipage}[b]{\linewidth}\raggedright
\textbf{Gösterge}
\end{minipage} & \begin{minipage}[b]{\linewidth}\raggedright
\textbf{Değer}
\end{minipage} \\
\midrule\noalign{}
\endhead
\bottomrule\noalign{}
\endlastfoot
1 Kimlik & Rev. 2.1 karşılığı; CPA & 14.10, 14.21, 14.22, 14.29 (1:4); CPA 14.19 \\
& Başlıca HS4 (ihracat, 2024) & 6114 \%21, 6211 \%17, 6111 \%14, 6112 \%12 \\
& BEC payları: ara / sermaye / tüketim & \%0 / \%0 / \%100 \\
2 Ölçek & Üretim; katma değer (GSYH payı) & \$1.645 mn; \$466 mn (\%0,03) \\
& Çalışan; girişim & 41.506; 4.599 \\
& Girişim ölçeği: 0 / 1--9 / 10--49 / 50--249 / 250+ çalışan & 208 / 4.251 / 521 / 164 / 9 \\
& 250+ girişimlerin istihdam payı; kişi başı KD (AB27 = 1) & \%5; 0,35 \\
& İstihdam değişimi, 2015--2024 & \%31 \\
& İlk üç bölge (İBBS-2) & İstanbul \%25; Bursa, Eskişehir, Bilecik \%15; Samsun, Tokat, Çorum, Amasya \%6 \\
& Çapa yatırımcılar & İstanbul ve Bursa'daki bebek giyimi ve aksesuar üreticileri \etiketH{} \\
3 Dış ticaret & İhracat 2024; 2020--2024 değişimi & \$1.461 mn; −\%37 \\
& AB payı; ilk beş AB pazarı & \%47; İspanya \%23; Almanya \%22; Fransa \%11; Hollanda \%9; Polonya \%7 \\
& İlk beş AB dışı pazar & Birleşik Krallık \%13; ABD \%9; Irak \%7; Kazakistan \%7; Suudi Arabistan \%6 \\
& İthalat 2024; AB ve Çin payı & \$507 mn; AB \%23, Çin \%33 \\
& İlk beş ithalat kaynağı & Çin \%33; Bangladeş \%7; İtalya \%7; İspanya \%6; Myanmar \%5 \\
& Dış ticaret dengesi & +\$954 mn \\
& İthal girdi içeriği (I-O); KD ÷ üretim & \%23; \%28 \\
& RCA: dünya 2020 → 2024; AB pazarı 2024 & 2,75 → 1,59; 0,97 \\
& Birim değer ÷ AB pazarı ortalaması & 1,16 \\
& AB pazarındaki pay (AB dışı tedarik içinde) 2020 → 2024; Türkiye'nin sırası & \%6,6 → \%3,9; 6. \\
4 AB tarafı & AB'ye ihracat 2025; rejim payları (GB / İTÜ / tarım) & €514,9 mn; \%100 / \%0 / \%0 \\
& S4'te AB tarifesi: 1/98 sürerse / askıda / tam MFN ortalaması & \%11,2 / \%11,2 / \%11,2 \\
& Tepe satırlar (S4 artışı ≥ \%10) & 70 satır, AB'ye ihracatın \%89'i: KN 6114 30 00 \%12,0 (€148,8 mn); KN 6111 20 90 \%12,0 (€41,0 mn); KN 6211 43 90 \%12,0 (€34,6 mn) \\
& AB pazarında AB dışı rakipler (AB tercihi) & Çin \%43 (MFN/GSP); Bangladeş \%12 (EBA); Vietnam \%7 (STA/EPA); Hindistan \%5 (MFN/GSP); Kamboçya \%4 (EBA) \\
& Tercihli rakiplerin payı; AME (puan; ham) & \%42; 96 (6,88) \\
& Menşe kuralı (PEM) & Fasıl 61--62 kuralları; şal ve atkıda baskı + konfeksiyon; 6217: dokuma + konfeksiyon \\
& Menşe riski S2 / S3; c\_RoO (\%); p\_uyum & O/O; 2,5 / 3,0; 0,85 / 0,80 \etiketH{} \\
& Değer zincirindeki rol; DİR bağımlılığı & nihai; orta \\
& EUDR ve SKDM kapsamındaki AB'ye ihracat payı & \%0; \%0 \\
& PEM ortaklarına ihracat; kümülasyon maruziyeti; κ & \$146 mn; \$73 mn; 0,0 \etiketH{} \\
5 Türkiye tarafı & Uygulanan tarife 2026 (ithalat ağırlıklı): MFN; İGV dahil & \%9,5; \%45,4 \\
& S4'te AB menşeline tarife (c2) & \%44,8 \\
& Kritik girdiler: 2024 ithalatı, AB payı, ilk kaynaklar & --- \\
6 Kabiliyet & KS; sıra (ağırlık duyarlılığında P10--P90) & 42,8; 183 / 230 (112--212) \\
& KS bileşenleri (yüzdelik): RCA, PCI, yurt içi KD, Ar-Ge, tedarikçi, verimlilik, firma & 46, 3, 58, 25, 96, 24, 85 \\
& PCI (sınıf) & −1,66 \\
& Atlanabilir ürünler (yoğunluk üst çeyrek, PCI \textgreater{} sınıf) & yok \\
7 Stratejik önem & SS; en yüksek boyut; ağırlıklı ortalama & 0,0; --- 0; 0,0 \etiketH{} \\
8 Makro duyarlılık & TÜFE ağırlığı (doğrudan); tüketimde ithal pay & \%0,90; \%15 \\
& ÜFE vekili (imalat üretimindeki pay) & \%0,25 \\
& EDE, 1 puan tarife: tüm menşeler / yalnız AB (bp; ρ = 1) & 0,160 / 0,037 \\
& CKE: tarife (\$ mn / bp); mali (\$ mn / mlr TL); puan & 152,8; 22,74; 35 \\
9 Senaryolar & Senaryo tablosu & Tablo C.15b (aşağıda) \\
\end{longtable}

\endgroup

\Needspace{7\baselineskip}

\textbf{Tablo C.15b --- 14.19 Diğer giyim ve aksesuar: senaryo etkileri (T+5'te yıllık; P50, parantez içinde P10 / P90)}

\nopagebreak[4]

\begingroup\gbtablo\small

\begin{longtable}[]{@{}
  >{\raggedright\arraybackslash}p{(\columnwidth - 6\tabcolsep) * \real{0.2727}}
  >{\raggedleft\arraybackslash}p{(\columnwidth - 6\tabcolsep) * \real{0.2325}}
  >{\raggedleft\arraybackslash}p{(\columnwidth - 6\tabcolsep) * \real{0.2364}}
  >{\raggedleft\arraybackslash}p{(\columnwidth - 6\tabcolsep) * \real{0.2584}}@{}}
\toprule\noalign{}
\begin{minipage}[b]{\linewidth}\raggedright
\textbf{Değişken}
\end{minipage} & \begin{minipage}[b]{\linewidth}\raggedleft
\textbf{S2}
\end{minipage} & \begin{minipage}[b]{\linewidth}\raggedleft
\textbf{S3}
\end{minipage} & \begin{minipage}[b]{\linewidth}\raggedleft
\textbf{S4}
\end{minipage} \\
\midrule\noalign{}
\endhead
\bottomrule\noalign{}
\endlastfoot
AB'ye ihracat (\$ mn) & −49 (−119 / +14) & −59 (−117 / −11) & −107 (−178 / −45) \\
AB'ye ihracat (\%) & −7,2 (−17,5 / +2,0) & −8,6 (−17,1 / −1,6) & −15,7 (−26,1 / −6,5) \\
Toplam ihracat (\$ mn) & −36 (−86 / +9) & −42 (−83 / −8) & −86 (−145 / −36) \\
Nihai mal ithalatı (\$ mn) & −3 (−10 / +4) & −4 (−11 / +2) & −12 (−21 / −2) \\
Üretim (\$ mn) & −32 (−77 / +5) & −37 (−71 / −10) & −74 (−125 / −33) \\
İstihdam (kişi) & −563 (−1.328 / +100) & −642 (−1.253 / −177) & −1.289 (−2.229 / −568) \\
İstihdam (\%) & −1,4 (−3,2 / +0,2) & −1,5 (−3,0 / −0,4) & −3,1 (−5,4 / −1,4) \\
Cari denge (\$ mn) & −25 (−59 / +4) & −29 (−55 / −8) & −57 (−97 / −25) \\
TÜFE (bp) & +0,14 (+0,12 / +0,16) & +0,26 (+0,22 / +0,30) & +1,37 (+1,18 / +1,55) \\
\end{longtable}

\endgroup

\emph{Kaynak:} \etiketT{} f3e; toplamlar Faz 2B genel denge dağılımına kalibre. Yüzdeler 2024 düzeyine göre.

\hypertarget{nace_kartlar_13-15-c16-1420-kuxfcrkten-eux15fya-imalatux131}{%
\subsection{C.16 14.20 Kürkten eşya imalatı}\label{nace_kartlar_13-15-c16-1420-kuxfcrkten-eux15fya-imalatux131}}

\Needspace{7\baselineskip}

\textbf{Tablo C.16 --- 14.20: veri sayfası}

\nopagebreak[4]

\begingroup\gbtablo\small

\begin{longtable}[]{@{}
  >{\raggedright\arraybackslash}p{(\columnwidth - 4\tabcolsep) * \real{0.1362}}
  >{\raggedright\arraybackslash}p{(\columnwidth - 4\tabcolsep) * \real{0.3651}}
  >{\raggedright\arraybackslash}p{(\columnwidth - 4\tabcolsep) * \real{0.4987}}@{}}
\toprule\noalign{}
\begin{minipage}[b]{\linewidth}\raggedright
\textbf{Alan}
\end{minipage} & \begin{minipage}[b]{\linewidth}\raggedright
\textbf{Gösterge}
\end{minipage} & \begin{minipage}[b]{\linewidth}\raggedright
\textbf{Değer}
\end{minipage} \\
\midrule\noalign{}
\endhead
\bottomrule\noalign{}
\endlastfoot
1 Kimlik & Rev. 2.1 karşılığı; CPA & 14.24 (n:1); CPA 14.20 \\
& Başlıca HS4 (ihracat, 2024) & 4303 \%100 \\
& BEC payları: ara / sermaye / tüketim & \%0 / \%0 / \%100 \\
2 Ölçek & Üretim; katma değer (GSYH payı) & \$127 mn; \$37 mn (\%0,00) \\
& Çalışan; girişim & 1.602; 117 \\
& Girişim ölçeği: 0 / 1--9 / 10--49 / 50--249 / 250+ çalışan & 5 / 87 / 31 / 6 / 0 \\
& 250+ girişimlerin istihdam payı; kişi başı KD (AB27 = 1) & \%0; 1,00 \\
& İstihdam değişimi, 2015--2024 & −\%35 \\
& İlk üç bölge (İBBS-2) & İstanbul \%95; İzmir \%2; Kocaeli, Sakarya, Düzce, Bolu, Yalova \%1 \\
& Çapa yatırımcılar & İstanbul (Kapalıçarşı ve Zeytinburnu) kürk atölyeleri \etiketH{} \\
3 Dış ticaret & İhracat 2024; 2020--2024 değişimi & \$94 mn; \%24 \\
& AB payı; ilk beş AB pazarı & \%37; Fransa \%37; İtalya \%22; Almanya \%14; Hollanda \%7; Polonya \%4 \\
& İlk beş AB dışı pazar & ABD \%27; Kazakistan \%21; Rusya \%18; Kırgızistan \%8; Birleşik Krallık \%6 \\
& İthalat 2024; AB ve Çin payı & \$5 mn; AB \%91, Çin \%3 \\
& İlk beş ithalat kaynağı & İtalya \%45; Fransa \%23; Yunanistan \%15; İspanya \%3; Çin \%3 \\
& Dış ticaret dengesi & +\$89 mn \\
& İthal girdi içeriği (I-O); KD ÷ üretim & \%23; \%29 \\
& RCA: dünya 2020 → 2024; AB pazarı 2024 & 5,05 → 7,07; 6,84 \\
& Birim değer ÷ AB pazarı ortalaması & 1,10 \\
& AB pazarındaki pay (AB dışı tedarik içinde) 2020 → 2024; Türkiye'nin sırası & \%20,3 → \%27,2; 2. \\
4 AB tarafı & AB'ye ihracat 2025; rejim payları (GB / İTÜ / tarım) & €29,6 mn; \%100 / \%0 / \%0 \\
& S4'te AB tarifesi: 1/98 sürerse / askıda / tam MFN ortalaması & \%3,7 / \%3,7 / \%3,7 \\
& Tepe satırlar (S4 artışı ≥ \%10) & yok \\
& AB pazarında AB dışı rakipler (AB tercihi) & Çin \%31 (MFN/GSP); Filipinler \%9 (GSP+); Birleşik Krallık \%8 (TCA); İsviçre \%6 (AEA/STA); ABD \%2 (MFN/GSP) \\
& Tercihli rakiplerin payı; AME (puan; ham) & \%30; 68 (2,31) \\
& Menşe kuralı (PEM) & 4303: birleştirilmemiş, tabaklanmış ya da işlenmiş kürkten imalat \\
& Menşe riski S2 / S3; c\_RoO (\%); p\_uyum & D/D; 1,5 / 2,0; 0,92 / 0,88 \etiketH{} \\
& Değer zincirindeki rol; DİR bağımlılığı & nihai; düşük \\
& EUDR ve SKDM kapsamındaki AB'ye ihracat payı & \%0; \%0 \\
& PEM ortaklarına ihracat; kümülasyon maruziyeti; κ & \$2 mn; \$1 mn; 0,0 \etiketH{} \\
5 Türkiye tarafı & Uygulanan tarife 2026 (ithalat ağırlıklı): MFN; İGV dahil & \%3,7; \%28,8 \\
& S4'te AB menşeline tarife (c2) & \%30,1 \\
& Kritik girdiler: 2024 ithalatı, AB payı, ilk kaynaklar & --- \\
6 Kabiliyet & KS; sıra (ağırlık duyarlılığında P10--P90) & 60,7; 24 / 230 (3--94) \\
& KS bileşenleri (yüzdelik): RCA, PCI, yurt içi KD, Ar-Ge, tedarikçi, verimlilik, firma & 93, 22, 60, 25, 96, 96, 59 \\
& PCI (sınıf) & −0,59 \\
& Atlanabilir ürünler (yoğunluk üst çeyrek, PCI \textgreater{} sınıf) & yok \\
7 Stratejik önem & SS; en yüksek boyut; ağırlıklı ortalama & 0,0; --- 0; 0,0 \etiketH{} \\
8 Makro duyarlılık & TÜFE ağırlığı (doğrudan); tüketimde ithal pay & \%0,00; \%3 \\
& ÜFE vekili (imalat üretimindeki pay) & \%0,02 \\
& EDE, 1 puan tarife: tüm menşeler / yalnız AB (bp; ρ = 1) & 0,000 / 0,000 \\
& CKE: tarife (\$ mn / bp); mali (\$ mn / mlr TL); puan & 1.461,0; 3,89; 32 \\
9 Senaryolar & Senaryo tablosu & Tablo C.16b (aşağıda) \\
\end{longtable}

\endgroup

\Needspace{7\baselineskip}

\textbf{Tablo C.16b --- 14.20 Kürk: senaryo etkileri (T+5'te yıllık; P50, parantez içinde P10 / P90)}

\nopagebreak[4]

\begingroup\gbtablo\small

\begin{longtable}[]{@{}lrrr@{}}
\toprule\noalign{}
\textbf{Değişken} & \textbf{S2} & \textbf{S3} & \textbf{S4} \\
\midrule\noalign{}
\endhead
\bottomrule\noalign{}
\endlastfoot
AB'ye ihracat (\$ mn) & −1 (−3 / +1) & −1 (−3 / 0) & −2 (−3 / −1) \\
AB'ye ihracat (\%) & −3,1 (−7,4 / +2,0) & −4,2 (−8,2 / −0,8) & −5,7 (−9,4 / −2,4) \\
Toplam ihracat (\$ mn) & −1 (−2 / 0) & −1 (−2 / 0) & −2 (−3 / −1) \\
Nihai mal ithalatı (\$ mn) & 0 (−1 / 0) & 0 (−1 / 0) & −1 (−1 / 0) \\
Üretim (\$ mn) & −1 (−1 / 0) & −1 (−1 / 0) & −1 (−1 / 0) \\
İstihdam (kişi) & −5 (−12 / +1) & −7 (−13 / −3) & −8 (−13 / −4) \\
İstihdam (\%) & −0,3 (−0,8 / +0,1) & −0,4 (−0,8 / −0,2) & −0,5 (−0,8 / −0,3) \\
Cari denge (\$ mn) & 0 (−1 / 0) & −1 (−1 / 0) & −1 (−1 / 0) \\
TÜFE (bp) & 0,00 (0,00 / 0,00) & 0,00 (0,00 / 0,00) & +0,01 (+0,01 / +0,01) \\
\end{longtable}

\endgroup

\emph{Kaynak:} \etiketT{} f3e; toplamlar Faz 2B genel denge dağılımına kalibre. Yüzdeler 2024 düzeyine göre.

\hypertarget{nace_kartlar_13-15-c17-1431-uxf6rme-trikotaj-ve-tux131ux11f-iux15fi-kroux15fe-uxe7orap-imalatux131}{%
\subsection{C.17 14.31 Örme (trikotaj) ve tığ işi (kroşe) çorap imalatı}\label{nace_kartlar_13-15-c17-1431-uxf6rme-trikotaj-ve-tux131ux11f-iux15fi-kroux15fe-uxe7orap-imalatux131}}

\Needspace{7\baselineskip}

\textbf{Tablo C.17 --- 14.31: veri sayfası}

\nopagebreak[4]

\begingroup\gbtablo\small

\begin{longtable}[]{@{}
  >{\raggedright\arraybackslash}p{(\columnwidth - 4\tabcolsep) * \real{0.1340}}
  >{\raggedright\arraybackslash}p{(\columnwidth - 4\tabcolsep) * \real{0.3433}}
  >{\raggedright\arraybackslash}p{(\columnwidth - 4\tabcolsep) * \real{0.5227}}@{}}
\toprule\noalign{}
\begin{minipage}[b]{\linewidth}\raggedright
\textbf{Alan}
\end{minipage} & \begin{minipage}[b]{\linewidth}\raggedright
\textbf{Gösterge}
\end{minipage} & \begin{minipage}[b]{\linewidth}\raggedright
\textbf{Değer}
\end{minipage} \\
\midrule\noalign{}
\endhead
\bottomrule\noalign{}
\endlastfoot
1 Kimlik & Rev. 2.1 karşılığı; CPA & 14.10 (n:1); CPA 14.31 \\
& Başlıca HS4 (ihracat, 2024) & 6115 \%100 \\
& BEC payları: ara / sermaye / tüketim & \%0 / \%0 / \%100 \\
2 Ölçek & Üretim; katma değer (GSYH payı) & \$1.642 mn; \$607 mn (\%0,04) \\
& Çalışan; girişim & 25.205; 1.367 \\
& Girişim ölçeği: 0 / 1--9 / 10--49 / 50--249 / 250+ çalışan & 45 / 1.202 / 150 / 66 / 18 \\
& 250+ girişimlerin istihdam payı; kişi başı KD (AB27 = 1) & \%45; 0,51 \\
& İstihdam değişimi, 2015--2024 & −\%4 \\
& İlk üç bölge (İBBS-2) & İstanbul \%39; Tekirdağ, Edirne, Kırklareli \%32; Kocaeli, Sakarya, Düzce, Bolu, Yalova \%15 \\
& Çapa yatırımcılar & Çorlu--Çerkezköy ve İstanbul çorap kümesi; Avrupa marka ve perakendecilerinin tedarikçileri \etiketH{} \\
3 Dış ticaret & İhracat 2024; 2020--2024 değişimi & \$1.296 mn; \%17 \\
& AB payı; ilk beş AB pazarı & \%65; Almanya \%32; Fransa \%20; Hollanda \%9; İspanya \%7; İsveç \%6 \\
& İlk beş AB dışı pazar & Birleşik Krallık \%42; ABD \%10; İsviçre \%6; Kazakistan \%3; Norveç \%2 \\
& İthalat 2024; AB ve Çin payı & \$59 mn; AB \%19, Çin \%50 \\
& İlk beş ithalat kaynağı & Çin \%50; İtalya \%7; Bangladeş \%6; Pakistan \%3; Gürcistan \%3 \\
& Dış ticaret dengesi & +\$1.237 mn \\
& İthal girdi içeriği (I-O); KD ÷ üretim & \%23; \%37 \\
& RCA: dünya 2020 → 2024; AB pazarı 2024 & 9,13 → 7,53; 6,47 \\
& Birim değer ÷ AB pazarı ortalaması & 1,29 \\
& AB pazarındaki pay (AB dışı tedarik içinde) 2020 → 2024; Türkiye'nin sırası & \%28,2 → \%25,8; 2. \\
4 AB tarafı & AB'ye ihracat 2025; rejim payları (GB / İTÜ / tarım) & €656,9 mn; \%100 / \%0 / \%0 \\
& S4'te AB tarifesi: 1/98 sürerse / askıda / tam MFN ortalaması & \%12,0 / \%12,0 / \%12,0 \\
& Tepe satırlar (S4 artışı ≥ \%10) & 13 satır, AB'ye ihracatın \%100'i: KN 6115 95 00 \%12,0 (€433,0 mn); KN 6115 96 99 \%12,0 (€123,3 mn); KN 6115 29 00 \%12,0 (€21,0 mn) \\
& AB pazarında AB dışı rakipler (AB tercihi) & Çin \%43 (MFN/GSP); Pakistan \%11 (GSP+); Sırbistan \%4 (MFN/GSP); Vietnam \%3 (STA/EPA); Endonezya \%2 (MFN/GSP) \\
& Tercihli rakiplerin payı; AME (puan; ham) & \%21; 99 (9,44) \\
& Menşe kuralı (PEM) & Fasıl 61 (biçimine göre örülmüş): eğirme ya da filament çekimi + örme veya örme ile konfeksiyon tek işlemde (ithal iplik serbest) \\
& Menşe riski S2 / S3; c\_RoO (\%); p\_uyum & D/D; 1,5 / 2,0; 0,92 / 0,88 \etiketH{} \\
& Değer zincirindeki rol; DİR bağımlılığı & nihai; orta \\
& EUDR ve SKDM kapsamındaki AB'ye ihracat payı & \%0; \%0 \\
& PEM ortaklarına ihracat; kümülasyon maruziyeti; κ & \$83 mn; \$43 mn; 0,0 \etiketH{} \\
5 Türkiye tarafı & Uygulanan tarife 2026 (ithalat ağırlıklı): MFN; İGV dahil & \%11,9; \%50,0 \\
& S4'te AB menşeline tarife (c2) & \%47,7 \\
& Kritik girdiler: 2024 ithalatı, AB payı, ilk kaynaklar & HS4 5402, 5205: \$1.653 mn (AB \%7; Çin \%36, Özbekistan \%18, Hindistan \%15) \\
6 Kabiliyet & KS; sıra (ağırlık duyarlılığında P10--P90) & 51,2; 104 / 230 (32--167) \\
& KS bileşenleri (yüzdelik): RCA, PCI, yurt içi KD, Ar-Ge, tedarikçi, verimlilik, firma & 71, 4, 77, 25, 96, 62, 51 \\
& PCI (sınıf) & −1,44 \\
& Atlanabilir ürünler (yoğunluk üst çeyrek, PCI \textgreater{} sınıf) & yok \\
7 Stratejik önem & SS; en yüksek boyut; ağırlıklı ortalama & 0,0; --- 0; 0,0 \etiketH{} \\
8 Makro duyarlılık & TÜFE ağırlığı (doğrudan); tüketimde ithal pay & \%0,00; \%3 \\
& ÜFE vekili (imalat üretimindeki pay) & \%0,25 \\
& EDE, 1 puan tarife: tüm menşeler / yalnız AB (bp; ρ = 1) & 0,003 / 0,001 \\
& CKE: tarife (\$ mn / bp); mali (\$ mn / mlr TL); puan & 4.089,9; 13,72; 56 \\
9 Senaryolar & Senaryo tablosu & Tablo 5.7 \\
\end{longtable}

\endgroup

\hypertarget{nace_kartlar_13-15-c18-1439-uxf6rme-trikotaj-ve-tux131ux11f-iux15fi-kroux15fe-diux11fer-giyim-eux15fyasux131-imalatux131}{%
\subsection{C.18 14.39 Örme (trikotaj) ve tığ işi (kroşe) diğer giyim eşyası imalatı}\label{nace_kartlar_13-15-c18-1439-uxf6rme-trikotaj-ve-tux131ux11f-iux15fi-kroux15fe-diux11fer-giyim-eux15fyasux131-imalatux131}}

\Needspace{7\baselineskip}

\textbf{Tablo C.18 --- 14.39: veri sayfası}

\nopagebreak[4]

\begingroup\gbtablo\small

\begin{longtable}[]{@{}
  >{\raggedright\arraybackslash}p{(\columnwidth - 4\tabcolsep) * \real{0.1332}}
  >{\raggedright\arraybackslash}p{(\columnwidth - 4\tabcolsep) * \real{0.3357}}
  >{\raggedright\arraybackslash}p{(\columnwidth - 4\tabcolsep) * \real{0.5311}}@{}}
\toprule\noalign{}
\begin{minipage}[b]{\linewidth}\raggedright
\textbf{Alan}
\end{minipage} & \begin{minipage}[b]{\linewidth}\raggedright
\textbf{Gösterge}
\end{minipage} & \begin{minipage}[b]{\linewidth}\raggedright
\textbf{Değer}
\end{minipage} \\
\midrule\noalign{}
\endhead
\bottomrule\noalign{}
\endlastfoot
1 Kimlik & Rev. 2.1 karşılığı; CPA & 14.10 (n:1); CPA 14.39 \\
& Başlıca HS4 (ihracat, 2024) & 6110 \%100 \\
& BEC payları: ara / sermaye / tüketim & \%0 / \%0 / \%100 \\
2 Ölçek & Üretim; katma değer (GSYH payı) & \$1.910 mn; \$552 mn (\%0,04) \\
& Çalışan; girişim & 32.975; 2.428 \\
& Girişim ölçeği: 0 / 1--9 / 10--49 / 50--249 / 250+ çalışan & 51 / 2.154 / 286 / 101 / 18 \\
& 250+ girişimlerin istihdam payı; kişi başı KD (AB27 = 1) & \%29; 0,38 \\
& İstihdam değişimi, 2015--2024 & \%34 \\
& İlk üç bölge (İBBS-2) & İstanbul \%59; Hatay, Kahramanmaraş, Osmaniye \%9; Mardin, Batman, Şırnak, Siirt \%8 \\
& Çapa yatırımcılar & İstanbul triko kümesi; Hatay ve Mardin--Batman'daki fason atölyeler \etiketH{} \\
3 Dış ticaret & İhracat 2024; 2020--2024 değişimi & \$2.719 mn; \%18 \\
& AB payı; ilk beş AB pazarı & \%66; Almanya \%30; İspanya \%14; Polonya \%14; Hollanda \%13; Fransa \%11 \\
& İlk beş AB dışı pazar & Birleşik Krallık \%27; ABD \%9; Kazakistan \%7; İsviçre \%7; Rusya \%5 \\
& İthalat 2024; AB ve Çin payı & \$459 mn; AB \%27, Çin \%22 \\
& İlk beş ithalat kaynağı & Çin \%22; Bangladeş \%20; İtalya \%11; İspanya \%8; Kamboçya \%6 \\
& Dış ticaret dengesi & +\$2.260 mn \\
& İthal girdi içeriği (I-O); KD ÷ üretim & \%23; \%29 \\
& RCA: dünya 2020 → 2024; AB pazarı 2024 & 4,35 → 3,60; 2,65 \\
& Birim değer ÷ AB pazarı ortalaması & 1,33 \\
& AB pazarındaki pay (AB dışı tedarik içinde) 2020 → 2024; Türkiye'nin sırası & \%12,0 → \%10,6; 3. \\
4 AB tarafı & AB'ye ihracat 2025; rejim payları (GB / İTÜ / tarım) & €1.306,1 mn; \%100 / \%0 / \%0 \\
& S4'te AB tarifesi: 1/98 sürerse / askıda / tam MFN ortalaması & \%12,0 / \%12,0 / \%12,0 \\
& Tepe satırlar (S4 artışı ≥ \%10) & 15 satır, AB'ye ihracatın \%100'i: KN 6110 30 99 \%12,0 (€416,1 mn); KN 6110 20 91 \%12,0 (€349,8 mn); KN 6110 20 99 \%12,0 (€331,2 mn) \\
& AB pazarında AB dışı rakipler (AB tercihi) & Çin \%33 (MFN/GSP); Bangladeş \%27 (EBA); Kamboçya \%7 (EBA); Pakistan \%4 (GSP+); Vietnam \%4 (STA/EPA) \\
& Tercihli rakiplerin payı; AME (puan; ham) & \%52; 99 (9,51) \\
& Menşe kuralı (PEM) & Fasıl 61: örme + konfeksiyon ya da tam kalıp örme (ithal iplik serbest) \\
& Menşe riski S2 / S3; c\_RoO (\%); p\_uyum & D/D; 1,5 / 2,0; 0,92 / 0,88 \etiketH{} \\
& Değer zincirindeki rol; DİR bağımlılığı & nihai; orta \\
& EUDR ve SKDM kapsamındaki AB'ye ihracat payı & \%0; \%0 \\
& PEM ortaklarına ihracat; kümülasyon maruziyeti; κ & \$186 mn; \$88 mn; 0,0 \etiketH{} \\
5 Türkiye tarafı & Uygulanan tarife 2026 (ithalat ağırlıklı): MFN; İGV dahil & \%12,0; \%51,0 \\
& S4'te AB menşeline tarife (c2) & \%50,9 \\
& Kritik girdiler: 2024 ithalatı, AB payı, ilk kaynaklar & HS4 5509, 5205, 5108: \$672 mn (AB \%4; Özbekistan \%45, Hindistan \%20, Türkmenistan \%9) \\
6 Kabiliyet & KS; sıra (ağırlık duyarlılığında P10--P90) & 45,7; 163 / 230 (80--199) \\
& KS bileşenleri (yüzdelik): RCA, PCI, yurt içi KD, Ar-Ge, tedarikçi, verimlilik, firma & 67, 1, 60, 25, 96, 32, 66 \\
& PCI (sınıf) & −1,98 \\
& Atlanabilir ürünler (yoğunluk üst çeyrek, PCI \textgreater{} sınıf) & yok \\
7 Stratejik önem & SS; en yüksek boyut; ağırlıklı ortalama & 0,0; --- 0; 0,0 \etiketH{} \\
8 Makro duyarlılık & TÜFE ağırlığı (doğrudan); tüketimde ithal pay & \%1,30; \%8 \\
& ÜFE vekili (imalat üretimindeki pay) & \%0,29 \\
& EDE, 1 puan tarife: tüm menşeler / yalnız AB (bp; ρ = 1) & 0,132 / 0,036 \\
& CKE: tarife (\$ mn / bp); mali (\$ mn / mlr TL); puan & 306,7; 32,51; 44 \\
9 Senaryolar & Senaryo tablosu & Tablo 5.8 \\
\end{longtable}

\endgroup

\hypertarget{nace_kartlar_13-15-c19-1511-derinin-tabaklanmasux131-ve-iux15flenmesi-kuxfcrkuxfcn-iux15flenmesi-ve-boyanmasux131}{%
\subsection{C.19 15.11 Derinin tabaklanması ve işlenmesi; kürkün işlenmesi ve boyanması}\label{nace_kartlar_13-15-c19-1511-derinin-tabaklanmasux131-ve-iux15flenmesi-kuxfcrkuxfcn-iux15flenmesi-ve-boyanmasux131}}

\Needspace{7\baselineskip}

\textbf{Tablo C.19 --- 15.11: veri sayfası}

\nopagebreak[4]

\begingroup\gbtablo\small

\begin{longtable}[]{@{}
  >{\raggedright\arraybackslash}p{(\columnwidth - 4\tabcolsep) * \real{0.1356}}
  >{\raggedright\arraybackslash}p{(\columnwidth - 4\tabcolsep) * \real{0.3591}}
  >{\raggedright\arraybackslash}p{(\columnwidth - 4\tabcolsep) * \real{0.5053}}@{}}
\toprule\noalign{}
\begin{minipage}[b]{\linewidth}\raggedright
\textbf{Alan}
\end{minipage} & \begin{minipage}[b]{\linewidth}\raggedright
\textbf{Gösterge}
\end{minipage} & \begin{minipage}[b]{\linewidth}\raggedright
\textbf{Değer}
\end{minipage} \\
\midrule\noalign{}
\endhead
\bottomrule\noalign{}
\endlastfoot
1 Kimlik & Rev. 2.1 karşılığı; CPA & 15.11 (1:1); CPA 15.11 \\
& Başlıca HS4 (ihracat, 2024) & 4107 \%38, 4104 \%21, 4302 \%17, 4112 \%13 \\
& BEC payları: ara / sermaye / tüketim & \%100 / \%0 / \%0 \\
2 Ölçek & Üretim; katma değer (GSYH payı) & \$903 mn; \$201 mn (\%0,01) \\
& Çalışan; girişim & 8.899; 879 \\
& Girişim ölçeği: 0 / 1--9 / 10--49 / 50--249 / 250+ çalışan & 71 / 653 / 171 / 39 / 0 \\
& 250+ girişimlerin istihdam payı; kişi başı KD (AB27 = 1) & \%0; 0,38 \\
& İstihdam değişimi, 2015--2024 & −\%18 \\
& İlk üç bölge (İBBS-2) & Kocaeli, Sakarya, Düzce, Bolu, Yalova \%28; İstanbul \%21; Tekirdağ, Edirne, Kırklareli \%9 \\
& Çapa yatırımcılar & Tuzla, Çorlu, Menemen ve Gerede organize deri sanayi bölgeleri; KOBİ ağırlıklı \etiketH{} \\
3 Dış ticaret & İhracat 2024; 2020--2024 değişimi & \$268 mn; \%19 \\
& AB payı; ilk beş AB pazarı & \%40; İtalya \%26; Almanya \%21; İspanya \%14; Portekiz \%12; Fransa \%4 \\
& İlk beş AB dışı pazar & Hindistan \%20; Güney Kore \%11; Çin \%11; Dominik Cumhuriyeti \%8; Rusya \%4 \\
& İthalat 2024; AB ve Çin payı & \$138 mn; AB \%51, Çin \%1 \\
& İlk beş ithalat kaynağı & İtalya \%36; Hindistan \%11; Tacikistan \%5; Azerbaycan \%5; İspanya \%4 \\
& Dış ticaret dengesi & +\$131 mn \\
& İthal girdi içeriği (I-O); KD ÷ üretim & \%23; \%22 \\
& RCA: dünya 2020 → 2024; AB pazarı 2024 & 1,58 → 1,58; 1,66 \\
& Birim değer ÷ AB pazarı ortalaması & 1,41 \\
& AB pazarındaki pay (AB dışı tedarik içinde) 2020 → 2024; Türkiye'nin sırası & \%5,6 → \%6,6; 5. \\
4 AB tarafı & AB'ye ihracat 2025; rejim payları (GB / İTÜ / tarım) & €89,4 mn; \%100 / \%0 / \%0 \\
& S4'te AB tarifesi: 1/98 sürerse / askıda / tam MFN ortalaması & \%3,7 / \%3,7 / \%3,7 \\
& Tepe satırlar (S4 artışı ≥ \%10) & yok \\
& AB pazarında AB dışı rakipler (AB tercihi) & Brezilya \%19 (MFN/GSP); Çin \%9 (MFN/GSP); ABD \%9 (MFN/GSP); Hindistan \%8 (MFN/GSP); Güney Afrika \%6 (STA/EPA) \\
& Tercihli rakiplerin payı; AME (puan; ham) & \%32; 73 (2,50) \\
& Menşe kuralı (PEM) & 4104--4106: yeniden tabaklama ya da başlık değişikliği; 4107: başlık değişikliği (ithal wet-blue'dan bitmiş deri menşe kazanır) \\
& Menşe riski S2 / S3; c\_RoO (\%); p\_uyum & D/D; 1,5 / 2,0; 0,92 / 0,88 \etiketH{} \\
& Değer zincirindeki rol; DİR bağımlılığı & ara; orta \\
& EUDR ve SKDM kapsamındaki AB'ye ihracat payı & \%0; \%0 \\
& PEM ortaklarına ihracat; kümülasyon maruziyeti; κ & \$14 mn; \$8 mn; 0,0 \etiketH{} \\
5 Türkiye tarafı & Uygulanan tarife 2026 (ithalat ağırlıklı): MFN; İGV dahil & \%3,9; \%10,8 \\
& S4'te AB menşeline tarife (c2) & \%11,8 \\
& Kritik girdiler: 2024 ithalatı, AB payı, ilk kaynaklar & HS4 4101, 4104, 4105: \$92 mn (AB \%15; Irak \%10, İtalya \%9, Türkmenistan \%8) \\
6 Kabiliyet & KS; sıra (ağırlık duyarlılığında P10--P90) & 45,1; 168 / 230 (93--194) \\
& KS bileşenleri (yüzdelik): RCA, PCI, yurt içi KD, Ar-Ge, tedarikçi, verimlilik, firma & 56, 14, 45, 25, 96, 32, 78 \\
& PCI (sınıf) & −0,86 \\
& Atlanabilir ürünler (yoğunluk üst çeyrek, PCI \textgreater{} sınıf) & yok \\
7 Stratejik önem & SS; en yüksek boyut; ağırlıklı ortalama & 0,0; --- 0; 0,0 \etiketH{} \\
8 Makro duyarlılık & TÜFE ağırlığı (doğrudan); tüketimde ithal pay & \%0,00; \%19 \\
& ÜFE vekili (imalat üretimindeki pay) & \%0,14 \\
& EDE, 1 puan tarife: tüm menşeler / yalnız AB (bp; ρ = 1) & 0,008 / 0,004 \\
& CKE: tarife (\$ mn / bp); mali (\$ mn / mlr TL); puan & 3.907,5; 15,64; 56 \\
9 Senaryolar & Senaryo tablosu & Tablo C.19b (aşağıda) \\
\end{longtable}

\endgroup

\Needspace{7\baselineskip}

\textbf{Tablo C.19b --- 15.11 Deri tabaklama: senaryo etkileri (T+5'te yıllık; P50, parantez içinde P10 / P90)}

\nopagebreak[4]

\begingroup\gbtablo\small

\begin{longtable}[]{@{}lrrr@{}}
\toprule\noalign{}
\textbf{Değişken} & \textbf{S2} & \textbf{S3} & \textbf{S4} \\
\midrule\noalign{}
\endhead
\bottomrule\noalign{}
\endlastfoot
AB'ye ihracat (\$ mn) & −9 (−22 / +2) & −12 (−24 / −2) & −16 (−27 / −7) \\
AB'ye ihracat (\%) & −8,8 (−20,9 / +2,0) & −11,5 (−22,1 / −2,1) & −15,2 (−25,0 / −6,4) \\
Toplam ihracat (\$ mn) & −7 (−16 / +1) & −9 (−17 / −2) & −13 (−22 / −5) \\
Nihai mal ithalatı (\$ mn) & −4 (−12 / +5) & −5 (−13 / +3) & −12 (−21 / −2) \\
Üretim (\$ mn) & −3 (−6 / −1) & −4 (−6 / −1) & −2 (−4 / +1) \\
İstihdam (kişi) & −22 (−41 / −10) & −26 (−41 / −9) & −13 (−26 / +6) \\
İstihdam (\%) & −0,2 (−0,5 / −0,1) & −0,3 (−0,5 / −0,1) & −0,1 (−0,3 / +0,1) \\
Cari denge (\$ mn) & −2 (−5 / −1) & −3 (−5 / −1) & −1 (−3 / +1) \\
TÜFE (bp) & +0,01 (+0,01 / +0,01) & +0,01 (+0,01 / +0,01) & +0,04 (+0,03 / +0,05) \\
\end{longtable}

\endgroup

\emph{Kaynak:} \etiketT{} f3e; toplamlar Faz 2B genel denge dağılımına kalibre. Yüzdeler 2024 düzeyine göre.

\hypertarget{nace_kartlar_13-15-c20-1512-bavul-el-uxe7antasux131-ve-benzerleri-ile-sarauxe7lux131k-ve-koux15fum-takux131mux131-imalatux131}{%
\subsection{C.20 15.12 Bavul, el çantası ve benzerleri ile saraçlık ve koşum takımı imalatı}\label{nace_kartlar_13-15-c20-1512-bavul-el-uxe7antasux131-ve-benzerleri-ile-sarauxe7lux131k-ve-koux15fum-takux131mux131-imalatux131}}

\Needspace{7\baselineskip}

\textbf{Tablo C.20 --- 15.12: veri sayfası}

\nopagebreak[4]

\begingroup\gbtablo\small

\begin{longtable}[]{@{}
  >{\raggedright\arraybackslash}p{(\columnwidth - 4\tabcolsep) * \real{0.1361}}
  >{\raggedright\arraybackslash}p{(\columnwidth - 4\tabcolsep) * \real{0.3638}}
  >{\raggedright\arraybackslash}p{(\columnwidth - 4\tabcolsep) * \real{0.5001}}@{}}
\toprule\noalign{}
\begin{minipage}[b]{\linewidth}\raggedright
\textbf{Alan}
\end{minipage} & \begin{minipage}[b]{\linewidth}\raggedright
\textbf{Gösterge}
\end{minipage} & \begin{minipage}[b]{\linewidth}\raggedright
\textbf{Değer}
\end{minipage} \\
\midrule\noalign{}
\endhead
\bottomrule\noalign{}
\endlastfoot
1 Kimlik & Rev. 2.1 karşılığı; CPA & 15.12 (1:1); CPA 15.12 \\
& Başlıca HS4 (ihracat, 2024) & 4202 \%88, 9605 \%6, 4205 \%5, 4201 \%1 \\
& BEC payları: ara / sermaye / tüketim & \%9 / \%0 / \%91 \\
2 Ölçek & Üretim; katma değer (GSYH payı) & \$731 mn; \$234 mn (\%0,02) \\
& Çalışan; girişim & 16.310; 3.071 \\
& Girişim ölçeği: 0 / 1--9 / 10--49 / 50--249 / 250+ çalışan & 81 / 3.140 / 206 / --- / --- \\
& 250+ girişimlerin istihdam payı; kişi başı KD (AB27 = 1) & \%7; 0,18 \\
& İstihdam değişimi, 2015--2024 & \%38 \\
& İlk üç bölge (İBBS-2) & İstanbul \%67; İzmir \%5; Malatya, Elazığ, Bingöl, Tunceli \%4 \\
& Çapa yatırımcılar & Derimod, Desa (deri çanta); İstanbul çanta kümesi \etiketH{} \\
3 Dış ticaret & İhracat 2024; 2020--2024 değişimi & \$377 mn; \%82 \\
& AB payı; ilk beş AB pazarı & \%37; İtalya \%22; Almanya \%19; Fransa \%16; Hollanda \%11; Danimarka \%4 \\
& İlk beş AB dışı pazar & ABD \%17; Suudi Arabistan \%12; Birleşik Krallık \%9; BAE \%5; Azerbaycan \%4 \\
& İthalat 2024; AB ve Çin payı & \$668 mn; AB \%51, Çin \%29 \\
& İlk beş ithalat kaynağı & Çin \%29; İtalya \%22; Fransa \%12; Malezya \%5; Kamboçya \%4 \\
& Dış ticaret dengesi & −\$292 mn \\
& İthal girdi içeriği (I-O); KD ÷ üretim & \%23; \%32 \\
& RCA: dünya 2020 → 2024; AB pazarı 2024 & 0,33 → 0,37; 0,25 \\
& Birim değer ÷ AB pazarı ortalaması & 1,36 \\
& AB pazarındaki pay (AB dışı tedarik içinde) 2020 → 2024; Türkiye'nin sırası & \%0,9 → \%1,0; 11. \\
4 AB tarafı & AB'ye ihracat 2025; rejim payları (GB / İTÜ / tarım) & €115,6 mn; \%100 / \%0 / \%0 \\
& S4'te AB tarifesi: 1/98 sürerse / askıda / tam MFN ortalaması & \%3,9 / \%3,9 / \%3,9 \\
& Tepe satırlar (S4 artışı ≥ \%10) & yok \\
& AB pazarında AB dışı rakipler (AB tercihi) & Çin \%62 (MFN/GSP); Vietnam \%9 (STA/EPA); Hindistan \%7 (MFN/GSP); Kamboçya \%3 (EBA); Endonezya \%3 (MFN/GSP) \\
& Tercihli rakiplerin payı; AME (puan; ham) & \%23; 69 (2,35) \\
& Menşe kuralı (PEM) & Fasıl 42: başlık değişikliği ya da menşesiz malzeme ≤ \%50 \\
& Menşe riski S2 / S3; c\_RoO (\%); p\_uyum & D/D; 1,5 / 2,0; 0,92 / 0,88 \etiketH{} \\
& Değer zincirindeki rol; DİR bağımlılığı & nihai; orta \\
& EUDR ve SKDM kapsamındaki AB'ye ihracat payı & \%0; \%0 \\
& PEM ortaklarına ihracat; kümülasyon maruziyeti; κ & \$46 mn; \$23 mn; 0,0 \etiketH{} \\
5 Türkiye tarafı & Uygulanan tarife 2026 (ithalat ağırlıklı): MFN; İGV dahil & \%4,3; \%36,6 \\
& S4'te AB menşeline tarife (c2) & \%15,4 \\
& Kritik girdiler: 2024 ithalatı, AB payı, ilk kaynaklar & HS4 4107, 3921, 5407: \$842 mn (AB \%36; Çin \%20, Mısır \%19, İtalya \%12) \\
6 Kabiliyet & KS; sıra (ağırlık duyarlılığında P10--P90) & 37,6; 215 / 230 (153--226) \\
& KS bileşenleri (yüzdelik): RCA, PCI, yurt içi KD, Ar-Ge, tedarikçi, verimlilik, firma & 29, 5, 68, 25, 96, 2, 69 \\
& PCI (sınıf) & −1,39 \\
& Atlanabilir ürünler (yoğunluk üst çeyrek, PCI \textgreater{} sınıf) & 420500 (PCI −0,63, RCA 0,70); 420291 (PCI −1,12, RCA 0,61) \\
7 Stratejik önem & SS; en yüksek boyut; ağırlıklı ortalama & 0,0; --- 0; 0,0 \etiketH{} \\
8 Makro duyarlılık & TÜFE ağırlığı (doğrudan); tüketimde ithal pay & \%0,65; \%36 \\
& ÜFE vekili (imalat üretimindeki pay) & \%0,11 \\
& EDE, 1 puan tarife: tüm menşeler / yalnız AB (bp; ρ = 1) & 0,271 / 0,137 \\
& CKE: tarife (\$ mn / bp); mali (\$ mn / mlr TL); puan & 153,5; 28,09; 39 \\
9 Senaryolar & Senaryo tablosu & Tablo C.20b (aşağıda) \\
\end{longtable}

\endgroup

\Needspace{7\baselineskip}

\textbf{Tablo C.20b --- 15.12 Bavul ve çanta: senaryo etkileri (T+5'te yıllık; P50, parantez içinde P10 / P90)}

\nopagebreak[4]

\begingroup\gbtablo\small

\begin{longtable}[]{@{}lrrr@{}}
\toprule\noalign{}
\textbf{Değişken} & \textbf{S2} & \textbf{S3} & \textbf{S4} \\
\midrule\noalign{}
\endhead
\bottomrule\noalign{}
\endlastfoot
AB'ye ihracat (\$ mn) & −6 (−15 / +3) & −8 (−16 / −2) & −12 (−19 / −5) \\
AB'ye ihracat (\%) & −4,5 (−10,9 / +2,0) & −5,9 (−11,9 / −1,1) & −8,4 (−13,8 / −3,5) \\
Toplam ihracat (\$ mn) & −4 (−11 / +2) & −6 (−12 / −1) & −9 (−15 / −4) \\
Nihai mal ithalatı (\$ mn) & −6 (−19 / +13) & −9 (−22 / +7) & −23 (−42 / −5) \\
Üretim (\$ mn) & +1 (−11 / +9) & +3 (−9 / +11) & +14 (+1 / +27) \\
İstihdam (kişi) & +14 (−164 / +140) & +41 (−135 / +175) & +214 (+14 / +422) \\
İstihdam (\%) & +0,1 (−1,0 / +0,9) & +0,2 (−0,8 / +1,1) & +1,3 (+0,1 / +2,6) \\
Cari denge (\$ mn) & +1 (−8 / +7) & +2 (−7 / +9) & +11 (+1 / +21) \\
TÜFE (bp) & +0,25 (+0,21 / +0,29) & +0,47 (+0,40 / +0,55) & +1,81 (+1,56 / +2,06) \\
\end{longtable}

\endgroup

\emph{Kaynak:} \etiketT{} f3e; toplamlar Faz 2B genel denge dağılımına kalibre. Yüzdeler 2024 düzeyine göre.

\hypertarget{nace_kartlar_13-15-c21-1520-ayakkabux131-bot-terlik-vb-imalatux131}{%
\subsection{C.21 15.20 Ayakkabı, bot, terlik vb. imalatı}\label{nace_kartlar_13-15-c21-1520-ayakkabux131-bot-terlik-vb-imalatux131}}

\Needspace{7\baselineskip}

\textbf{Tablo C.21 --- 15.20: veri sayfası}

\nopagebreak[4]

\begingroup\gbtablo\small

\begin{longtable}[]{@{}
  >{\raggedright\arraybackslash}p{(\columnwidth - 4\tabcolsep) * \real{0.1353}}
  >{\raggedright\arraybackslash}p{(\columnwidth - 4\tabcolsep) * \real{0.3564}}
  >{\raggedright\arraybackslash}p{(\columnwidth - 4\tabcolsep) * \real{0.5083}}@{}}
\toprule\noalign{}
\begin{minipage}[b]{\linewidth}\raggedright
\textbf{Alan}
\end{minipage} & \begin{minipage}[b]{\linewidth}\raggedright
\textbf{Gösterge}
\end{minipage} & \begin{minipage}[b]{\linewidth}\raggedright
\textbf{Değer}
\end{minipage} \\
\midrule\noalign{}
\endhead
\bottomrule\noalign{}
\endlastfoot
1 Kimlik & Rev. 2.1 karşılığı; CPA & 15.20 (n:1); CPA 15.20 \\
& Başlıca HS4 (ihracat, 2024) & 6403 \%41, 6402 \%29, 6404 \%19, 6406 \%5 \\
& BEC payları: ara / sermaye / tüketim & \%0 / \%0 / \%100 \\
2 Ölçek & Üretim; katma değer (GSYH payı) & \$3.034 mn; \$765 mn (\%0,06) \\
& Çalışan; girişim & 50.596; 7.162 \\
& Girişim ölçeği: 0 / 1--9 / 10--49 / 50--249 / 250+ çalışan & 263 / 6.978 / 743 / 127 / 12 \\
& 250+ girişimlerin istihdam payı; kişi başı KD (AB27 = 1) & \%6; 0,38 \\
& İstihdam değişimi, 2015--2024 & \%5 \\
& İlk üç bölge (İBBS-2) & İstanbul \%42; Gaziantep, Adıyaman, Kilis \%9; İzmir \%8 \\
& Çapa yatırımcılar & Flo, Hotiç, Greyder (marka); İstanbul, Konya ve İzmir ayakkabı kümeleri \etiketH{} \\
3 Dış ticaret & İhracat 2024; 2020--2024 değişimi & \$1.176 mn; \%33 \\
& AB payı; ilk beş AB pazarı & \%33; Almanya \%21; İtalya \%11; Romanya \%11; Yunanistan \%9; Polonya \%8 \\
& İlk beş AB dışı pazar & Suudi Arabistan \%15; Irak \%13; Rusya \%6; Azerbaycan \%5; BAE \%4 \\
& İthalat 2024; AB ve Çin payı & \$1.720 mn; AB \%21, Çin \%32 \\
& İlk beş ithalat kaynağı & Çin \%32; Vietnam \%27; İtalya \%12; Endonezya \%11; Kamboçya \%2 \\
& Dış ticaret dengesi & −\$544 mn \\
& İthal girdi içeriği (I-O); KD ÷ üretim & \%23; \%25 \\
& RCA: dünya 2020 → 2024; AB pazarı 2024 & 0,67 → 0,62; 0,29 \\
& Birim değer ÷ AB pazarı ortalaması & 0,86 \\
& AB pazarındaki pay (AB dışı tedarik içinde) 2020 → 2024; Türkiye'nin sırası & \%1,3 → \%1,2; 14. \\
4 AB tarafı & AB'ye ihracat 2025; rejim payları (GB / İTÜ / tarım) & €303,4 mn; \%100 / \%0 / \%0 \\
& S4'te AB tarifesi: 1/98 sürerse / askıda / tam MFN ortalaması & \%9,9 / \%9,9 / \%9,9 \\
& Tepe satırlar (S4 artışı ≥ \%10) & 23 satır, AB'ye ihracatın \%32'i: KN 6404 19 90 \%16,9 (€28,1 mn); KN 6402 99 39 \%16,8 (€14,0 mn); KN 6402 99 50 \%16,8 (€9,3 mn) \\
& AB pazarında AB dışı rakipler (AB tercihi) & Çin \%35 (MFN/GSP); Vietnam \%29 (STA/EPA); Endonezya \%7 (MFN/GSP); Hindistan \%5 (MFN/GSP); Kamboçya \%3 (EBA) \\
& Tercihli rakiplerin payı; AME (puan; ham) & \%45; 92 (4,13) \\
& Menşe kuralı (PEM) & Fasıl 64: her başlıktan; yalnız 6406'daki saya + iç taban birleşimleri kullanılamaz \\
& Menşe riski S2 / S3; c\_RoO (\%); p\_uyum & D/D; 1,5 / 2,0; 0,92 / 0,88 \etiketH{} \\
& Değer zincirindeki rol; DİR bağımlılığı & nihai; orta \\
& EUDR ve SKDM kapsamındaki AB'ye ihracat payı & \%0; \%0 \\
& PEM ortaklarına ihracat; kümülasyon maruziyeti; κ & \$132 mn; \$63 mn; 0,0 \etiketH{} \\
5 Türkiye tarafı & Uygulanan tarife 2026 (ithalat ağırlıklı): MFN; İGV dahil & \%13,2; \%51,9 \\
& S4'te AB menşeline tarife (c2) & \%48,7 \\
& Kritik girdiler: 2024 ithalatı, AB payı, ilk kaynaklar & HS4 6406, 4107, 3921: \$470 mn (AB \%50; Çin \%31, İtalya \%17, Almanya \%13) \\
6 Kabiliyet & KS; sıra (ağırlık duyarlılığında P10--P90) & 39,4; 207 / 230 (142--222) \\
& KS bileşenleri (yüzdelik): RCA, PCI, yurt içi KD, Ar-Ge, tedarikçi, verimlilik, firma & 28, 4, 52, 25, 96, 33, 86 \\
& PCI (sınıf) & −1,46 \\
& Atlanabilir ürünler (yoğunluk üst çeyrek, PCI \textgreater{} sınıf) & 640610 (PCI −1,33, RCA 0,39) \\
7 Stratejik önem & SS; en yüksek boyut; ağırlıklı ortalama & 7,6; savunma 10; 2,0 \etiketH{} \\
8 Makro duyarlılık & TÜFE ağırlığı (doğrudan); tüketimde ithal pay & \%1,11; \%34 \\
& ÜFE vekili (imalat üretimindeki pay) & \%0,46 \\
& EDE, 1 puan tarife: tüm menşeler / yalnız AB (bp; ρ = 1) & 0,469 / 0,099 \\
& CKE: tarife (\$ mn / bp); mali (\$ mn / mlr TL); puan & 654,8; 54,65; 57 \\
9 Senaryolar & Senaryo tablosu & Tablo 5.9 \\
\end{longtable}

\endgroup

\hypertarget{nace_kartlar_13-15-ek-d--uxf6nceki-paruxe7alara-ve-fazlara-duxfczeltmeler}{%
\section{Ek D --- Önceki parçalara ve fazlara düzeltmeler}\label{nace_kartlar_13-15-ek-d--uxf6nceki-paruxe7alara-ve-fazlara-duxfczeltmeler}}

\textbf{D-4. EUDR kapsamı.} AB'nin 2026/2102 sayılı Yetki Devrine Dayanan Tüzüğü 17 Eylül 2026'da Resmî Gazete'de yayımlandı ve ertesi gün yürürlüğe girdi \etiketV{}. Tüzük, EUDR Ek I'den sığır derisini (ex 4101, ex 4104, ex 4107) ve bazı vulkanize kauçuk eşyayı (ex 4010, ex 4016) çıkardı. Kullanılmış lastiği ve tohumluk soyayı kapsam dışı bıraktı, ahşap oturma yerlerini daralttı. Hazır kahveyi (2101 11 00), dondurulmuş sığır dilini ve on dört palm türevini 30 Aralık 2027'den itibaren kapsama aldı. İlk parçada kullanılan EUDR listesi bu değişikliği içermiyordu. Liste güncellendi. İlk parçanın metnindeki EUDR paragrafları ve EUDR uyum masasının (E-10-03) kapsamı buna göre düzeltildi. Sınıf payları da değişti: tabaklamada (15.11) \%62'den sıfıra, kahvede (10.83) hazır kahvenin eklenmesiyle \%22 düzeyine. Tabaklamanın kural tabanlı arketipi ``Yeşil Mutabakat'a maruz'' (G) olmaktan çıktı. Bu parçada uzman düzeltmesiyle ikincil arketip olarak korundu, ama gerekçesi artık EUDR değil, alıcıların çevre özen yükümlülükleri (Bölüm 5.13).

\textbf{D-5. İlk parçanın S2 sonuçları ve imalat toplamları.} S2'de AB'ye ihracattaki toplam kayıp genel denge sonucuna bağlı ve sınıflara kısmi denge kayıp paylarıyla dağıtılıyor. Bu parçada giyimin menşe parametreleri varsayılan düşük riskten orta riske çekildi; bu, S2 kaybının bir kısmını giyime kaydırdı ve diğer sınıfların S2 kaybını çok az küçülttü. İlk parçanın S2 istihdam etkisi −286 kişi oldu (önceki sürümde −299). S4'te ise PEM kümülasyonu kanalı imalatın toplam istihdam kaybını büyüttü: medyan −46,1 bin kişi (önceki sürümde −41,2 bin), P10 −70,1 bin kişi. İlk parçanın belgesi bu sayılarla yeniden üretildi; metnindeki bütün sayılar veriden geldiği için ayrıca düzeltme gerekmedi.

\textbf{D-6. Faz 2'nin giyim istihdam tahmini.} Faz 2A ve 2C, aktif S4'te giyimde yaklaşık 59 bin kişilik net istihdam kaybı raporlamıştı. Bu tahmin tek bir esneklik değeriyle (σ = 5), genel denge kalibrasyonu olmadan, durağan bir yöntemle yapılmıştı. Faz 3'ün genel dengeye kalibre modelinde giyim bölümünün (14) S4 kaybı medyanda 29.021, dağılımın alt ucunda 47.388 kişi. Faz 2'nin tahmini bu dağılımın alt ucundan da büyük bir kayıp gösteriyordu ve Faz 3 kartlarıyla geçersiz. Faz 2'nin vergi takası hesabı (02a, Tablo 1.5a) ise geçerli ve bu parçanın defterinde kullanıldı.

\hypertarget{nace_kartlar_13-15-ek-e--varsayux131mlar-ve-veri-boux15fluklarux131}{%
\section{Ek E --- Varsayımlar ve veri boşlukları}\label{nace_kartlar_13-15-ek-e--varsayux131mlar-ve-veri-boux15fluklarux131}}

Bu parçada varsayım kaydına dört yeni madde eklendi ve beş veri boşluğu açıldı ya da güncellendi. VB-32, yani DTÖ bağlı oranları, bu parçada bağlayıcı olmadı. c2 uygulanan MFN oranlarını değiştirmiyor; giyimde de tarife satırlarının yalnızca \%3,7'si bağlı. Bağlı oranlar, MFN artışı önerilen ilk parçada yeniden ele alınacak.

\Needspace{7\baselineskip}

\textbf{Tablo E.1 --- Varsayım kaydına eklenen maddeler}

\nopagebreak[4]

\begingroup\gbtablo\small

\begin{longtable}[]{@{}
  >{\raggedright\arraybackslash}p{(\columnwidth - 2\tabcolsep) * \real{0.1445}}
  >{\raggedright\arraybackslash}p{(\columnwidth - 2\tabcolsep) * \real{0.8555}}@{}}
\toprule\noalign{}
\begin{minipage}[b]{\linewidth}\raggedright
\textbf{Kimlik}
\end{minipage} & \begin{minipage}[b]{\linewidth}\raggedright
\textbf{Varsayım}
\end{minipage} \\
\midrule\noalign{}
\endhead
\bottomrule\noalign{}
\endlastfoot
V-35 (güncellendi) & SS boyut puanları ve ağırlıkları \etiketH{} (veri/f3\_ss\_boyutlar.csv; 10--15 gözden geçirildi, diğerleri ilk tahmin) \\
V-40 & PEM kümülasyonu kanalı yalnız S4'te; kayıp = λ · κ · maruziyet, λ \textasciitilde{} üçgen(0,3; 0,6; 0,9), κ: kumaş 1, teknik ve diğer tekstil 0,3, diğerleri 0 \etiketH{}; genel denge toplamına eklenir \\
V-41 & tekstil, giyim ve deride menşe parametreleri revize PEM kurallarından (Karar 1/2023): giyimde p\_uyum 0,85--0,88 (Asya kumaşı payı \%10--15), kumaşta 0,92 \etiketH{} \\
V-42 & geri ödeme yasağı (PEM Ek I md. 16) bugün AB'ye A.TR ihracatında telafi edici vergi olarak uygulandığından S2'de ek maliyet sayılmadı \\
V-43 & KDV takası (E-13-04) giyim ve ayakkabının S4 TÜFE şokunu tasarım gereği siler; net mali maliyet Faz 2A merkez değeri (laspeyres, c2) \\
\end{longtable}

\endgroup

\Needspace{7\baselineskip}

\textbf{Tablo E.2 --- Açık veri boşlukları}

\nopagebreak[4]

\begingroup\gbtablo\small

\begin{longtable}[]{@{}
  >{\raggedright\arraybackslash}p{(\columnwidth - 2\tabcolsep) * \real{0.1360}}
  >{\raggedright\arraybackslash}p{(\columnwidth - 2\tabcolsep) * \real{0.8640}}@{}}
\toprule\noalign{}
\begin{minipage}[b]{\linewidth}\raggedright
\textbf{Kimlik}
\end{minipage} & \begin{minipage}[b]{\linewidth}\raggedright
\textbf{Veri boşluğu}
\end{minipage} \\
\midrule\noalign{}
\endhead
\bottomrule\noalign{}
\endlastfoot
VB-36 & firma düzeyinde kumaş menşei (giyim ihracatında PEM dışı kumaşın payı; p\_uyum'un kalibrasyonu) \\
VB-37 & PEM ortaklarında Türk kumaşının AB'ye giden üretimdeki gerçek payı (kümülasyon maruziyeti üst sınır; κ ve λ \etiketH{}) \\
VB-38 & tekstil ve giyimde 2025 istihdam daralmasının sınıf düzeyinde dağılımı (SGK yalnız bölüm düzeyinde) \\
VB-39 & deri eşya ve çantada KDV oranı (takasın kapsamı); hızlı modada sınır gecikmesinin maliyeti \\
VB-32 (sürüyor) & DTÖ bağlı oranları, satır düzeyi. 13--15'te bağlayıcı değil: c2 uygulanan MFN oranlarını değiştirmiyor, yalnız AB tercihini kaldırıyor (giyimde satırların \%3,7'si bağlı) \\
\end{longtable}

\endgroup

\begin{gbilerleme}

İLERLEME: Faz 3 parça 2 / 16 tamamlandı (54 / 230 C sınıfı kartı) --- SONRAKİ ADIM: Faz 3 parça 3 (16--18: ağaç, kağıt, basım); DEVAM bekleniyor.

\end{gbilerleme}
